% Edición española — Seitz 1903, II. Disputatio Luther-Eck (impresas 55-218)
% Compilar: xelatex edicion-espanola.tex (dos pasadas)
\documentclass[11pt,a4paper,twoside]{article}
\usepackage{fontspec}
\setmainfont{Libertinus Serif}
\usepackage[lmargin=3cm,rmargin=3cm,top=2.8cm,bottom=2.5cm,headheight=14pt]{geometry}
\usepackage{setspace}
\setstretch{1.08}
\setlength{\parskip}{0.4em}
\setlength{\parindent}{1.2em}
\usepackage[hang,flushmargin]{footmisc}
\usepackage{fancyhdr}
\usepackage[hidelinks]{hyperref}
% llamada con la etiqueta ORIGINAL del impreso + nota con etiqueta repetida
\newcommand{\apnota}[2]{\textsuperscript{#1)}\,\footnote{[#1] #2}}
\pagestyle{fancy}
\fancyhf{}
\fancyhead[LE]{II.\ Disputa Lutero--Eck}
\fancyhead[RO]{Seitz 1903}
\fancyfoot[LE,RO]{\thepage}
\begin{document}
\begin{titlepage}
\centering
{\LARGE El texto aut\'entico\\de la Disputa de Leipzig\par}
\vspace{1em}
{\large II.\ Disputa Lutero--Eck (4--14 de julio de 1519)\par}
\vspace{2em}
{\normalsize Traducci\'on al espa\~nol de Seitz 1903 \textperiodcentered{} impresas 55--218\par}
{\small (aparato cr\'itico en lat\'in/alem\'an)\par}
\vfill
{\small Edici\'on de trabajo \textperiodcentered{} \today\par}
\end{titlepage}
\noindent\textit{Fuente latina: Seitz, Der authentische Text der Leipziger Disputation (Berl\'in 1903), ejemplar Internet Archive. Traducci\'on cuerpo a cuerpo con la transcripci\'on; el aparato cr\'itico se conserva en su lengua original (siglas M\,U\,Uv\,Mc\,L\"o\,P\,R.). Anomal\'ias documentadas en \texttt{lotes/RESOLUCIONES.md}.}

\bigskip
II.

En el nombre del Señor. Amén\apnota{1}{Diese Vorrede findet sich nur in P, vgl. oben p. 4.}. En el año de la encarnación del mismo M. D. XIX, indicción séptima, el lunes día cuatro del mes de julio, en el año séptimo del pontificado del santísimo padre en Cristo y señor nuestro, el señor León por la divina providencia Papa décimo, a la hora séptima antes del mediodía, después de que durante algunos días, a saber desde el veintisiete de junio, se hubiese celebrado la disputa entre los egregios y preclaros varones y señores Juan Eck, canónigo de Eichstätt y procanciller de Ingolstadt, y Andrés Rodolfo Bodenstein Carlstadt, canónigo y archidiácono de Wittenberg, doctores en artes, en decretos y en teología, sobre la materia del libre albedrío: se inició la disputa entre los magníficos señores doctores D. Juan Eck, ya mencionado, y el señor Martín Lutero, agustino de Wittenberg, en presencia de nosotros, notarios públicos, y de los testigos abajo firmantes, del modo que sigue. Pero, como el egregio señor doctor Martín al comienzo de su disputa se remitió a las protestaciones anteriores de los mencionados doctores, por eso las mismas, tomadas de la disputa precedente, se transcriben a la letra e insertan al inicio de esta disputa\apnota{2}{Hier folgen in P die Protestationen Karlstadts und Ecks, vgl. oben p. 14 f. [W. A. 254]}.

El señor D. Martín comenzó de este modo\apnota{5}{U Mc haben hier: Incipit disputatio excellentium theologorum, Johannis Eccii et Martini Lutheri Augustiniani, quae coepta fuit quarta die Julii M. D. XIX. hora septima. Protestatio patris M. Lutheri [Mc doctoris Martini].}.

«En el nombre del Señor. Amén\apnota{6}{U ἀρχήν.}. Abrazo y sigo la protestación de ambos egregios señores, tanto\apnota{7}{fehlt in M} de Andrés Carlstadt como de Juan Eck. Añado solo esto: que por reverencia al sumo pontífice y a la Iglesia romana de buen grado\apnota{10}{M videns.} habría omitido esta materia, innecesaria y extrañamente odiosa, si no hubiese sido arrastrado a ella por la proposición\apnota{12}{M protestationem.} del egregio señor doctor Juan Eck.

Lamento también que no estén presentes aquellos a quienes más convenía\apnota{14}{M oportet.}; ellos, que en privado y en público tantas veces me han difamado con la acusación de herejía\apnota{15}{M in haereseos.}, ahora, cuando urge el conocimiento de la causa, se han sustraído; me refiero a los inquisidores de la depravación herética, que pospusieron la amonestación fraterna\apnota{17}{M Lücke, Mc paternam admonitionem.} y la doctrina a sus acusaciones.

[255] El señor doctor Eck\apnota{e3}{U Scopus. De potestate, immo de primatu Romani pontificis (Uv a. R. Additio). Eccii protestatio, Mc [rot] = U.}.

En tu nombre, dulce Jesús\apnota{e4}{fehlt in M.}.

Antes de descender a la arena, protesto ante vosotros, ilustrísimos, nobles, magníficos y excelentísimos señores, que todo lo que he de decir y he dicho está sometido, en primer lugar, al juicio de la primera sede y del señor que en ella se sienta, y después\apnota{e8}{M dehinc.} al de cualesquiera otros a quienes corresponde corregir a los que yerran y reducirlos al conocimiento de la verdad. Y, como el reverendo padre en su prefación, a modo de excusa, atestigua que de buen grado habría omitido\apnota{e10}{U omisisse.} «esta» materia por reverencia al sumo pontífice, si no hubiese sido arrastrado por mi proposición, recuerde el reverendo padre: si él mismo no hubiese negado antes, en el resolutorio anterior a los tiempos de Silvestre, que el pontífice romano\apnota{e13a}{M Romanam.} es superior a los demás\apnota{e13b}{M aliis esse.}, no habría sido necesario que yo formulase aquella decimotercera proposición. Y en cuanto a que en las actas ante el legado\apnota{e14}{M legatis.} de la sede apostólica se queja de que el bienaventurado pontífice Pelagio torció los escritos evangélicos\apnota{e15}{M dicta.} —él, sin embargo, más que ningún otro, entendió las palabras de Cristo según la intención de los santos padres—. En vano, pues, el reverendo padre se esfuerza en volver contra mí la causa, a la que él mismo dio ocasión no una sola vez; pero, cortados ya los rodeos\apnota{e18}{P reiectis.}, acometamos, con Dios por guía, lo principal de nuestro propósito. Reverendo padre: vuestra conclusión XIII, que contradice a la mía, confiesa que la Iglesia romana «es» superior a las demás según los frigidísimos decretos de los pontífices romanos, nacidos dentro de los últimos cuatrocientos años\apnota{e22}{M quadringenta annis.}; a los cuales decís que se oponen el texto de la sagrada Escritura y las historias aprobadas de mil y\apnota{e23}{et fehlt in M.} cien «años». Contra ella opongo así.

La monarquía y el principado único en la Iglesia de Dios son de derecho divino e instituidos por Cristo. Por lo cual ni el texto de la sagrada Escritura ni las historias aprobadas se le oponen, pues aquella Iglesia militante, que es como un solo cuerpo según la sentencia del divino Pablo, está instituida y hecha a imagen de la Iglesia triunfante, en la cual hay una sola monarquía

estando todos dispuestos por orden hasta una sola cabeza, a saber, Dios. Por lo cual también tal orden ha sido constituido por Cristo en la tierra, pues confiesa en Jn 5 [v. 19] que el Hijo no hace nada sino lo que viere hacer al Padre. Por lo cual no es del cielo quien rehúsa someterse a la cabeza, así como no es del cielo sino de Lucifer quien no quiere someterse a Dios. Todo lo cual puedo confirmar amplísimamente\apnota{34}{U confirmari possunt.}, sobre todo por aquella alma santa, el bienaventurado Dionisio Areopagita, en el libro sobre la jerarquía eclesiástica, donde dice: «Nuestra jerarquía, santamente dispuesta en los órdenes transmitidos por Dios, es conforme a las santas [256] y celestes jerarquías». Así Gregorio Nacianceno en el Apologético dice que los sacrosantos misterios se celebran a\apnota{2a}{fehlt in M.} imagen de la imagen celeste\apnota{2b}{P maginis.}; por ellos\apnota{2c}{M terque [gestrichen], Mc per quae.}, en efecto, en la tierra nos asociamos a los órdenes celestes. Pues qué monstruo sería que la Iglesia fuese acéfala\apnota{4}{M acevolam, U a. R. Id est, sine capite.}, cosa que casi todos los herejes intentaron, como insinúa san Cipriano a Rogaciano y Pupiano\apnota{5}{M Rogatianum et Cupianum, U rogationum et Puppianum.}, para así, debilitada la cabeza, poder inculcar impunemente sus errores y su veneno en las mentes de los hombres. Y esta fue la razón principal, junto con otras conexas, por la que el feliz\apnota{7}{M Parisiensis velut.} estudio parisiense condenó a Juan de Tournai, que negaba el primado de la Iglesia romana. Así también fue error de Wiclef\apnota{9}{M wiewerbiticus, Mc wicleffiticus.} que la Iglesia romana no es superior a las demás por ley evangélica\apnota{10}{U evangelii.}.

El doctor D. Martín.

Cuando el señor doctor argumenta que hay absolutamente\apnota{12}{M U omnino esse.} una cabeza universal de la Iglesia, hace muy bien. Y si hay alguien que por pacto privado\apnota{13a}{M a. R. praefato.} haya convenido con\apnota{13b}{fehlt in M.} su señoría defender lo contrario, salga al medio; a mí nada me toca.

El doctor D. Eck.

Puesto que el reverendo padre dice que nada le toca acerca de lo contrario de aquello que yo pretendía probar —que la monarquía en la Iglesia militante es de derecho divino, como también en la triunfante—, en lo cual yo\apnota{18}{fehlt in U.} le alabo, porque también en esto concuerda con el divino Juan en el Apocalipsis: «Vi la ciudad santa nueva, que descendía», etc. Pero\apnota{20a}{M ad, Mc at.}, acercándonos más al asunto: si la Iglesia militante no estuvo sin monarquía\apnota{20b}{U monarcha.}, querría oír quién fue o habría sido alguna vez ese otro monarca\apnota{21}{M ista monarchia.} sino el pontífice romano, o cuál\apnota{22}{M atque (statt ant quae).} otra primera sede sino la sede de Pedro y de sus sucesores, según aquello del bienaventurado Cipriano, quien en la epístola tercera\apnota{23}{M P U secunda.} a Cornelio, pontífice romano, escribiendo contra los novacianos que se acercaban a Roma con astucia, dice: «Después de esto, además, osan navegar, constituido para sí un pseudoobispo\apnota{25a}{P U cedo episcopo, M sedo episcopo.} por los herejes\apnota{25b}{W. A. setzt fälschlich hinzu Romam und Anm. „Romam fehlt"; corrigiert am Ende des 2. Bandes.}, y llevar cartas de los cismáticos y profanos a la cátedra de Pedro\apnota{26}{U. a. R. non ad Petris cabellum.} y a la Iglesia principal, de donde ha brotado la unidad sacerdotal, y no\apnota{27}{U ne.} piensan que esos son romanos, cuya fe fue alabada por el Apóstol, a los cuales la perfidia no puede tener acceso». Y lo mismo atestigua Jerónimo contra los luciferianos: «La salud de la Iglesia —dice— pende de la dignidad del sumo sacerdote; a quien, si no se le da una potestad exenta\apnota{31a}{M exsors.} y eminente sobre todos, se producirán\apnota{31b}{M tot in ecclesiis tot efficienter (so!)} en las iglesias tantos cismas cuantos sacerdotes». Y que aquel sumo sacerdote sea el romano consta en el mismo bienaventurado Jerónimo, en dos epístolas al papa Dámaso, de las cuales casi cada palabra hace a nuestro propósito\apnota{34b}{M a. R. facio.}. Pero, por brevedad, señalemos aquello: «Hablo con el sucesor del pescador y discípulo\apnota{35}{U discipuli.} de Cristo. Sin seguir ningún premio sino a Cristo, me asocio a tu beatitud, esto es, a la cátedra [257] de Pedro\apnota{1b}{M consocio.}; sé que sobre aquella piedra está fundada la Iglesia»; y más abajo: «Quienquiera que no recoge contigo, ese dispersa». De lo cual cualquier buen cristiano fácilmente colige y entiende\apnota{3b}{et intelligit fehlt in U.} que la unidad sacerdotal fluye del pontífice romano, y que aquella fue siempre la cátedra principal y antepuesta a todas las demás\apnota{4b}{M praelatis.}, y que ella es la piedra sobre la que, como dice Jerónimo que sabe\apnota{5b}{M se dicit.}, está fundada la Iglesia. O bien asigne el reverendo padre otro monarca\apnota{6a}{M aliam monarchiam.} de la Iglesia en los tiempos antiguos\apnota{6b}{M U assignet (M assignat) alium monarcham etc.}.

El señor D. Martín.

Confieso por entero la monarquía de la Iglesia militante, y que su cabeza no es un hombre sino Cristo mismo, y esto por autoridad divina: primero, 1 Cor 15 [v. 25]: «Conviene, pues\apnota{11}{fehlt in M U.}, que él reine hasta que ponga a todos los enemigos bajo sus pies». Y poco antes: «Después el fin, cuando entregue el reino a Dios Padre y destruya todo principado». Lo cual el bienaventurado\apnota{13ab}{fehlt in U.} Agustín, en el libro I\apnota{13bb}{libro I fehlt in M.} sobre la Trinidad, capítulo último\apnota{13cb}{fehlt in M, ultimo cap U.}, expone acerca del reino de Cristo del tiempo presente; a saber: que Cristo, cabeza de la Iglesia, nos transferirá por la fe —nosotros, que somos su reino— a\apnota{15}{ad P M Uv, per U.} la visión. Así en Mt último: «He aquí que yo estoy con vosotros hasta la consumación del siglo». Asimismo en Hch 9 [v. 4]: Pablo oyó del cielo: «Saulo, Saulo,

¿por qué me persigues?», donde de nuevo Agustín dice que la cabeza habla por sus miembros\apnota{18}{fehlt in M.}. Por lo cual de ningún modo hay que escuchar a quienes empujan a Cristo fuera de la Iglesia militante hacia la triunfante\apnota{20a}{Mc ecclesiam triumphantem.}, siendo así que es reino de la fe\apnota{20b}{M filii.}; esto es: que no vemos a nuestra cabeza y, sin embargo\apnota{21a}{M attamen.}, la tenemos, según aquello del Sal 122\apnota{21b}{P CXIX, M U 119}: «Allí se sentaron sillas en juicio sobre la casa de David»; muchas sillas, ciertamente\apnota{22a}{Uv sunt.}, en las que se sienta un solo Cristo\apnota{22b}{Mc scilicet Christus.}. Vemos las sillas, no al que se sienta o al rey. Después\apnota{24a}{M U proinde.}, viniendo a las autoridades del señor\apnota{24b}{M U egregii domini.} doctor\apnota{24c}{M veniendo ad auctoritates etc.}, cuando afirma que en la Iglesia militante\apnota{25a}{M triumphante, Mc militante.} hay un principado único por derecho divino e instituido por Cristo\apnota{25b}{M instituto.}, dice lo suyo, pero nada prueba. Pues su primera autoridad, la de Pablo, a saber\apnota{26}{fehlt in P U.} Ef 4 [v. 15], que dice que Cristo es cabeza de la Iglesia, hace por mí y contra él\apnota{27}{U valet.}, pues habla con toda certeza allí\apnota{28a}{M ibi loquatur.} de la Iglesia\apnota{28b}{M ecclesiae.} militante y llama\apnota{28c}{M U appellat.} a Cristo su cabeza. Y es también\apnota{29a}{M Est est.} la misma sentencia, contraria a esta\apnota{29b}{M U contraria huic.}, en 1 Cor 3 [v. 5; 1 v. 13]: «¿Qué es Apolo? ¿Qué Cefas? ¿Qué Pablo? ¿Está dividido Cristo?\apnota{30}{U add. etc.}», donde manifiestamente prohíbe otra cabeza que Cristo.

Su segunda autoridad\apnota{32a}{fehlt in M.}, la de Jn 5: «No puede\apnota{32b}{M poterit.} el Hijo hacer nada\apnota{32c}{M quicquam facere.} sino lo que\apnota{33}{fehlt in M.} viere hacer al Padre», no habla ni de la Iglesia militante ni de la triunfante, sino, según sentencia de todos los doctores, de su igualdad con el Padre, a saber\apnota{35a}{U sed, Uv scilicet.}: que nada puede hacer el Padre que no pueda igualmente el Hijo\apnota{35b}{M quod nihil faciat nec facere possit, quod non idem facit filius, U quod nihil facit aut possit facere pater, quin idem possit et filius.}. Omito\apnota{36a}{M 'Omitto' omitto.} lo que dijo\apnota{36b}{fehlt in M.}: que no es del cielo quien rehúsa someterse a la cabeza\apnota{36c}{U recusat.}, y que es de Lucifer «quien\apnota{37a}{M est, qui.} no quiera\apnota{37b}{U non vult.} someterse a Dios»; porque, así como lo precedente está mal\apnota{37c}{U male sunt.} asumido, así\apnota{38a}{M sic, Mc ita.} esto\apnota{38b}{U et haec.} está mal inferido.

La tercera, «la de Dionisio», nada dice contra nosotros, pues no negamos la jerarquía eclesiástica\apnota{40a}{M coelestem, Mc ecclesiae.}, sino que disputamos sobre la cabeza de la monarquía, no de la jerarquía\apnota{40b}{M monarchia, non hierarchia.}.

[258] La cuarta\apnota{1a}{U Quarto.}, que es\apnota{1b}{fehlt in M U.} la de Gregorio Nacianceno, se aduce\apnota{1c}{M dicitur Mc additur (adducitur U), al. inducitur.} porque por los sacrosantos misterios nos asociamos\apnota{2a}{M sociamur.} a los órdenes celestes\apnota{2b}{Mc coelestibus ordinibus, U sanctis ordinibus coelestibus.}, etc.; «cualquier gramático entiende» que nada suena ni de la monarquía ni «de la cabeza»\apnota{3}{M capite hierarchia, letztes Wort gestrichen.}. En cuanto a lo que añadió\apnota{4a}{M nam quod adducit}: que sería un monstruo si\apnota{4b}{M quod, Mc si.} la Iglesia fuese acéfala\apnota{4c}{M A scevola.}, lo confieso\apnota{4d}{M (Mc) 'fateor' fateor.}; pero esta cabeza ni el mismo\apnota{5a}{fehlt in M.} señor doctor puede dar\apnota{5b}{M dari.} otra que Cristo, lo cual pruebo evidentemente: porque si su cabeza, a la que llaman pontífice romano\apnota{7a}{M vocat romanum pontificem, U Romanum pontificem appellat.}, muere, como\apnota{7b}{Mc quia.} es hombre, ya la Iglesia queda acéfala\apnota{7c}{M acephola.}; pero si entretanto Cristo es cabeza de la Iglesia hasta que se elija otro, no es menor monstruo\apnota{8b}{M U est.} que Cristo «ceda» ante el pontífice vivo y suceda al muerto.

La quinta, la del bienaventurado Cipriano, que persiguió a los herejes\apnota{10b}{M insectatur, Mc insectatus.} porque, debilitada la cabeza, intentaron\apnota{11}{U moliti.} inculcar impunemente sus errores a los hombres, nada hace en absoluto al propósito, pues habla\apnota{12b}{M U de quolibet.} no de la cabeza romana, sino de cualquier\apnota{12c}{M U de quolibet.} cabeza de cualquier\apnota{13a2}{U cuiusque.} «episcopado». Y\apnota{13b2}{M Additur et, Mc a. R. Atamen (so!) et.} si el egregio\apnota{13c}{fehlt in M.} señor doctor quisiere atenerse a la autoridad «de Cipriano», en esta hora terminaremos la disputa\apnota{14a2}{fehlt in U.}, pues él nunca saludó a Cornelio, pontífice romano\apnota{14b}{U a. R. Eccius dixit non fiet.}, de otro modo que\apnota{15b}{U salutat nisi.} como queridísimo hermano. Además, al describir las elecciones y confirmaciones\apnota{16a}{M electionem et confirmationem.} de los obispos, prueba por muchas epístolas, eficacísimamente\apnota{16b}{M longissime, Mc efficacissime.} desde las divinas letras, que pertenecen a la plebe y a dos o tres obispos vecinos\apnota{17b}{M suos episcopos duos.}, como también fue sancionado en el sacratísimo concilio de Nicea\apnota{18b}{U Niceno concilio; M add. ut beatus Augustinus allegat, gestrichen, weil vom Disputator widerrufen.}. Más aún, el mismo bienaventurado mártir, como alega el bienaventurado Agustín en el libro 2, cap. 2, sobre el bautismo\apnota{19a}{Mc add. parvulorum.}, dice así\apnota{19b}{M dicens.}: «Ninguno de nosotros se constituye obispo de los obispos\apnota{20c}{M constituat.} ni con terror tiránico\apnota{21a2}{U errore.} obliga a sus colegas a la “necesidad de obedecerle"\apnota{21b2}{Mc obsequendam statt ad obsequendi.}\apnota{21c}{M addidit.}, ya que cada obispo tiene, por la licencia de su libertad y potestad, su propio arbitrio, como si no pudiera ser juzgado por otro\apnota{23b}{M posse.}, del mismo modo que “él tampoco puede juzgar a otro”. Sino esperemos todos\apnota{24a2}{M omnes.} el juicio de nuestro Señor Jesucristo\apnota{24b2}{iudicium domini nostri Jesu Christi fehlt in M, dafür eine Zeile frei.}».

En cuanto a lo que dice de que de la cátedra romana y de Pedro ha nacido\apnota{25b2}{fehlt in M.} la unidad sacerdotal\apnota{26b}{M dignitatem.}, lo admito de muy buen grado en cuanto atañe a la Iglesia occidental, pues en verdad también la Iglesia romana ha nacido de la jerosolimitana: esta es propiamente la matriz\apnota{28d}{U matrix proprie, M nur matrix.} de todas las Iglesias.

Pero tampoco vale la consecuencia: de la Iglesia romana ha nacido la unidad sacerdotal; luego ella misma es cabeza y señora de todas\apnota{30}{U add prima.}; de otro modo, concluiría insuperablemente que la jerosolimitana es cabeza y señora\apnota{31}{et dominam fehlt in M.} de todas.

La última autoridad, la del divino Jerónimo, no la aduce rectamente el egregio señor doctor, aun si en todo fuese\apnota{33a}{M esse, Mc esset.} verdadera autoridad de Jerónimo\apnota{33b}{M beati Hieronymi.}, porque el señor doctor pretende probar la potestad monárquica de la Iglesia romana por derecho divino e\apnota{34}{fehlt in U.} instituida por Cristo. Esto no lo contienen las palabras de Jerónimo, pues dice: a quien\apnota{35a}{fehlt in M}, si no se le da una potestad exenta y\apnota{35b}{fehlt in U.} eminente sobre todos\apnota{36a}{M eminus, Mc eminuens.}, nacerán\apnota{36b}{M efficenter (so!), Mc efficerentur, U in ecclesiis efficientur.} en la Iglesia tantos cismas cuantos sacerdotes\apnota{37a}{M add sunt.}. «Désele» (dice)\apnota{37b}{M add non habet.}, esto es: por derecho humano puede\apnota{37c}{M U posset.} hacerse, consintiéndolo todos los demás fieles. Pues tampoco yo niego esto\apnota{38}{M et ego non nego.}: si todos los fieles de todo el orbe consintieren\apnota{39a}{M omnes fideles, a. R. omnes totius orbis.} en que el romano\apnota{39b}{M Romanam.} o el parisiense\apnota{39c}{U pariensem.} o el magdeburgués o cualquier otro [259] fuese el primer pontífice «y sumo», que este, por reverencia a toda la Iglesia de los fieles que así lo consiente\apnota{2a}{M consentientes.}, haya de tenerse por sumo monarca\apnota{2b}{M unam monarchiam, a. R. summum monarchum (so!)}. Pero esto ni se ha hecho nunca, ni se hace ni se hará\apnota{3a}{M factum est nunquam neque fit neque fiet, a. R. = P U.}, pues hasta nuestros\apnota{3b}{M U usque ad nostra.} tiempos la Iglesia griega no ha consentido, y sin embargo no ha sido tenida\apnota{4a}{M U sit habita.} por herética. Y que esta es\apnota{4b}{M video esse.} la sentencia del divino Jerónimo\apnota{5a}{U sententiam divi [fehlt in M] Hieronymi.} lo pruebo por la epístola a Evagrio\apnota{5b}{M ephasum.}, donde dice: «Dondequiera que hubiere obispo, sea en Roma o en Gubbio, en Constantinopla o en Regio, en Alejandría o en Tanis, es del mismo mérito y del mismo sacerdocio. La potencia de las riquezas y la humildad de la pobreza lo hacen o más sublime o más inferior. Por lo demás, todos son sucesores de los apóstoles\apnota{9}{U a. R. Martinus: Das sie den Hiero: nicht lang verbrant haben.}». Esta epístola se recita en los decretos no fríos\apnota{10a}{M vidi.}, dist. XCIII, c. legimus. Lo mismo sobre la epístola a Tito\apnota{10b}{M deditum, Mc Titum.}: «Es, pues, lo mismo\apnota{11a}{M Sicut ergo est vor Idem, gestrichen.}\apnota{11b}{M ergo est.} el presbítero que el obispo. Y antes de que por instinto del diablo se hicieran banderías en la religión y se dijera entre los pueblos\apnota{12a}{M populo.}: “Yo soy de Pablo, yo de Apolo, yo de Cefas”\apnota{12b}{U Ego sum Pauli, ego sum Cephae.}, las Iglesias se gobernaban por el consejo común de los presbíteros\apnota{13}{M ecclesia gubernabatur.}». Pero después que cada uno juzgaba suyos a «aquellos» que bautizaba, en todo el orbe se decretó que uno elegido de entre los presbíteros fuese antepuesto\apnota{15}{M unus vir super presbyteros superponeretur, U = P, nur superponeretur.}. Y, aducidas las autoridades de la Escritura, dice al final: «Así como los presbíteros saben que por costumbre de la Iglesia están sujetos a aquel que les ha sido antepuesto, así los obispos sepan que son mayores que los presbíteros más por costumbre que por verdad de disposición divina\apnota{18a}{U dispositionis dominicae, W. A. Anm. 2 vermutet die richtige L. A.}\apnota{18b}{M presbyteres.}». Por eso lo que añadió el señor doctor —que por Jerónimo ha de entenderse el pontífice «sumo» romano, cuando dice: «Hablo con el sucesor del pescador y discípulo\apnota{20d}{U discipuli.} de Cristo, y a tu beatitud, esto es, a tu cátedra\apnota{21a}{U cathedrae Petri, M capite Petro, a. R. catedrae.}, me asocio\apnota{21b}{fehlt in M, U consotius.}; sé que sobre\apnota{21c}{M et super.} aquella piedra está fundada la Iglesia»— nada hace. No se sigue: «me asocio a esta Iglesia, luego esta es la única primera». Ni se sigue: «está fundada sobre\apnota{23}{M super.} la piedra, luego solo ella está fundada».

Se añade a esto el decreto del concilio africano, dist. XCIX, c. Primae: «El obispo de la primera sede —dice— no se llame príncipe de los sacerdotes ni sumo sacerdote ni nada semejante, sino solo obispo de la primera sede. Y ni siquiera el romano se llame pontífice universal». Por lo cual, si la monarquía del pontífice romano se sostuviese por derecho divino\apnota{28}{M Romana ierarchia summorum pontificum.}, todo esto sería herético, lo cual es temerario afirmar. Y al final oigamos al Señor mismo, que en Lc 22 [v. 24-26] dice: «Se suscitó entre ellos una contienda\apnota{62}{U contentio discipulorum.} sobre quién parecía ser el mayor. Y les dijo: Los reyes de las naciones se enseñorean de ellas\apnota{32}{U earum.}, y los que tienen potestad sobre ellas son llamados bienhechores\apnota{33}{M heretici.}; vosotros, en cambio, no así, sino que el mayor entre vosotros hágase como el menor\apnota{62}{M add. Ecce hoc (haec?) est mea sententia.}».

El D. D. Eck\apnota{35}{In M beginnt mit diesem Abschnitt eine neue Seite Bl. E 3a. Auf der vorangehenden findet sich ein andrer (von Eck widerrufener?) Anfang: „Reverendus pater debens;" dann folgt ein leerer Raum von etwa 4 Zeilen.}.

El reverendo padre desciende a la arena bastante bien provisto, teniendo su materia tejida en buen orden en un libelo impreso y por él compuesto [260]. Por lo cual vuestras ilustres\apnota{1a}{U illustrissimae.} dominaciones, magnificencias y dominaciones\apnota{1b}{M et magnificentiae statt magnificentiae et dominationes.} perdonarán a Eck, ocupado ya desde hace mucho tiempo en otros negocios, si no pudiere acumular al instante\apnota{62}{PU evestigio} de modo tan redondo y preciso cuantas\apnota{62}{P quantam.} cosas ya ha amontonado el reverendo padre. Pues ha venido a disputar, no a publicar un libro. Veamos, pues, por orden qué ha dicho el reverendo padre.

En primer lugar pretende probar que Cristo es cabeza de la Iglesia, lo cual, sin embargo, era superfluo, pues nadie presume negar esto sino quien sea Anticristo\apnota{8a}{fehlt in M.}. Me admiro\apnota{8b}{MU miror.}, con todo, vehementemente de que esto

no advierte, como promete en las actas ante el legado de la sede apostólica poder presentar\apnota{9a}{M se posse.} cierto jurista-teólogo\apnota{9b}{M iuristam theologum, U iuristatheologum.}, que puede haber varias cabezas subordinadas\apnota{10}{U esse posse, M om. posse.}; en lo cual la razón de la cabeza mística o también simbólica\apnota{11}{M et symbolici, U aut symbolici.} difiere de la razón de la cabeza natural. Por lo cual enseguida se probará que, además de la cabeza que es Cristo, debe buscarse otra cabeza en la Iglesia. Y no le favorece que aduzca al Apóstol en 1 Cor 3 [v 5; 1 v. 13]: «¿Está dividido Cristo?», etc.\apnota{13}{fehlt in U.}, pues, aunque Pablo mencione allí a Pedro, no dijo falsedad el bienaventurado Jerónimo en el libro 1 contra Joviniano, para mí columna 18: «Uno es elegido (habla de Pedro) para que, constituida la cabeza, se quite la ocasión del cisma\apnota{16}{M ut schismatis.}». Llama claramente a Pedro cabeza constituida en la Iglesia\apnota{17a}{U esse constitutum.}. Pero dejemos esto\apnota{17b}{M U facimus.}, para rechazar las soluciones a lo por nosotros aducido\apnota{18}{M refusuri.}.

En primer lugar, cuando respondió\apnota{19a}{M U respondet.} a aquello de Jn 5: «Pues el Hijo no puede\apnota{19b}{M U filius potest.} hacer nada sino lo que viere hacer al Padre», dijo que según los santos padres se expresa aquí\apnota{21}{fehlt in M.} la igualdad del Padre y del Hijo. Pero lea con más atención, se lo ruego, el reverendo padre al bienaventurado e inengañable\apnota{22}{M inadulabilem, Mc adulabilem.} padre Bernardo, en el libro 3 de la Consideración a Eugenio, donde, hablando de la forma de la Iglesia y probando que es de derecho divino, formó nuestro argumento\apnota{24a}{M formavit probationem, U rationem formavit.}, para mí columna 7\apnota{24b}{mihi col. VII fehlt in M.}: «Ni tengas por vil\apnota{25}{U reputans, Uv reputes.} esta forma porque está en la tierra: tiene su ejemplar en el cielo. Pues el Hijo no puede hacer nada sino lo que viere hacer al Padre. Sobre todo porque a él, bajo el nombre de Moisés, se le dijo: “Mira, hazlo todo\apnota{27}{M ne omnia.} según el ejemplar que te ha sido mostrado en el monte”. Esto lo había visto quien decía\apnota{28}{U dicebar.}: “Vi la ciudad santa\apnota{29}{U add. etc.}”. Pues creo que se dijo por la semejanza: que así como allí los serafines y querubines\apnota{30a}{M cherubim et seraphim.} y todos “los demás”\apnota{30b}{U ceteri qui; M ceteri, constitutique, gestrichen constituti.}, hasta los ángeles y arcángeles, se ordenan bajo una sola cabeza que es Dios, así también aquí, bajo un solo sumo pontífice, primados o patriarcas, obispos, presbíteros\apnota{32}{M Archipresbyteri statt episcopi, presbyteri; U Archiepiscopi, episcopi, presbyteri.} o abades y los restantes de este modo». Entonces\apnota{33}{fehlt in M.} «añade Bernardo»: «No ha de tenerse en poco lo que tiene a Dios\apnota{34}{U habeat.} por autor y trae del cielo su origen. ¿Quién ya no entiende\apnota{35}{U intelligat.} que esta jerarquía eclesiástica, según Bernardo, ha sido instituida por Cristo? ¿Y que, como Cristo\apnota{36a}{M U deus.} es cabeza en el cielo, así el sumo pontífice es cabeza\apnota{36b}{M U sit.} en la Iglesia militante, sin excluir\apnota{37}{M nec dicam.} en modo alguno a Cristo, de quien se confiesa vicario?». Pero a la racioncilla plebeya\apnota{38a}{U pleiam, Uv plebeiam.} que aduce\apnota{38b}{U induxit.} —que la Iglesia queda acéfala\apnota{38c}{M ad Scepolam.}, [261] muerto el papa, a no ser que ridículamente dijésemos que Cristo cede ante el papa vivo y sucede al papa muerto—, «racioncilla del todo ridícula», la cual\apnota{2}{M et quae.} en asunto tan serio, entre tantos varones insignes, apenas es digna de ser repetida\apnota{3}{M U enarrari.}, pues desde el principio dije que aquella cabeza simbólica difiere de la razón de la verdadera cabeza «natural»\apnota{4}{M naturalis 2 mal.} en algunas conveniencias. Y Cristo, cuyo reino permanece para siempre y cuyo sacerdocio es sempiterno\apnota{5}{U et cuius sempiternum est.}, ni cede ante el papa ni le sucede, pues a él le ha sido dada toda potestad en el cielo y en la tierra (Mt último), y, muerto el papa, del mismo modo el colegio de los cardenales, como, muerto el obispo, el cabildo, mantienen\apnota{8}{M U teneat, P Uv teneant.} aquellos derechos hasta que se elija un nuevo pontífice\apnota{9}{M delegatur.}.

En cuanto a lo que en segundo lugar dice el reverendo padre —que Cipriano habla de cualquier obispo, no del pontífice romano—, me admiro sobremanera, pues la inteligencia de lo dicho ha de tomarse de las causas de decirlo\apnota{12a}{M dicendis.}, y Cipriano en los lugares por mí\apnota{12b}{M a me.} aducidos increpa a quienes se apartaron de Cornelio, que ciertamente fue pontífice romano. Por eso sepa el reverendo padre que no me sacio con palabras desnudas\apnota{14a}{M nullis, Mc nudis.}\apnota{14b}{M satiare, Mc satiari.}, como solemos apacentar a los sofistas. Lo que aduce de Cipriano en contra, oponiéndolo, creo que lo hará con mayor abundancia. Pues que Cipriano llame hermano a Cornelio, y que los apóstoles fuesen hermanos, nadie lo ignora. Sin embargo, Pedro, como también su sucesor Cornelio, fue cabeza de los apóstoles, ápice y vértice, según sentencia del bienaventurado Dionisio, cap. 3 o 7, sobre los nombres divinos\apnota{64}{fehlt in M (Lücke).}. Lo escrito sobre la elección, tanto por\apnota{64}{fehlt in U (ist in Uv zugefügt).} el concilio de Nicea como por Cipriano, ni favorece ni estorba al asunto. Y mucho menos debió aducirse, después de Cipriano, a Agustín, lib. 2, cap. 2, sobre el bautismo de los párvulos, pues Agustín reprende la arrogancia y temeridad de quienes por ambición y soberbia se introducen en las prelaturas eclesiásticas: que no deben constituirse a sí mismos ni obligar a otros a constituirse, sino que cada prelado debe esperar a ser llamado, como Aarón.

En tercer lugar, para diluir\apnota{64}{M dilaturus.} otra sentencia de Cipriano, de la epístola tercera\apnota{26}{P M U ex secunda epistola.} a Cornelio, respondió que la unidad sacerdotal nació de la romana en la Iglesia occidental, no en la oriental. Calló aquí el reverendo padre que Cipriano antes había llamado a la Iglesia romana cátedra de Pedro\apnota{64}{M U petri cathedram.} y Iglesia principal. Pero lo que pueda su solución es manifiesto para quien mire la médula de las palabras, pues el reverendo padre, a modo de gramático, entendió el nacimiento de la unidad sacerdotal\apnota{32a}{M U intelligit.} en cuanto a los auspicios\apnota{32b}{U auspicium.} y al inicio\apnota{32c}{M initia.}, siendo así que Cipriano quiso explicar el nacimiento de la comisión,

de subordinación o de influjo: de modo que de un solo Pedro, como de cabeza, se haya derivado la jurisdicción a todos los demás. De otro modo no dará un solo sacerdote ni siquiera en Jerusalén. Callo que no le sirve la pequeña glosa añadida sobre la iglesia occidental, pues el bienaventurado Jerónimo, escribiendo desde oriente, llama al inicio de su carta por ello cismática a la iglesia oriental, la cual desgarra en menudos trozos\apnota{39a}{W. A. notiert: „per frustra alle Gesammtausgaben und Löscher"} la túnica inconsútil del Señor y tejida de arriba abajo, y\apnota{39b}{fehlt in U.} «las zorras (dice) exterminan la viña de Cristo», dando a entender Jerónimo, sin duda, aquello de que se quejaba la esposa en el Cantar [2 v. 15] [262]: «Cazadnos las raposillas, que destrozan\apnota{1a}{M demoliantur.} la viña». Calle, pues, os lo ruego\apnota{1b}{M quaero.}, el reverendo padre, y no nos insulte con griegos y orientales, los cuales, al apartarse de la iglesia Romana, se han hecho también desterrados de la fe cristiana. De ello se sigue de qué modo ha de entenderse la inferencia\apnota{4}{M accipiendo sit dilatio.}: «Es raíz, luego señora». Pues no hablamos de raíz de inicio o de tiempo, sino de raíz de influjo y de primacía.

En cuarto lugar se esfuerza el reverendo padre por zafarse de las palabras de Jerónimo y evitarlas: pues al sumo sacerdote se le atribuye\apnota{8}{M U tribuit.}, ciertamente, y se le da la suma dignidad, pero esto se haría por derecho humano. Pero ¿por qué entonces el bienaventurado Jerónimo llama a Dámaso sucesor del pescador y quiere ser asociado a la cátedra\apnota{10}{M cathedra.} de Pedro, citando aquella palabra divina de Mt 16: «Sé que sobre aquella piedra está fundada la Iglesia»? Lo cual no puede apropiarse a las demás iglesias del modo que razona Bernardo. Y ¡ay dolor!, con suma injuria de los cristianos\apnota{13}{M christiana.} hemos experimentado que las puertas del infierno han prevalecido contra la iglesia de Jerusalén, la de Antioquía, la de Alejandría, y añado también la de Bohemia. Lo cual, sin embargo, la verdad incorrupta de Cristo no lo permite respecto de la iglesia que está fundada sobre\apnota{15}{M super.} la piedra. Pero es verísimo en lo que toca a la fe\apnota{16}{U iis.} aquello que Jerónimo afirma en la misma carta: «dondequiera que estuviere el cuerpo, allí se congregarán también las águilas». Desechado el patrimonio por la mala prole, entre vosotros solos se conserva del todo incorrupta la autoridad de los padres. Pero esto lo probaremos\apnota{20a}{U probabimus, Uv probabimus.} aún con evidencia en lo principal: que la iglesia Romana\apnota{20b}{M Romanam ecclesiam.} obtiene el primado no por derecho humano, sino divino. Conviene, sin embargo, refutar lo que el venerable padre adujo en su favor a partir de Jerónimo. Primero, a Evagrio, donde dice que el obispo de Roma y el de Eugubio\apnota{23}{M U egubii.}, el de Constantinopla y el de Regio son del mismo mérito y sacerdocio. Esto lo sabíamos antes de que naciera Teognis\apnota{25a}{M diagonos.}. Y tampoco\apnota{25b}{Uv a. R. (von moderner Hand?) atqui, namque, sed.} el papado es un orden por encima del episcopado. Pues así también en otro lugar dice que los apóstoles fueron iguales, donde sin embargo no quitó al santo Pedro el primado. Pero lo que el reverendo padre me exigía con tanto ahínco, que no

buscara rodeos, cosa que nunca suelo hacer, él mismo\apnota{28}{U ipsum.} lo entona. Aduciendo el capítulo Legimus\apnota{29a}{M legem.} de la distinción 93 me lleva a un brocardo\apnota{29b}{M U procardum; W. A. erinnert an das französ. brocard.}, cosa que los señores canonistas y teólogos controvierten\apnota{30a}{fehlt in U, W. A. ergänzt aus Löscher = M commentantur, a. R. alii antevertunt.}: si el episcopado es un orden y\apnota{30b}{fehlt in M U.} añadido al sacerdocio, insigne por su propio carácter; sobre lo cual, por ser impertinente, nada decido de momento. Esto digo, sin embargo: me parece, salvo\apnota{32}{M solvo.} siempre mejor juicio, que no hubo tal confusión en la primitiva\apnota{33}{M primativa.} iglesia que el obispo no se distinguiera del sacerdote, pues los doce apóstoles superan a los setenta y dos discípulos; en testimonio de lo cual aduzco al bienaventurado Dionisio, más antiguo que Jerónimo y jerarca de la primitiva iglesia, el cual en su libro sobre la jerarquía eclesiástica\apnota{36}{U ecclesiae.} pone el episcopado entre los órdenes sagrados, y al supremo\apnota{37}{M supremam.} jerarca, y de qué modo debe ser consagrado; asintiendo\apnota{38a}{M assentient.} a él, opino que desde el inicio\apnota{38b}{M indicio.} mismo de la iglesia los obispos superaron al sacerdocio común.

[263] En quinto\apnota{1a}{P M U Tertio, also kommt dieser Fehler auf Ecks Rechnung.} lugar adujo el capítulo del concilio\apnota{1b}{M consilii.} Africano, distinción\apnota{1c}{U add. can. 1., in Uv gestrichen.} XCIX, donde el concilio\apnota{2}{M consilium.} prohíbe que el pontífice Romano deba llamarse universal, y que Cristo lo prohibió, Lc 22 [v. 24ss.]: «Los reyes de las gentes\apnota{3}{U add. etc.}». Respondo que es verdad que se prohibió el nombre fastuoso de obispo universal\apnota{4}{M episcopi universali (so!).}; no porque el pontífice Romano no haya sido tenido nunca por un verdadero cristiano como primero y sumo pontífice, sino porque el obispo propio peculiarmente de la iglesia Romana no es obispo propio de cada iglesia, sino primero, pues de otro modo\apnota{7}{M alioqui.} no se tributaría a los obispos inferiores el honor debido. Pero nada de crimen hay en que alguien llame universal al pontífice Romano, en el sentido de primero. Pero más rectamente se dirá que él no es obispo universal, sino obispo de la iglesia universal\apnota{10}{U dicetur non episcopus universalis, sed universalis ecclesiae episcopus.}, como vicario de Cristo. Que el Señor reprendiera la ambiciosa contienda de los apóstoles, cual se busca entre seglares, no quita el primado a la iglesia Romana. Pero que el bienaventurado Gregorio fue el primero en hacerlo y reconoció que por ello era prelado de la iglesia católica, para reconocerse siervo de los siervos\apnota{14}{fehlt in M.}, en lo cual procuremos no zaherir con injurias\apnota{15}{M cum vicii.} a los pontífices siguientes, sino impetrar de Dios con oraciones que lleguen a ser tales.

A la hora segunda de la tarde se continuará la disputa, proclamándolo el bedel público de la universidad\apnota{17a}{Die Zeitangabe fehlt in M U.}.

A la misma hora segunda de la tarde\apnota{17b}{U Hora secunda pomeridiana quarta die Julii 1519.}.

El D. D. Martín.

Mi primera respuesta, que mostré\apnota{21a}{M U ostendi ab apostolo.} por la autoridad de Pablo en 1 Cor 3, que estaba prohibido que los fieles se arrogaran a Cefas\apnota{21b}{M Sceve.} o a Pablo\apnota{22a}{M Pauli.} o a Apolo, la refutó así el egregio señor doctor\apnota{22b}{M Egregius doctor, U dominus doctor egregius.}: que aunque\apnota{22c}{M si.} Pablo mencione allí a Pedro\apnota{23}{M Petri ibidem.}, no por ello dice falsamente Jerónimo contra Joviniano: «uno solo es elegido, para que\apnota{24}{U quod.}, constituida la cabeza, se quite la ocasión de cisma». Claramente\apnota{25a}{M quare.} llama a Pedro cabeza\apnota{25b}{M caput Petrum.} constituida en la iglesia. Añadió: «pero dejemos pasar estas cosas».

Respondo. No tolero ser apartado, por una autoridad menor aducida\apnota{26}{M adductam.}, de una mayor\apnota{27}{M afelli.}. Ni es tan grande Jerónimo que por él abandonemos\apnota{28}{M U deseram.} a Pablo. No es, pues, que Pablo mencione allí a Pedro sin más, como intenta atenuar esta refutación; sino que con plena autoridad enseña y prohíbe que nadie diga que es\apnota{29}{M se dicat.} de Pedro. Por eso también concluye el mismo capítulo\apnota{30}{M capitulo.}: «Todo es vuestro, sea Pablo, sea Cefas\apnota{31}{M scevela (?).}, sea muerte, sea vida; vosotros, empero, de Cristo, y Cristo de Dios». Permanece, pues, esta respuesta\apnota{32}{responsio ista.} mía aún invicta. Y si no fuere refutada con más fuerza, la opongo a todos los argumentos pasados y futuros del señor doctor. Pues la palabra de Dios está sobre todas las palabras de los hombres. — En cuanto a Jerónimo, digo lo que también yo dejo pasar\apnota{35}{M missefacia.} estas cosas, pues es un pasaje ambigüísimo, como bien advirtió el señor doctor.

Mi segunda respuesta a la autoridad de Jn 5\apnota{37}{M 15.}, donde dije que Cristo habla de la igualdad con el poder paterno: me mandó el señor doctor leer con más atención [264] al bienaventurado Bernardo, que aduce esta autoridad en favor de la iglesia militante, como se ha oído. Respondo: venero al santo Bernardo y no desprecio su sentencia. Pero en la contienda ha de tomarse el sentido genuino y propio de la Escritura, que pueda sostenerse en la línea de batalla. Del cual los santos padres se apartan a veces para enriquecer el discurso, y sin culpa\apnota{5a}{et sine culpa fehlt in M (Lücke).}. Ahora\apnota{5b}{M Nunc, Mc setzt dafür Nunquam.}, en cambio, por el texto precedente y siguiente resulta claro que Cristo habla de su igual poder\apnota{7a}{M potentia Christi, U omnipotentia sua.} con el Padre; (el texto es)\apnota{7b}{textus est fehlt in M U.}: «Por eso perseguían los judíos\apnota{7c}{M Judaei persequebantur.} a Jesús, porque hacía\apnota{8a}{M U faciebat haec.} estas cosas en sábado\apnota{8b}{M sabbata.}». «Por eso\apnota{8c}{fehlt in U.}, pues, con más ahínco buscaban los judíos matarle\apnota{8d}{M Judaei interficere eum.}, porque no solo quebrantaba el sábado, sino que llamaba a Dios padre suyo, haciéndose igual a Dios». Respondió, pues,

Jesús y les dijo: «Amén, amén, os digo: no puede el Hijo hacer por sí nada sino lo que viere hacer al Padre». Y así queda claro que Bernardo trata esta palabra de Cristo en otro sentido\apnota{12}{fehlt in M.}.

A la tercera, a saber, la plebeya\apnota{13}{M blebeiam.}, «como dijo», y ridícula razoncilla, con la que dije que también sin papa la iglesia tiene su cabeza, opuso que no era digna de ser recontada\apnota{15a}{M narrari, U enarrari.} ante tantos varones y en asunto tan serio; respondo: Sea plebeya\apnota{15b}{M blebeia.} y ridícula, con tal de que sea invicta. Y aún no la veo refutada. Pues no entiendo: si la iglesia, muerto el papa, durante tres o cuatro meses no está acéfala\apnota{17}{M a scevola statt non est acephala.} mientras haya otros obispos, cómo va a estar acéfala si no hay ningún papa\apnota{18}{M modo quomodo non sit a sceevola si nullus est papa, U quomodo etiam non [in Uv gestrichen] sit acephala, cum nullus sit papa.}. Pues lo que adujo de los cardenales, en poder de los cuales está el derecho\apnota{19}{U vis.} de elegir etc., confirma mi respuesta, pues de ahí se sigue que en el tiempo en que aún no había cardenales, «como en tiempo de Jerónimo», no pudo haber papa\apnota{21}{M caput.}.

A la cuarta, la autoridad de Cipriano, donde dije que él habla de cualquier pontífice, objetó que por el texto se demuestra que habló del pontífice Romano Cornelio\apnota{24a}{U Cornelio Romano.}\apnota{24b}{fehlt in M.} contra Novaciano\apnota{24c}{M U Novatianos.}; respondo: Nada me importa, ni tengo en la memoria esa carta. Esto sé, en cambio: que el santo Cipriano en muchas cartas se ocupa todo en afirmar que a cada iglesia\apnota{26}{M episcopo (gestrichen) ecclesiae.} se le ordene su cabeza y su obispo por sufragio del pueblo y juicio de los obispos vecinos. Y por eso, si es como alega el señor doctor, que se dijo contra\apnota{28}{M de (gestrichen) contra.} los novacianos acerca de Cornelio\apnota{29a}{U dico.}, es cierto que habló de la «cabeza» de la iglesia Romana\apnota{29b}{M de Romana ecclesia, Mc = P.}, no de la iglesia universal. Asimismo refutó también aquello de que Cipriano llama siempre a Cornelio hermano\apnota{31}{M appellat.}, nunca en cambio señor, como ahora hacen los obispos, poniendo el relativo sin correlativo. Objetó que también Pedro tuvo por hermanos\apnota{32}{P apostolus.} a los apóstoles y fue sin embargo cabeza y vértice de los apóstoles, según Dionisio en los Nombres divinos\apnota{33}{U a. R. de divinis nominibus.}. Respondo: Si el egregio señor doctor pudiere probar que Pedro ordenó alguna vez a uno de los apóstoles, más aún, que envió a uno\apnota{35}{M illum illorum unquam, U ullum illorum unquam.} de los setenta discípulos, cedo\apnota{36a}{M concedo.} en todo y quiero\apnota{36b}{M volo volo st. et volo.} darme por vencido. Pero si yo probare que ni siquiera todos los apóstoles pudieron\apnota{37a}{M U potuerint.} enviar u ordenar a uno\apnota{37b}{nec unum.} de los apóstoles, ruego me conceda que Pedro no tuvo potestad alguna sobre los demás apóstoles. De lo cual se sigue que mucho menos potestad tiene el obispo, sucesor de Pedro\apnota{39}{M U Petro.}, sobre [265]

sobre los obispos, sucesores\apnota{1}{fehlt in M.} de los demás apóstoles. Ahora bien, clarísimo es el texto de Hch 1 [v. 23ss.], que Matías apóstol no pudo\apnota{3b}{M poterat.} ser ordenado por todo el coeto de los apóstoles y discípulos\apnota{3a}{et discipulorum fehlt in M.}, sino que fue elegido del cielo, como todos los demás, y ordenado\apnota{4}{M ordinatus et electus [est fehlt].} por Cristo. Así también en el capítulo 13 Saulo y Bernabé fueron asumidos\apnota{5}{U sint.} a la obra, segregándolos el Espíritu Santo. Es, pues, error apertísimo que Pedro tuviera potestad sobre los apóstoles. Esto, ciertamente, lo confieso: que el apóstol Pedro fue el primero en el número de los apóstoles y que se le debe prerrogativa de honor, pero no de potestad. Igualmente fueron elegidos y recibieron igual potestad. Así siento también del pontífice Romano: que por prerrogativa de honor debe ser antepuesto a los demás, salva la igual potestad de cada uno, y no como infiere Pelagio en su frigidísimo decreto: «Donde mayor autoridad, allí mayor potestad, y permanece para los demás\apnota{12a}{M ceteris.} la necesidad\apnota{12b}{M potestas.} de obedecer».

A la quinta\apnota{13}{M quintum.}, donde aporté la elección del obispo a partir de Cipriano y del concilio Niceno, el egregio señor doctor la despreció con bastante retórica, diciendo que ni promueven ni impiden el negocio. Pero con ello no se disuelve mi respuesta. Permanece, pues, aún el decreto Niceno; o si no permanece y establecieron contra el derecho divino, no será concilio católico, sino diabólico conciliábulo\apnota{17}{U appellandum erit.}. Así, como juzgaba\apnota{18a}{M confutaret statt cum putaret.} que no debía aducirse\apnota{18b}{M dicendum.} a Agustín, e interpretó con su hermosísima glosa a Cipriano, aducido por medio de Agustín, en el sentido de que Cipriano reprendía la ambición y soberbia de quienes se entrometían por sí mismos antes de ser llamados como Aarón: esto lo dijo por mera presunción. Por lo demás\apnota{21}{M veterum.}, el texto es bastante claro: que ningún obispo, que ya sea obispo, debe constituirse obispo de otros obispos\apnota{23a}{M se aliorum se constituere episcoporum (so!)}. Permanece, pues, aún\apnota{23b}{M U adhuc mea.} mi respuesta.

A la sexta me reprendió gravemente\apnota{25a}{M Sextam graviter, Mc ad sextam ubi graviter.} el egregio señor doctor\apnota{25b}{M U idem egregius dominus (fehlt in M) doctor.}, porque en la otra autoridad de Cipriano\apnota{26}{fehlt in M.} callé este vocablo («principal»); después, como a gramático, me colgó de la nariz\apnota{27a}{P sub naso, M me nasum.}, porque dije\apnota{27b}{M dixerit, Mc dixerim.} que la unidad sacerdotal\apnota{28a}{M unitatem sacerdotalem ortam esse.} había nacido de la sede de Pedro. Por eso este nuevo dialéctico\apnota{28b}{M dialectius.}, o más bien filósofo, interpreta este nacimiento como de «comisión, subordinación»\apnota{29}{M ordinationis.} o influjo. De otro modo, dice, ni desde Jerusalén dará un solo sacerdote. Respondo: Ya callara\apnota{31a}{U docuerim, W. A. obticuerim.}, ya dijera este vocablo («principal»)\apnota{31b}{M U principale.}, lo mismo

es. Pues no puede llamarse principal respecto de la iglesia oriental, como bastante se ha dicho.

Después, su invención\apnota{33a}{fehlt in M.} sobre el nacimiento por influjo la desprecio con la misma facilidad\apnota{33b}{fehlt in M (Lücke), Mc füllt diese Lücke durch contemptu aus.} con que\apnota{33c}{M quae.} él la fingió. Y no me es difícil dar\apnota{34}{mihi dare fehlt in M.} un solo sacerdote desde Jerusalén, a saber, el mismo Cristo, de quien comenzó\apnota{35}{M U add et hinc orta, in Uv gestrichen.} y procedió la iglesia, según aquello de Is 2 [v. 3]: «De Sión saldrá\apnota{36}{M ex Sion exivit.} la ley, y la palabra del Señor de Jerusalén». Ya lo que añadió con la autoridad de Jerónimo, que la iglesia oriental fue cismática y desgarró en menudos trozos\apnota{38b}{M minutatim domini tunicam.} la inconsútil\apnota{38a}{fehlt in M.} túnica del Señor, no sé en verdad qué quiere. Pues no puede decir que toda la iglesia oriental y siempre fue [266] cismática, como\apnota{1a}{U hinc, Uv sicut.} tampoco puede negar\apnota{1b}{M negari.} que también la iglesia latina tuvo alguna vez sus cismas y sin embargo permaneció iglesia. Por tanto, nada hay que me mande callar y no insultar por medio de la iglesia griega, porque al apartarse de la iglesia Romana se hicieron a la vez desterrados de la fe de Cristo. Yo más bien ruego al señor doctor Eck que, por la eckiana modestia\apnota{5}{M modestiam.} de que se jacta, perdone a tantos miles de santos, pues hasta nuestros tiempos duró la iglesia griega y sin duda dura hasta hoy y durará. Pues Cristo no recibió el centro de la tierra romana, sino todos los confines de la tierra en posesión\apnota{8a}{M possionem.} y heredad del Padre, Sal 2 [v. 8]\apnota{8b}{fehlt in M (Lücke).}.

Al séptimo\apnota{9}{M Septimum, U Ad septimam.}, que opuso desde Jerónimo sobre el sumo sacerdote, dijo que mi respuesta era una evasión. Por eso, para confirmar su dicho, añadió por qué el bienaventurado Jerónimo llama a Dámaso sucesor del pescador\apnota{11}{successorem dicat fehlt in M, Mc hat appellat successorem.} y quiere ser asociado a la cátedra de Pedro, citando aquella palabra divina de Mt 16: «Sé que sobre aquella piedra está fundada la Iglesia», lo cual\apnota{13}{fehlt in M.} no puede apropiarse a las demás iglesias. Y después llora la «caída»\apnota{14a}{M U comploravit.} de Jerusalén, de Antioquía\apnota{14b}{fehlt in M.}, de Alejandría y por fin también de Bohemia, y que, siendo testigo el mismo Jerónimo, solo entre los romanos se conservaba incorrupta la autoridad de los padres\apnota{15}{M et apud Romanos Pauli servaretur auctoritas (z. t. von Mc aufgefüllt); in Uv ist patrum gestrichen.}. Respondo, y ruego al egregio señor doctor quiera alegar los dichos de los padres según conciencia, no sea que en vez de teólogos «parezcamos» sofistas\apnota{18}{M sophisticare.}. Pues a sumo sacerdote en ese lugar Jerónimo llama a cualquier obispo, aun como elevado\apnota{19}{M assumptus.} del número de los demás sacerdotes. Por eso nada atañe propiamente al

pontífice Romano\apnota{20a}{proprie ad Romanum pontificem fehlt in M (Lücke); Mc setzt proprie ad reverendum patrem.}. Después, aquella\apnota{20b}{fehlt in M.} autoridad de Mt 16 no se apropia solo a la iglesia Romana\apnota{21}{M U Romanae ecclesiae.}, lo cual indican expresamente las palabras de Cristo, cuando\apnota{22a}{U cum.} dice: «mi iglesia». Cualquiera que sea, pues, iglesia de Cristo\apnota{22b}{fehlt in U, Uv fügt es bei.}, está edificada sobre la piedra, no solo\apnota{23}{M et non sola, U et non solum.} la Romana. O si esta palabra no conviene a las demás iglesias, la iglesia Romana\apnota{24}{M ecclesia Romana.} será la única y por consiguiente no la primera. Por lo cual la unidad de la iglesia no pende de la unidad del primado Romano, sino mucho mejor, según el Apóstol en Ef 4 [v. 5], de la unidad de la fe, del bautismo, del Señor, como también siente con frecuencia Cipriano en sus cartas. Ni permaneció solo entre los romanos\apnota{28a}{M sola apud 'Romanos' (das letzte Wort steht über der Zeile), U apud Romanos solos.} incorrupta la autoridad de los padres\apnota{28b}{M integritas, Mc auctoritas.}, salvo quizá en el tiempo en que escribía Jerónimo; más aún, cuentan las historias que Liberio, pontífice Romano, cedió ante los arrianos, y aquello escribe Jerónimo en los Varones ilustres: que Acacio de Cesarea, obispo arriano, discípulo\apnota{31}{M episcopus discipulos (so!)} del arriano Eusebio, por autoridad del emperador Constancio\apnota{32}{Constantii imperatoris fehlt in M.} ordenó a Félix como pontífice Romano.

En octavo lugar, refutando la autoridad de Jerónimo aducida por mí a Evagrio, dice que sabía\apnota{34}{M U omnes eiusdem.} que todos los obispos fueron\apnota{35}{M fuisse vel esse, U fuisse et esse, Uv fuisse.} del mismo mérito y del mismo sacerdocio, pero que el papado es un orden por encima del episcopado. Pero no ha vencido mi respuesta, porque Jerónimo atribuye la sublimidad e inferioridad de los obispos no al derecho divino, sino a la costumbre y al poder de las riquezas; por eso aún permanezco con Jerónimo.

Al noveno: Sobre el canon Legimus de la distinción 93\apnota{39}{P M U can. 99. dist., W. A. verbessert Ad nonam: Super etc.}, dice que no le parece [267] que en la iglesia primitiva hubiese tal confusión que el obispo no se distinguiera del sacerdote. Respondo: ¿Qué me importa eso? Lídiese con Jerónimo y con los cánones. Que aduzca a Dionisio, que enumera el episcopado entre los órdenes sagrados, me admira que no aduzca también del mismo autor la monarquía de la iglesia Romana, pues tiene tanto peso para la jerarquía que sin ella no podría subsistir la semejanza con la iglesia triunfante. Convenía, empero, al profesor\apnota{6}{M scriptorem.}, emprendido el argumento de escribir la jerarquía, satisfacer sobre todo en su parte más excelente. Pero Dionisio solo desciende hasta el obispo.

En décimo\apnota{9a}{U ad decimam.} lugar, al capítulo «Primae»\apnota{9b}{fehlt in allen Quellen, ist nach 263, 1 ergänzt; cf. W. A.} de la distinción 99, donde respondí que se prohibió que el pontífice Romano fuese llamado obispo universal\apnota{10}{U universalis ecclesiae, Uv = P.}, lo refutó de este modo: No porque el pontífice Romano no sea

primero y sumo, sino porque el obispo propio peculiarmente de la iglesia Romana no es obispo propio de cualquier iglesia. Respondo: Como si a mortal alguno pudiera caberle tan necio antojo de querer una sola persona presidir todas y cada una de las iglesias, que fuese necesario prohibir tan gran locura\apnota{15a}{M pro tanta insania.}\apnota{15b}{M providere.}. Sin embargo\apnota{15c}{W. A. tum gegen alle Quellen.}, dejando esta refutación, dio otra más recta, a saber, llamando al pontífice Romano\apnota{17a}{M U eum Romanum, W. A. esse Romanum.} no obispo universal, sino obispo de la iglesia universal\apnota{17b}{U dicendum.}. Si yo no le perdonara\apnota{18a}{so W. A., et ego alle Quellen.}, suspendería también esta solución. Pero júzguenlo los jueces\apnota{18b}{fehlt in M.} y auditores.

Por último, a la autoridad de Cristo, donde dijo: «vosotros, empero, no así», dijo que se había reprendido la ambición, pero no el primado. Respondo: Esta es petición de principio, como si ya hubiese obtenido que hay, a saber, principado\apnota{22}{M iam principatum.}. Además, claro es el texto: que prohibió no solo la ambición, sino la misma mayoría sobre los demás.

El D. D. Eck.

A las defensas de las soluciones dadas por el reverendo padre digo, sobre lo primero: que Pablo no reprendió a los que decían\apnota{27a}{M dixerint.} ser de Pedro, si\apnota{27b}{U quia, Uv si.} hubiesen tenido intención del primado\apnota{28}{M U principatus.} de los apóstoles, sino porque miraban la singularidad de las personas; lo cual indican claramente las mismas palabras de división y cisma. Y\apnota{29}{fehlt in M.} aunque con razón antepone a Pablo a Jerónimo, piadosamente ha de creerse, sin embargo, que Jerónimo entendió bien en ese lugar la sentencia de Pablo. Ni es ambiguo el pasaje, por ello constituida la cabeza en la iglesia, para que se quitara la ocasión de cisma, lo cual «al lector gramático\apnota{32}{M quod grammaticus, Mc aliquod quo grammaticus.} satisface\apnota{33}{U satisfecit.}», disciplina que, por encima de las demás partes de la filosofía, dijo el reverendo padre en cierta disputa que sirve\apnota{34}{M theologiae sufficientem, Mc add alii subvenientem, U theologiae subservientem.} a la teología\apnota{35}{von dixit an fehlen die Worte in M, Mc dixit d. doctor in quadam disputatione.}.

Segundo. Nadie negó que Cristo en Jn 5 afirmara\apnota{37a}{M confirmasse.} la igualdad con el Padre, salvo el arriano\apnota{37b}{U Arrianos, W. A. Arriani.}. Ni en otro sentido lo aduce Bernardo, sino que lo que el reverendo padre opina, que los santos padres aducían la santa Escritura [268] para enriquecer el discurso\apnota{1a}{M gratia arbitratur.}, lo rechazamos; pues no es lícito sospechar de ellos tal\apnota{1b}{M pauli.} ostentación.

Sobre lo tercero, de muerto el papa, que entonces la iglesia está acéfala\apnota{3}{M ascevolam.}, digo: Nunca se negó que Cristo es cabeza de la iglesia, según también la glosa al Cantar 5 [v. 11]: «Su cabeza, oro óptimo», glosa: Cabeza, esto es

Cristo, y el papa es su vicario. En consistorio, empero, el obispo y el vicario se computan por la misma\apnota{7a}{U una.} persona; por eso apelar del vicario al obispo está prohibido\apnota{7b}{M U non licet.}. Sobre los cardenales dije que ahora, según la iglesia ya instituida\apnota{8}{M U iam institutam.}, la elección ha venido\apnota{9}{M venit.} a los cardenales, ordenándolo Nicolás, «cierto pontífice», pero que en tiempo de Jerónimo hubo cardenales lo creo, a no ser que Jerónimo no fuera presbítero\apnota{10}{fehlt in M.} cardenal\apnota{11}{U a. R. Martinus respondit, Hieronymum non fuisse cardinalem.}.

Cuarto, sobre Cipriano dice\apnota{12}{fehlt in U, M dicit quod iam non viderit, zitiert bei Löscher p. 346 [eam statt iam].}: ciertamente no puede ser que las palabras del santo mártir se constriñan a las estrecheces de la diócesis Romana\apnota{13a}{M coarctat.}, porque al obispo máximo\apnota{13b}{ad fehlt in P M U, W. A. vermutet episcopi Novatiani.} desde Numidia, que es región de África más allá del Atlante\apnota{14}{U a. R. D. Martinus dixit citra Atlantem.}, como son autores Ptolomeo y Estrabón, llegaron\apnota{15}{M adventavit.} a Roma. Que Cipriano llamara hermano a Cornelio, pienso que fue cosa de los recopiladores\apnota{16}{M U collectoris.}, no mente de Cipriano. Pues si leyéremos las cartas de los santos obispos, hallaremos claramente\apnota{17}{U clare.} que en aquel tiempo fueron más usuales aquellas laudables y magníficas intitulaciones que lo que ahora se hace\apnota{18}{M U fiat.} con el pontífice Romano; como consta de Ambrosio, Agustín, Hilario y los demás\apnota{19}{M U ceterisque.} padres. Pues mutuamente se llaman beatísimos, santísimos, amadísimos de Dios\apnota{20}{U beatissimus, sanctissimus, deo amabilissimus.} etc. A lo añadido se queja de que busco rodeos y cosas no pertinentes al asunto. Con buena venia\apnota{21}{U vinea, Uv venia.}, de grado diría: «Vergonzoso\apnota{22}{M turpi.} es para el doctor etc.». Me pide que pruebe que Pedro ordenó a algún apóstol. Pero\apnota{23a}{U sed.} esto es impertinente para nuestro instituto\apnota{23b}{U ad institutum nostrum.}. Pues no preguntamos quién ordenó\apnota{24}{U aut alterum ordinavit.} a uno u otro, sino quién recibió del Señor Jesús el primado sobre los demás\apnota{25}{M susceperit.}. Lo siguiente lo rechazo por completo, pues hizo\apnota{26}{W. A. gegen alle Quellen facit.} esta inferencia: Pedro no pudo ordenar a ningún apóstol; por lo cual tampoco el sucesor de Pedro pudo\apnota{27}{M poterit.} ordenar al sucesor de otro ni tener potestad\apnota{28a}{P patestatem.} sobre él. Pues lo asumido es verdadero, lo inferido es manifiestísimamente\apnota{28b}{M manifestarie, U manifeste.} falso, pues el sumo pontífice ahora tiene potestad y ordena a otros obispos. Pero\apnota{30a}{M vero.} la razón para desenredar el nudo será aquella\apnota{30b}{M erat ista, U erit ista.}: que el oficio del apostolado dice más fundamento de la iglesia que ser obispo. Por lo cual León X es sucesor del apóstol Pedro\apnota{32}{fehlt in M.}, no apóstol. Pero lo que confiesa, que Pedro fue bien el primer apóstol\apnota{33}{fehlt in M.} en número y en prerrogativa de honor y no de potestad, no satisface al instituto. Primero, porque los evangelistas no comenzaron igualmente el número de los apóstoles, como

Crisóstomo anotó\apnota{35}{MU super Matthaeo annotat.} sobre Mateo. Y véase la glosa a Mt 10. Después\apnota{36}{M U Deinde.}, lo que distingue entre potestad y prerrogativa de honor [269] es expresamente contra\apnota{1}{est contra fehlt in P.} el santo mártir Cipriano, el cual en el tratado Sobre la sencillez de los prelados, contra Novaciano, hablando de la astucia del diablo, increpa a los que, como ministros de justicia, afirman la noche por día, la perdición por salud, la desesperación «bajo apariencia\apnota{4a}{U obtutu (falsch, cf. W. A. Anm. 2), Uv obtentu.} de esperanza», la perfidia bajo pretexto de fe. Y\apnota{4b}{fehlt in M.} más abajo: «Y aunque a todos los apóstoles\apnota{5}{M apostolis, Mc Christus.} después de su resurrección atribuye igual potestad y dice: como me envió el Padre etc., sin embargo, para manifestar la unidad, dispuso con su autoridad que el origen de la misma unidad comenzara desde uno. Esto eran ciertamente también\apnota{8a}{utique et fehlt in M.} los demás apóstoles\apnota{8b}{U discipuli.} lo que\apnota{8c}{M quo.} Pedro: dotados de igual “consorcio”\apnota{8d}{Mc praediti doppelt.} de honor y de potestad (esto ha de ponderarse\apnota{9}{M ponderando.}), pero el comienzo procede de la unidad, para que se muestre\apnota{10}{M demonstretur.} que la iglesia es una». Y más abajo: «Esta unidad quien no la guarda, no guarda la ley de Dios, no guarda la fe del Padre y del Hijo, no guarda la vida y la salud». Estas palabras de Cipriano, dignísimas de notarse\apnota{12}{P M notata.}, el cual entre prerrogativa de honor y de potestad no distingue a los apóstoles.

Al quinto. Sobre la elección digo, como antes, que nada disputamos sobre el modo de elegir, sino sobre la cualidad o más bien cantidad del elegido. Por lo cual el concilio Niceno\apnota{16}{M Niceni.} no fue conciliábulo. Pero lo que mira a la jactancia y al trato puede\apnota{17}{M U possunt.} mudarse según la calidad del tiempo, de las personas y de los lugares, como vemos en muchos cánones.

Al sexto. Lo que el observante señor padre\apnota{19}{U doctor.} me opone como dialéctico, que fingí la distinción del doble nacimiento, ya hemos oído a Cipriano, y habla\apnota{20}{M U et loquatur, W. A. quod loquatur.} en favor de Eck, al cual no le sobra tanto el ingenio para fingir cosas nuevas, sino que explica\apnota{22a}{M U explanat.} según sus fuerzas los viejos\apnota{21}{M cetera.} dichos de los santos. Pero cuando me da\apnota{22c}{M U dat mihi.} un\apnota{22b}{U omnium, Uv unum.} solo sacerdote, Cristo, primero no alcanza la intención\apnota{23}{U attentionem, Lö Anm. richtig.} de Cipriano ni la de Jerónimo, pues quieren que Pedro sea constituido primero de los apóstoles, de donde derive la potestad de los demás sacerdotes, no ciertamente confiriendo interiormente, lo cual es de Cristo como cabeza, sino comunicando\apnota{26}{M communicatum.} la potestad eclesiástica.

Séptimo, lo que desea la eckiana modestia, que juzga a griegos y orientales por condenados, respondo: que los griegos durante largo tiempo no solo fueron cismáticos, sino\apnota{29a}{M et.} herejísimos, como lo atestiguan tan gran cúmulo\apnota{29b}{M cumulos.} de errores y la pertinaz obstinación\apnota{30}{U add Clementina de summa trinitate, in Uv gestrichen; M add de summa trinitate.}, como sobre el Espíritu Santo, sobre la confesión,

de la falsedad de los tres evangelistas y de otras innumerables, aunque a menudo fingieron obediencia a la Iglesia Romana\apnota{32a}{U Romanae ecclesiae.}, como también\apnota{32b}{fehlt in M U.} se hizo en el concilio Florentino en tiempos de Eugenio IV\apnota{33}{M factam statt factum est.}. Por lo cual, si pocos cristianos entre nosotros deben salvarse según sentencia de muchos\apnota{34a}{M sententias.}\apnota{34b}{M debeat, U debent.}, con cuánta mayor razón poquísimos o ninguno se salvan en Turquía, salvo algunos religiosos y los adheridos a ellos, que mantienen la obediencia romana.

De octavo, el reverendo padre me pide que alegue autores según conciencia, en lo que nada dude\apnota{37}{M haesito.}. Quisiera poder aducirlos también [270] según ciencia. Pero\apnota{1}{M ad.} que Jerónimo reconoce a Dámaso como sumo pontífice\apnota{2a}{U pontifice, Uv pontificem.} no puede dudarlo nadie\apnota{2b}{P dubimu (so!)}. Así también que la Iglesia universal\apnota{2c}{fehlt in P.} está fundada sobre la piedra nadie lo duda\apnota{3}{U ambigat, Uv ambigit.}. Pero que aquella piedra sea Pedro y su sucesor\apnota{4a}{U successores.}, lo probaré en otro momento\apnota{4b}{M probamus.}.

Sin embargo zahiere no poco a Jerónimo\apnota{5}{M Hieronymo.}, cuando dijo: Solo entre vosotros se conserva íntegra la autoridad de los Padres\apnota{6}{U prima, Uv primum statt patrum.}; como si también los pontífices romanos hubieran estado sin mancha. Si lo refiere al tiempo de Jerónimo, pues Liberio\apnota{7}{fehlt in M U.} y Anastasio le precedieron. Esto digo, lo cual con razón es admirable en las mentes de los fieles, que ningún pontífice romano, por malo o errante que fuese\apnota{9}{U quantum cuique, Uv quantumcumque.}, haya definido\apnota{10a}{M diffinivisse.} judicialmente y por constitución, que yo recuerde, cosa alguna que fuese\apnota{10b}{U est.} contra los mandatos de la fe cristiana, aunque de hecho erraron a menudo, más aún, intentaron hacer tales juicios erróneos\apnota{12a}{M attemptaverant.}, donde\apnota{12b}{W. A. setzt ubi vor talia.} también fueron corregidos por juicio divino, como de León, arriano adversario de Hilario, como de Anastasio can. Anastasio dist. 19 o quizá 20\apnota{14}{U 20. dist., Uv add forte vel 19.; can. Anastasius etc. fehlt in M.}.

En cuanto al noveno, sobre el canon Legimus 93 dist., quizá el reverendo padre no me entendió. Nunca me vino a la mente que el papado sea un orden por encima del episcopado, sino una dignidad. En cuanto a que dice que estoy en guerra con Jerónimo y los cánones, dije lo que pienso\apnota{18}{M sentiunt.}, prefiero en este lugar la autoridad de Dionisio como más antiguo. Pero cuando el reverendo padre discursea retóricamente sobre Dionisio, por qué no describió al monarca\apnota{20a}{M monarchem.} de la Iglesia\apnota{20b}{P U descripsit.}, sino que no avanza más allá del episcopado, fácil es lo que he de decir. Pues Dionisio trata sumariamente los sacramentos de la Iglesia\apnota{21}{P M perstringens.}, y como yo he negado que el papado sea un orden, el episcopado\apnota{22}{U episcopatum.} obtendrá por consenso de todos el lugar supremo entre los órdenes.

En décimo lugar, en cuanto a que juzga que nadie es tan insensato como para haber creído que haya un obispo propio para cada iglesia\apnota{25a}{M per ecclesiis, Mc per ecclesias.}: ¿qué hay\apnota{25b}{M U quid.} de lo que yo me lamento, sino que es infinito el número de los necios y de los que se complacen en la singularidad? Lea el señor\apnota{27a}{U reverendus.} padre a Álvaro en el Llanto de la Iglesia, a Juan de Torquemada\apnota{27b}{M clamato.} en la Suma de la Iglesia, a Guillermo de Ockham\apnota{28}{U de Ocken.} en el Diálogo, y hallará que alguna vez hubo hombres\apnota{29}{U homines quandoque.} de aquella insania.

En cuanto a que quiere suspender nuestra solución sobre aquello que dije, que el papa es mejor llamado pontífice de la Iglesia universal que obispo universal, lo hice por autoridad del divino Bernardo y por costumbre de los sumos pontífices. Dice en efecto Bernardo, lib. 2 del De consideratione a Eugenio\apnota{33a}{M Eusegenium, Mc Eugegenium (so!)} col. 7: Ciertamente señal del singular pontificado de Pedro. Y más abajo\apnota{33b}{M U etc. statt et; infra fehlt in M.} así: Como cada uno de los demás tiene su iglesia\apnota{34}{U habet ecclesiam, M habet (ecclesiam fehlt).}, a ti te ha sido encomendada una sola, gran nave hecha de todas, la misma Iglesia universal difundida por todo el orbe.

Al undécimo\apnota{36}{U Ad decimum.}, no se ha satisfecho a la palabra de Cristo en Lc 22, como si fuera invención mía. Pero para no hablar sin autoridad\apnota{37}{U non sine auctore loquor.}, le aduzco a Ricardo de Armagh, lib. 7 cap. 3 sobre las cuestiones de los armenios, que entiende así aquel lugar [271] por autoridad del bienaventurado León. Y que esto sea verdad lo indican las palabras: Y el que es mayor entre vosotros etc. Luego Cristo presuponía que alguno era mayor. Quién fuese el mayor, no lo indicó entonces, sino después, cuando dijo a Pedro acerca del demonio que los reclamaba\apnota{4}{M expotente.}, y cómo rogaba por él para que, una vez convertido, confirmase a sus hermanos, donde declaró su mayoría.

Mañana a las siete de la mañana se continuará la disputa según el edicto público como siempre suele hacerse.

El día martes cinco del mes de julio a las siete de la mañana.

El D. D. Martín comenzó como sigue\apnota{7}{Die Zeitangabe fehlt in M; U Martinus Lutherus mane hora 7. die Martis, quae erat 5. Julii.}.

Habiendo de refutar el egregio D. D. mi respuesta\apnota{8a}{U responsionem egre. d. d. (meam fehlt, in Uv nachgetragen), M ad responsionem meam.}\apnota{8b}{U add in causa, in Uv gestrichen.} a aquello\apnota{8c}{fehlt in M.} de 1 Cor 3: ¿Qué es Pablo? ¿Qué es Cefas? etc.\apnota{9}{fehlt in U.}? dijo allí que no se trata de principalidad, sino de singularidad de la persona, y que esto se indica por las palabras del cisma; respondo: Más me mueve el texto mismo de Pablo que

tan violenta y forzada distinción, que en nada se apoya\apnota{12}{M U nititur.} en autoridad alguna, aunque quiera argumentar desde el derecho divino. Es claro, en cambio, que la contienda de ellos fue sobre la preeminencia\apnota{14}{U add vel principalitate, in Uv gestrichen.} de las personas, como lo indica la propia refutación de Pablo, que insiste con atenuaciones y comparaciones diciendo: ¿Qué es Cefas? ¿Qué es Pablo? Ministros por quienes creísteis. Así, el que planta y el que\apnota{16}{U et qui.} riega son uno. Ni el que planta\apnota{17a}{U add est aliquid, in Uv gestrichen.} ni el que riega es algo\apnota{17b}{M aliquid est.}, sino Dios que da el crecimiento. En lo cual manifiestamente rebaja a la persona, esto es, la ocasión del cisma. Como también en Gál 2 [v. 6], habiendo sido seducidos los gálatas con el título de preeminencia de Pedro y de los demás apóstoles, se atreve a decir: Los que parecían ser algo, cuáles hayan sido alguna vez\apnota{21a}{M fuerit.}, esto es, cuán grandes\apnota{21b}{id est quam magni in M a. R.}, nada me importa, pues Dios no hace acepción de personas. Como si dijera: Sea Pedro\apnota{22a}{M add sive Paulus.} u\apnota{22b}{U vel.} otro cualquier apóstol el primero o el último, 'nada importa'. Quisiera, pues, que el egregio D. D. leyera mejor las Escrituras antes de alegarlas 'así'\apnota{24}{M quam.}. Queda, pues, en pie que aquel primado o persona en nada atañe a la Iglesia, al menos por derecho divino.

Paso por alto aquello de que dice que Jerónimo\apnota{26a}{U Hieronymum dicit, M dicit Hieronymus (so!)} entendió\apnota{26c}{M concessisse.} la sentencia de Pablo\apnota{26b}{fehlt in M.}. Asimismo\apnota{27}{Der Satz Item et illud etc. fehlt in M} también aquello que concedió sobre Jn 5. Asimismo paso por alto aquello de la plebeya razoncilla, porque no ha sido refutado\apnota{28a}{M U confutata, Uv. a. R. confutatum habent alia.}. Pues lo que\apnota{28b}{fehlt in U.} adujo en el consistorio sobre el vicario y el obispo\apnota{28c}{M U episcopo et vicario.} nada viene al caso. Mi razón quiso esto: Si la Iglesia no queda acéfala\apnota{29}{M a cephola.} muerto el papa, tampoco quedará acéfala\apnota{30}{fehlt in U, in Uv nachgetragen.} sin ningún papa. Paso por alto aquello de los cardenales, porque es notorio a todos cuándo comenzaron\apnota{31a}{quando coeperint fehlt in M.}. Y paso por alto\apnota{31b}{M transeat statt transeo et.} aquello de Cipriano, porque el D. D.\apnota{32a}{doctor fehlt in M.} dijo no sé\apnota{32b}{M nescit.} qué de su propia saliva sobre Numidia más allá del Atlante, cosa que no existe en la naturaleza. Y aquello de que Cornelio es llamado hermano por Cipriano\apnota{33}{U Damaso, Uv wie oben, M a. R. de Cypriano.} lo paso por alto\apnota{34a}{fehlt in P.}, porque no ha sido refutado, aunque se dijo que se honraron con títulos honorabilísimos\apnota{34b}{M dignissimis, Mc honestissimis.}, llamándose 'beatísimos', santísimos\apnota{35a}{M add etc.}, pues también yo\apnota{35b}{M ego hoc.} digo esto. Pero no se escribía así solo al pontífice romano, como hoy. Pues esto era lo que había que probar.

Y paso por alto también aquello que dijo\apnota{37}{M quod impertinens esse dixerat.} que es impertinente al propósito si el apóstol Pedro ordenó\apnota{38}{U an Petrus apostolum aliquem ordinaret.}, porque este es el único nervio [272] e invicto contra D. Eck, por eso lo pasa sin advertirlo\apnota{1}{Mc transeo.}, para no tropezar irremediablemente

impugnar. Asimismo admito aquello\apnota{2}{fehlt in U, in Uv nachgetragen.} que adujo del mártir Cipriano sobre el honor igual y la potestad igual de los apóstoles, y lo recibo como una gracia que se me ha hecho\apnota{3}{M factam mihi.} por el D. D. Con igual cuidado, que haya pasado\apnota{4a}{M add electionem.} por la elección\apnota{5}{U ecclesiam.} de los obispos descrita por el concilio Niceno y por Cipriano\apnota{4b}{M U a Cypriano.} me place, porque no puede refutarse. Pero lo que repitió 'sobre la distinción' del origen sacerdotal, a saber\apnota{6}{U unde patet.}, que la intención de Cipriano trata del flujo\apnota{7}{M de fluxu, Mc de defluxu.} de la potestad sacerdotal, como nada dijo sino sus propias palabras, ante mí nada concluye\apnota{8a}{M quae apud me nihil concludunt.}.

Aquello\apnota{8b}{U illud quoque.} debe atenderse más\apnota{8c}{fehlt in U.}, que se atrevió a llamar a los griegos herejísimos con bastante modestia, cuando en toda la Iglesia ninguna parte dio\apnota{10a}{P dedit.} más\apnota{10b}{fehlt in M.} excelentes escritores que la griega. Pues en cuanto a que tantas veces, mordiéndome por lo bajo, me reprocha Bohemia y entreteje otras\apnota{11}{M alias quam, Mc alias que.} calumnias, lo remito a los sofistas\apnota{12}{M zophistos.}, son cosas más indignas que para ser mencionadas en una grave y teológica disputa, mucho menos para ser objetadas\apnota{13}{M obicerentur.}.

Por lo cual veamos aquello de Mt 16: Sobre esta piedra etc. Donde digo: o la piedra significa en este lugar la potestad o la fe (confío en efecto que el señor D. nunca satisfará su promesa de que la piedra significa al papa o sucesor de Pedro); si la potestad, superfluamente se dice\apnota{17}{P dicitut (so!)} abajo: Te daré las llaves del reino de los cielos\apnota{18}{regni caelorum fehlt in M.}, esto es, la potestad, a no ser que digas que la potestad fue entregada a la potestad. Además, si en absoluto significa la potestad\apnota{19}{P significet.}, la misma será de todas las Iglesias, porque dice: Sobre esta piedra edificaré mi Iglesia, no la romana\apnota{20}{U add tantum.}. O significa la fe (lo cual es verdad), de nuevo la fe es la misma de todas las Iglesias. Así es patente que\apnota{22}{M Jam hoc facit, Mc ita patet quod.} este solo pronombre «mía» hace común la piedra, signifique lo que signifique por piedra.

Por eso frigidísimamente\apnota{23}{M frigidissimam.} aducen esta 'autoridad' los decretos de los pontífices en favor de la singularidad del principado, la cual defiende tan\apnota{24}{U tamen.} constantemente\apnota{25a}{U defendunt.} la comunidad de la piedra\apnota{25b}{M U petri.}, y así concuerda con el Apóstol a los Efesios: Una fe, un bautismo\apnota{26}{M unum baptisma, una fides.}, un Señor. Pues no, como dice el bienaventurado Jerónimo\apnota{27a}{fehlt in M.} a Evagrio, es una la fe de la Iglesia romana y otra\apnota{27b}{M et alia.} la de Britania y de todo el orbe, así tampoco es otro Cristo ni otra piedra. Y si la misma fe y el mismo Señor y el mismo bautismo es de todas las Iglesias, se sigue que también todo lo demás, que sigue a la fe y a Cristo y al bautismo\apnota{30}{U et ad baptisma et ad Christum.}, es lo mismo.

Por último\apnota{31a}{M Ultimo quod erat auctoritas Matthaei.}, a partir de Ricardo de Armagh\apnota{31b}{P U Amarcano, M a marcano.} interpretó la palabra de Cristo sobre la ambición, no sobre la mayoría, porque dice que Cristo presuponía que había uno mayor entre ellos, aunque no indicaba\apnota{33a}{M U indicarit.} quién\apnota{33b}{M quid.} fuese. Respondo: Sea lo que sea de Armagh\apnota{34}{P M Amarcano.}, prevalece el clarísimo texto sobre el oscuro autor, porque está escrito que la contienda se hizo 'entre' ellos sobre quién de ellos parecía\apnota{35}{M videtur.} ser mayor. Y Cristo: Quien quiera ser mayor entre vosotros, indica bastante que no se suponía que hubiese\apnota{36}{fehlt in M.} ningún mayor, y por la consecuencia del texto es patente\apnota{37a}{M patuit.} que no quiso que hubiese ningún 'mayor etc.'\apnota{37b}{Mc ceteris statt etc.}. Y en cuanto a que\apnota{38}{U quod.} para su confirmación adujo aquello de Lc: Rogué por ti, Pedro, y tú una vez convertido confirma etc.\apnota{39}{conversus fehlt in M U, Uv add a. R. confirma; et tu confirma fehlt in M.}, vale en mi favor, porque manda que los hermanos de Pedro sean confirmados por Pedro, no la mayoría.

[273] El doctor D. Eck.

Sobre los apóstoles, lo que adujo el reverendo padre, por decirlo de una vez, oímos por Cipriano y Jerónimo que también el santo papa y mártir\apnota{3}{M U martyr et papa.} Anacleto atestigua en el cap.\apnota{4a}{M U canone.} In novo\apnota{4b}{Mc nova.}, dist. 21: Que fueron iguales en el apostolado, lo cual nadie niega\apnota{4c}{M negarvit.}. Pero aquella igualdad del apostolado no prejuzga al primado y principado, a no ser que se diga impudentemente que el santo mártir Anacleto se contradijo a sí mismo en el mismo decreto. En cuanto a que pone tanto peso y apoyo en la ordenación de los apóstoles, sé que Pablo escribió ingenuamente: Después que subió a Jerusalén, nada recibió de aquellos\apnota{9a}{U iis.} que parecían\apnota{9b}{Mc sibi videbantur.} ser algo. Pero si ese nervio debe atar tan invenciblemente a Eck, úselo mientras\apnota{10}{U cum, Uv dum.} tenga ocasión de replicar. Pero en cuanto a que dice que se le ha hecho una gracia a partir de Cipriano, que los apóstoles están dotados de igual consorcio de honor y potestad; si a esto llama gracia, destruir sus soluciones, estoy dispuesto a mostrarle a menudo tal gracia. Pues había distinguido el reverendo padre que Pedro había sido primero por prerrogativa de honor, no de potestad; Cipriano los equipara en ambas cosas.

Sobre Numidia y Mauritania\apnota{16a}{M U De Numidia, quam ait esse inter Mauritaniam et; Mc Uv wie oben.}, 'lo que adujo', es impertinente al negocio\apnota{16b}{U negotium, Uv negotio.}, y como hace tiempo que no trato negocio cosmográfico, tampoco me acuerdo de todo. Esto sé: que Mauritania Tingitana\apnota{18}{U Turganicam, M Tinganicam (so!)} y Cesariense se extienden desde el Atlante hasta el mar, y que el Atlante no se pone\apnota{19}{M ponit.} como término de África, sino que divide\apnota{20}{M dividerit.} a África en gran parte suya en muchos grados; baste que los que vinieron a Cornelio desde Numidia no fueron de la diócesis romana, sino

acudieron a la primera Iglesia. Por eso permanece invicta hasta ahora aquella sentencia de Cipriano aún no resuelta. Sobre el influjo me admiro de que aduzca que fueron palabras mías, habiendo yo aducido a Cipriano contra Novaciano\apnota{24a}{M induxerim, Mc adduxerim.}, a quien\apnota{24b}{M quam.} el reverendo padre nada respondió.

Sobre los griegos, confieso que antaño fueron cristianísimos y doctísimos, cuando llamaban a la Iglesia romana primera sede. Pero ensoberbecidos por la soberbia e infectados de envidia, sustrayéndose a la obediencia de la sede romana, cayeron en pésimos errores, y a la vez perdieron la fe con el imperio. En cuanto a que adivina que yo no exhibiré lo prometido, e hizo argumento a partir de Mt 16, me admiro de que el venerable padre, tan hostil a los sofistas, use sin embargo tan bellamente las cautelas de los sofistas\apnota{30}{U et . . . utitur statt ut . . . utatur.}. Y siendo él quien responde, se traslada a la parte del oponente. Por lo cual nada respondo ya, sino que probaré lo que pide que sea probado por mí\apnota{32}{M probare.}, para que no siempre venga a sus partes el tiempo de deliberar.

Por último rechaza al oscuro doctor Ricardo de Armagh\apnota{35}{M U Richardum dixi.}, olvidando aquello que dije, que Ricardo hace eso por autoridad del bienaventurado León. Además no demuestra por la letra lo que pide, que Cristo reprendió con razón la contienda hecha entre los discípulos. Por lo cual también las palabras de Cristo han de tomarse de modo que repriman la contienda, sin quitar el primado. Y el texto aducido después cree [274] el reverendo padre que hace en su favor, porque\apnota{1}{M quo.} Cristo llamó hermanos a los apóstoles y así no lo hizo mayor. Pero si ponderase las palabras del texto, como suele hacer en otras ocasiones, ya vería que Pedro es superior a los demás apóstoles, pues mayor es el que confirma que el confirmado. Esto dicho sobre las soluciones\apnota{4}{Hoc de solutionibus dictum fehlt in M.}.

Vengo, pues, a lo principal que pide, habiendo de probar que el primado de la Iglesia romana es de derecho divino y por constitución de Cristo, de modo que Pedro fue monarca de la Iglesia instituido por Cristo con sus sucesores. Para lo cual repito la razón de Bernardo aducida y aún no resuelta. Repito el dicho de Cipriano igualmente, y en tercer lugar pruebo por aquellas palabras de Cristo, Mt 16: Tú eres Pedro, y sobre esta piedra edificaré mi Iglesia. Donde la glosa ordinaria: especialmente le concedió a Pedro aquella potestad para invitarnos a la unidad, por eso lo constituyó príncipe de los apóstoles, para que la Iglesia\apnota{12}{M ecclesiam.} tuviese un solo principal vicario de Cristo, al cual recurriesen\apnota{13}{M concurrerent.} los diversos miembros si acaso disintiesen entre sí\apnota{14}{U dissentirent (so!).}, porque, si hubiese diversas cabezas, se rompería el vínculo de la unidad.

Así expuso el bienaventurado Agustín en la epístola contra Donato\apnota{15}{W. A. Donatistas.}: Tú eres Pedro, y sobre esta piedra, esto es, Pedro, edificaré mi Iglesia, y aunque el mismo Agustín en otro lugar expuso sobre esta piedra, esto

es, Cristo, según la mente del Apóstol: La piedra era Cristo, sin embargo en el lib. de las Retractaciones no se arrepintió de su primera exposición, así también el bienaventurado Jerónimo en el libro primero contra los pelagianos col. 5 dice: ¿Qué tiene que ver Platón con Pedro? pues\apnota{20}{fehlt in M.} como aquél fue príncipe de los filósofos, así también éste lo fue de los apóstoles, sobre el cual la Iglesia del Señor está fundada con estable mole, la cual ni por ímpetu del río\apnota{22}{U fluminum.} ni por tempestad alguna es conmovida. Así el bienaventurado Ambrosio en el Serm. 47 dice que Pedro fue la piedra.

A esta sentencia asiente también Crisóstomo al inicio del mismo capítulo. ¿Qué pues, dice, Pedro, boca de todos los apóstoles, vértice de todo el consorcio. Y más abajo: Le hizo saber cosas más sublimes, y lo constituyó pastor de la futura Iglesia, y después\apnota{27a}{M subdit statt post.}: a éste lo prepuso Cristo a todo el orbe de la tierra\apnota{27b}{M proposuit.}. Confiesa esto el santo mártir Cipriano a Cornelio papa diciendo: Pedro, sobre el cual había sido edificada\apnota{29a}{M fuerit.} por el mismo Señor la Iglesia\apnota{29b}{M ecclesiae}, uno hablando por todos y respondiendo con la voz de la Iglesia, dice: Señor, ¿a quién iremos\apnota{30a}{M add tu habes verba vitae, U statt dessen etc.}? Por eso san León 24\apnota{30b}{M 34.} q. 1 cap.\apnota{31}{fehlt in U.} Cum beatissimus confiesa ingenuamente que el apóstol Pedro recibió el primado de la Iglesia del Señor. Así todos los santos concuerdan en que Pedro obtuvo de Cristo el primado de toda la Iglesia, omito a los más recientes, Beda, Bernardo y semejantes. Hable el santo papa y mártir Anacleto, que no compuso frigidísimos decretos hace cuatrocientos años, sino hace mil cuatrocientos\apnota{35}{fehlt in P.} así intona: La sacrosanta Iglesia romana y apostólica no obtuvo el primado de los apóstoles, sino del mismo Señor y Salvador nuestro, como dijo al bienaventurado apóstol Pedro\apnota{38}{M dixerit.}: Tú eres Pedro, y sobre esta piedra etc. Y más abajo: Esta sede apostólica, cabeza y quicio, como se ha dicho antes, fue constituida por el Señor y no por otros. Se transcriben\apnota{40a}{U transferibuntur.} estas palabras en el can.\apnota{40b}{M in canone.} Sacrosancta 22 dist.

[275] Así san Marcelo, cuarto antes de Silvestre, can. Rogamus, 24 q. 1, escribe a los obispos antioquenos: aunque la primera sede estuvo en Antioquía, después fue trasladada a Roma por mandato del Señor. Así el santo\apnota{3a}{U beatus.} Julio 3 q. 6 can. Dudum\apnota{3b}{P M U can. b. Petrus.} atestigua que la Iglesia romana tiene por Cristo el primado y es cabeza de las demás Iglesias.

Así el papa Pelagio, no hace cuatrocientos años, sino hace novecientos veintiocho, siguiendo la sentencia de los santos padres, tomó del mismo modo las palabras de Cristo: Y sobre esta piedra etc.\apnota{6}{fehlt in P.} De aquí que entre los condenados y pestíferos errores de Juan Wiclef fue condenado también aquél: No es de necesidad de salvación creer que la Iglesia romana es la suprema entre las demás. Así entre los pestilentes errores de Juan Hus también se cuenta\apnota{11a}{M enumerantur.} aquél: Pedro no es ni

fue cabeza de la santa Iglesia\apnota{11b}{U Romanae ecclesiae.} católica. Y otro: No hay centella de apariencia de que convenga que haya una sola cabeza en 'lo espiritual rigiendo la Iglesia'\apnota{12}{Mc rei ecclesiae statt regens ecclesiam.}, que esté siempre con la militante\apnota{13}{M conmilitante, U in militante.} Iglesia conversando.

Y éste\apnota{14a}{U ista.}: la dignidad papal nació del César. Y: La prefectura\apnota{14b}{P M perfectio.} e institución del papa emanó del César. Así Bonifacio VIII, condenando la herejía de los lioneses, decreta contra su error que es de necesidad de salvación que toda criatura humana\apnota{17}{M U humanam creaturam.} esté sometida al pontífice romano, como se contiene sobre esto en la decretal Unam sanctam. También Juan XXII, condenando la herejía de Marsilio de Padua\apnota{18}{fehlt in M (Lücke).}, reprobó asimismo ésta, que el bienaventurado apóstol Pedro no fue más cabeza de la Iglesia que lo fue otro apóstol\apnota{20a}{M U quam alios apostolos.}. Así en larga serie, comenzando por la primitiva\apnota{20b}{M primativa.} Iglesia, siempre fue confesado entre los buenos cristianos que la Iglesia romana obtuvo el primado\apnota{22}{M U primatum obtinuisse} de Cristo, no por derecho humano y consenso del pueblo.

Confieso que los bohemios en la pertinaz defensa de sus errores\apnota{23a}{M erroro (so!)} conmemoran aquello y con estas\apnota{23b}{M illis.} armas virulentas se defienden, como puede verse en su oblación hecha en el concilio de Basilea, y en otra disputa hecha ante el rey y los magnates\apnota{25}{M aptimatibus.} del reino entre Ragusa\apnota{26a}{M in terra — statt inter Ragusium.} y aquel ministro de iniquidad Juan Rokycana\apnota{26b}{P Rockezanum}. Por lo cual pido venia al venerable padre si soy hostil a los bohemios (no hablo de los cristianos, sino de\apnota{28a}{fehlt in U.} los cismáticos) como enemigos de la Iglesia, y porque los recuerdo en\apnota{28b}{fehlt in M U.} la presente disputa, porque tanto la conclusión misma como lo que ayer fue aducido\apnota{29}{M illata.}, que el primado de la Iglesia fue constituido por derecho humano\apnota{30}{M constituto.}, a mi pequeño y exiguo juicio favorecen muchísimo a los errores de ellos, y, como es fama, de esto se congratulan muchísimo. Esto baste aducir por ahora en presencia, habiendo de oír\apnota{32}{M U auditurus.} del reverendo padre tanto la sentencia como la solución de lo alegado.

El doctor D. Martín.

En cuanto a que\apnota{35}{U Primum diluam contumeliam quod.} el egregio D. doctor me acusa de partidario y claramente patrono de la facción bohemia, perdónelo el Señor, especialmente en tan\apnota{37a}{fehlt in M.} grande corona de tantos varones. Nunca me plugo ni me placerá jamás\apnota{37b}{in aeternum fehlt in M.} cisma alguno\apnota{38}{M scissa.}; inicuamente obran los bohemios, que por propia autoridad se separan de nuestra unidad, aunque el derecho divino estuviese por ellos\apnota{39}{M stat.}, pues el supremo [276] derecho 'divino' es la 'caridad' y la unidad del espíritu\apnota{1}{fehlt in M.}.

Esto solo busqué yo, y ruego a cualquier cristiano\apnota{2}{M U bonum Christianum.} que se digne ponderar con cristiana caridad si no es\apnota{3}{M U sit longe.} de impudentísima iniquidad expulsar de la Iglesia a tantos millares de mártires y santos tenidos durante mil cuatrocientos años en la Iglesia griega y querer ahora derribar\apnota{5b}{M reprobare.} también\apnota{5a}{fehlt in M.} a los que reinan en el cielo. Pues aunque\apnota{6a}{M etiam si.} enloquezcan todos los aduladores del pontífice romano\apnota{6b}{M Romani adulatores pontificis.}, no pueden negar que la Iglesia de Cristo estuvo 20 años fundada, coronada\apnota{7}{M et coronatam.} por mucha parte del orbe de las tierras\apnota{8a}{M terrarum orbis.}, 'antes que la Iglesia romana'\apnota{8b}{M U Romana ecclesia.} fuese hecha de Pedro\apnota{8c}{M a. R. fieret petra.}, como clarísimamente consta por la epístola a los Gálatas [1 vv. 18 ss.], donde escribe Pablo\apnota{9}{fehlt in M.} que después de tres años vino a Pedro, después, pasados catorce\apnota{10}{M a. R. quadraginta.} años, subió de nuevo a Pedro. Los cuales si se computan, se hallarán\apnota{11}{M invenienter.} casi dieciocho años después de la ascensión de Cristo, cuando Pedro estaba aún en Jerusalén, por callar los años que residió en Antioquía, de modo que no puede decirse que la Iglesia romana sea la primera y cabeza por derecho divino. 'Pero\apnota{14a}{M U Jam.} aquello\apnota{14b}{M illud quod.} urge más, que la Iglesia griega' hasta nuestros\apnota{14c}{M ista.} tiempos nunca recibió sus obispos confirmados desde Roma\apnota{15}{U, W. A. Romana.}. Por eso, si hubiese sido derecho divino, durante tanto tiempo todos los obispos de Alejandría, de Constantinopla, santísimos algunos, como Gregorio Nacianceno y otros muchísimos\apnota{17}{M. plures.}, habrían sido condenados, herejes y bohemios, blasfemia de la cual nada puede decirse más detestable.

A los argumentos\apnota{20}{U Ad argumenta solutionum.}.

Cuando el egregio D. D. dice: La igualdad del apostolado no prejuzga al primado\apnota{22}{M primatum.}, y aduce al santo mártir Anacleto dist. 21 ca. in Novo test.\apnota{23}{fehlt in M.}, digo brevemente: este es uno de los más frigidísimos decretos, que impugno, y nadie me persuadirá que este decreto sea del santísimo 'pontífice\apnota{24}{M illius sanctissimi pontificis (Lücke), U huius sancti pontificis et martyris.}.

Al segundo'\apnota{26}{M ad secundum dico.}: No puse fuerza en la ordenación de los apóstoles, de lo cual 'en la oposición veremos'.

En tercer lugar, en cuanto a que adujo a Cipriano equiparando a los apóstoles tanto en honor como en potestad, habiendo yo dicho\apnota{29}{M dixisset.} que Pedro fue primero en honor, lo admito gustosísimamente, y si es preciso quiero errar gustoso\apnota{30a}{fehlt in M.}\apnota{30b}{M volo errare.}, con tal de que el D. D. mantenga\apnota{30c}{U obtineat, in M fehlen die Worte modo d. d. bis principatus.} esto por verdad y me diga dónde permanece después\apnota{31}{U postea (statt mihi), ubi postea.} el principado. Digo, sin embargo, que igual es el honor de los apóstoles

respecto a los otros, pero entre 'sí con razón' dieron a Pedro el primer lugar. Pues cada uno de los apóstoles en su suerte, como\apnota{33}{U et sicut.} cada obispo en su diócesis, tiene igual honor.

En cuanto a que reprocha que nada respondí a Cipriano, busqué la epístola y no la hallé. Sin embargo, bastante se ha respondido por lo precedente, que Cipriano atribuye a cada Iglesia su cabeza por muchas epístolas. Estos\apnota{37}{M hic, Mc hii.} que fueron de Numidia no prueban que sea derecho divino, sino que muestran el hecho, como en caso semejante Acacio de hecho constituyó a Félix pontífice romano\apnota{39}{M Romanum pontificem felicem.}, 'como se dijo ayer'. Más aún, Epifanio obispo de Chipre depuso al mayor obispo constantinopolitano, a saber, a Juan Crisóstomo, como está en la historia tripartita. Pero de ese hecho no se hizo derecho, y muchos otros ejemplos.

[277] Bastante me admiro de que el egregio\apnota{1}{U d. egregium d. (im Druck als 2) gezählt; 1) ohne eigene Definition)} D. D. se haya propuesto probar el derecho divino, y hasta hoy no aduzca ni una sílaba\apnota{2}{M ne quidem unam sillabam, U ne unam quidem syllabam.} de la Escritura, sino solo dichos y hechos de los Padres y además repugnantes entre sí mismos\apnota{3}{M U sibi ipsi, Uv sibi ipsis.}.

En cuanto a que confiesa\apnota{4a}{U fatetur.} que los griegos antaño fueron cristianísimos\apnota{4b}{M fuisse Christianissimos, U Christianissimos fuisse.}, y que después, apartándose de la obediencia de la Iglesia romana\apnota{5}{U obedientia Romanae ecclesiae, in M fehlt Romanae ecclesiae (Lücke).}, perdieron la fe con el imperio, es lo mismo que ayer dijo con la misma autoridad ecciana, que prevalecieron\apnota{6a}{M praefalnuisse.} contra ellos las puertas del infierno\apnota{6b}{M eam adversus inferos.}, 'tratando las Escrituras' así, que entiende por la prevalencia de los infiernos contra la Iglesia\apnota{7}{Mc inferiorem ecclesiam; in U fehlt adversus ecclesiam, in Uv nachgetragen.} la pérdida de la vida terrena y de las cosas temporales\apnota{8a}{M U vel rerum amissionem.}. 'Bellísima glosa ciertamente', como si no hubiera podido permanecer\apnota{8b}{M U manere.} la fe perdido el imperio y así fuera lícito adivinar que no hay cristianos en 'Grecia\apnota{9}{M U sint Christiani in Graecia.}, porque no hay imperio'. Por la misma razón 'diría\apnota{10}{U dicet.} que también los mártires fueron superados por los infiernos'.

También reprocha\apnota{12a}{U etiam dicens.} que, siendo yo quien responde, me he hecho oponente\apnota{12b}{U sum, Uv sim.}\apnota{12c}{M quod fecerim me ex respondente opponentem.}, recibo grato la monición, y no lo hice porque\apnota{13}{U at hoc ideo fecerim ut.} tuviese tiempo 'de deliberar'. No son tan grandes\apnota{14a}{M non tanta sunt gestrichen, a. R. non sunt tanta.} los fulmines eccianos ('para que también yo me gloríe en algo'\apnota{14b}{Mc 'ut ego mihi quid glorier'.}) que me sea necesaria\apnota{15a}{U sit mihi, in M fehlt mihi.} la 'deliberación'.

También aquello de que Ricardo de Armagh\apnota{15b}{P M U Amareanus.} trató la palabra de Cristo no por la suya, sino por autoridad de León\apnota{16}{M U tractaverit.}, poco me mueve. Muestre él mismo por el texto mismo\apnota{17}{M U ipso textu.} (esto es, por derecho divino) lo propuesto, y quedaré contento.

Por último, donde razona así: El que confirma es mayor que el confirmado,

luego Pedro es\apnota{19a}{fehlt in M U, in Uv nachgetragen.} mayor que los apóstoles\apnota{19b}{U apostolus, Uv apostolis.}, quizá porque piensa,\apnota{19c}{M putes.} con Aristóteles, que el agente es más excelente que su paciente\apnota{20a}{M praestantius esse suo passo, U suo passo esse praestantius.}. Pero argumentaré\apnota{20b}{P U argutabor, M argumentabar.} también yo de modo plenamente aristotélico: el que confirma es en verdad mayor por sí, pero por accidente bien puede ser\apnota{21a}{fehlt in M (Lücke).} menor, a no ser que entienda aquí\apnota{21b}{M confortasset.} por confirmación el sacramento de la confirmación, lo cual no creo; por lo demás, no es raro que el superior sea exhortado, consolado, confirmado por\apnota{23}{M vel. statt per.} el inferior. Esto en cuanto a sus refutaciones\apnota{24}{hoc ad confutationes eius fehlt in M.}.

A lo principal\apnota{25}{Ad principale fehlt in P Uv.}.

Al pasar a lo principal, probó que Pedro es monarca por derecho divino, esto es\apnota{27}{U et hoc ex, Uv hoc est.} por la autoridad de Bernardo aducida arriba, de Cipriano, Jerónimo, Ambrosio, Crisóstomo, Agustín, y después por el consenso de todos los santos, asimismo por muchos decretos y decretales de diversos pontífices. Respondo: el D. D. Eck quiso probar por derecho divino, y enseguida, olvidado de sí mismo, cayó en autoridades de los Padres, de las que ya hemos tratado la mayor parte, y hemos visto que en diversos lugares sintieron a veces cosas diversas, y que son mucho\apnota{32}{M multum.} más numerosos y con mayor frecuencia están a mi favor que a favor del D. D. Eck.

En segundo lugar, veámoslas sin embargo por orden.

La primera, la de Bernardo, ya quedó arriba suficientemente disuelta, a saber, que no prueba, sino que tan solo persuade, por cuanto se apoya en un sentido ajeno de la Escritura, como el propio señor doctor concedió\apnota{36}{doctor. fehlt in M; U concessit doctor.}. De igual modo, también\apnota{37a}{fehlt in M.} la autoridad de Cipriano sobre la primacía y sobre el origen de la unidad sacerdotal\apnota{37b}{M dignitatis.} ya se ha visto suficientemente.

En tercer lugar, probó por aquellas palabras: Tú eres Pedro, y sobre esta piedra edificaré\apnota{39}{fehlt in M U.} etc., que Agustín expuso así: Sobre esta piedra, esto es, Pedro, y no lo retractó. Respondo: ¿qué me importa a mí? Si quiere contender contra mí [278], concilie él mismo primero los dichos contrarios. Pues es cierto que Agustín expuso las más de las veces que la piedra es Cristo, y quizá apenas una sola vez que es Pedro; por eso está más a mi favor que en mi contra.

Pero aunque Agustín y todos los Padres hubieran entendido\apnota{4a}{U intellexerunt.} a Pedro por la piedra, yo solo les resistiré\apnota{4b}{M U eis ego.} con la autoridad del apóstol Pablo\apnota{4c}{fehlt in M U, in Uv nachgetragen.}, esto es, por derecho divino, que escribe en 1 Cor 3 [v. 11]: Otro fundamento nadie puede poner, fuera del que está puesto, que es Jesucristo; y con la autoridad de Pedro, en su primera\apnota{7}{U Petri 1. 2.} 2 [v. 4 ss.], donde llama a Cristo piedra viva y angular, enseñando que seamos sobreedificados como casa espiritual. De otro modo, si Pedro fuese el fundamento de la Iglesia, la Iglesia habría caído a la voz de una sola esclava portera, a la que sin embargo ni las puertas del infierno podrán vencer.

Se sigue, pues, que los santos Padres, cuando en este lugar llaman a Pedro piedra, o bien padecen cosas humanas o bien tienen algún otro sentido, sobre el cual no me pronuncio. Por eso admito sin dificultad aquello de Ambrosio cuando dice que Pedro es piedra, puesto que también cualquier cristiano es piedra por Cristo, en cuya solidez se afianza y se hace uno con él.

Cuando Crisóstomo llama a Pedro\apnota{15}{M vocat.} pastor de la futura Iglesia y constituido prepósito de todo el orbe, estoy de acuerdo, con tal de que esa pastura no se entienda de toda la Iglesia, no sea que excomulguemos al apóstol Pablo, que apacentó muchas más iglesias que Pedro; y confieso que Pedro es el primero en honor en todo el orbe, y también esto toca Crisóstomo cuando dice que es el vértice de toda la sagrada\apnota{20a}{fehlt in M U, in Uv nachgetragen.} compañía apostólica. El vértice no es «cabeza del vértice», «sino que es\apnota{20b}{fehlt in U.} parte de la cabeza». Más aún, más abiertamente lo llama boca de los apóstoles, lo que también afirman Jerónimo y Cipriano, porque no en su sola persona, sino en la de todos los apóstoles y de toda la Iglesia, escuchó: Te daré las llaves etc.\apnota{23}{fehlt in M.}.

A los decretos nada digo, que dije ser frigidísimos, y sobre todo al de aquel\apnota{24}{fehlt in M.} Anacleto tan jactado en esta hora, pues ningún buen cristiano cree\apnota{25a}{Mc tenet credat.} que sea de Anacleto\apnota{25b}{U Anacleti esse.} mártir, el cual interpreta Cefas como cabeza y «llama a la Iglesia romana quicio».

En fin, puesto que tanto desagradan los bohemios al egregio D. D., muestre su memoria y su ingenio y escriba contra ellos. Bastante me admiro\apnota{28}{U miror.} de que se hallen tantos acusadores y enemigos de los bohemios y, sin embargo, no haya ninguno que se digne refutar su error con caridad fraterna para gloria de la Iglesia romana.

[El D. Martín pidió a Eck que no quisiera inferirle tan gran contumelia como hacerlo bohemio, porque siempre le habían sido odiosos, por disentir de la unidad]\apnota{34}{Die eingeklammerten Worte fehlen in P M, sind in Uv als „Additio" bezeichnet. Da sie nur den Inhalt des Gesagten zusammenfassen, sind sie wohl irrtümlich in den Text von U aufgenommen.}.

[279] A la hora segunda se continuará la disputa después del mediodía, conforme al edicto público del bedel de la universidad.

A la misma hora segunda de la tarde

el señor D. Martín prosiguió lo hoy comenzado de este modo\apnota{3}{Die Zeitangabe fehlt in M; U hat Hora secunda continuata est disputatio eadem 5. die Julii Martinus Lutther.}.

Añadió\apnota{4a}{M U obiecit, Uv = P.} el egregio D. D. al final los artículos de Wiclef y\apnota{4b}{Vicleff et fehlt in M.} de Juan Hus condenados, y a Bonifacio octavo\apnota{5}{fehlt in U.} como condenador «de los mismos».

Respondo, como antes, que no quiero ni puedo defender el cisma de los bohemios\apnota{6}{M me non velle Bohemorum defendere, U me non velle nec posse defendere Bohemorum.}, sino a la Iglesia griega de mil cuatrocientos años, hayan sentido con ella los bohemios o no\apnota{7}{sive non fehlt in U, in Uv nachgetragen.}; nada me importa; tengo por cierto que ni el pontífice romano ni todos sus aduladores pueden\apnota{8}{M possunt.} derribar del\apnota{10a}{M non possint de (so!).} cielo a tan gran número de santos que nunca estuvieron bajo la potestad del pontífice romano\apnota{9}{M Romanae.}\apnota{10b}{M probari, Mc deturbari.}.

En segundo lugar, y esto es cierto, entre los artículos de Juan Hus o de los bohemios hay muchos que son llanamente cristianísimos y evangélicos, que la Iglesia universal no puede condenar, como es aquel\apnota{13}{M illa.} y semejante: que solo hay una Iglesia universal. Pues este, por obra de impíos aduladores, ha sido inicuamente condenado, cuando reza la universidad de la Iglesia: Creo en el Espíritu Santo, en la santa Iglesia católica, en la comunión de los santos. Este nobilísimo artículo de la fe lo cuentan entre los artículos de Juan\apnota{17a}{fehlt in M.} Hus.

Luego aquel\apnota{17b}{M illam.}: No es de necesidad de salvación creer que la Iglesia romana es «superior a las demás»\apnota{18}{U esse aliis.}, sea de Wiclef o de Hus, nada\apnota{19}{U non.} me importa. Sé que se salvaron Basilio Magno, Gregorio Nacianceno\apnota{20a}{M U Greg. Naz., Basil. Magn.}, Epifanio de Chipre\apnota{20b}{Mc Cypriae.} e innumerables otros obispos griegos\apnota{20c}{M episcopi, U Graeciae episcopi.}, y sin embargo no sostuvieron este artículo. Ni está «en potestad del pontífice romano» ni del inquisidor\apnota{22}{M inquisitorem, U inquisitorum.} de la depravación herética «crear nuevos artículos de fe», sino juzgar según los ya creados. Ni puede el fiel cristiano\apnota{23}{M Christianus fidelis.} ser obligado más allá de la sagrada Escritura, que es propiamente el derecho divino, a no ser que se añada una revelación nueva y probada. Más aún, por derecho divino se nos prohíbe creer sino lo que esté probado, o por la Escritura divina o por manifiesta revelación, como también afirma Gersón,\apnota{26}{M enim.} «aunque más reciente», en muchos lugares, y como observa el divino Agustín, más antiguo, a modo de canon singular, diciendo al divino Jerónimo: Solo a aquellos libros he aprendido a tributar este honor, que se llaman canónicos; los demás los leo de tal modo que, por grande que sea\apnota{29}{U quantolibet, Uv quantalibet.} su doctrina y santidad, no por eso juzgo verdadero algo porque ellos así lo sintieron, sino si pudieron persuadírmelo por los libros canónicos o por razón probable.

Es más, hasta los propios juristas, de quienes menos se esperaría, en el can. Significasti, de electione et elect. pote.\apnota{33a}{et elect. pote. fehlt in M U.} establecen\apnota{33b}{U statuerint, W. A. statuerunt.} que prevalece la sentencia de un solo hombre privado\apnota{33c}{M sententiae.} tanto sobre el pontífice romano como sobre el concilio y la Iglesia, si se apoya en mejor autoridad o razón. Por eso nada es,

lo que el egregio señor doctor, queriendo argumentar contra mí desde el derecho divino, dejando de lado el derecho divino argumenta contra mí desde las compilaciones de los inquisidores de la depravación herética\apnota{37a}{M inquisitoribus.}.

Por tanto, también\apnota{37b}{fehlt in U.} esta proposición de Juan Hus\apnota{37c}{M add vel Wickleff.}: «la dignidad papal» brotó del César\apnota{38}{a Caesare inolevit fehlt in M.}, si es falsa, sea borrado Platina en la vida de Benedicto [280] segundo, donde escribe que Constantino IV, emperador de los griegos\apnota{1}{M imperatorum (so!) Graecorum, U imperatorem Graecorum.}, sancionó que el pontífice romano\apnota{2}{U pontificem Romanum.} es vicario general de Cristo, aunque ni aun así fue observado\apnota{3a}{M servatum.} por los obispos griegos\apnota{3b}{U Graeciae.}.

Por lo cual, cuanto me apremia el egregio D. doctor con los bohemios de aún no cien años, tanto le apremio yo a él con la Iglesia oriental, la mejor parte de la Iglesia universal y de mil cuatrocientos años. Si aquellos\apnota{6}{M ille.} son herejes porque no reconocieron al pontífice romano, acusaré de hereje al adversario, que se atreve a afirmar que están condenados tantos santos celebrados por la Iglesia universal. Lo mismo digo respecto a Bonifacio VIII, de qué clase de pontífice fue y con qué fe deben recibirse sus gestas; bastante lo prueban las historias.

Por tanto, concluyo y ruego al D. doctor que quiera conceder que los pontífices romanos fueron hombres y no hacerlos dioses, sobre todo cuantas veces juzgaron en causa propia; y además, no por sí mismos, sino por medio de indoctísimos aduladores, cuando el divino Gregorio en muchas cartas, aunque pontífice romano, rechazó\apnota{13}{M reiicit.} de sí el primado de todo el orbe, alegando para ello a su predecesor\apnota{14}{praedecessorem suum fehlt in M, die Lücke von Mc ausgefüllt mit dicens (falsch: doppelt).} Pelagio, diciendo entre otras cosas que el venerable sínodo de Calcedonia ofreció este honor del primado\apnota{15a}{fehlt in M.} al pontífice romano\apnota{15b}{M Romani.}, y sin embargo ninguno\apnota{16}{U tam.} se atrevió a aceptarlo. Si, pues, yerro yo, yerra conmigo Gregorio primero con sus predecesores\apnota{17}{M praecessoribus.}, y pecaron damnablemente por no haber asumido el primado ofrecido.

Con esto quiero que quede probado que con los decretos, condenaciones y aprobaciones más recientes\apnota{19}{M recensoribus.} de la Iglesia romana nada se actúa contra mí, puesto que todas son sospechosísimas y en todo contrarias a la antigua verdad y costumbre; no obstante, por reverencia y para evitar el cisma, de muy buena gana las tolero y persuado que se toleren, con tal de que no condenemos como por\apnota{22}{U tantum.} derecho divino a tantos santos anteriores. Esto tuve que decir sobre los artículos\apnota{24}{Haec habui etc. fehlt in M.}.

El D. D. Eck.

Que el reverendo padre, para excusar su honor, niega ser patrono de los bohemios, si los hechos respondieran a las palabras, lo magnificarían

a él\apnota{27}{M magnificarem deum, U magnificarem eum.}, pero lo último no concuerda con lo primero, pues llama cristianísimos a los pestilentísimos «errores de los husitas», de modo no cristiano, dice\apnota{28}{M dico, U dicit, Uv dicat.}. Pero\apnota{29}{M ad.} de «esto» después.

Odio\apnota{30}{M odeo.}, sin embargo, aquella condición que los bohemios cismáticos y los picardos podrían aceptar como dada\apnota{31a}{U deum, Uv datum.}\apnota{31b}{M accipere possunt.}: «si por ellos estuviera el derecho divino». A los griegos laudatísimos y santos mártires siempre los alabé. Pero el reverendo padre, poco instruido en el arte culinaria, mezcla a los santos\apnota{33}{M sanctis, Mc sanctos.} griegos con los cismáticos y herejes, para defender con el afeite de la santidad de los Padres la «perfidia»\apnota{34}{Mc perfidiam tueatur.} de los herejes; porque inculca a la vez mil cuatrocientos\apnota{35}{M quagrigenta annos (so!).} años,

El D. D. Martín\apnota{36a}{Die Zwischenbemerkung in M a. R.; U Überschrift Carolostadius, gestrichen und Martinus überschrieben; U D. Martinus protestabatur dicens.}. Protesto (dice)\apnota{36b}{fehlt in M U.} ante todos vosotros\apnota{37a}{[M U add et publice.} que el egregio D. D. dice esto\apnota{37b}{U hoc.} de mí impúdica y mendazmente\apnota{37c}{M falso et mendaciter dixerit, U mendaciter et impudenter de me loquitur.}. Y el egreg. D. D. Eck protesta a su vez que quiere probar esto\apnota{38a}{rursus fehlt in M U; M se probare, U se velle probare.} con escritos y dichos\apnota{38b}{scriptis et dictis fehlt in M; hier ist der Einschub zu Ende, nicht wie W. A. nach U angibt hinter confessores.},

en [281] los cuales durante largo tiempo la mayor parte de los griegos y la Iglesia\apnota{1}{et ecclesia fehlt in M.} fue herética y cismática, pero no hay concierto de la luz con Belial, ni de los cismáticos con los santos mártires y confesores.

Sobre la Iglesia existente veinte años antes que la romana diré cuando\apnota{4}{U quam, Uv cum.} el reverendo padre me lo oponga\apnota{5}{U componit, Uv opponet.}; no me mueve que los obispos griegos no fueran confirmados por el pontífice romano, pues tampoco los curas o párrocos\apnota{6}{seu parochiani fehlt in M (Lücke, von Mc ausgefüllt mit plebei).} sacerdotes son ya confirmados por el papa. Pero sería de extrema demencia decir que por esto el sumo pontífice no tiene primado\apnota{8a}{M potestatem.} sobre los curas\apnota{8b}{M plebeiis, Mc. a. R plebanos, U plebeianos.} sacerdotes.

Sobre el frigidísimo decreto de Anacleto diré abajo, y defenderé también los demás decretos. Sobre Numidia dijo el reverendo padre que fue hecho, no derecho, lo cual va contra la sentencia de Cipriano, que los acusa como cismáticos, a los que se atreven a recurrir a la cátedra de Pedro y a la Iglesia principal, de donde nació la unidad. Que Epifanio de Chipre expulsó del episcopado a Juan Crisóstomo, obispo superior, véalo él mismo; esto, digo, juzgo que fue hecho, no derecho\apnota{14}{U iustum.}\apnota{15}{M esse.}.

Además\apnota{16}{M praetereo.}, lo que contra mí retuerce cavilaciones, como si yo tratara las sagradas letras sin idoneidad, porque por la pérdida\apnota{17}{M omissionem.} de las cosas temporales y de las riquezas habría dicho\apnota{18a}{M dixerim Graecos (so!).} que contra los griegos prevalecieron las puertas del infierno, me imputa un pensamiento nunca\apnota{18b}{M mihi cogitatum.} tenido por mí, pues dije que las herejías, los cismas y los errores prevalecieron contra la Iglesia de los griegos. Es verdad que también siguió la pérdida del imperio, con gran vergüenza de la cristiandad.

Después\apnota{22}{M U deinde.}, en cuanto a que el reverendo padre se gloría de hablar por derecho divino desde Lc 22, yo aduciré a un varón oscuro, Ricardo de Armagh, apoyado también en la autoridad de León, digo\apnota{24a}{U dicat; Uv = P.} que también yo creo la misma\apnota{24b}{U eundem, Uv = P.} Escritura del Evangelio y el derecho divino. Pero el reverendo padre, apoyado en su propio entendimiento, me rechaza a mí, que sigo\apnota{25}{M me sequentem (so!).} la inteligencia de los antiguos. Es indudable que tanto Arrio hereje como Atanasio tuvieron el Evangelio (el Padre es mayor que yo)\apnota{26}{in U steht der Zwischensatz Pater maior me est hinter Arium, in M fehlt me; Uv verbessert nach P.}, pero Arrio lo entendió erróneamente, y Atanasio como exigía el Espíritu Santo. Sobre la confirmación añadida\apnota{28}{M de obiecta confutatione.} nada importa, pues cualquier prudente entiende que quien confirma a otros con la autoridad concedida por un superior aventaja en esto a los confirmados\apnota{30}{U confirmatos.}. Sobre lo de por sí y por accidente no sé a qué viene\apnota{31a}{U a. R. nisi forte ad currum Carolostadii.}. Esto por lo primero\apnota{31b}{Hoc pro primo fehlt in M.}.

Al asumir los argumentos del reverendo padre contra la intención principal, me reprocha desmedidamente que mi propósito fue probar por derecho divino que la Iglesia romana es «prelada» a las demás\apnota{34}{M U esse aliis.}, y sin embargo solo habría caído en dichos de Padres y santos, como si, a punto de hacer un ánfora, tomara un jarrito\apnota{35}{M urceolam accipiam (W. A. incipiam).}. Perdóneme el reverendo padre si no quiso o no pudo alcanzar mi intención. Pues casi\apnota{37}{alle Quellen fore, W. A. fere.} se satisface bastante lo del derecho divino cuando tantos santos Padres dijeron que es de derecho divino, aunque no callé el derecho divino\apnota{38}{U obticuerimus.} de Mt 16 [v. 18]: Sobre la cual\apnota{39a}{U hanc.} piedra; después aducí las autoridades\apnota{39b}{U ubi auctoritates.} de los santos Padres.

[282] Sobre Bernardo, cuya autoridad y a la vez razón es invencible, no sé cómo el reverendo padre, olvidado de sí, dijo que yo había admitido que Bernardo habló en otro sentido, lo que con ninguna arma podría demostrar\apnota{3}{M possit efficere, U possit devincere, Uv = P.}; pues, conservado el verdadero y genuino sentido de la Escritura, de la igualdad del Hijo con el Padre, procede óptimamente la razón de san Bernardo.

Sobre Agustín y los demás, que dijeron que Pedro es piedra, insinúa que no los recibe, como si dijeran cosas contrarias. Contra él digo yo: ¿cómo se atreve a creer que tan santo y tan docto Padre sintió

cosas contrarias en el mismo libro, en el mismo capítulo, en la colación de una y otra sentencia, lib. 1. retract. cap. 21? Pero cuán modesta y humildemente respondió el reverendo padre agustiniano, júzguenlo otros, cuando él solo prometió oponerse a tantos santos Padres. Esto es lo verdaderamente\apnota{12}{W. A. verbessert virus.} bohemico: querer entender la sagrada Escritura más que los sumos pontífices\apnota{13}{U pontificis.}, los concilios, los doctores y las universidades en gran vigor existentes, cuando, sin embargo, el Espíritu Santo no abandonó a su Iglesia. Y sería maravilla si Dios hubiese ocultado aquella verdad a tantos santos y mártires hasta la venida del reverendo padre. Ni prueba\apnota{16}{M evinci.} nada la inducción del bienaventurado Pablo, pues sus dichos de ningún modo son adversos a los santos Padres y doctores. Cristo es piedra, es fundamento, es piedra angular, es cabeza indudable\apnota{19}{M indubitato.} de la Iglesia. Pero que aquellas cosas no deban atribuirse a su vicario, contra la atestación de los santos Padres y la nuestra\apnota{20}{M U martyrum.}, no ha de aceptarse.

Añadió también hace poco una razoncilla diseminada en el sermón vulgar\apnota{21}{M sermone vulgari.}: Si sobre Pedro\apnota{22a}{U petram.}, ¿cómo pudo edificarse sobre Pedro\apnota{22b}{M U Petrum.} la Iglesia\apnota{22c}{U potuit, Uv = P.}, él que a la voz de una sola esclavilla\apnota{23a}{M aucilluae.} negó a Cristo y la fe cristiana\apnota{23b}{fehlt in M.}\apnota{23c}{U abnegaret.}? Perdóneme el reverendo padre: cuando desprecia\apnota{24}{M contemnit.} la filosofía aristotélica y tanto estima la gramática, ¿por qué aquí con sus ojos de lince no miró\apnota{26}{M U perspexit.} la palabra «edificaré», de tiempo futuro? Pues no, presente el esposo, cuando los hijos se alegraban, era menester vicario, ni entonces Cristo le dio la potestad de las llaves, sino más bien se la prometió. Antes, pues, de la dación de las llaves y de la potestad fue interpelado san Pedro\apnota{29a}{M pater.} por la portera, como atestiguan\apnota{29c}{M attestantur.} Ambrosio y después\apnota{29b}{M U dehinc.} Gregorio.

A Crisóstomo\apnota{30}{M U Ad Chrysostomum.} respondió que Pedro fue superior en honor, pues también Pablo apacentó más iglesias. Como si el reverendo padre no hubiese predicado\apnota{31}{M iam praedicaverit.} más que ya el sumo pontífice, y por esto no\apnota{32}{W. A. verbessert ob hoc sit dicendus.} debiera llamarse mayor pastor que el papa\apnota{33}{M pater.}. Así se retuerce con demasiada sutileza en la dicción «vértice», como si Crisóstomo no hubiese querido\apnota{34}{U voluit, Uv = P.} que se entendiera el primado. ¿Quién soportará esto, siendo simbólico, y usurpándose vértice por cima «en tal traslación» y metáfora? Así también se le llama boca de los apóstoles, porque a menudo habló por todos los apóstoles, lo que también pondera\apnota{37}{M perpenderat.} el bienaventurado Crisóstomo.

Sobre los decretos\apnota{38}{M decreto.} dice que son frigidísimos decretos, de lo que muchísimo me maravillo, pues en su disputa puso\apnota{39}{M posuerat.}: que la Iglesia romana es

superior a todas las demás [283], se prueba por los frigidísimos decretos de los pontífices romanos nacidos dentro de cuatrocientos años, y yo le habría aducido otros mucho\apnota{2}{M multa.} más vetustos, antes de que empezara a enfriarse la caridad de muchos\apnota{3a}{M sanctorum.}.

Su subterfugio no lo acepto, cuando\apnota{3b}{M U quo.} niega que aquella constitución sea de Anacleto, pues así se eludirían todos los decretos de los concilios\apnota{4}{M consiliorum.} y de los sumos pontífices, a no ser sellados\apnota{5}{U mihi sigillatim, Uv = P.}, y así toda la facultad jurídica es tachada de falsedad\apnota{6a}{W. A. verbessert insimularetur.}, como si estuviera apoyada en falsos fundamentos en las lecciones\apnota{6b}{M illectionibus (statt in lectionibus).}, decisiones, juicios y cosas semejantes. Y así sus derechos pontificios serán de vil haba negra, a no ser que hayan sido aprobados por consenso\apnota{8}{M cum consensu.} de toda la cristiandad. De donde de ningún modo ha de recibirse la sentencia del reverendo padre, puesto que los decretos de todos los sumos pontífices y de los concilios se hallan en muchísimos lugares, y como es audaz el género humano, hace tiempo se habrían hallado quienes despedazaran a aquella mona de Anacleto que avanza con piel de león. A no ser, pues, que muestre que en los originales no se contienen, no le doy fe.

Por último, sobre los bohemios dijo, ciertamente no sin contumelia de los doctores cristianos, que hay muchísimos acusadores de los bohemios, ¿dónde están los que por caridad y buen celo escriban contra ellos y los exhorten?; ahí ejerciten el ingenio\apnota{16}{M et ingenium.} y la memoria. ¿Qué fábula contaría a un sordo? Pero ningún cristiano debe negar que muchísimos, movidos por óptimo celo, escribieron contra los bohemios, cuales fueron los padres diputados en el concilio de Constanza, cual fue el doctor que respondió a los bohemios\apnota{20a}{M prohemis.} en el concilio de Basilea\apnota{20b}{U Herbipolensi, Uv = P.}, cual Raguza, cual Juan Capistrano, de óptima probidad y doctrina, sacerdote del divino Francisco\apnota{21}{M divi s. francisci.}, notísimo a todos los fieles, Nicolás de Cusa, el doctísimo de los germanos\apnota{22}{U Nicolaus Cusanus Germanorum doctissimus.}, con otros muchos, que omito\apnota{23}{W. A. verbessert quos praetereo; P M haben eine andere Konstruktion, s. o.} —como me objeta el reverendo padre— a los inquisidores. Por lo cual no faltaron a los bohemios quienes escribieran cosas buenas, sino que ellos mismos faltaron, obstinados en la herejía, al no seguir las cosas buenas.\apnota{25}{W. A. verbessert bona sequerentur.}

Sobre los santos griegos ya dijimos muchas veces, pero juzgo\apnota{27}{U esse arbitror.} horrendo para todos los fieles de Cristo que el reverendo padre contra tan santo y laudable concilio de Constanza, congregado con tan gran consenso de toda la cristiandad, no tema decir que algunos artículos husitas y wiclefitas fueron cristianísimos y evangélicos\apnota{30}{M evangelica.}.

Protestó en este lugar el D. D. Martín que no es verdad haber dicho contra el concilio de Constanza. Eck, por el contrario, se ofrece a probarlo por escritos y dichos\apnota{32}{M. a. R. 'protestor' Martinus (ursprünglich Carolostadius) non est verum quod contra Constantiense concilium dixerim. Eccius offert se probaturum ex dictis et scriptis [Mc rot fügt bei: ‚per notarios'], U a. R. Protestatur Martinus, im Text Non est verum, quod contra Constantiense concilium locutus sim. Eccius contra offert se probaturum ex dictis et scriptis.}.

a los que la Iglesia universal no puede condenar; así como suena pésimamente que haya sido inicuamente condenado aquel\apnota{33}{M istum.} artículo husita de que no\apnota{34a}{fehlt in P.} es de necesidad de salvación creer que la Iglesia romana es superior a las demás\apnota{34b}{M U de necessitate salutis Rom. ecclesiam esse superiorem aliis ecclesiis [U aliis superiorem], Mc (rot) 'non est' de . . . Rom. 'credere' ecc.}, dicho del cual no sin motivo se regocijan\apnota{35}{M gaudent.} los bohemios, y esto pidieron en sus oraciones ante\apnota{36}{fehlt in M.} Dios, pero con gran quebranto de la Iglesia. Porque (como en otro lugar\apnota{37}{M alibi Augustinus.} raciocina el bienaventurado Agustín), si se admitiese alguna mentira en las sagradas Escrituras, «toda quedará» sospechosa de la verdad. Así también los husitas condenados, apoyados\apnota{39a}{M vulti, Mc fulti.} en el patrocinio del reverendo padre,

El D. D. Martín, interrumpiendo, dijo que es un impudente mendacio, con protesta\apnota{39b}{Die Zwischenbemerkung fehlt in M; U a. R. Martinus Luther protestatur, im Text Mendacium est impudentissimum.}

sin duda [284] dirán: Si el concilio\apnota{1}{M U consilium.} erró en estos dos artículos cristianísimos, así vacilará\apnota{2}{M vacillabit.} entre nosotros su autoridad en los demás artículos.

Por eso, en cosa ya condenada, no quiero prodigar más palabras; sobre a qué debe ser obligado o admitido el cristiano, digo esto: que se hace por autoridad del concilio o del pontífice romano, de modo que una sentencia no puede defenderse sin sospecha de herejía, la cual por lo demás podría\apnota{6}{M possit.} defenderse impunemente sin lesión de la fe. A mano está el ejemplo: si la esencia en lo divino engendra. Sentencia que Ricardo abrazó en su lib. de «trinitate» sin ser inculpado\apnota{8a}{M inculpavit.}. Pero hecha la decisión por el concilio\apnota{8b}{M consilium.}, ya nadie diría sin sospecha de herejía que la esencia engendra, quizá en\apnota{9}{U credo in (statt forte in) fehlt in M.} el capítulo: Damnamus o firmiter\apnota{10a}{vel firmiter fehlt in M U.} de summa «trinitate et fide catholica»\apnota{10b}{et fide catholica fehlt in U.}.

A lo aducido de Agustín sobre leer\apnota{11a}{P legdis (so!).} las Escrituras canónicas nada me detengo\apnota{11b}{M modo, Mc movet, U moveor, Uv = P.}, porque no excluye\apnota{12}{M excludo.} los decretos de los concilios y de los sumos pontífices. Qué ponderen mis señores los jurisconsultos en el c. Significasti, de electione, ya no lo recuerdo. Pero, atentos a su profesión, como cultivadores de la justicia no concederán, según creo, que las leyes pontificias sean arruinadas. En cuanto a que aduce a Platina, como si él fuese más que el sumo pontífice o Agustín o Cipriano, el cual en la vida de Benedicto II cuenta que Constantino IV sancionó que a la Iglesia romana se debe el primado, lea\apnota{18}{M lege.}, por favor,

el reverendo padre lo que mis señores los jurisconsultos anotan en el canon Ecclesia s. Mariae\apnota{19}{M in canone S. Mariae, U in canone ecclesiae S. Mariae.} de constitut., qué, cuánto y cómo valen las constituciones de los seculares en las cosas eclesiásticas, y verá en qué sentido deben\apnota{22b}{U debent.} tomarse las palabras de Platina (por no decir del emperador romano) sobre las iglesias y los herejes\apnota{22a}{M) ne dicam Romanae ecclesiis [Mc rot ecclesiae] et de haereticis.}.

En último lugar propone que yo apremio con los bohemios, apremio\apnota{23}{fehlt in U.} con los griegos; confieso que a los cismáticos. Pues a los santos Padres griegos canonizados no los condeno, como los gigantes querían en las fábulas expulsar a Júpiter del cielo. Pero entonces verdaderamente\apnota{25a}{M vero.} nos saldrá\apnota{25b}{U occurit.} al encuentro el reverendo padre\apnota{25c}{fehlt in M.} cuando dé un griego desobediente y rebelde a la Iglesia romana que haya sido canonizado. ¿Quién\apnota{27}{M U et quis.} es este? ¿Y no lo alabaremos? En vano, pues, asume tolerar la monarquía, como si no fuese de derecho divino, sino inducida por cierta tolerancia del pueblo y consenso de los fieles.

Lo cual clarísimamente consta por el Gregorio aducido, que rechazó\apnota{30}{M reiicedit, Mc (rot) reiecit.} el primado ofrecido a él entre los fieles. Pero siga el reverendo padre el sentido, no la corteza de las palabras, y hallará en las cartas la suma humildad del óptimo padre, de modo que con su mansedumbre y humildad quebrantaba más la soberbia del patriarca de Constantinopla\apnota{33}{M Constantinopolitanae.} que si de él se tuviera sospecha de elevación o soberbia. Por eso también fue el primero\apnota{34}{M primum.} en escribirse siervo de los siervos\apnota{35}{fehlt in M.}. Y que esto es verdad, si tan diligente lector de Gregorio es el reverendo padre, pudo leerlo en su misma\apnota{36a}{M eiusmodi.} carta (como se transpone\apnota{36b}{M transsumitur.} en los decret. canónicos 2. q. 6. Decreto\apnota{37}{P M U de cetero statt decreto.}), donde, memorioso de su potestad el humildísimo padre, despreciador de riquezas y gloria, escribe: A los demás obispos, fuera del romano, llamados en parte a la solicitud, no a la plenitud de la potestad. Nadie hay que aquí pueda dudar que a los demás obispos quiere dar la parte, y al pontífice romano la plenitud [285] según el sentido gramatical. Por eso ruego, reverendo padre, que lo dicho, sea por otros o por mí, no quiera siempre arrojarlo al afán de adulación, pues adular no aprendí, adular tampoco sé.

Mañana a la hora séptima se continuará la disputa.

Actuado como arriba, presentes allí mismo los honorables varones Miguel Gabler de Weissenburg y Nicolás Meyer de Leibnitz, laicos de las diócesis de Eichstätt y Augsburgo, requeridos como testigos de lo antedicho.

Juan Graumann, notario, suscribió lo antedicho\apnota{4a}{Die Worte von Crastina luce hora an fehlen in M U.}.

Hora séptima antes del mediodía del día miércoles 6 del mes de julio.

Antes del exordio el esforzado D. Cesario Pflugk pronunció un prefacio, como se señalará al final después de las pruebas aducidas del D. D. Martín.\apnota{4b}{Zeitbestimmung und Zwischenbemerkung fehlen in M, letztere auch in U, die Zeitbestimmung lautet in U: Sexta die Julii hora septima mane}

El D. D. Martín.\apnota{5}{U Martinus exorsus est.}

Después de que ayer el egregio D. D., usurpando el oficio no de parte sino de juez, contra lo pactado en la convención y la voluntad del ilustrísimo príncipe Jorge\apnota{7}{M Gerogii.}, nuestro patrono, me proclamara tantas veces «definido y» hereje, siendo propio de su oficio dejar, con solas razones y autoridades aducidas, a los jueces si yo era hereje o no; en lo cual, si se ha violado\apnota{10a}{M vigolata, Mc violata.} la fe pública, véanlo aquellos\apnota{10b}{U ii.} a quienes incumbe.

A la causa, primero objetó que yo habría llamado cristianísimos a los pestilentísimos errores de Hus, sobre lo cual atestiguo mi inocencia. Y esto no podrá probarlo nunca, y exijo que señale esos mismos artículos que, pestilentísimos, yo habría llamado cristianísimos, o que revoque su palabra.

En segundo lugar me acusa de haber mezclado\apnota{16}{M U commiscuerim.} a los santos griegos con los cismáticos. ¿Qué otra cosa haría quien no tuviera qué decir, siendo certísimo que en el tiempo de la recientísima devastación de Constantinopla\apnota{18}{M fastatae Constantinopolitanae.} hubo integérrimos cristianos en Grecia, después acogidos\apnota{19b}{M receptus.} en Italia\apnota{19a}{M alia, Mc Italia.}? Y aun si\apnota{19c}{M U etiam si.} esto no apremiara, permanece todavía que\apnota{20a}{adhuc stat fehlt in M, quo (statt quod) M.} al menos hasta el concilio de Nicea\apnota{20b}{M Niceni.} hubo cristianos íntegros por todo\apnota{21}{M totam.} el Oriente, no sujetos al Romano pontífice, como clarísimamente atestigua el decreto del mismo concilio niceno\apnota{22a}{M U concilii Niceni.} que dice lib. 10. ca. 3\apnota{22b}{P can. (Druckfehler für can. iii), can. 3 fehlt in U, in Uv nachgetragen, M can 6; dicentis und historiae ecclesiasticae fehlen in M.} hist. Eccl.: Y para que tanto en Alejandría como en la «Urbe Roma\apnota{23a}{M Romano (so!)} vetusta» se guarde\apnota{23b}{U conservetur.} la costumbre, a saber que aquél cuide de Egipto y éste\apnota{24}{P M U hoc, Uv hic.} de las «iglesias» suburbicarias. El mismo concilio atribuyó el primado de honor\apnota{25a}{U honorum.} no al Romano, sino al de Jerusalén, diciendo allí mismo\apnota{25b}{fehlt in M.}: Y para que\apnota{26a}{M quod, Mc ut.} se guarde al obispo de Jerusalén la prerrogativa de honor\apnota{26b}{U honorum.} transmitida desde antiguo.

Pero si ni siquiera esto basta, aún\apnota{27}{U quod.} concluye esto, que el egregio doctor con bastante astucia calló, que la Iglesia de Cristo existió veinte años\apnota{28}{M annus, Mc annos,} al menos antes de que naciera\apnota{29a}{M fuerit.} la Iglesia Romana: por eso\apnota{29b}{M igitur.} es ridícula

es su refutación, pues finge que yo mezclo a los santos con los cismáticos, sin poder señalar quiénes sean sus cismáticos.

En tercer lugar, habiendo refutado mi razón\apnota{31a}{M confutationem.}, a saber\apnota{31b}{M U scilicet.} que los obispos de Grecia no son confirmados por el Romano pontífice, dijo que esto nada vale, pues tampoco los párrocos por doquier son confirmados\apnota{33}{M confirmarentur.} por el Romano pontífice. ¿Quién no ve que esto se dice para perder tiempo?; cuando, sin embargo, los párrocos\apnota{34}{U plebani tamen.} son ordenados por los obispos. Todo esto a la vez lo remito a lo ya dicho sobre Jerónimo a Evagrio. Por la misma causa repite tantas veces «aquel» asunto de Numidia, cuando de aquellos se escribe que vinieron a Roma no llamados y se les acusa de haber venido. Y todavía el D. D. atribuye al Romano pontífice por este\apnota{37}{fehlt in P.} hecho un derecho divino\apnota{38}{U divinum tribuit, M stellt divinum hinter ius.}. Pero también aquello de que Epifanio de Chipre depuso a Juan Crisóstomo, no lo toma como derecho, sino como\apnota{39a}{fehlt in M U.} hecho\apnota{39b}{U accepit.}, y así pretende haber refutado, [286] cuando aquél obró según lo estatuido por el concilio Niceno\apnota{1}{M consilii.} y la autoridad de Cipriano sobre ordenar obispos sin autoridad del Romano pontífice.

Permanece, pues, todavía invicto que o el primado Romano no está afianzado\apnota{4a}{U formatus.} por derecho divino, o los santos de toda la Iglesia oriental han de ser condenados eternamente\apnota{4b}{M in aeterum esse.}. Paso por alto que interpreta la prevalencia de las puertas del infierno como invalencia de las herejías, lo que no me importa; sin embargo no demuestra que los griegos fueran herejes, siendo cosa divulgada aun por los derechos que los griegos no son tenidos por herejes. Paso también aquello de León\apnota{8a}{M leoni, Mc leonis.} y del Armacano\apnota{8b}{P U Amarcani, M domini marcani.}, donde dice seguir la inteligencia de los antiguos, tomando número por número, a saber de un singular León. Asimismo nada digo\apnota{10a}{U dicam, Uv dico.} de la confirmación de los hermanos encomendada a Pedro, bastante se ha dicho\apnota{10b}{fehlt in P M.} de cómo\apnota{11a}{U cum.} nada prueba, ni ha sido aún refutado\apnota{11b}{U confutatum sit.}.

A la cuestión principal dice no haber callado\apnota{12}{U non tenuisse, W. A. tenuisse.} el derecho divino, porque siguió aquello de Mt 16 [v. 18]: Tú eres Pedro y sobre etc.\apnota{13a}{et super etc. fehlt in M U.}, con las sentencias de los Padres, a las que siguió\apnota{13b}{M sic.}, sobre todo de Ambrosio y Agustín, quienes dijeron\apnota{14}{M U dixerint.} que Pedro es la piedra, habiendo osado añadir además que Agustín no lo retractó\apnota{15a}{M retractavit.}; yo después\apnota{15b}{M U postea.}, consultando el libro de las Retractaciones, hallo lo contrario. Pues verdaderamente lo retracta y dice que Pedro no es la piedra, sino que lo confesado es la piedra. Esto mismo lo dice en la homilía que todos los sacerdotes rezan en el día de Pedro y\apnota{18}{fehlt in M.} Pablo, donde dice: Sobre la piedra, no sobre ti, sino sobre la piedra que confesaste. Lo mismo hallo en Ambrosio, aunque también él varía. A esto se añade a la vez aquella áurea glosa, como se glorían, sobre el capítulo: «Así

el Señor»\apnota{21a}{M innovo, a. R. ita dominus.}, que dice: Y sobre esta piedra\apnota{21b}{U petram etc.}, por esta dicción no creo que el Señor mostrara otra cosa que estas palabras que Pedro respondió al Señor cuando dijo\apnota{23a}{domino cum dixit fehlt in M.}: Tú eres Cristo, Hijo de Dios vivo, porque sobre aquel artículo\apnota{23b}{M U super illo articulo.} de la fe fue fundada la Iglesia; luego sobre sí mismo\apnota{24}{M U ipso.} fundó Cristo la Iglesia.

Nada hay, pues,\apnota{25}{M ergo est.} de que se gloríe el egregio D. D. de que los dichos de los Padres estén de su parte, cuando con mucho mayor fuerza se hallan de la mía. Podía, pues, haber ahorrado a su lengua y a nuestros oídos cuando con retórico bramido exclamó: que yo querría\apnota{27}{M vellem.} ser uno solo más docto que todos, querer entender la Escritura más que los concilios, universidades, doctores, que los Romanos pontífices\apnota{29}{M quam consilium (so!), doctores, pontifices, universitates etc.. U quam doctores, universitates, concilia et Romanum pontificem.}, y que sería maravilla si solo a mí, tras tan largo tiempo oculta, se me hubiese revelado la verdad. Pues esto no fue disputar, sino mover envidia.

En cuanto a que se ríe de mi raciocinio, cuando dije que, si sobre él hubiese\apnota{32}{M super eam esse.} sido edificada, la Iglesia habría caído al negar Pedro, aduciendo gramaticalmente el verbo en tiempo futuro «edificaré», como si después de la muerte\apnota{33}{W. A. ergänzt Christi.} hubiese sido edificada la Iglesia sobre Pedro, lo paso por alto, pues cualquiera ve fácilmente qué vale. No obstante, también después de enviado el Espíritu Santo cayó Pedro en gravísimo escándalo de la fe, en Gal 2, cuando fue reprendido por Pablo, en quien permaneció íntegra la fe y la confesión; en Pedro, ciertamente la fe, pero a la vez simulación contra la verdad del Evangelio.

Admito, sin embargo, que el egregio D. D., apoyado sobre todo en esta autoridad, ahora se excluye a sí mismo de ella, de modo que ningún derecho divino le queda a él. Pues si solo se hizo\apnota{40a}{M tanta.} la promesa, hay que ir\apnota{40b}{M eundem, a. R. eundum.} al lugar donde la promesa fue cumplida [287], y se hallará aquello de Jn últ. [20 v. 22 f.], donde no dice: Pedro, recibe el Espíritu Santo, sino igualmente a todos: Recibid el Espíritu Santo, y\apnota{2}{fehlt in M U.} a quienes perdonareis etc. Paso también aquello del Crisóstomo, donde dije que Pedro fue superior en honor\apnota{4}{fehlt in M.}, y el D. D., para refutarlo, opuso: como Pablo predicó más que Pedro, también tendrá más honor, como si yo hubiese hablado de la multitud del trabajo y no\apnota{6}{M U non potius.} de la primacía del orden. Paso lo del vértice y su\apnota{7}{fehlt in M.} simbólica cabeza, que son meras palabras.

Sobre los decretos, se admira de que yo hable de frigidísimos decretos y nacidos dentro de cuatrocientos años. Sobre los cuatrocientos años me abstengo, para decirlo después. Pues no era menester que se me enseñara que mil doscientos años antes\apnota{11}{fehlt in M U.}, más aún cuatrocientos años antes, hubo decretos, más aún hubo también contiendas sobre el primado. Admírese entre tanto el mismo D. D., que entiende ser la autoridad de Mateo palabra de promesa, no de cumplimiento, cuando, sin embargo, los decretos mismos sobre

la misma palabra como de cumplimiento se apoyan por entero. Es necesario\apnota{15a}{U necessarie.}, pues, o que los decretos traten impropiamente, aun por autor\apnota{15b}{M doctore, Mc auctore.} del D. D., o que él mismo yerre.

En cuanto a que no admite el decreto de Anacleto negado por mí, y lo que allí bromea sobre decretos sellados, lo paso por alto, pues aún no ha refutado que ese mismo decreto interpreta Cefas como cabeza, ignorancia que no ha de atribuirse a tan gran pontífice, sobre todo en aquel tiempo en que florecieron las lenguas y se disponía de abundancia de judíos. Pero también\apnota{21}{fehlt in U.} consta esto, que el libro de los decretos aún no ha sido aprobado.

También aquello atacó gravemente, que dije que son más los que acriminan a los bohemios que los que los instruyen. Ojalá hubiese dicho mentira. Veo, ciertamente, que mucho se dice y escribe contra ellos, pero con bastante infelicidad, pues no se les nombra con nombre fraterno, con el que\apnota{25a}{M U quod.} sin embargo Pablo no desdeña a los Gálatas\apnota{25b}{M Paulos Galatos.} caídos en la perfidia de los judíos. Creo yo que los bohemios\apnota{26}{U Bohemus.} son hombres y que con blandos sermones y concesiones pueden ser atraídos, mientras que con acriminaciones y oprobios del nombre de hereje más se endurecen. No hay, pues,\apnota{28}{M ergo est.} por qué excusarnos con que no ha de contarse fábula al sordo, cuando, mandándolo Pablo, hay que instar a tiempo, a destiempo etc.\apnota{29}{U importune, opportune; etc. fehlt in MU.} [2 Tim 4 v. 2].

De los santos griegos ya hemos hablado repetidas veces\apnota{30a}{M saepe.}. Ahora ha de tratarse\apnota{30b}{U autem [statt nunc] tractandum [Uv tractandum].} aquello que para concitar envidia ponderó mucho\apnota{31}{M me elevavit, U olenavit, W. A. clamavit.}, que es horrendo para todos los fieles de Cristo que yo no temiera\apnota{33}{U veretur.} decir, contra tan santo\apnota{32a}{M sanctam.} y laudable concilio de Constanza\apnota{32b}{M consilium.}, que algunos artículos de Hus fueron cristianísimos y evangélicos, los cuales no puede condenar la Iglesia universal. Respondo: Entre los artículos de Hus hay\apnota{35a}{M ut est.} también aquel: Única\apnota{35b}{U unita, W. A. una} es la santa Iglesia universal, que es la universidad de los predestinados. Asimismo otro: La universal santa Iglesia solo es una, como solo es uno el número de los predestinados\apnota{37}{M praedestinatorum sanctorum, U omnium praedestinatorum.}. Estos dos no son de Hus, sino de Agustín sobre Juan 6\apnota{38a}{„6“. fehlt in [M] U, in Uv nachgetragen.} al pie de la letra propiamente\apnota{38b}{M sed beati Augustini ad verbum proprie [P prope].}, y se repiten por el Maestro de las Sentencias en el 4.\apnota{39a}{M U magistrum 4. sententiarum.} sobre el sacramento de la Eucaristía. El tercero es\apnota{39b}{de sacram. eucharistiae und est fehlen in M.}: Las dos naturalezas, divinidad y humanidad, son un solo Cristo. Estos artículos creo que los confiesa conmigo el D. [288] Juan\apnota{1a}{fehlt in M.} Eck. El cuarto: La división inmediata de las obras humanas\apnota{1b}{M humanorum immediata, Mc notiert die Umstellung.} es\apnota{1c}{fehlt in M.} que o son virtuosas o viciosas, porque si el hombre es vicioso y hace\apnota{2}{M U agat.} algo, entonces obra viciosamente, y si es virtuoso y hace algo, entonces obra virtuosamente\apnota{3}{U add etc.}.

Este artículo, como oímos la semana precedente\apnota{4}{P hebdomoda.}, fue triunfalmente defendido por el egregio D. D. Carlstadt, de modo que el egregio D. D. Juan\apnota{5}{fehlt in M.} Eck se vio obligado a repudiar a Escoto y a los escotistas, a Capreolo y a los tomistas con todos los doctores más célebres de su facción y a consentir con él.

Por tanto, lo que ayer vomitó contra mí, que yo era patrono de los bohemios\apnota{8}{M haereticorum Bohemorum.}, hereje pestilentísimo, y otras cosas semejantes si las hay\apnota{9}{Die Worte von haereticus an fehlen in M, a. R. steht 'haereticus pestilentissimus'.}, lo devuelvo a su seno. Éste siente con Juan Hus; cuanto dijere en su favor, dicho queda en mi favor contra sus acriminaciones\apnota{10}{M criminationem.}. Añado aquello que rectamente dije, cristianísimos y evangélicos, sobre todo aquel\apnota{12a}{M ille.}: Las dos naturalezas, divinidad y humanidad, son un solo Cristo\apnota{12b}{U etc. statt der Worte von divinitas an.}. Por lo cual debía admitirme «esto» el D.\apnota{12c}{U hoc admittere egreg. dominus.} D., que por reverencia al concilio de Constanza creyera que estos y semejantes artículos no fueron allí\apnota{14}{U ibi, M ibi non fuisse.} condenados, sino interpolados por algún impostor. Pero más allá de esto\apnota{15}{M ultra haec, U contra haec.}, cuando el mismo concilio dice que algunos son heréticos, algunos erróneos, algunos blasfemos, algunos temerarios, algunos sediciosos, algunos ofensivos de oídos piadosos, debía primero el egregio D. D.\apnota{17a}{doctor fehlt in M.}, por su prudencia\apnota{17b}{M reverentia.}, reducir cada uno a su propio orden\apnota{18a}{U suum ordinem.}, y no condenar contra la determinación del concilio\apnota{18b}{M consilii.} como herejes a todos en general\apnota{19a}{M omnes in universo omnes (so!)}, a quienes el concilio apenas quiso quizá\apnota{19b}{fehlt in M U.} que fuesen temerarios. Pues también la purísima verdad misma puede ser acusada de temeridad, escándalo, sedición, ofensa de los oídos, como aconteció a Cristo\apnota{21a}{M Christum [contigit fehlt].}. Y no por eso es falso\apnota{21b}{U falsus est.} el artículo, mucho menos herético, porque sea temerario u ofensivo. Y así consta que con demasiada precipitación y muy por debajo de la modestia eckiana fui acusado de hereje, a quien apenas quizá\apnota{24}{M U fortassis vix.} podrá probar ofensivo; más aún, como a él en modo alguno le compete juzgar sobre esos artículos, cuáles erróneos, cuáles heréticos, cuáles temerarios, claramente\apnota{26}{M clarum.} consta cuán inigua y temerariamente me acusó de hereje y con pésimo nombre\apnota{27}{U accusat.}.

Por tanto, que no sea de necesidad de salvación que la Iglesia Romana sea superior a las demás, aún no se convence de que sea artículo herético, aunque se cuente\apnota{29}{M enumeretur.} entre heréticos.

En cuanto a lo que razona con el ejemplo de Agustín: Si se admite alguna mentira en el concilio, vacilará\apnota{31}{M facillabit.} toda la autoridad del concilio, es una semejanza infeliz; Agustín razona sobre las divinas Escrituras, que son palabra infalible de Dios, mientras que el concilio es criatura de esa palabra. Por eso se hace injuria a la palabra de Dios\apnota{33}{U verbo dei fit.} con esta comparación, siendo concedido que

el concilio puede errar, como\apnota{34}{M quod.} nota el Panormitano en el cap.\apnota{35a}{M U in cap.} Significasti\apnota{35b}{fehlt in M.}. Paso aquello de la generación de la esencia\apnota{36}{M substantiae.} divina, porque nada hace al propósito.

A Agustín, que manda leer con juicio los escritos de todos exceptuada la Escritura, dice el egregio\apnota{38}{fehlt in M.} D. D. que por él no quedan excluidos los decretos del Romano pontífice y de los concilios. Esto se dice, pero no se prueba; más aún, confirmo mi solución [289] con la autoridad de Pablo a los Tesalonicenses: [1 Tes 5 v. 21] Examinadlo todo, retened lo bueno. El Romano pontífice y los concilios son hombres, luego han de ser examinados y así retenidos, y no han de eximirse de esta regla apostólica.

En cuanto a que del cap. Significasti por mí aducido dice no acordarse, y amonesta a los señores juristas a no arruinar\apnota{5}{M pessumdant, Mc pessumdarent.} las leyes pontificias, lo cual también vale para concitar\apnota{6}{M concitandum.} envidia contra mí. Y no por esto se arruinan las leyes pontificias, si se les prefieren\apnota{7a}{U eius.} las leyes divinas; y cuando esto hacen\apnota{7b}{M U faciunt.} los señores juristas, pésimamente hacen los teólogos que en esta materia resisten aun a los óptimos juristas.

A Platina por mí alegado dice no ser más que\apnota{9}{P plusquam.} el Romano pontífice, Agustín y\apnota{10a}{fehlt in U.} Cipriano, como si por Cipriano, Agustín\apnota{10b}{M U Augustinum et Cyprianum.} o el sumo pontífice\apnota{11}{aut summum pontificem fehlt in M.} ya hubiese probado el primado, él que tantas veces prometió probarlo por derecho divino\apnota{12}{M se divino iure probaturum, U iure divino se probaturum.} y solo exhibió frigidísimos decretos y algunas autoridades de los Padres falsamente\apnota{13}{M falsas [Mc false] 'auctoritates intellectas.'} entendidas. A Platina nada atribuyo, sino a la historia, que es madre de la verdad, la cual escribe Platina.

Lo que dice de las constituciones, lo paso por alto\apnota{15}{de constitutionibus dicit transeo fehlt in M.}.

Finalmente confiesa que los santos griegos no están condenados, sino solo los cismáticos. Esto nada hace al propósito, pues yo\apnota{17}{U ego enim} nunca respondí sobre los griegos cismáticos, sino sobre los santos, sobre todo los del concilio Niceno\apnota{18}{M consilii.}, como tampoco creo que él mismo sienta sobre los latinos cismáticos cuando habla de la Iglesia Romana.

A Gregorio por mí aducido dice que ha de seguirse según el sentido, no\apnota{20}{M et non.} según la corteza de las palabras, y así queda bastante refutada\apnota{21}{M confortata est, Mc confutata und confundata.} mi respuesta; pues basta la sola «palabra eckiana»\apnota{22}{M 'Sed' sufficiunt 'verba Eccii', U sufficiunt enim sola verba Ecciana.}.

En cuanto a que adujo 2. q. 6 cap. Decreto\apnota{23}{2 q. 6 fehlt in P, M 22 (statt 2), P M U de cetero.}, donde escribe el mismo Gregorio que los demás obispos fuera del Romano han sido llamados a participar\apnota{24}{M vocatus.} de la solicitud, no de la plenitud de la potestad, debía el D. D. mostrar que

entendió a todos los obispos por el orbe y no solo a los obispos\apnota{26}{M episcopum.} de la Iglesia occidental. Y aunque esto hiciera en sumo grado, no por eso se prueba por derecho divino. Gregorio es hombre, y tanto menos se prueba cuanto en muchísimos\apnota{28a}{U multo in.} más lugares siente diversamente, de modo que mayor\apnota{28b}{U magis.} derecho tengo yo a decir al D. doctor que\apnota{29}{M quod.} él mismo siga en sus aducidos el sentido y no la corteza, o primero concilie consigo al mismo.

Por último\apnota{31a}{fehlt in M, U ultimum.}, cuando yo contra el egregio D. D. aducía\apnota{31b}{M induxi contra doctorem.} a Pablo, 1 Cor 3 [v. 11]: Que nadie puede poner otro fundamento\apnota{32}{U fundamentum aliud.} fuera del puesto, que es Jesucristo, dijo que también ha de atribuirse al vicario de Cristo; luego habría que decir así: Nadie puede poner otro fundamento fuera del puesto, que es el Romano pontífice, y así habrá que resistir a la evidentísima autoridad de Pablo, o revocar que esto no ha de\apnota{36}{U sint.} atribuirse al vicario de Cristo.

Semejantemente también\apnota{37}{fehlt in M U.} aquello a Gal 2, [v. 6], donde dijo Pablo: los que parecían ser algo (habla de Pedro, Santiago y Juan)\apnota{38a}{M Johanne et Jacobo.}, cómo fueron alguna vez, nada me\apnota{38b}{M ad me.} importa, pues Dios no acepta la persona del hombre; allí Pablo dice evidentísimamente que la grandeza y cualidad de Pedro nada importa, ni es aceptada por Dios. Pero si [290] es de derecho divino, es certísimo que Pablo aquí\apnota{1}{M hic Paulum.} miente, pues lo que es de derecho divino en todo es aceptado por Dios, y mucho importa. Y por eso esa autoridad de Pablo obliga a que el primado de ningún\apnota{3a}{M U unius.} hombre obispo sea de derecho divino\apnota{3b}{M est iuris divini.}, y por eso ha de revocarse lo que el D. D. dijo ser\apnota{4a}{fehlt in U.} de derecho divino\apnota{4b}{M divino iure.}, pues patentemente\apnota{4c}{U apparenter.} contradice al Apóstol. Esto\apnota{5}{U hic [zum vorigen Satze gezogen!]} lo digo de modo que, si el egregio D. D. no mudare su sentencia, como evidentísimamente contraria a Pablo, no\apnota{6}{M U sim.} lo tendré por católico. Baste esto de mi respuesta a lo objetado por el D. D. Juan Eck sobre la proposición decimotercera\apnota{8}{M super propositionem 13.} acerca de la potestad y primado del sumo pontífice, lo cual dejo al juicio de los jueces y de todos los lectores.

Hora primera después del mediodía continuará la disputa por mandato de los señores\apnota{9a}{Die Zeitbestimmung fehlt in M U.}.

Sigue el prefacio del esforzado señor Cesario Pflug, del que arriba se hace mención\apnota{9b}{vgl. oben S. 95. Dieser ganze Zwischenfall [Rede Pflugks] fehlt in M U.}.

Por las acriminaciones surgidas en la disputa de ayer entre los egregios señores disputantes, como consta por las protestas hechas, y

referidas al consejo de los señores presidentes, al día siguiente, a saber el miércoles 6 del mes de julio, al principio de la hora séptima, antes del exordio de la disputa, celebrado primero consejo, el esforzado señor Cesario Pflugk en nombre del ilustrísimo príncipe y de los señores del consejo, en quienes sobre los defectos emergentes de este género comprometieron las partes, en vigor de tal compromiso exhortó y quiso que los señores disputantes en adelante se abstuvieran de mutuas acriminaciones y lesión de la fama; y asimismo que no trataran temerariamente de la santa Iglesia y de sus concilios, y que terminen en el día de hoy la disputa de esta materia, no sea que la disputa, prorrogada más de lo justo, resulte más costosa para los forasteros y obligue a los nuestros a posponer por más tiempo sus estudios; añadiendo aquello de que las partes con sus adherentes no intentaran nada contra la paz y el salvoconducto, como lo llaman.

Juan Graumann, notario, lo suscribió.

Hora primera después del mediodía del mismo día\apnota{9c}{Unterschrift und Zeitbestimmung fehlen in M U.}.

El D. D. Eck.

Aunque el reverendo padre me reprocha que usurpé\apnota{11}{U occupaverim, Uv = P.} el oficio de juez, nada hace a la cosa; soy disputante. Ni dije\apnota{12}{M dicta dixi (so!)} hereje, sino que sus dichos favorecen y patrocinan a los herejes\apnota{13}{M haeretica.} y sobre todo a los bohemios, principalmente cuando osó decir aquello horrendo, que algunos artículos de Juan Hus reprobados por el sacro concilio de Constanza son cristianísimos y evangélicos. Descendamos más bien a la cosa. Dos artículos de este género adujo\apnota{16}{M heri protulit heri, Mc streicht das erste heri.} ayer, hoy añadió algunos, de los cuales\apnota{17}{M quo.} abajo.

En cuanto a lo que toca a los griegos, es admirable\apnota{18}{M miror, Mc mirum.} con qué empeño los defiende el reverendo padre, de quienes afirma que fueron buenos\apnota{19a}{M bonos esse.} aun en la ruina\apnota{19b}{M concilio.} constantinopolitana, los cuales después\apnota{20}{M dehinc, U deinde.} vinieron a Italia; que yo debo mostrar y nombrar quiénes fueron los griegos cismáticos y herejes. Me admiro de que el reverendo padre en su carta me objete ignorancia de las historias, y él mismo en este lugar no quiera saber lo que toda la Iglesia conoce. ¿Acaso no fue cismático\apnota{23a}{M fuit haeretícus.} Macedonio, Nestorio\apnota{23b}{M Nestorius nonne Macedonius, U Nestorius non Macedonius.}, Eutiques\apnota{23c}{M nonne Eutyches}, Acacio, Juan de Constantinopla, y los que perseverando largo tiempo en el cisma bajo Eugenio IV hicieron obediencia a la Iglesia en el concilio de Florencia, aunque por pertinaz perfidia al punto volvieron al vómito\apnota{26}{M fomitum.}? De otro modo, si los griegos no hubiesen sido cismáticos y herejes\apnota{27}{M haeretici et schismatici.}, no tendríamos «decisión sobre la suma Trinidad y\apnota{28a}{trinitate et fehlt in U} la fe» católica contra los griegos\apnota{28b}{contra Graecos fehlt in M.}. O quizá

el reverendo padre aún\apnota{28c}{fehlt in M U.} ignora que «Tomás editó un libro» sobre los errores de los griegos.

Añado, porque ya repetidas veces adujo que la Iglesia oriental no estuvo sujeta al Romano pontífice, y que Epifanio de Chipre\apnota{31}{fehlt in P M.} repelió del episcopado a san Juan Crisóstomo, lo que el reverendo padre dice que estuvo bien hecho por autoridad del concilio Niceno\apnota{32}{M consilii.} y de Cipriano. Pero cayó en la fosa\apnota{33}{M faveam.} que hizo, pues san Juan Crisóstomo fue depuesto por odio de la emperatriz, y\apnota{34}{fehlt in M.} sustituido un hereje arriano, al que el sumo pontífice Julio mandó retirarse, y restituyó al santo Crisóstomo\apnota{35}{M sanctum Johannem Chrys., U Johannem Chrys.}. Y como esto no se hubiese hecho, el papa Julio preceptuó que Juan Crisóstomo fuese restituido y excomulgó al emperador. Así también procedió contra otros obispos herejes, de modo que, a saber,\apnota{38}{fehlt in M U.} no una vez, sino diez veces las historias transmiten que los griegos se apartaron de la obediencia de la Iglesia Romana [291]. Vaya\apnota{1a}{M Ea.} ahora el reverendo padre y\apnota{1b}{fehlt in M.} diga que el Romano pontífice no dispuso sobre las Iglesias orientales, o justifique la repulsa\apnota{3}{M repulsum.} de tan santo padre\apnota{2}{U patres.} Crisóstomo.

En tercer lugar, sobre el sínodo Niceno\apnota{4a}{M Niceno.} que adujo de la Historia escolástica\apnota{4b}{W. A. verbessert ecclesiastica.} lib. 10, si esa es la constitución que quiso en su conclusión, es buena, digo, pero fría para el propósito. Pues\apnota{6a}{U iam, Uv nam.} ningún\apnota{6b}{nullo.} sínodo aun en aquel tiempo se juzgaba legítimo si no hubiese sido congregado por autoridad del Romano pontífice. Así León, así Marcelo, así Julio, santos, no fríos papas, lo sancionaron, viviendo en torno al mismo tiempo. Sus constituciones se leen dist. 17.\apnota{9}{P M U dist. XIX.} cap. Synodum, y siguientes; de donde lo que concedieron al patriarca de Jerusalén la prerrogativa del honor, pero no el primado de toda la Iglesia, esto debía probar el señor padre. La sexta sínodo, en cambio, da expresamente la primera sede a la Iglesia Romana, lo que se transmite XXII dist. Lo cual también obtuvo el papa León contra el emperador Miguel y san Gregorio contra Juan de Constantinopla y el emperador Mauricio\apnota{15}{fehlt in M.}.

De donde en este lugar también discutiré aquello de san Gregorio, a quien así aduce como si hubiese rechazado el primado, lo que consta ser falsísimo, aunque tanto atribuya a Platina, quien así escribe sobre el santo Gregorio: Además, como Juan de Constantinopla, celebrado un sínodo de griegos, se hubiese creado a sí mismo ecuménico\apnota{19}{M economicum, U icomenicum.}, esto es patriarca universal, y Mauricio hubiese amonestado\apnota{20}{que fehlt in U.} a Gregorio a obedecer a Juan, respondió el hombre de constante ingenio y fe que la potestad de atar y desatar fue entregada a Pedro y a sus sucesores, no a los obispos constantinopolitanos; por tanto, dejase de concitar contra sí la ira de Dios.

De donde consta que Gregorio, sin duda santo y humilde padre, no solo no\apnota{25a}{fehlt in M.} rechazó el primado, sino que también lo afirmó contra el emperador constantinopolitano\apnota{25b}{M contra Const., U et Const., W. A. et episcopum Const.}. Por eso no quiero imponer a tan santo padre lo que el reverendo padre, como poco consecuente consigo mismo, que Gregorio dijera cosas contrarias\apnota{27a}{W. A. quod contraria.} y repugnantes\apnota{27b}{M U. pugnantia.}, sino que negó ser obispo universal, en el sentido dado en la disputa anterior, y sin embargo mantuvo el primado.

En cuanto a que siempre fue esta la costumbre, aun en tiempo del concilio Niceno, observada, consta por la epíst. 92 de Agustín y sus coobispos a Inocencio I, donde dicen: Porque el Señor por singular don de su gracia te colocó en la sede apostólica y tal te mostró en nuestros tiempos, pedimos que te dignes aplicar pastoral diligencia a los grandes peligros de los enfermos miembros de Cristo. Responde el papa\apnota{35}{U add in epistola sequenti.}: Diligente y congruamente consultáis a la cumbre del honor apostólico, del honor, digo, al que\apnota{36}{M U quoniam.}, fuera de lo extrínseco, le incumbe la solicitud de todas las Iglesias, sobremanera solícito por las cosas, qué sentencia haya de tenerse, habiendo seguido la forma de la antigua regla. ¿Qué\apnota{38a}{M quis.} podría decirse más expreso\apnota{38b}{M U possit.}, [292] que el papa en tiempo de Agustín\apnota{1}{M U tempore Augustini papam.} tuvo solicitud de todas las Iglesias, y que entonces no comenzó, sino que es forma de la antigua regla?

Después viene\apnota{3}{M Dein tenet; U Deinde venit.} el reverendo padre con su Hércules, esto es con aquella razón de los veinte años, donde estuvo la Iglesia antes que la Romana. Respondo: Ser verdad que Cristo en el Evangelio no expresó a la Iglesia Romana, sino que constituyó a Pedro príncipe de los apóstoles. De donde, porque Pedro, cuyo era el principado, mandándolo el Señor trasladó la sede de Antioquía a Roma, por lo cual entonces no por derecho humano, sino por mandato de Dios fue hecha primera, con derecho antes habido por san Pedro, como\apnota{9}{fehlt in P M.} atestigua el santo Marcelo mártir y papa ca. Roga. 24. q. 1: Aunque la primera sede estuvo en Antioquía, después, mandándolo el Señor, fue trasladada a Roma. Pues ¿quién tan desatinado que quiera fijar el primado al lugar? Pero esta es la sentencia de los santos Padres, porque por derecho divino Pedro fue constituido vicario de Cristo, y sus\apnota{13a}{M omnes [statt eius], U et omnes.} sucesores, cuales son los Romanos pontífices, dondequiera\apnota{13b}{M U ubi.} que estuvieren, son vicarios de Cristo. El rey o emperador de los Romanos\apnota{14a}{U Roma: W. A. Romanus.} no lo es por\apnota{14b}{U vel ob.} esto en Roma. Así parece haber sentido san Agustín contra los maniqueos. Manifiesto es, dice, que en cosa dudosa vale para la certeza de la fe la autoridad de la Iglesia católica, la cual desde las fundadísimas sedes de los apóstoles hasta el hoy\apnota{17}{M U hodiernam.} de hoy se afianza por la serie de los obispos sucediéndose entre sí y por el consenso del pueblo; por lo cual en las cartas refiere san Agustín el catálogo de los Romanos pontífices. Esto sobre los preámbulos\apnota{20}{Haec de praeambulis fehlt in M.}.

Ahora\apnota{21}{fehlt in M} a lo principal.

Habiendo de responder el reverendo padre a lo aducido por mí —dije, en efecto, que Pedro fue antepuesto a los demás apóstoles por derecho divino, Mt 16; alegué a Jerónimo, Bernardo, León y Cipriano, a los cuales nada respondió, aunque de su propia mano confiesa en su libro que el bienaventurado Cipriano sintió que la «iglesia» fue fundada sobre Pedro\apnota{26}{M U petram.}, pero se atrevió a añadir que el bienaventurado Cipriano «ahí» se equivoca—. Intentó responder a Agustín, pues respecto a Cipriano es indudable, ya que en otra carta a Cornelio\apnota{28}{U Cornelium 8, Mc hat 8 vor epistola eingeschoben.} llama a la iglesia Romana matriz y raíz de las demás.

Me achaca que habría citado un dicho retractado por Agustín. El reverendo padre no puede oscurecer el juicio del lector\apnota{30}{M nequid.}, pues tan poco se arrepiente [Agustín] de su primera sentencia que incluso introduce al bienaventurado Ambrosio como apoyo; sino que, como aducía yo desde el principio\apnota{32}{M induxit, Mc induxi.}, atestigua Agustín que también expuso a Cristo por medio de la piedra, sin retractar sin embargo ninguna\apnota{33}{tamen fehlt in U, M nec tamen.} de las dos sentencias y sin preferir tampoco ninguna\apnota{34a}{in M Lücke, Mc 'neutram sententiam praeeligendam' (so!).} añade: «De estas dos\apnota{34b}{fehlt in U, autem in M gestrichen.} sentencias, cuál sea más probable, elíjalo el lector». Agustín no se atreve a definir\apnota{35}{M U diffinire; U audit.}, y el reverendo padre quiere rechazar una sentencia «atestiguada» por tantos santos padres y aceptar la otra a\apnota{36}{U acceptatam et cum [statt attestatam et aliam pro], Uv = P.} su arbitrio. Por eso no admito que Agustín haya dicho en un\apnota{37}{M U novo.} capítulo cosas contrarias u opuestas, sino diversas. Pues abrazo ambas\apnota{38a}{U amplectatur.} sentencias, y Cristo es piedra\apnota{38b}{U quoniam petra est Christus; M fuit, Mc est.}, y [293] Pedro fue piedra. Por lo cual permanezco con Ambrosio, Jerónimo, Cipriano, Bernardo y los demás sagrados concilios\apnota{2}{M alii [Mc aliis] consiliis.} y decretos.

En segundo lugar, como el reverendo padre hubiese tomado el argumento de la portera, para impugnar la sentencia común\apnota{4}{U meam.}, le mandé que ponderase mejor con razón gramatical las palabras de Cristo, pues por consenso de todos en Mt 16: «sobre esta piedra edificaré y a ti daré etc.» Cristo prometió a Pedro, no le dio entonces. Pero dejando de lado su gramática, de la que sin embargo dijo\apnota{7a}{M tamen dicit, U dixit tamen.} que vale más para la teología\apnota{7b}{ad theologiam fehlt U. M valere, Mc pertinere.} que las restantes partes de la filosofía\apnota{8a}{Mc theologiae.}, objeta por qué entonces se fundan los decretos\apnota{8b}{M fundant.} en aquel\apnota{8c}{M hoc.} lugar de Mt 16 y yo mismo me fundé\apnota{9a}{M ipse met fundaret.}, «del cual ya me aparto»\apnota{9b}{M 'quo tamen iam recedo', U qua iam recedo.}. Respondo: porque Cristo es camino, verdad y vida, por eso ha de creerse indudablemente que él

cumplió a Pedro lo que le había prometido. Por lo cual los decretos argumentan rectamente a partir de aquel lugar, donde Cristo prometió, pero lo cumplió por primera vez después de la resurrección.

Pero tácitamente objeta que Cristo no solo cumplió esto a Pedro, sino a todos los apóstoles, diciendo: «Recibid el Espíritu Santo etc.»\apnota{14}{M U quorum remiseritis etc.}. No quiso expresar el reverendo padre dónde dio a Pedro el principado prometido, pues en esto convienen los doctores: que en la cena Cristo hizo sacerdotes a los discípulos, dándoles potestad sobre\apnota{17a}{U supra.} el cuerpo verdadero de Cristo, diciendo\apnota{17b}{M U dicens.}: «Haced esto en memoria mía», Lc 22 [v. 19]\apnota{18}{U in meam Luc. 20 (so!) Uv 22.}, y después, en el día de la resurrección, les dio potestad sobre el cuerpo místico de Cristo diciendo\apnota{19}{fehlt in M, eis . . . Christi und dicendo fehlen in U.}: «Recibid el Espíritu Santo». Pero el primado y la prelatura de toda la iglesia los confió\apnota{20}{U promisit.} a Pedro: «Apacienta mis ovejas», Jn últ. [21 vv. 15 ss.] Esto\apnota{21a}{U hinc, Uv = P.} lo atestiguan Gregorio, Crisóstomo, Jerón.\apnota{21b}{fehlt in U.} y los demás santos padres\apnota{21c}{M attestantur.}. Dice, en efecto, Crisóstomo: «El primero de los apóstoles era Pedro, boca de los discípulos y vértice del colegio. Por eso, borrada la negación, le confía la\apnota{23a}{U promittit [ei fehlt].} prelación sobre los hermanos». Y en la homilía 80.\apnota{23b}{U 8.} sobre la penitencia: «Pero aquella misma noche Pedro cayó y se levantó; después de tan\apnota{25a}{fehlt in M.} grave caída de nuevo la penitencia lo condujo al grado anterior y le\apnota{25b}{M egit totum orbem terrarum ei, U ei per totum orbem terrarum.} entregó todo el orbe de la tierra, la presidencia de la iglesia»\apnota{26}{M dedit.}. Así la glosa a 1 Pe 5 [v. 2]: «Apacentad la grey que está entre vosotros»\apnota{27}{fehlt in U.}, dice: «como el Señor confió al solo Pedro el cuidado de toda la grey etc.»\apnota{28}{fehlt in M.}. Por lo cual lo antes prometido Cristo lo cumplió entonces. Como también el bienaventurado\apnota{29}{fehlt in P.} Gregorio concuerda: «El cuidado de toda la iglesia y el principado se confía a Pedro, a saber: “Apacienta mis ovejas”». Y aunque no hubiese mostrado\apnota{30}{U coram, Uv = P.} dónde lo dio Cristo, la verdad no pudo mentir\apnota{31}{U poterit.}. Y muchas son las cosas hechas que no están escritas, como enseña elegantemente\apnota{32a}{M U eleganter docet.} Alejandro III\apnota{32b}{3. fehlt in M.} en el cap. Con Marthae, sobre la celebración de las misas\apnota{32c}{U etc. [statt de celebratione missarum].}.

Pero de nuevo objeta el reverendo padre que, después de la misión del Espíritu Santo\apnota{33}{M missionem spiritus sancti.}, Pedro pecó todavía\apnota{34}{U peccaverit.}. Y magnificó bastante el pecado de Pedro, como se escribe en Ga 2. Aunque primero hubo disensión entre los apóstoles Pedro y Pablo, después\apnota{36}{M dein.} entre los doctores de la iglesia, Jerónimo y Agustín, sin embargo no puedo ser inducido a esta sentencia: que alguno de los apóstoles, después de la misión del Espíritu Santo, cayese en algún pecado mortal, como tampoco se admite aquello de los santificados en [294] el útero, y

no se admite. Pero esto es ajeno al propósito. Sin embargo, dado y en modo alguno concedido que san Pedro hubiese pecado mortalmente en su simulación, aun así habría permanecido piedra y cabeza de la iglesia\apnota{3a}{petra et caput ecclesiae fehlt in M, Mc füllt die Lücke aus petra et ecclesia.}, a no ser que «el reverendo padre quiera» defender también este\apnota{3b}{Mc vellet [et hunc fehlt].} artículo husita\apnota{4}{M hussiticum articulum.}, cosa que no creo: «Ninguno es señor civil, ningún prelado, ningún obispo en pecado mortal»; lo cual produciría la suma\apnota{5}{U summum.} incertidumbre de la religión cristiana.

En cuarto lugar, sobre los decretos, para él frigidísimos, digo esto: tantas veces han sido recibidas en la iglesia las constituciones de los sumos pontífices en materia de fe y, lo que es más, antaño el pontífice Rom. no confirmaba al obispo si no enviaba la cédula\apnota{9}{M cedam.} de protestación de la fe, en la cual profesaba creer los evangelios, los cuatro concilios, los sínodos legítimos\apnota{11}{M legitimos.} y los decretos de los sumos pontífices. Por esta razón el papa Juan no quiso confirmar a Gilberto de Colonia\apnota{12}{M gibrechtum, Mc Gilibrechtum [Coloniensem fehlt].} como obispo, porque no había ofrecido perfectamente esta cédula de la fe\apnota{13a}{M cedam.}, como leemos en el can. Optatum, dist. 100\apnota{13b}{U 20. dist.}.

En cuanto a que rechaza el decreto de Anacleto, ayer dijimos con qué autoridad o con qué prueba hace esto. Pues estando los originales a mano, nadie ha dicho aún que este decreto del santo mártir Anacleto no sea suyo. Pero Nicolás de Cusa, el más docto de los alemanes, que confiesa haber escrito la concordancia de la fe a partir de los originales, alega aquel decreto de Anacleto. Añadió un afeite el reverendo padre, porque en tan santo mártir no habría tal impericia como para exponer Cefas, esto es, cabeza, como si para la santidad de vida se requiriese tanta ciencia. Diré\apnota{20}{U dicamus.} sin embargo una cosa: que Cefas\apnota{21}{U Cephe aut Cephas.} debe ser siríaco y familiar también para los hebreos, valiendo tanto como Pedro o sólido\apnota{22}{M solidum vel validum.}, como refiere Erasmo siguiendo a Jerónimo. Añadamos\apnota{23a}{M addamus, gestrichen, Mc. a. R. ad damasum} también esto singular\apnota{23b}{et singulare fehlt in M.}: que Cusa, docto en caldeo y hebreo, atestigua en el libro\apnota{24a}{U libris.} de las Excitaciones que Pedro vale también\apnota{24b}{M U etiam valere.} tanto\apnota{24c}{fehlt in P.} como cabeza de la casa. Por lo cual Anacleto no padecía tanta impericia si interpretó Cefas por cabeza\apnota{26}{U interpretatur.}. Pero sea como fuere, no veo con qué afeite pueda negar los demás decretos de los santos mártires y confesores, como los de Clemente\apnota{27}{M add 22. dist., U 20. dist. W. A. vermutet 80. dist.}, Marcelo, Julio, Pelagio, Nicolás, Agatón, Símaco\apnota{28a}{U Simonis.}, Benito\apnota{28b}{fehlt in P. U Vigilli Benedicti.} etc. Por lo cual prefiero condescender con tantos padres y con el sacro concilio.

Pero\apnota{30}{U als Überschrift ad scrupos.} con grandísimo disgusto oí que el reverendo padre, contra el mandato del ilustrísimo príncipe —con el cual quiso el ilustrísimo príncipe que permaneciese intacto lo decidido por los sacros concilios—, el reverendo padre, sin embargo, en la de ayer

se endureció en su sentencia y propuso cuatro artículos condenados entre los husitas, los cuales él mismo\apnota{34}{M ipse reputat.} reputa católicos y evangélicos, hermoso patrocinio ciertamente para los bohemios\apnota{35a}{U pulchre certe Bohemis patrocinatur.}. Sin embargo, por reverencia al concilio\apnota{35b}{M conciliis}, quiso que hubiesen sido añadidos por algún impostor\apnota{36}{M aliquot.}. Ante todo, el egregio doctor y noble Jerónimo de Croaria\apnota{37a}{fehlt in M, Mc U Croatia.}, que cuidó de imprimir las actas del concilio\apnota{37b}{M consilii.}, hizo sacar traslado en forma auténtica, y habiéndose actuado la cosa tan manifiestamente\apnota{38}{M manifesta.} en Constanza\apnota{39a}{M Constantia.}, los husitas «no habrían callado»\apnota{39b}{Mc reticnisse, U reticuissent.} tan largo tiempo que algunos artículos fuesen espurios [295]. Ni «podrá» sospechar el reverendo padre\apnota{1a}{reverendus pater fehlt in M.}, como ayer había asumido, que se hizo por la pernición de aduladores\apnota{1b}{U adulatoris.} contra el pontífice Rom., pues Juan Hus fue quemado en el año 15\apnota{2}{U a. R. in Julio.}, Jerónimo, colega suyo\apnota{3a}{M U suus.}, en el año 16\apnota{3b}{M 16 Julii, U a. R. in Maio.}, electo Martín V por primera vez\apnota{3c}{M primo.} en el año 17, en el mismo día de Martín. Por lo cual los santos padres y varones incapaces de adulación, diputados para esto de entre todas las naciones principales, discutieron aquellos artículos, y la sacrosanta sinodo los condenó, reprobó y quemó al autor. Y por eso han de ser tenidos por todo buen cristiano «como condenados y reprobados». Ni me imponga el reverendo padre que quiero juzgar sobre aquellos artículos, porque ya han sido juzgados.

No aprovecha que el sínodo\apnota{9}{M U quod cum synodus.} recuerde que algunos son heréticos, otros temerarios, sediciosos y ofensivos a los oídos piadosos\apnota{10}{M animarum, Mc a. R. aurium.}; pues a cualquier orden en que haya reducido aquellos artículos, no pueden decirse cristianísimos y evangélicos, y aunque no haya asumido este «trabajo de defender» a toda\apnota{12}{M totum.} la sinodo en la condenación de los artículos\apnota{13a}{M condemnationem.}, digamos sin embargo algunas cosas\apnota{13b}{U doceamus, (Uv dicamus) tamen aliquid.}, y brevemente.

El primer artículo lo juzga católico y que es de Agustín, cap. 6 sobre Juan\apnota{15}{U Johanne.}. Digo: acaso el reverendo padre interpreta benignamente el artículo, pero porque la cosa no está sujeta al sermón, sino el sermón a la cosa, no fueron estúpidos los definidores del concilio\apnota{17}{M consilii.}. Es verdad que una es la santa y universal iglesia, pero que sea solo una, como es uno el número de los predestinados, según la inteligencia husita\apnota{19a}{M Hussiticum intellectum.}, es heretiquísimo\apnota{19b}{Mc (haereticissimum) velut amissa fide, die drei letzten Worte gestrichen.}, por lo cual él mismo quiso que los que están en pecado mortal, como perdida\apnota{20}{M a missa.} la fe, no estén en la iglesia, aunque el reino de los cielos es comparado\apnota{21}{M U a Christo comparetur.} a diez vírgenes, donde cinco eran prudentes y cinco fatuas, con las demás parábolas que sirven al mismo propósito.

Ni Agustín en el trat. 26 sobre Juan aporta una jota\apnota{23}{M U iotam.} al propósito husita, al recomendar la comunión de la «eucaristía»\apnota{24}{M ecclesiae gestrichen, Mc eucharistiae.}, sobre lo cual deseo el juicio del lector.

Así dice del otro artículo: «Dos naturalezas, divinidad y humanidad\apnota{25}{U divina et humana.}, uno es Cristo». Nada adujo en favor de esto, sino que es de fe. En Atanasio\apnota{26}{U Athana.; W. A. Athanasiano.} leemos de otro modo: «Dios y hombre uno es Cristo, no deidad y humanidad».

En otro artículo: «la división inmediata de las obras humanas»\apnota{28}{fehlt in U, in Uv nachgetragen.} etc., la cual él admirablemente retuerce contra mí y, como si yo le hubiese dado un juez\apnota{29a}{U citatus.}, proclamó a su\apnota{29b}{M suam.} colega triunfador sobre mí\apnota{29c}{de me fehlt in U.}. Sobre lo cual todos vosotros, los que estuvisteis presentes, podéis dar testimonio cuán verdaderamente «dijo el reverendo padre»\apnota{30}{Mc 'dixit'.} y qué encomio «cantó» antes de la victoria\apnota{31}{ante victoriam fehlt in M.} en Wittenberg, ya cantó también otro\apnota{32}{M alium; cecinit fehlt in M U.}, como perdida la victoria.

A la cosa.

Ese artículo nunca lo reputé\apnota{34}{M U illum; nunquam reputavi fehlt in U}, nunca lo reputaré cristiano, y por esto lo rechacé siguiendo a Gregorio de Arim. 28 dist. 2. Ni se trató\apnota{35a}{M fuerit.} entre nosotros si «puede»\apnota{35b}{M fuerit, gestrichen, Mc possit.} darse obra indiferente, o si toda vida de los infieles\apnota{36a}{P fidelium.} es pecado, y\apnota{36b}{M aut, U vel, Uv = P.} cosas semejantes que miran a este artículo. Por eso, lo que insinúa paliando los errores husitas: que alguna vez las palabras se entienden siniestramente\apnota{38a}{M sinistra.}, como le aconteció al Señor Cristo, las cuales\apnota{38b}{Mc et quae.} —pregunto—, ¿qué comparación de la luz con Belial? que\apnota{39}{fehlt in U.} compara en esto a Juan Hus con Cristo, no habiendo sufrido que yo, en cuanto a la sugilación de la verdad parcialmente admitida, comparase\apnota{40}{M admissa ne compares.} los sacros [296] concilios\apnota{1}{U consilia.} con las sagradas escrituras, en los cuales, en unos y otras, hay verdad indudable e infalible. Pues a ningún cristiano debe mover\apnota{2a}{M moveri.} lo que\apnota{2b}{U et quod.} dice: que los concilios\apnota{2c}{M U consilia.} son hombres y así criaturas, y por tanto pueden pecar, pues si yerran —como fue el concilio\apnota{3}{M consilium.} Efesino condenado por el papa León, creo, como el Ariminense, el «Aquisgranense»—, entonces no son concilios, sino conciliábulos. Por lo cual más bien debemos sostener con constantísima fe que cuanto los concilios legítimamente congregados determinaron y definieron en aquello\apnota{6}{M U his.} que es de fe\apnota{7a}{U et diffiniverint; et fehlt in P M.} es certísimo\apnota{7b}{P certissima, M certissimam.}. Pues así Cristo permanece con nosotros hasta la consumación del siglo, y «si dos se congregaren en mi nombre etc.»\apnota{9a}{M Christi (statt meo); etc. fehlt in M U.}. Es rancio decir que, porque son hombres\apnota{9b}{M U sunt.}, en

\noindent\textit{[test-p03: impresas 59-60 (test pp. 5-6). Claves: serie p.59 simples; serie p.60 con sufijo 2/b donde la numeración del original colisiona (18, 24a, 25a, 25b, 26).]}
\noindent\textit{[test-p04: impresas 61-62 (test pp. 7-8). Serie p.62 con sufijo -62 donde colisiona (31, 34, 3a, 3b). El fragmento cruza el salto 61->62 a mitad de frase ("scripturae / dicit in fine"), sin guion. Termina a mitad de frase ("quod hoc" -> sigue p.63).]}
\noindent\textit{[test-p05: impresas 63-64 (test pp. 9-10). Serie p.64 con sufijo -64 donde colisiona (19a, 19b, 27, 29). Cierra las 10 pp. de prueba (impresas 55-64). Termina a mitad de frase (sigue p.65, fuera del test).]}
\noindent\textit{[nota 32: venía partida entre pp. 92-93; texto completo fusionado en [\textasciicircum{}092-32] de pp-092.md]}

\noindent\textit{Notas del aparato sin llamada en el cuerpo:} [14] P U Mc baptizabat, M baptizabas, W. A. baptizarat.; [18] M sancto Gregorio.

concilio, y que así puede pecar o errar\apnota{9c}{M et ita peccare (Mc add errare) posse, U et ita errare possunt, Uv = P.}; pues aunque como hombres sean falibles, ha de creerse que el verdadero concilio legítimamente reunido\apnota{10}{M in concilio legitime congregato, U in consilium legitime congregatum.} no se rige por sentido humano, sino por espíritu divino, de donde\apnota{11}{regi credendum est fehlt in U; U ut (statt unde); Uv = P.} están suficientemente probadas\apnota{12a}{M probata sunt, U probata sint (Uv ? = M).} las cosas que el concilio aprueba\apnota{12b}{U probavit.}, y no han\apnota{12c}{fehlt in U} de ser discutidas ulteriormente por la singularidad y testarudez de cualquiera\apnota{12d}{M U cuiusque.}\apnota{13}{M singularitate et capitositate, Mc a. R. capciositate, U singularitatem vel capitositatem.}, sino que debemos cautivar el entendimiento en obediencia de la fe. Por lo cual me maravillo\apnota{14a}{M minor, Mc miror.} de que, estando por negar el texto de los derechos\apnota{14b}{U Siculam (statt iurium), Uv trägt iurium nach.}, admita como comentarista\apnota{15}{M admittit.} al Panormitano.

Así, en cuanto a lo que por último adujo de las palabras del apóstol Pablo\apnota{16}{U patet} sobre el fundamento, siempre he dicho que es sentir de todos los católicos que Cristo es el fundamento de la Iglesia\apnota{17}{Christum esse fundamentum ecclesiae fehlt in U, ecclesiae fehlt in M.}, que Cristo es la cabeza de la Iglesia, pues él es\apnota{18}{MU enim est.} el Señor principal. Pero porque constituyó a Pedro como vicario y le confirió potestad sobre toda la Iglesia, como\apnota{19}{U ubi.} atestigua\apnota{20b}{M testavit.} san León en el can. Ita Dominus 19. dist.\apnota{20a}{alle Quellen 21. dist.}. Y me maravillo sobremanera de que, dejado a un lado el texto de san León, el reverendo padre recoja allí una\apnota{21}{fehlt in M.} glosa contraria al texto. Convenía en todo caso, al ascender Cristo\apnota{22a}{U Immo tantum confert Christo ascendenti (W. A. Christus ascendens).} a los cielos, que en su lugar\apnota{22b}{M suo.} constituyese aquí un vicario como cabeza, al que se recurriese en las dudas y en los demás casos emergentes, para que toda la Iglesia de Cristo no fuese un «uno» confuso, un caos\apnota{24}{M caus, Mc chaos.} anaxagórico. Esta sentencia la sostiene toda la Iglesia, como está en el responsorio\apnota{25}{M P responso.} con que la Iglesia habla de san Pedro: Tú eres el pastor de las ovejas, príncipe de los apóstoles. A ti te entregó Dios todos los reinos del mundo. Así oramos en la Iglesia. Más aún, el divino\apnota{27}{M U divinus.} Francisco, aprobado con los sellos de los estigmas, enseñó a sus hermanos y a todo el orbe de la tierra a obedecer al sumo pontífice, y esto lo preceptuó en la regla. Por eso canta de él la Iglesia: Francisco, varón católico y todo apostólico, enseñó a guardar la fe de la Rom. Iglesia. Así muchísimos sumos\apnota{31a}{M romani.} pontífices, así los sagrados concilios lo definieron\apnota{31b}{M U diffiniverunt.}, cosas que por brevedad no aduzco. Muchísimas sentencias de los santos padres suenan y aprueban lo mismo. Todo lo cual no descuidaré\apnota{34}{M negligabo, Mc negligam.} aducir a su tiempo, cuando vea el tratado del reverendo padre sobre este asunto.

Si, sin embargo, el pri-[297]mado del Rom. pontífice se apoya sólo en derecho humano y es elevado por consenso del pueblo, ¿de dónde entonces tiene el reverendo padre\apnota{2}{U reverendi patris.} el privilegio de la mendicidad? ¿De dónde ese hábito de la religión? ¿De dónde la potestad\apnota{3}{M potestas praedicandi (also doppelt!)}

de oír confesiones, de predicar, con otros innumerables\apnota{4a}{M innumerandis.} privilegios, «con los que son abrumados por la sede apostólica\apnota{4b}{M U a sede apostolica.}», contradiciéndolo con frecuencia arzobispos, obispos\apnota{5}{U episcopis archiepiscopis.} y curas? Los cuales, sin embargo, como hijos de la obediencia escuchan al Rom. pontífice y en él a Cristo, para alcanzar así por el bien de la paciencia la vida eterna.

Estas breves cosas he querido añadir\apnota{8}{M adducere.} en el presente asunto de la disputa, según me lo dio el tiempo, habiendo de aducir por lo demás más, si no me arredrara el fastidio y el mandato del ilustrísimo príncipe\apnota{10}{Mc U iussus.}.

El D. D. Martín.

Dos cosas me incumben hacer: primero\apnota{12}{M U primum.} responder a las confutaciones del egregio D. D., segundo objetar contra él según el pacto de la convención. Y porque ya durante tres días he respondido y se ha consumido\apnota{14}{M U consumptum.} el tiempo, de modo que ni una sola puedo cumplir, digo brevemente:

Que todo lo aducido por el egregio D. D. ha sido dicho ciertamente con grandísima copiosidad, pero del todo nada\apnota{16}{M nihil prorsus.} al propósito, pues durante tres días prometió que actuaría\apnota{17a}{U acturum se.} por derecho divino contra mí, y hasta ahora no se ha oído ni una\apnota{17b}{fehlt in M.} autoridad de derecho divino, fuera de aquella de Mt 16, que todos oímos, de cómo diversamente es tratada por los santos padres, y de que la mayor y más sana parte de ellos sostiene la sentencia en mi favor. Él mismo adujo en su favor un solo lugar de Agustín y otro de Ambrosio; todo lo demás fueron palabras de Anacleto, Marcelo, Pelagio etc.\apnota{21}{M U Pelagii, Anacleti etc. [etc. fehlt in U].}, de quienes escribieron derechos humanos. Cipriano, en cambio, casi\apnota{22a}{M fero.} todo siente\apnota{22b}{M sit.} conmigo, lo cual encomiendo al juicio del prudente lector. Igualmente «el clarísimo\apnota{23}{M clarus.} texto» de Jerónimo a Evagrio, y sobre Tito. Luego Gregorio en el registro por casi seis epístolas seguidas, después de esto la historia y el uso de la iglesia oriental. Pues nada hace al caso que el D. D. haya alegado a los griegos cismáticos, «a Nestorio, Macedonio, Acacio, Eutiques y demás»\apnota{26}{Macedonium, Achatium, Eutychem fehlt in M U, Uv = P.}, pues aquéllos no eran la iglesia oriental. Podría «también» yo describir así a la iglesia latina por medio de «pelagianos, maniqueos, jovinianos», vigilantianos\apnota{28}{M U possum et ego [M ergo] latinam ecclesiam sic describere per Pelagios [M Pelagianos], Manichaeos [fehlt in M], Jovinianos, Vigilantianos [U add Helpidianos].} y monstruos semejantes. Pero líbreme de llamar cismática a toda una iglesia por unos pocos malos y cismáticos\apnota{30a}{M appellam.}.

Lo demás «no podré proseguirlo». Y\apnota{30b}{fehlt in M.} me reservo especialmente los artículos de Juan Hus a mí objetados «y entre ellos el novísimo sobre los actos humanos» para el cálamo y el papel, y todo lo demás que tenía\apnota{32}{M haberem.}

que objetar, para prestar por escrito lo que aquí no puedo hacer\apnota{33}{U possim.}, excluido por los límites de espacio, y quiero en esto haber cedido y obedecido a los mandatos del ilustrísimo príncipe\apnota{34}{U principi.}, clementísimo patrono de los estudios\apnota{35}{M patronum [clementissimi fehlt]; P clementissi.}, de lo cual os requiero a vosotros, señores notarios, y llamo por testigos a los oyentes.

Hecho como arriba, presentes allí los honorables Olavo Roters de Osvaberga y Celestino Goritio de Gorlitz, clérigos de las diócesis de Colonia y de Meissen, testigos especialmente llamados y legítimamente requeridos.

Juan Graumann, notario, lo suscribió\apnota{36}{Die Worte von Actum an fehlen in M U.}.

[298] El jueves, día siete del mes de julio, a la hora séptima de la mañana\apnota{2}{Die Zeitbestimmung fehlt in M, in U lautet sie: Die septima Julii anno 1519. mane hora septima.}.

El D. D. Eck.

Porque el reverendo padre tuvo ayer la hora de responder a nuestras cosas, para que así hubiésemos ejecutado los mandatos del ilustrísimo príncipe, donde al responder apeló muchísimo por la palabra\apnota{6}{per verbum fehlt in M.} «pase» al cálamo, sin embargo, habida deliberación, propuso que ya respondería más plenamente. Yo no me opuse, por amor a la verdad y para su esclarecimiento. Me admiro de cómo el reverendo padre con gran estruendo se atreve a tronar ante tantos varones\apnota{9}{M U doctis viris.} que yo habría dicho que probaría por derecho divino la monarquía y el primado en la Iglesia, cuando sin embargo sólo alegué dichos de los santos padres y\apnota{11a}{fehlt in M U.} decretos de los pontífices\apnota{11b}{M pontificum et sic iura humana.}, cuando sin embargo yo\apnota{11c}{fehlt in U.} le aduje el derecho divino de Mt 16 sobre\apnota{12a}{M pro.} la promesa de la edificación sobre\apnota{12b}{U supra.} la piedra, de que\apnota{12c}{M quo.} Cristo rogó especialmente por la fe de Pedro, de que especialmente le encomendó\apnota{13}{U promisit.} confirmar a los hermanos, de que especialmente le encomendó en presencia de Juan y de Santiago apacentar las ovejas, por tres veces. Y que este derecho divino ha de entenderse así, lo probé entonces con los santos padres, los santos\apnota{16}{fehlt in U.} sumos pontífices y mártires, el consenso de los concilios y de todas las universidades. Por lo cual por derecho divino quedó probado que hay primado en la Iglesia\apnota{18a}{U ecclesia dei.}, a no ser que el reverendo padre anteponga\apnota{18b}{M U proponit.} su propia inteligencia al entendimiento de tantos santos padres, sumos pontífices y del sacro concilio\apnota{19}{U consilii.}, pues el laudable concilio Constanciense\apnota{20a}{U consilium.} condenó\apnota{21a}{M damnavit.} tal artículo entre los pestilentes errores\apnota{20b}{M articulos; errores fehlt in U.} de Juan Hus. A lo cual el reverendo padre calla. Pero sin embargo\apnota{21b}{Sed nihilominus fehlt in U.}, si quiere ofrecer\apnota{21c}{M U affere.} cosas mejores, para esclarecimiento de la verdad, lo que ayer habría podido

haber hecho según el mandato del ilustrísimo príncipe\apnota{23a}{M iussu.}, estoy dispuesto a oírlo, con el pacto de que, mientras\apnota{23b}{M tum.} desempeña el oficio de respondiente, ya al fin del acto no se traslade taimadamente a la persona del oponente, sino que cualquier cosa que haya de oponer contra esta ineludible verdad, estoy dispuesto a responder a estas cosas aducidas y a disolverlas, de lo cual protesto.

El D. D. Martín.

Como ayer se había establecido\apnota{28}{M esse, Mc esset.} en nombre del ilustrísimo príncipe, nuestro patrono, que esta materia se consumase\apnota{29}{M consumeretur; ista fehlt in M.} aquel día, y el egregio D. D. hubiese consumido sin necesidad y más odiosamente de lo que convenía a este espléndido auditorio la mayor parte del tiempo, a mí para una respuesta de tres días sólo me quedó\apnota{31}{M mihi.} una horilla, en la que a la vez respondiese a tanta farrago\apnota{32}{M varragini.} suya y objetase. Plugó después por gratuito afecto al mismo ilustrísimo\apnota{33a}{M illustrissimi.} príncipe nuestro que\apnota{33b}{U ut.} se me diese a mí\apnota{33c}{M mihi daretur.} copia de responder y de objetar\apnota{34a}{M et opponendi et respondendi.}, por lo cual doy a su ilustrísima\apnota{34b}{M illustrissimi.} gracia las máximas gracias. Y por eso, viniendo al propósito, antes de continuar la comenzada respuesta de ayer, disiparé primero esta niebla matutina del egregio D. D.\apnota{36}{U d. d Eccii.}, donde a su modo con bastante odiosidad siempre me objeta las pravidades husitas, no ob-[299]stante que tantas veces me he purgado de ellas. Y por eso prologaré un poco en alemán, porque entiendo que soy muy mal oído entre el vulgo.

Decía\apnota{3}{U Declarabat.} pues que no impugnaba, ni podía ser cristianamente reprobado\apnota{4}{U impugnari.} por nadie el primado y la obediencia de la Romana iglesia. Sino que había sido llevado al punto de tener que declarar que ese primado no es de derecho divino, así como tampoco puede reprobarse la potestad imperial entre los germanos\apnota{6}{U potestatem apud Germanos.}, aunque no fundada en la sagrada escritura.

El D. D. Eck.\apnota{7}{U quamvis d. Eccius.}

dijo que era verdad que ésa era la controversia, si aquel primado\apnota{8}{U primatus ecclesiae Romanae.} es de derecho divino, cosa que juzgaba\apnota{9}{U constat.} reprobar como husítica y artículo condenado\apnota{10}{U condemnatur.} por el concilio Constanciense, y que habría que dolerse si los corazones de los cristianos fuesen tan fríos que nadie defendiese esto.

Después el D. D. Martín prosiguió en latín diciendo\apnota{11}{U post hoc}.

Que en cuanto aduce esta autoridad de Mt 16, contendiendo que allí fue prometido el primado, bastante se ha respondido en lo precedente, que

el egregio D. apenas tiene dos autores y esos mismos ambiguos para su sentencia, cuando sin embargo la mayor parte establece que Pedro fue allí\apnota{15}{fehlt in M.} la persona de todos los apóstoles y de los fieles, cosa que también urge el mismo tenor consecuente del texto\apnota{16}{W A. continentia.}, pues Cristo interroga a todos, y no sólo a Pedro, diciendo: Mas vosotros ¿quién decís que soy yo? Y allí Pedro, boca de los apóstoles, como dice Crisóstomo, y persona de los apóstoles, como dicen Jerónimo y Orígenes y Beda y Agustín y Ambrosio, responde por todos. Por eso dije que nada se prueba de esta autoridad sobre el primado\apnota{21a}{M primatui.}. Igualmente, que el texto mismo\apnota{21b}{M U textus ipse.} no dice así: Tú eres Pedro, y sobre ti, sino sobre esta piedra, indicando manifiestamente otra piedra, de la que Pedro fue dicho. Por lo demás, como son derechos divulgados y los mismos Rom. pontífices atestiguan que la interpretación de la escritura vale más en los doctores que en ellos mismos, pero en las causas por decidir\apnota{25}{U descendendis (W. A. verbessert nach Löscher.)} al contrario.

Sobre aquel artículo husítico dije bastante ayer, que aún no se ha probado que sea herético. Esto por respuesta matutina.

Volviendo a lo de ayer, cuando el egregio D. D. en segundo lugar opuso que la deposición de Crisóstomo fue hecha ilegítimamente porque se hizo por odio de la emperatriz, niego la consecuencia; no fue por eso ilegítima la potestad de ordenar y deponer establecida por el sacro concilio Niceno\apnota{31}{M concilium tragens.}, porque por acaso tuvo mal uso. Pues en la misma historia\apnota{32}{M U etiam historia.} se escribe que el Rom. pontífice Víctor\apnota{33}{fehlt in M.} mártir quiso excomulgar a los obispos de Asia. Pero ellos a su vez como superiores les mandaron silencio, y les ordenaron\apnota{34b}{M perturbarent iussarent [Mc iusserunt].} no perturbar a la iglesia de Dios\apnota{34a}{fehlt in U, in Uv nachge-Nicenum.}, lib. 5 de la historia eclesiástica\apnota{35}{libro 5. eccles. hist. fehlt in M.}. Pero también Ireneo de Lyon de la Galia\apnota{36a}{M Tirenaeus Luctunens Galli (so!)}, obispo, coerció al mismo Rom. pontífice\apnota{36b}{U episcopum.}, y eran entonces tem-[300]porales en Grecia Epifanio, celebradísimamente alabado por Jerónimo, Gregorio Nacianceno, Basilio\apnota{2}{M Basilius etc.}, cristianísimos varones, y sin embargo nunca estuvieron bajo el Rom. pontífice, sino que según los estatutos\apnota{3b}{M constituta.} del concilio Niceno\apnota{3a}{M Nicelmi.} eran ordenados por los obispos de las provincias.

En cuanto a que el egregio D. doctor aduce que Julio I excomulgó al emperador Arcadio, nada hace al caso. Pues también Bonifacio VIII se atrevió a querer arrojar de la sede al rey de los francos. No prueba\apnota{6}{U probatur.} que sea derecho todo lo que los Rom. pontífices por humana fragilidad alguna vez intentaron. Ni es maravilla que aquellos santos\apnota{8}{M U istos.} varones fuesen vejados por la tentación de la ambición

cuando más de una vez los santísimos apóstoles aun en presencia de Cristo padecieron lo mismo.

En cuanto a que dice que al obispo «jerosolimitano» le fue dada prerrogativa de honor, no primado\apnota{11}{M praerogativam datam (Mc a. R. gratiam) non primatum, U datam honoris praerogativam [non primatum fehlt].}, rectamente lo dice, pero no prueba que haya primado\apnota{12}{U esse hinter iure divino, M hat esse an beiden Stellen.} del Rom. pontífice, al menos por derecho divino.

También allí dice óptimamente el egregio D. D. que el sexto sínodo dio a la Romana iglesia el primado, pero no los sínodos precedentes. Esto es pues lo que yo quise, que por estatutos sinodales y por derecho humano, con el consenso\apnota{15}{M consensum, Mc consensu.} de los fieles de Cristo, al que no es lícito resistir, fue dado ese primado. Como también el divino Gregorio dice en el registro con estas palabras\apnota{17}{U dixit.}: Sabéis que el venerable sínodo Calcedonense ofreció a esta sede apostólica el primado y sin embargo ninguno de mis predecesores presumió\apnota{19}{M U praesumpserit.} jamás aceptarlo. Por lo demás nada hace que el egregio D. D. diga que Gregorio no rechazó\apnota{20}{M reiecisse.} el primado, sino que lo afirmó. Lea el piadoso lector las epístolas\apnota{21}{M [epistolam.} de Gregorio, y hallará que Gregorio prohibió el primado y no vindicó el pontificado ecuménico\apnota{22a}{M icumenicum [et fehlt], Mc universalem.}, como arriba por lo relatado en la dist. 99.\apnota{22b}{P M U dist. 95.} bastante probamos.

A la autoridad de Agustín, a Inocencio Rom. pontífice y a la de éste a aquél respondo: Lea el lector las palabras de uno y otro y juzgue si el egregio D. D. impuso nuevas significaciones a los vocablos.

Sobre los veinte\apnota{27}{U „22".} años de la iglesia antes de la Romana dijo, habiendo sido trasladado con Pedro el primado a Roma, queriendo quizá con este proverbio: Donde está el papa, allí está Roma. Respondo: había que probar el primado que se trasladaría, no la traslación, que de buen grado concedo.

Adujo también la autoridad de Agustín contra los maniqueos, que sospecho es del libro 2 de\apnota{32a}{M beati Augustini 2. lib.} la doctrina cristiana\apnota{32b}{U doctrinae christianae.} c. 8, que en cosa dudosa para certeza de la fe vale la autoridad de la iglesia católica\apnota{33}{U ecclesiae catholicae [catholicae fehlt in M].}, la cual es firmada\apnota{34}{M firmati.} por las sucesiones de los obispos que desde las mismas sedes de los apóstoles se suceden entre sí y por el consenso del pueblo. Nada dice de la Rom. iglesia\apnota{35a}{U loquimur.}, sino de la universal, y hace por mí contra el D.\apnota{35b}{fehlt in M.} D., porque Agustín describe las sucesiones de las diversas sedes apostólicas entre sí y el consenso\apnota{36}{U consensu.} del pueblo, pero no la confirmación de las otras sedes por una sola sede.

Al\apnota{39}{U nunc ad.} asunto mismo.

[301] Sobre la autoridad de Mt 16: Tú eres Pedro, y sobre esta piedra etc.\apnota{1}{et super hanc petram etc. fehlt in M U.} y a los autores aducidos dice que no respondí, cosa que encomiendo

a vosotros oyentes y lectores, si es verdad. Lo de Agustín en las retractaciones\apnota{3a}{M retractatione.}, dejando la elección al lector\apnota{3b}{fehlt in P.} sobre dos sentencias suyas de la misma autoridad, lo opuso diciendo\apnota{4}{M dicit.} que abraza ambas sentencias, lo cual me place\apnota{5a}{M Uv et ego facio [statt mihi placet].}, y así nada prueban\apnota{5b}{M U probat.} contra mí.

La segunda autoridad de derecho divino aducida por su dignación\apnota{6}{M U dominatione.} es ésta: Apacienta mis ovejas, Jn último [21 v. 15 ff.], y\apnota{7}{fehlt in M.} me gozo de oír por fin tras tres días alguna vez una autoridad de derecho divino, en que firmísimamente se apoye su sentencia, pero veamos también\apnota{9}{fehlt in M.} ella misma.

Primero esta autoridad parece entenderse de dos modos, de un modo\apnota{10}{fehlt in U.} por el verbo amar, esto que es, fiarse de sí mismo\apnota{11}{M videre, P a. R. lectum fuit videre media longa.}, y que todo le sea lícito, y por el verbo apacentar, lo mismo que ser el primero y reinar. Y así será el sentido: Pedro, si me amas, esto es, si buscaras lo tuyo e hicieres todo lo que place\apnota{13}{M placet.} a tus aduladores, apacienta mis ovejas, esto es\apnota{14}{id est fehlt in M.}, sé el primero y señor de todos. Este sentido no lo tengo en mi códice.

De otro modo: Si me amas, esto es, si\apnota{15a}{fehlt in U.} te niegas a ti mismo\apnota{15b}{U te ipsum.}, si pones tu alma por mí, si aun toda dignidad desdeñares y no\apnota{16}{M U nihil.} amares sino a mí, como egregiamente en este lugar expone Agustín\apnota{17}{hoc loco fehlt in M; U exponit Augustinus in eodem loco.}: Apacienta mis ovejas, esto es, enseña, predica la palabra\apnota{18}{M verbo.}, exhorta, ora, precede con buen ejemplo. Pues la palabra griega\apnota{19}{M U verbum Graecum.} en este lugar no significa simplemente regir y apacentar, sino cuidar suave y blandamente y hacer todo para que nada falte a las ovejas\apnota{20}{M omnibus.}, y creo que este sentido es el evangélico.

Por lo demás ruego al D. D.\apnota{21}{M et proinde rogo doctorem.} Eck que ore\apnota{22}{M U roget.} conmigo al Señor Jesucristo, para que quiera inspirar no sólo al sumo, sino a todos\apnota{23a}{U etiam omnibus.} los pontífices\apnota{23b}{non solum summo, sed omnibus pontificibus fehlt in M.}, que crean que estas palabras les atañen. No hay duda de que todo el orbe recibiría\apnota{24}{M accepturus.} con manos abiertas y lágrimas derramadas a tal varón que quisiese regirse según estas palabras. Por lo demás, si el egregio D.\apnota{25}{U egregius dominus.} D. inspecciona diligentemente el texto, no se concede allí derecho ni privilegio, como se piensa, sino que se impone precepto, y se injunge un trabajo hoy intolerable, y se encomienda un oficio a quien de hecho está constituido ya en el primero ya en el medio lugar, y no se constituye primero. Que si urge que el oficio no puede administrarse si primero no se tiene el lugar de la potestad, place, pero entonces con otra palabra\apnota{31}{U alio modo.} ha de construirse ese lugar, pues aquí esta palabra no encomienda sino el oficio.

Por lo demás rectamente hicieron los santos padres, que dicen que este lugar atañe a todos

pues esto nadie puede negarlo, si a Pedro le fueron aquí encomendadas\apnota{33}{M sunt hic.} todas las ovejas, porque él no apacentó a todas, como dice el bienaventurado Agustín y el libro de los Hch evidentemente demuestra\apnota{35}{M evidentissime ostendit.}, y así fue inobediente a Cristo. Ni puede decirse que, aunque por sí no apacentó, sin embargo lo hizo por otros subordinados. Primero, por conceder lo que podemos devolver, no fue dicho a Pedro: Apacienta por medio de otro, sino a él mismo en propia persona: Apacienta tú\apnota{38}{W. A. verbessert nach Walch non esse dictum Petro ipsimet in persona propria „pasce tu", sed „pasce per alium".}, y sin embargo esto convence de que ni ningún apóstol ni al menos Pablo fueron subordinados a Pedro, quienes apacentaron muchas más\apnota{39}{M multos, Mc multo.} ovejas [302] que Pedro. Por eso la inteligencia del verbo ha de pedirse\apnota{2}{M esset.} de toda la escritura y de la circunstancia de los hechos, porque esta autoridad o no atañe sólo a Pedro o no a todos\apnota{3a}{M U omnes oves.}.

En tercer lugar digo\apnota{3b}{M tertium; dico fehlt in M U.}, si por estas palabras se probara el primado, será incierto y nulo el primado, porque no se pone el primado sino bajo condición de amor, pues no conviene romper la escritura en diversas partes, sino más bien concertarla. Y por eso, como es incierto quién ama a Cristo, también nos será incierto\apnota{7}{M U erit nobis} quién sea nuestro pastor. Y si a él no le es necesario el amor, tampoco nos será necesario a nosotros reconocer al pastor. Más rectamente pues se dice que esta palabra es un precepto dado a todos los prelados de la iglesia, para que, despreciadas las riquezas, la dignidad, aun el mismo primado, por fin con vida y muerte se gasten por las ovejas de Cristo. ¿Quién es éste, y lo alabaremos?

Paso lo que adujo de Crisóstomo sobre el vértice del colegio, y cómo a Pedro fue entregada la presidencia del orbe\apnota{14}{M tradita praesentia (so!) totius orbis.} de la tierra, porque todo lo concedo, que sea el primero en honor, pero no solo en la administración, al menos\apnota{15}{in fehlt in P, saltem in M.} por derecho divino.

Sobre el caso de Pedro, Gál 2, donde de nuevo el egregio D. D. más odiosamente me objetó el artículo husítico de que ningún señor espiritual o civil que\apnota{17}{M qui esset sine, die beiden letzten Worte gestrichen.} esté en pecado mortal\apnota{18a}{M U peccato mortali.} —si dijese que Pedro por eso\apnota{18b}{M ideo Petrum.} no es prelado porque pecó mortalmente—, respondo: Esto no se trataba, y esta disputa está fuera del punto [ad rhombum]\apnota{19}{M romgum.}. Sé\apnota{20}{M U scio et optime scio, Uv = P.} que el prelado malo no ha de ser rechazado. Por eso también yo condeno aquel\apnota{21a}{P artieu- (so!) Hussiticum illum, U Hussiticum articulum; illum fehlt in M U.} artículo husítico. Pero esto quise\apnota{21b}{M vult, U voluit.}, que Pedro, cuando en causa de fe dio escándalo, si no hubiese sido enmendado por Pablo, con mérito debió ser removido de la prelatura, pues el pastor hereje o el que simula\apnota{24}{M simulum, Mc simulat.} en peligro

de la fe, peca gravemente. Pues por esta simulación de Pedro perecía desde el fundamento la fe de Cristo, como dice Pablo, por eso antes habría de hacerse, si Pedro no hubiese sido enmendado\apnota{26}{fehlt in M.}, que Pedro fuese tenido por nada, mucho menos por sumo pontífice. El pecado mortal daña a la propia persona, pero la herejía daña a la persona común y\apnota{28}{personae communi et fehlt in M.} a la comunidad. Por eso doy gracias al señor D. de que al menos esto aprendí de esta disputa, que edificar sobre la piedra significa\apnota{29}{M U super petram significat.} estar bajo el prelado, sea bueno, sea malo. Pues esta significación de la edificación y de la piedra confieso que antes ignoraba\apnota{31}{M me ignorasse.}, porque dijo: aun si Pedro hubiese caído mortalmente, aún habría permanecido piedra, esto es primero\apnota{32}{M primus praelatus, U primus et praelatus.}.

En cuanto a que dice que antaño los obispos no eran confirmados si no prestada cédula\apnota{33}{M ceda.} de fe, lo admito, estoy contento, pero\apnota{34}{M U verum.} nada al derecho divino, pues nada derogo al pontífice, cualquier cosa que se le atribuya.

Después de esto adujo a Nicolás Cusano y excusó a Anacleto, varón santo, porque no le fue necesaria tanta ciencia\apnota{37}{U tanta scientia.} a tan santo varón. Lo confieso. Pero al sumo pontífice, sea santo, sea no santo, le es torpe ignorar los evangelios, siendo pastor de las ovejas y doctor del evangelio. Como\apnota{39}{M dum.} Cefas es interpretado por Juan evangelista cap. 1 [v. 42] diciendo Pedro\apnota{40}{M Petrum, fehlt in U.}: Tú eres Simón, tú [303] serás llamado Cefas\apnota{1}{M cephas.}, lo que se interpreta Pedro, cuando entonces el hijo del trueno así truena, el que conoce\apnota{2}{P novat.} todas las lenguas, con mérito debe callar la vanidad del terrenal cordel, sea Nicolás Cusano, caldeo o hebreo. No porque reprenda al egregio D. D., que dijo que Pedro vale lo mismo que cabeza de casa en caldeo, por autoridad de su Cusano, sino porque alabo al impertérrito disputador, que osa meterse en la arena teológica\apnota{6a}{U theologorum.} y ni siquiera estudió antes tanto en los evangelios\apnota{6b}{U evangelio.} que\apnota{7a}{M quod.} leyese\apnota{7b}{Johannis fehlt in M; U in Johanne legit Evangelista.} el primer capítulo del evangelio de Juan. Igualmente no tiene\apnota{7c}{M habeo.} en su libro qué es Cefas, lo cual perdono, pero lea el capítulo primero de 1 Cor [v. 12]\apnota{8}{M U cap. 1 ad Cor. 1 (Die letzte Zahl fehlt in M).} y lo hallará.

Igualmente en último lugar sobre los artículos bohemios\apnota{10}{U [item fehlt] Ultimum de articulis Bohe: [W. A. Bohemorum]} no quiere juzgar, aunque nunca desiste\apnota{11a}{M dicit [Mc a. R. licet] nunquam desistere.} de denigrarme, y yo sobre ellos paso\apnota{11b}{M ei conscio [statt de eis transeo].}, sin decir nada, sino que a Gregorio de Rímini en la dist. XXVIII, reprobado por el egregio D. D., yo lo apruebo. Pues es todo él nada más que Agustín y la divina escritura, resistiendo ciertamente a todos los doctores escolásticos, entonces\apnota{14}{M eum, Mc tum.} máximamente a Aristóteles, pero aún por ninguno confutado\apnota{15}{U confutatus est.}.

Convengo con el señor D. en que los estatutos de los concilios en estas\apnota{16}{U iis.} cosas que son de fe han de abrazarse de todo modo, sólo esto me reservo, que también ha de guardarse que el concilio alguna vez erró y puede\apnota{18}{U aliquando posse.} errar, sobre todo en estas\apnota{19a}{U iis.} cosas que no son de fe, ni tiene el concilio autoridad\apnota{19b}{fehlt in M.} de forjar nuevos artículos en la fe, de otro modo tendremos al fin tantos artículos cuantas\apnota{20}{M quod.} opiniones de hombres.

En cuanto a que al fin adujo a san Francisco, que debemos obedecer al sumo pontífice, me maravillo contra quién lo dijo\apnota{23a}{M U dixit.}. Igualmente paso estas criminaciones\apnota{23b}{M [item fehlt] taceo istam criminationem.} sobre los hermanos\apnota{24a}{fehlt in M.} mendicantes, este miserable hábito también atormenta al D. doctor\apnota{24b}{U miser [M nisi] iste etiam habitus cruciat egregium d. doct.}. Yo por mi sentencia pronuncio que quiero que no haya\apnota{25}{M U vellem me esse nullum.} ninguna orden de mendicantes. Esto\apnota{26}{U Ecce [M etc.?] haec.} tuve que responder a lo objetado por el egregio D. D. sobre la proposición 13. Por lo cual, acabada ahora la respuesta, resta que también yo mueva contra el egregio D. D. las cosas que me mueven\apnota{28}{quae me movent fehlt in M.}.

Aquí el señor doctor Martín quiso pasar a objetar por las partes de oponente a él devueltas por pacto, como decía. Pero el señor doctor Eck pretendía que entre lo soluble ya se habían dicho algunas cosas a modo de oponente, las cuales primero querría disolver según la protestación sobre esto hecha. No obstante permitió que proceda ya la oposición del señor doctor Martín. Comenzó pues\apnota{34}{in U lautet der Bericht über den Zwischenfall: Hic intonuit Eccius, se obiectis velle respondere, quasi et ipse respondens fuisset. Respondit Martinus, hoc pacto nunquam finiendam disputationem, accusans Eccianam petulantiam, quae solet singula etiam verba ancupari. Contigit prorsus virgula divina, ut clementissimus Dux Georgius interesset, cuius nutu et pio sceptro concessum est Martino, impetere Eccianam conclusionem, quae contraria est huic iam discussae. In M fehlt der Bericht ganz.}

El D. D. Martín\apnota{35}{M Martinus opponens Egregie domine doctor.}.

Contra la sentencia del egregio D. D. parece militar lo que con bastante fuerza saltó al confutar, y se escribe en 1 Cor 3 [v. 5], donde Pablo, rebajando [304] la persona de todos los apóstoles, dice: ¿Qué es Apolo? ¿qué Pablo? ¿qué Cefas\apnota{2}{M Scephe.}? ministros por quienes creísteis. Donde Pablo quiere del todo quitar la ocasión de contención\apnota{3}{M contentionem, Mc contentionis.} y de cisma; niega que ninguna iglesia pertenezca a ningún apóstol, sino que todo es común a todos. De donde al fin del tercer capítulo [3. v. 21 ff.] dice: Todo es vuestro, sea Cefas, sea Apolo, sea Pablo\apnota{5}{M sive Pauli (so!) sive Apollo.}, mas vosotros de Cristo.

No vale la evasión con que se dice que contendieron no sobre la principalidad, sino sobre la singularidad de las personas\apnota{7}{M eos contendisse.}, siendo manifiesto el texto que contendieron sobre las dignidades de las personas, pues unos preferían a Pedro sobre todos, otros a Pablo, otros a Apolo, otros en cambio solo a Cristo. Se confirma por aquello de Gal 1 y\apnota{9}{M Gal. 1. 2., U Gala. 2.} 2, donde los Gálatas, seducidos por semejante contienda a causa de la primacía de Pedro que se les encomiaba, abandonaron a Pablo y a su doctrina como si fuese más indigna. Contra lo cual prueba con largo texto que nada importa al caso la majestad de Pedro o de los otros apóstoles, diciendo que él no fue enviado a ellos ni por hombre ni mediante hombre\apnota{13}{M a deo (statt ad eos).}, ni siquiera vio a Pedro ni aprendió de él\apnota{14}{M a Petro.}, sino que todo lo tuvo y lo transmitió sin Pedro. Si, pues, la autoridad de Pedro fue\apnota{15a}{M fuerit, Mc fuit.} necesaria como\apnota{15b}{U et.} derecho divino, Pablo será en este lugar manifiestamente impío y blasfemo, como quien no quiere que se vea que ha sido enviado ni siquiera por Dios mediante hombre, y rechaza por completo\apnota{17}{M prorsum.} la autoridad de Pedro. En tercer lugar, y\apnota{18}{fehlt in M.} más abajo con mayor claridad, donde dice: A mí los que parecían ser algo nada me comunicaron, y cómo hayan sido alguna vez nada me importa, pues Dios no acepta la persona del hombre. He aquí que dice claramente que la cualidad de Pedro y de los otros apóstoles nada le importa, lo cual sería impíísimo decir si la cualidad de Pedro hubiese de guardarse por derecho divino. Con la misma impiedad diría: Dios no acepta la persona del hombre, cuando el derecho divino y lo que es de derecho divino lo preceptúa incluso bajo pena de eterna maldición. Por lo cual parece que esta primacía y majestad, o con cualquier nombre que se designe la persona o cualidad de Pedro, no se sostiene por derecho divino.

El D. D. Eck quiso responder ya repitiendo al aire las objeciones. Entonces, por mandato de los señores transmitido por el alguacil público, fue interrumpida la disputa y suspendida hasta la hora segunda de la tarde.

Hora segunda de la tarde del mismo día\apnota{26}{Zwischenbemerkung und Zeitangabe fehlen in M U.}.

El D. D. Eck.

Cuando el reverendo padre aduce contra mí tres autoridades del bienaventurado Pablo, ventiladas tiempo ha por los adversarios de la sede apostólica y aclaradas por los defensores de la sede apostólica\apnota{29}{M ab adversariis sedis [gestrichen, U sedes] apostolicae sedis defensores enodatos [offenbar in Unordnung].}, fácilmente le respondo, añadiendo los malos afeites con que se esfuerza por paliar su dogma y untar a los hombres la boca y los ojos. Digo lo primero\apnota{30}{M primo.}, que no pasé por alto, como me imputa, el lugar del Apóstol 1 Cor 3 [v. 5.], aunque él mismo lo adujo falsamente con Pablo y Apolo a Cefas, porque el texto desde el principio no lo contiene, sino que al final hace mención de Cefas, 'sino

al punto' la glosa interlineal confiesa que él es mayor que todos. No reciba, pues, el reverendo padre la Sagrada Escritura por su propia cabeza, sino según los santos.

Y digo que mi solución dada\apnota{35}{meam fehlt in M U, datam fehlt in M.} fue óptima e irrefutable\apnota{36}{M infallibilem.}, y nada vale el afeite del reverendo padre, con que se burla de los sencillos; testigo es el bienaventurado Jerónimo en el libro 1 contra Joviniano\apnota{37}{P Jovinianum 1.}, col. 18 para mí, donde dice expresamente que sobre Pedro fue edificada la Iglesia. Allí pasó por alto con fuerza\apnota{38}{M resiliit.} el reverendo padre y nada dijo\apnota{39a}{U dicit.}. Y más abajo\apnota{39b}{et infra fehlt in M.}: para que, constituida la cabeza, se quitase [305] la ocasión del cisma. Luego lo que el apóstol Pablo quiso evitar\apnota{1a}{M vitare.} —los cismas y\apnota{1b}{fehlt in M U.} las contiendas—, con ello no quitó\apnota{2}{M tollit.} el primado de la cabeza, pues los cismas 'no pueden quitarse mejor'\apnota{3a}{Mc possint.} sino por la cabeza. De donde ni una palabra hace\apnota{3b}{M dicit.} ni un ápice toca al propósito sobre la negación\apnota{4a}{de negatione fehlt in M.} del primado. Por lo cual falsamente ha comenzado hoy diciendo que apenas\apnota{4b}{M vix in me.} aduje dos s.\apnota{5a}{fehlt in M.} padres y ellos ambiguos, cuando aduje\apnota{5b}{M induxerim.} padres ciertos e indudables: a Agustín, a Jerónimo\apnota{6a}{M U Ambrosium, Hieronymum.}, a Crisóstomo, a s.\apnota{6b}{fehlt in U.} León, a Bernardo etc.\apnota{6c}{M et Bernhardum (statt Bernh. etc.)}. Y en cuanto a que me alaba por abrazar ambas sentencias de Agustín en el libro 1 de las retractaciones cap. 21, y afirma\apnota{8a}{W. A. verbessert affirmat et, wahrscheinlicher ist, dass mit facere der Satz abbricht.} que él mismo hace\apnota{8b}{U facere idem, Uv = P.}. . . . .

El D. D. Martín niega haberlo hecho etc., salvo por desliz de lengua y de modo poco premeditado etc.\apnota{8c}{Die Zwischenbemerkung fehlt in M; U. a. R. Negavit d. Martinus id se fecisse, ut supra protestatus est.}.

Y en cuanto a lo que en segundo lugar aduce a Pablo a los Gálatas, como si aquél hubiese escrito\apnota{9}{M U scripsit, Uv = P.} contra Pedro, queriendo confirmar a los Gálatas en la fe, y mostrase en todo aquel proceso que él fue igual a Pedro y a los demás, porque no lo recibió ni por hombre ni mediante hombre etc.\apnota{12}{M acceperat; etc. fehlt in M U.}, digo que, si el reverendo padre fuese tan diligente inquisidor de la verdad como promete, habría aprendido los verdaderos fundamentos y la inteligencia de las Escrituras, de modo que no sería preciso que tan jactanciosamente dijese en medio de tan grande asamblea que él solo quiere resistir a mil\apnota{15a}{M millis.}. Así\apnota{15b}{fehlt in U.} están la verdad y el asunto. En el apostolado y en la dignidad del orden todos los apóstoles fueron iguales; esto quiso la Sagrada Escritura, esto quisieron los santos padres. A Pablo\apnota{17a}{U palam Uv. = P.} dice Jerónimo que no habría reprendido a Pedro si se supiera desigual\apnota{17b}{M esse imparem.}. Así Anacleto en el cap. In novo dijo\apnota{18a}{U dicit.} que los apóstoles fueron dotados de igual\apnota{18b}{U inaequali.} potestad. Así Cipriano, así también los demás. Pues cuando

Pablo fue\apnota{19}{U affuerit.} designado apóstol por Cristo, ingenuamente dijo que no lo recibió ni por hombre ni mediante hombre, y que Santiago y Pedro nada le comunicaron, pues todo esto ha de referirse al apostolado.

Otra cosa, en cambio, es la potestad de régimen y de administración, donde es concorde la sentencia de los santos padres y del sagrado concilio\apnota{23}{M sacrum consilium, U et concilium.} contra Juan Hus, que tal primado perteneció al mismo Pedro. No usa, pues, rectamente los testimonios de la Sagrada Escritura, sino que los arrastra a otra parte\apnota{25a}{W. A. verbessert ea; M aliorsim.} a su arbitrio, contra lo que el Espíritu Santo\apnota{25b}{M U sanctus spiritus.} exige. Pues lo que hoy ponderaba —que Pablo dijo que nada le importaba cómo hubiesen sido— no significa esto que quería el reverendo padre, sino que significaba\apnota{27}{U significat, Uv = P.} que Pedro y Santiago y los demás\apnota{28}{et ceteros fehlt in M.} fueron idiotas y poco peritos antes de andar con el Señor, y así Dios no acepta personas.

Pero dirá\apnota{29}{M hinc (statt inquiet).}, como es mordazuelo\apnota{30}{U cordaculus, W. A. audaculus.}, que esto es un invento eckiano. Lea a Ambrosio en el mismo lugar, quien sin corromper la sentencia de Pablo atestigua ingenuamente\apnota{31}{M fatetur.} que la Iglesia fue fundada sobre Pedro\apnota{32}{M super petram, U supra petram.}. Por lo cual, si no adujeré cosas más firmes, ruego que desista del intento. De donde hoy, cuando yo había dicho que es artículo husita que la obediencia eclesiástica no esté expresada en\apnota{34}{fehlt in M U.} la autoridad de la Sagrada Escritura, respondió con brevedad que el artículo no es herético. Preferiría tratar punto por punto 'con el reverendo padre' y preguntar\apnota{36a}{U dicere.}: Este es un artículo condenado, y él mismo lo sostiene, luego sostiene un artículo condenado y reprobado\apnota{36b}{M U tenet articulum. P reprobum.}.

Asimismo le pregunto\apnota{37}{ab eo fehlt in U, in Uv nachgetragen.}: si el artículo no es herético, entonces o\apnota{38}{fehlt in M U.} es temerario o sedicioso u ofensivo de los oídos piadosos; comoquiera que [306] lo conceda el mismo padre, será o temerario o 'sedicioso' o sospechoso en la herejía\apnota{1}{U de haeresi suspectus [vel fehlt]; M vel haeresis suspectus.}, u ofende los oídos piadosos\apnota{2}{Mc vel piarum aurium offensivus [Lücke in M ausgefüllt], U [Uv aut] pias offen: (W. A. offendens) aures.} contra la 'institución' teológica.

Maravilla con qué color pavonino adorna\apnota{3}{M Miror, quo [Mc colore] pavonico colore ornat.} las palabras de Cristo en Mt 16. Que Cristo interrogó\apnota{4a}{U interrogaverat.} a todos los discípulos y que así Pedro\apnota{4b}{M Christus, Mc Petrus.} respondió por todos, ¿quién lo niega? Pero de su propia harina añadió\apnota{5a}{M varina hoc addit.} que por eso a Pedro\apnota{5b}{U Petrum, W. A. ad Petrum.} 'por sí\apnota{6a}{pro se fehlt; in M.} y' por los demás apóstoles le fue dicho: Tú eres Pedro, y sobre esta piedra edificaré\apnota{6b}{fehlt in M U.} etc., lo cual, como invento luterano\apnota{7a}{P Luterianum}, no recibo\apnota{7b}{M accipio.}, pues se opone a los santos padres\apnota{7c}{M est adversum adversetur fehlt in U, in Uv nachgetragen.}.

Sobre Crisóstomo, maravilla qué injuria hace\apnota{8a}{M U facit.} al santo padre\apnota{8b}{M patrem, Uv streicht quam bis patri, a. R. cum sanctis patribus adversetur (vgl. 7), nachher gestrichen.}, como si hubiese sido depuesto legítimamente y según los decretos del concilio\apnota{9}{M consilii.} Niceno, siendo cosa averiguada por las historias\apnota{10a}{M ab historia, U ab historiis.} que después de Eusebio 'de Antioquía', cismático\apnota{10b}{U credo schismaticum.}, fueron expulsados con otros rebeldes al Romano pontífice los obispos cristianos que, buscando refugio en la primera sede, padecieron repulsa\apnota{12}{U repulsum.} de parte de los impíos cismáticos, hasta los 'tiempos de Teodosio, emperador cristiano\apnota{13a}{U christianissimi.}. Cuando\apnota{13b}{U tunc.} los cismáticos hubieron ocupado la sede treinta años\apnota{14}{U 30. annis sedem.}, san Gregorio Nacianceno recibió en Constantinopla el 'episcopado'\apnota{15a}{U accepit.}, no 'por ambición', sino para que no le sucediera otro cismático\apnota{15b}{U alii schismatici succederent.}.

Al 'santo' Crisóstomo, expulsado injustamente por medio de Epifanio, sustituido por el cismático Arsacio, lo restituyó la autoridad del Romano pontífice. Y no vale lo que entonces\apnota{18}{M hunc.} calumnia el reverendo padre contra el Romano pontífice y contra los santos que se le adhirieron, como si hubiesen procedido contra aquellos cismáticos de hecho, no de derecho, cosa que dice con sus solas palabras desnudas\apnota{20a}{M U suis verbis.} y que\apnota{20b}{fehlt in M U.} tantos santos, aun encomiados por el señor padre, no consintieron\apnota{21}{U tot sanctos . . . commendatos . . . consensisse, Uv = P.} con el Romano pontífice; de donde sin causa empieza ya el reverendo padre a acusar de ambición a aquellos santos padres que vivieron hace mil años, vicio del que durante\apnota{23}{fehlt in M.} tantos siglos han sido tenidos por inmunes por toda la Iglesia.

Así\apnota{24}{M Ideo.} es capcioso y se burla de los sencillos aquello de que el bienaventurado Gregorio 'no quiso llamarse ecuménico' y pastor universal\apnota{25}{M appellari.}, pues ayer tanto apreció a Platina, siguiendo, 'como decía', la verdad de la historia. Ya\apnota{26}{M nam.} rechaza a Platina a su arbitrio, por juzgar que en todas partes le es lícito.

Alega el decreto en su favor; si se le alega, lo rechaza, y teniendo contra sí textos abiertísimos\apnota{28}{U apertissimus.} aun de León, cap. Ita dominus 19. dist., con otros semejantes, acepta la glosa\apnota{30a}{M habet sibi contrarias glossas [Mc contrariam glossam] accipit [Mc acceptat].}, como 'los pretendientes de Penélope', por decirlo en broma\apnota{30b}{M iocam.}, que no pudiendo poseer\apnota{31}{M potire.} a la señora seguían a las esclavas. Cierto es que también los santos pontífices escribieron llamándose obispos universales, como Sixto, como Víctor, y como en el concilio\apnota{33a}{M consilio.} Calcedonense fue aclamado a León\apnota{33b}{fehlt in M.}: Viva el santísimo León, patriarca ecuménico, aunque los pontífices, siguiendo la forma de la humildad y de Cristo, se abstuvieron de aquel nombre\apnota{35a}{M U temperaverunt.}, confesándose mejor siervos de los siervos\apnota{35b}{M servus servorum 'praestare'.}, y para que no se creyese que querían suprimir la autoridad de los obispos\apnota{36}{P auctoritate.}, como tiempo ha 'respondió' a aquello santo Tomás en el libro sobre la impugnación de la fe.

Sobre la carta de Agustín os plantea a vosotros, señores auditores, un embrollo, para que leáis las palabras; las palabras oídas fueron: Inocencio dijo hace mil cien años que es antigua costumbre consultar al sumo pontífice en aquello\apnota{40}{U veterum more consulendum pontificem in iis.} que es de la fe [307]. Esto dicho contra los maniqueos quiso el reverendo padre trasplantarlo a su huerto\apnota{2a}{M ortum.}. Pero no lo hará\apnota{2b}{M faciat.}, si lee la carta de Agustín, en la cual no enumera la serie de la universidad de los obispos, sino la de los Romanos pontífices desde el mismo Pedro hasta sus tiempos.

Dice que la solución dada por mí\apnota{5}{Solutionem per me in M doppelt, das erste Mal gestrichen.} sobre los veinte años no está probada por mí; agudo es el disputador 'reverendo padre; esto aún no sabía, que el respondiente debe\apnota{6}{M probare, Mc debere.} probar su solución; siempre juzgué que el respondiente debe\apnota{7}{Mc U suam solutionem semper arbitratus sum respondentem [debere fehlt].} ofrecer la solución para que el oponente la impugne. Débil fue su Hércules', si se apoya en Marte ajeno\apnota{9}{M utitur.}.

Pero que con tanta impudencia se atreva a sugerir a vuestras ilustrísimas dominaciones y reverendas paternidades que yo no leí el primer capítulo de Juan\apnota{12}{M U caput Johannis.}, es insigne temeridad en verdad, cuando siendo niño aún no decenal, quitados los profetas, leí toda\apnota{13}{M totam bibliam.} la Biblia; pero esto nada importa al caso, cuánto haya leído cada uno.

Dije\apnota{14a}{M dixit.} defendiendo a Anacleto que no\apnota{14b}{M defendendum [non fehlt].} 'fue tan ignorante aquella interpretación'\apnota{15a}{U interpretationem eius.} de Cefas, esto es, cabeza, no ignorando que κεφαλη en griego significa cabeza\apnota{16a}{fehlt in U.} en latín\apnota{16b}{M U significat.}, sino porque en\apnota{16c}{M et apud.} Juan, a quien ayer vi, en Agustín y en la glosa ordinaria\apnota{17}{et glossam ordin. fehlt in M.} Cefas se interpreta Pedro, y Pedro, con el Cusano por autor, también significa cabeza de la casa; todo lo cual arroja a risa\apnota{18}{M tradit.} el reverendo padre, sobre lo cual juzguen los jueces que han de ordenar.

Pero aquello es más excelente\apnota{20}{M quaestantius.}, que la autoridad «Apacienta mis ovejas» la envuelve en tantos rodeos\apnota{21a}{M imbalgibus, Mc ambalgibus (so!)}, solo para hacer que los hombres viendo no vean; paso por alto\apnota{21b}{M mitto, Mc omitto.} que aportó una exposición tan ridícula inventada por él. Pero vendré\apnota{22}{M venio.} a otro sentido suyo, donde dijo muchas cosas; mas yo, no atendiendo a la singularidad, que es madre de los errores, sino siguiendo la sentencia de los santos padres y de la santa madre Iglesia recibí la palabra «Apacienta», porque\apnota{25}{M pascere quia.} allí lo constituyó pastor y le dio el primado de la Iglesia universal. Así lo entendieron Gregorio, así Crisóstomo, así Ambrosio sobre el último de Lc, ayer por mí no alegado\apnota{27a}{M heri non allegatis.}. A los cuales pasando bellamente 'de largo'\apnota{27b}{P praegrediens, M Lücke, Mc progrediens.} nada dijo el reverendo 'padre, sino que aportó la verdadera\apnota{28a}{Mc varam attulit, U [veram fehlt] attuli, Uv attulit.} sentencia de Agustín, de cómo debe ser el prelado o\apnota{28b}{praelatus vel fehlt in U, in Uv nachgetragen.} pastor, la cual también yo admito, y creo que cualquier fiel\apnota{29a}{U fidelium.} ruega que también\apnota{29b}{M et et.} el sumo pontífice, 'cualquier príncipe', cualquier prelado buscase el honor de Dios y la salud de los súbditos, pues este\apnota{31}{U hoc, Uv = P, aber gestrichen.} es el verdadero apacentamiento, por el cual, así como aquí tienen el principado en la tierra, también son elevados para que reinen\apnota{32}{M regnant.} en los cielos. Por lo cual Agustín no dice cosas adversas ni pugnan­tes, sino tales que por ellas debe entender el prelado cuál sea su institución.

Y no vale lo que el reverendo padre adujo de que le fueron encomendadas las ovejas bajo condición, no habiendo allí condición alguna, salvo que su paternidad tenga\apnota{36a}{M habet.} otro evangelio que yo. Mas significando cómo debe ser\apnota{36b}{M U debeat esse.} el pastor, Cristo busca en Pedro la dilección y el amor de Dios, y así lo hizo pastor de toda la Iglesia, como canta la santa madre\apnota{38}{M mater sancta.} Iglesia, y el reverendo [308] padre ora\apnota{1a}{U orat vel legit.}: 'Tú' cabeza de la Iglesia, pastor de las ovejas\apnota{1b}{M U ovium etc.}. Por lo cual tenemos por derecho divino en este lugar que el prelado debe amar a Dios, tenemos a Pedro constituido pastor de las ovejas, como lo entendieron san Gregorio, san Ambrosio y Crisóstomo, aun el griego\apnota{4a}{M Graecus, Mc Graecos.}.

Nada importa la tenue\apnota{4b}{M praetenuit. Mc quae tenuit.} racioncilla aducida en contrario por el reverendo padre\apnota{5a}{M inductum, Mc inducta.}, que Pedro no apacentó a todos; lo que\apnota{5b}{M quo, Mc quod.} se hace en nuestro nombre se dice que lo hicimos nosotros, como atestigua Clemente\apnota{6}{U Clemens 21. ult. dist.} que Pedro ordenó obispos, presbíteros\apnota{7}{U presbyteros, episcopos.} y diáconos que sembrasen el Evangelio por todo el orbe de la tierra, como se lee en el can. In illis, dist. 80 o 81\apnota{8}{M can. in illis 81. dist. vel 80. dist., a. R. in c. in illis 81. dist, U in can. illis 80. vel 81. dist; richtig: 80. dist.}. De donde el enviado por el principal, como fue san Pablo, no está obligado a ser enviado por el vicario del mismo; sin embargo Pablo\apnota{10a}{M U Petrus, Mc Paulus.} subió con Tito, a Gal 2, y confirió con Pedro el Evangelio\apnota{10b}{M evangelia.}, quienes le dieron la diestra en señal de compañía.

Además, en cuanto a que forja contra mí una glosa jamás por mí pensada, como si yo hubiese dicho: Edificar sobre\apnota{13}{M super.} la piedra es estar sometido al Romano pontífice, ¡oh sueños, oh portentos de palabras!; jamás fue esto de mi mente. Pues Cristo edifica la Iglesia, no los que están sometidos, ya que los súbditos deben estar sometidos\apnota{15}{U esse, Uv subesse, vor aedificatae.} a la Iglesia edificada.

Sobre Gregorio de Rímini, cuya sentencia me desplace y a él le place, nada ya al propósito. No quiero, pues, descender a estos desvíos. Sino

como hemos recibido\apnota{18a}{P occepimus, M suscepimus.}, defendamos que el artículo del concilio\apnota{18b}{U consilii.} de Constanza es verdadero, y que quienquiera que lo contradiga yerra y contradice a la verdad y al honor del concilio\apnota{19}{U consilii.}.

El D. D. Martín.

Antes de venir a las autoridades por mí opuestas, recorreré brevemente lo que el egregio D. D.\apnota{22a}{M dom. egr. doctor.} ha mezclado sobre mi respuesta de ayer y de hoy\apnota{22b}{M meus, Mc mea.}, donde entre otras cosas, como desesperado tras consumir todas las armas\apnota{23}{M armis omnibus.}, coloca en un solo artículo del concilio de Constanza el postrer presidio de su confianza, y odiosísimamente lo repite e inculca, como si arriba no se hubiese respondido a él bastante sobradamente. Diré, sin embargo, una vez más.

Cuando el artículo sobre las obras de los hombres se cuenta entre los condenados y es afirmado por\apnota{28}{fehlt in M.} san Pablo y Agustín, defendido después por Gregorio de Rímini y por todas las universidades hasta hoy, no me moverán las odiosísimas inculcaciones de este artículo, hasta que el egregio D. D. pruebe que el concilio no puede errar, que no erró o aun que no yerra, pues el concilio no puede hacer derecho divino de lo que por su naturaleza no es derecho divino; por eso tampoco es herético sino lo que es contra el derecho divino. Esto por este artículo\apnota{33}{haec pro isto articulo fehlt in M.}.

En cuanto a que zahiere que Crisóstomo fue legítimamente\apnota{34}{M a me legitime.} depuesto, no me entendió, pues claramente dije hoy que la potestad de ordenar y deponer al obispo fue legítima en Epifanio de Chipre\apnota{36a}{P M U in beato Cypriano.} y en Teófilo, aunque nada disputé\apnota{37}{P potestare; M disputavi, U disputaverit.} sobre si entonces usaron bien o mal\apnota{36b}{M an tunc bene vel male usi sint, U an bene tum usi sint an male.} esta potestad legítima. Me basta que, sin [309] la autoridad del Romano pontífice, tantos santos y celebérrimos varones de Grecia\apnota{1}{M Graeci.}, catolicísimos, no cismáticos, tuvieron potestad de ordenar y deponer obispos.

También pondera egregiamente que yo habría acusado de ambición a los santos antiguos, y no atiende a que dije que también en los apóstoles estuvo esta fragilidad en presencia de Cristo\apnota{6}{U Christiano, Uv = P.}.

En cuanto a que ahora\apnota{7}{M U ita nunc.} engrandece tanto a Platina contra san Gregorio, por mi autoridad consta bastante claro a cualquiera aun medianamente ingenioso qué concluye. Pues yo de Platina tomo las historias y nada más. Y que el mismo Platina aduce ciertas razones, lo rechazo por la autoridad contraria de Gregorio.

Ya paso por alto aquello de que Sixto y Víctor\apnota{12}{W. A. et Victor.} escribieron llamándose pontífices universales y que a León fue aclamado en el concilio Calcedonense, sobre lo cual

hoy se ha dicho bastante. Hay que observar aquello de que, según dice el egregio D. D., los pontífices siguientes, siguiendo en Cristo la forma de la humildad\apnota{15}{M Christi formam secutus (so!)}, se abstuvieron de este nombre; lo cual no puedo entender de otro modo sino que por causa de humildad detrajeron al derecho\apnota{16}{U iure.} divino, y se hicieron mendaces y desobedientes, si es derecho divino, esto es precepto, que ellos fuesen pontífices universales.

Repitiendo la carta responsiva\apnota{19a}{U responsivam epistolam.} de Inocencio a Agustín dice\apnota{19b}{M dicit dicit.} que fue antigua costumbre consultar al Romano pontífice\apnota{20}{M ro: pontificis, U Romanos pontifices.} en lo que es de la fe. Óptimamente dice, y a mí\apnota{21}{fehlt in M U.} me place, sobre todo en la Iglesia latina; pero ¿qué tiene que ver con la fe y el primado? Esto es, ¿a qué\apnota{22}{U quid ad.} propósito?

También se quejó de que yo exigí de él la probación\apnota{23}{fehlt in M} de la solución sobre los veinte años\apnota{24a}{M 20 annorum solutionem.}; doblemente falta el egregio D. D.: primero\apnota{24b}{U primo}, porque él mismo fue oponente durante tres días y medio y solo a esta hora empezó a ser respondiente. Después, no exigí que probase la solución, sino el primado trasladado con Pedro.

Excusándose también\apnota{27a}{U se etiam.} sobre Cefas —que Pedro se interpreta en\apnota{27b}{M ad.} Juan—, aún persiste en que Pedro también significa\apnota{28}{M significat etiam.} cabeza de la casa, como si, a saber, Juan, para interpretar el Evangelio a los griegos, hubiese hablado\apnota{29}{sit in M gestrichen, Mc locutus locutus', das zweite gestrichen.} no en griego sino en caldeo, para no ser entendido acaso.

La palabra «apacienta» la toma de modo que quiere que signifique\apnota{32}{M U significare.} la constitución del pastor universal, y para ello adujo\apnota{33a}{fehlt in P M U, aus der ed. Viteb. zu ergänzen, vgl. W. A.} a Gregorio, a Ambrosio y a Crisóstomo, ayer no alegados, sobre los cuales nada diría yo. Diré, pues, ahora, si es que fui\apnota{33b}{M fuero.} olvidadizo, y a la vez opongo la regla de Agustín\apnota{34}{M U divi Augustini.}: Que los dichos de todos los escritores han de juzgarse por la divina Escritura, cuya autoridad es mayor que la capacidad de todo el género humano; no porque condene la sentencia de los ilustrísimos padres\apnota{36}{fehlt in P.}, sino porque\apnota{37a}{M qui.} imito a los que más de cerca se acercan a la Escritura\apnota{37b}{M ad scripturas.}, y ante todo, si la Escritura fuese clara, abrazo la Escritura misma. Siendo, pues, claro que apacentar, por fuerza del vocablo y por uso de la Escritura, no significa sino enseñar y cuidar [310] las ovejas, como Pedro lo recibe al final: Apacentad la grey que\apnota{1}{U Sic Petrus accipit ult. Pascite [eum fehlt].} hay en vosotros [1 Pe 5 v. 2]. No conviene equivocar el vocablo y dejar sin necesidad el sentido simple.

También disuelve lo que dije que hay condición en la palabra de Cristo, Si me amas, cuando yo juzgo\apnota{4}{M U arbitrer.} que no solo hay condición, sino también pacto

entre Cristo y Pedro, cuando aquél dice: ¿Me amas? y éste responde: Te amo.

Omito la dialéctica, en que ejercitadísimo es el señor D. D., y la conjunción\apnota{6}{M conditionem.} condicional «si». Querría oír qué diría a la decretal de elect. c. Significasti, donde el pontífice Pascual dice manifiestamente que aquí hay condición, si al menos tanto vale\apnota{9}{M valet, verbessert valeat.} la autoridad del Romano pontífice como hasta aquí ha contendido; dice en efecto\apnota{10a}{U enim ad sensum}. [Va a citar del libro esta decretal.]\apnota{10b}{Die eingeklammerten Worte in P a. R., M statt dessen Si conscientiarum cognitor et factor [aus d. Text d. Decr.].}

A mi razón, que\apnota{11a}{Ad fehlt in M, U at; quam fehlt in M U.} así la confutó —que Pedro no apacentó\apnota{11b}{M U paverit.} a todos, porque, dice, lo que se hace en nuestro nombre se dice que lo hicimos nosotros—. Aunque no creo que con esto haya querido el señor D. satisfacer a mi razón, pues no puede decir que Pablo predicó en nombre de Pedro, quien\apnota{14}{M si, Mc qui.} en cambio a Rom 1 [v. 1ss.] se dice apóstol en nombre de Cristo; lo cual colijo de que dice que Pedro, con Clemente por autor, ordenó obispos, presbíteros\apnota{16}{M presbyteros episcopos.} y diáconos que sembrasen el Evangelio por todo el orbe, entre los cuales no cuenta\apnota{17}{M enumerat.} a Pablo.

Pero esto parece hacer al caso, que el enviado por el principal no está obligado a ser enviado por el vicario del mismo; esto admito, si primero fuese probado que Pedro es vicario. Ahora, en cambio, no menos es Pablo vicario de Cristo\apnota{20}{fehlt in U.} que Pedro, como ya se ha dicho, porque es apóstol en nombre de Cristo, y así tantos vicarios cuantos apóstoles.

A\apnota{22}{M U Nunc ad.} propósito.

A la autoridad de Pablo 1 Cor 3 [v. 5] dice primero que Cefas no es del texto desde el principio, sino que se recuerda al final, donde la glosa confiesa que Pedro es mayor que todos; cosa levísima es esto, pase\apnota{25}{U transeo.}, pues al principio del primer capítulo [v. 12] ciertamente se pone\apnota{26}{P ponitur.} Cefas.

Segundo, nada dice sino que ni palabra ni ápice\apnota{27}{M add de sit, gestrichen.} hay en el texto sobre la negación del primado, y así quiere que esta autoridad quede disuelta. Añadió, sin embargo, la autoridad de Jerónimo, Agustín, Ambrosio, Crisóstomo, Bernardo y León\apnota{30}{Ambrosii und Bernhardi fehlt in M; M U Leonis etc.}, que dicen que Pedro es la piedra. Contra lo cual 'digo y' replico que este segundo dicho nada es al propósito y arriba está bastante disputado si Pedro es la piedra. Igualmente tampoco disuelve aquello de que dice que no hay palabra ni ápice sobre la negación del primado\apnota{33}{de negatione primatus fehlt in M, Mc füllt die Lücke durch 'primatum'.}; pues si aun Jerónimo establece\apnota{34a}{U constituat.} la cabeza para que se quite la ocasión del cisma,

bastante robustamente 'lucha en contra'\apnota{34b}{contra pugnat fehlt in M, Mc 'oppugnat'} el texto abierto, que quita la cabeza para que se quite la ocasión del cisma. Dice en efecto: ¿Qué es Pablo? ¿Qué Apolo?, queriendo que solo Cristo sea\apnota{36}{M est.} la cabeza, con lo cual verdaderamente se quita todo cisma. Y así\apnota{37a}{Et sic fehlt in U.} no se satisface con esta respuesta\apnota{37b}{U sic ad responsionem.}.

[311] A la segunda autoridad, Gal 2, dice que ciertamente todos los apóstoles fueron\apnota{1}{secundam und fuisse fehlt in U; letzteres in Uv nachgetragen, M omnes apostolos [U apostolus].} iguales, como atestiguan Jerónimo y el can. In novo\apnota{2}{M ut Hieronymus . . . attestantur.}. Pero disolver tan gran autoridad con estas tres palabras —que otra cosa es la potestad de régimen y de administración— me es bastante admirable. Pero tampoco\apnota{4}{M sed non, U sed et nec.} entiendo qué quiere, cómo distingue el régimen\apnota{5}{P rem.} y la administración.

Y en cuanto a lo que dijo, que Pablo 'con Tito' subió a Pedro para conferir\apnota{6a}{U ascendisse cum Tito . . . ut conferat.} el Evangelio con ellos, claro es por el texto precedente\apnota{6b}{fehlt in U.} y siguiente\apnota{6c}{M excedenti, Mc 'praecedenti'.}, como el mismo Jerónimo expone\apnota{7}{ut et H. i. exponit fehlt in M. [Lücke].}, que Pablo no hizo esto para consultar a Pedro sobre el Evangelio\apnota{8a}{U evangelium.}, que ya por más de 14 años\apnota{8b}{M annos} había predicado y del que se había jactado de no\apnota{9a}{fehlt in U.} haberlo recibido de hombre, sino para cerrar la boca a los pseudoapóstoles\apnota{9b}{pseudapostolorum.}, para confirmar a los Gálatas que eran lo mismo lo que él y los demás apóstoles predicasen.

El D. D. Eck.

En cuanto a que el reverendo padre mezcla el artículo de Juan Hus sobre la obediencia eclesiástica con otro artículo sobre las obras de los actos humanos\apnota{13}{M bonorum.}, por autoridad de Gregorio de Rímini, que sería defendido por todas las universidades, al cual hasta ahora en ese lugar en ninguna universidad, aunque en muchísimas estuve, oí defender\apnota{15}{M fui . . . defensam.}; pide, sin embargo, de mí que pruebe que el concilio\apnota{16}{M U ei concilium (M consilium).} no puede errar; no sé qué quiere\apnota{17a}{M velit sibi.} esta petición, si tácitamente quiere tener por sospechoso al laudable y glorioso\apnota{17b}{M an tacite 'an' laudabile et nobile, a. R. 'an nobile', U an tacite laudabile et [gloriosum fehlt].} concilio de Constanza. Esto os digo, reverendo padre: si creéis que el concilio legítimamente congregado erró o yerra\apnota{19a}{M esse [gestrichen] errasse et errare, U errare et errasse.}, sois para mí como étnico y publicano\apnota{19b}{M Lücke, Mc 'haereticus et ethnicus', a. R. 'et publicanus'.}. Qué sea hereje, por ahora no lo discutamos\apnota{20}{M U discutiam.}.

Sobre Crisóstomo, no quiere haber justificado\apnota{21}{M U se iustificasse.} su expulsión, sino haber alabado la ordenación de obispos instituida por el concilio Niceno\apnota{22}{M Ciceno, Mc 'Niceno'.}. Me place por reverencia a san Crisóstomo, y siempre los Romanos pontífices en diversas

permitieron a las naciones diversos modos de elegir a los obispos\apnota{24}{M episcopos [Mc episcoporum] eligendos tamen, U episcopos eligendi tamen.}, consta además con suficiente claridad, a partir de Gregorio Nacianceno y de Crisóstomo restituido, que intervino la autoridad del pontífice romano, como atestiguan san Atanasio de Alejandría, san Pablo y otros obispos de la Iglesia oriental, que recurrieron al pontífice romano para su restitución\apnota{28}{M se contulerunt (so!)}. Que también en los apóstoles hubo algo de humana fragilidad\apnota{29}{M signat.}, lo sabemos. Pero mientras desempeñaban públicamente su oficio, sin indicio\apnota{30b}{U iudicio. M officio fungerentur absque initio.} ni causa no ha de detraerse a tan santos varones\apnota{30a}{U cum.}.

Acerca de Platina dice que acepta la autoridad de su historia\apnota{32}{eius historiae auctoritatem fehlt in M, Mc eius historiam.}, y yo tampoco aduje razones tomadas de Platina\apnota{33}{U auctoritatem ex Platina adduxi.}, sino la historia.

A la carta de Inocencio responde que nada atañe al propósito, porque prueba acerca de la Iglesia latina, no acerca de otra\apnota{35}{Mc probat . . . . aliis, U de latina illum probat ecclesia et non de alia.}. Pero como esta es una glosa luterana sin apoyo en autoridad ni persuasión alguna, la desprecio con la misma facilidad\apnota{36}{M facilitatem, Mc facilitate.} con que se aduce.

A lo de los veinte años respondió que yo había sido oponente ya por tres días\apnota{38}{M per tribuum [Mc triduum], iam fehlt.}; y ciertamente debería haberlo sido\apnota{39}{M debuisse.}, si el reverendo padre con sus cautelas no hubiese convertido alguna vez de opo[312]nente en respondiente\apnota{1a}{M ex respondente opponentem, gestrichen, Mc a. R. wie oben.}. Pues este argumento no lo puse yo\apnota{1b}{M posui hoc, U hoc opposui.}, sino que lo puso él mismo\apnota{1c}{fehlt in M.}, y ciertamente aquella cátedra Pedro no se la llevó consigo a Roma\apnota{2}{secum und ad Romam fehlt in M, U sumpsit, Uv = P.}, sino que trasladó el primado.

Sobre la palabra «apacienta»\apnota{4a}{U = P.} quiere dejar de lado a Gregorio, a Ambrosio, a Crisóstomo, porque\apnota{4b}{M remittere [Mc dimittere] Hieronymum, [a. R. Ambrosium] quia.} una palabra simple no debe tomarse en sentido equívoco. Digo\apnota{5}{M debet aequivocari 'dico' dico.}: ¿quién ignora que la palabra pastor es nombre de oficio, de honor y también de carga? Por eso recibió a la vez el oficio de la prelacía y juntamente la carga de apacentar\apnota{7}{M oves.}.

La condición aducida por el reverendo padre no la acepto, porque\apnota{8}{M accipio allegata quia.} tampoco Cristo, como él mismo siempre\apnota{9a}{fehlt in M.} se adhiere a la letra, unió condición alguna\apnota{9b}{U iunxit.} al apacentar. Pues\apnota{9c}{M non.} así se confirmaría el error husita de que quien no está en caridad no sería\apnota{10}{fehlt in U.} obispo, prelado, pastor\apnota{11}{M praelatus etc., U praelatus pastor etc.}.

Sobre el capítulo Significasti, de elect., sé que Gerson y otros doctores discuten en la materia si el concilio\apnota{13}{Mc consilium.} está por encima del papa. Pero como no se han aducido las palabras de la decretal, nada puedo decir, sino que conjeturo

«sin querer ligarme a esto»\apnota{14}{M Lücke, Mc a. R. volens (so!) me huc (U hoc) allegare.}, que aceptó la condición pascual por cualidad\apnota{15}{U pastalem conditionem accepisse pro qualitatem (Uv qualitate).}.

Sobre Pablo, que no reconoció el primado de Pedro, sino que fue apóstol de Cristo: todos somos de Cristo, creo; sin embargo, el primado\apnota{17}{M U ut arbitror tamen Petri (Mc U Petro) 'primatum'.} de régimen a Pedro en ninguna parte se lo quitó, sino que incluso\apnota{18}{fehlt in M.} al morir, como se escribe en la carta sobre el tránsito de los apóstoles, se despidió de Pedro como de cabeza de la Iglesia; carta de la que no tengo certeza si es del bienaventurado Dionisio o no.

Dijo que Pablo fue vicario de Cristo, pero no general, pues también\apnota{21}{nam etiam fehlt in M, Mc 'ita etiam'.} en el Apocalipsis [21, v. 14] se ponen doce fundamentos de la Iglesia, y sin embargo otro fundamento\apnota{22}{M alium, Mc aliud.}, a saber, el principal, nadie puede ponerlo\apnota{23}{M aliud ponere.} sino el que ha sido puesto, es decir, Cristo.

A lo segundo dijo que no viene al propósito si Pedro es la piedra, y ese es todo nuestro asunto. Sin embargo, silenció\apnota{27}{M U ipse reticuit.} la solución dada por mí anteayer\apnota{26}{Mc a me.}, a la que me remití, aunque está expreso en la letra que Pablo reprueba las contiendas\apnota{28}{Mc increpare contentionem, Ui mprobare contentionem.}. Pero no sé con qué ingenio el reverendo padre juzga\apnota{29}{M U tolleret.} que Pablo negó por ello el primado, para que se quitara el cisma. Quizá no estuvo en una provincia que careció de príncipe\apnota{30a}{M fuit in principe [Lücke], Mc fuit in provincia quae caruit principatu.}, pues habría experimentado lo contrario: cuántas\apnota{30b}{P quod.} disensiones, contiendas y discordias surgen donde no hay príncipe ni\apnota{31}{Mc vel.} cabeza que juzgue y conozca entre las causas\apnota{32}{U agnoscens.}.

Y le parece admirable que con tan pocas palabras, esto es, apostolado y régimen, quiera yo resolver tan grande\apnota{33}{U tantum, Uv tantam.} autoridad. Los principios son mínimos en cantidad, máximos en virtud. Si el reverendo padre hubiese atendido a aquellas dos menudas palabras, apostolado y régimen, de la cabeza principal\apnota{35}{M Primarii, Mc principalis.} y de la secundaria, hace tiempo habría reconciliado las sentencias de la Sagrada Escritura y de los santos padres que le parecen pugnar entre sí; descuidadas estas, es necesario seguir dogmas exóticos\apnota{37}{M est ex hoc scil.} y ajenos.

Mañana a las siete de la mañana se continuará la disputa por edicto público.

[313] El viernes ocho del mes de julio a las siete de la mañana\apnota{1}{Die Zeitbestimmung von Crastina an fehlt in M, U Die Julii octva (Uv octava).}.

El D. D. Martín.

A las autoridades aducidas por mí, de Pablo a los Corintios y a los Gálatas, el egregio D. D. nada responde a su modo, sino que repite lo del concilio de Constanza y los husitas, con lo cual no se me resuelven\apnota{5}{U solvantur, Uv solvuntur.} las autoridades. Pues, como dije, el concilio de Nicea estableció otras cosas mucho más sagradas y célebres. Por eso, cuanto él se apoya en el suyo, tanto yo en el mío, y digo, como dije ayer: por mucho concilio que se alegue, aún no se tiene el derecho divino, del cual solo se trata; luego el concilio nada viene al propósito.

En segundo lugar, para eludir las autoridades más que\apnota{10}{M alluderet potius auctoritates.} para resolverlas, trajo la distinción entre apostolado y administración, la cual, puesto que ha sido hallada por arbitrio humano y no se encuentra ni en las sagradas letras ni en los padres eclesiásticos\apnota{12}{M neque ecclesiasticis, letzteres gestrichen; M U in sanctis patribus.} que las interpretan, no me doy por satisfecho con que se me haya devuelto\apnota{13}{M ea . . . . reddita.}, sobre todo cuando apostolado no es, como él piensa con los suyos, nombre de dignidad, sino de oficio, como\apnota{14}{M quod.} es claro por el apóstol a los Romanos [1, v. 5] y en todas partes, donde dice\apnota{15}{M dicit se.} que recibió la gracia y el apostolado para la obediencia de la fe. Pues apóstol es mensajero de la palabra, y apostolado o ἀποστολή en griego\apnota{17a}{M apostoli graece officium vergit, U apostole, (a. R. apostolē) officium verbi, Uv add graece.} es oficio de la palabra, y así en las sagradas letras apostolado y\apnota{17b}{U sive.} administración serán casi lo mismo\apnota{18}{M vere idem erit.}; a no ser que entienda por administración las contiendas forenses sobre bienes temporales eclesiásticos, la cual no es propiamente administración apostólica. Confieso, ciertamente, que si se permite esta libertad de fingir, fingiré\apnota{20}{U fingens, Uv fingam, Lö fingerem.} también yo que san Juan es canciller\apnota{21}{W. A. Anm. „nach S. 316 Z. 9 f. hat Luther beispielsweise auch noch Marcum sigilliferum hinzugefügt."} y que otro apóstol es otra cosa. Pero todo esto no es de derecho divino. Se desvanece, pues, aquel principio que jactó de mínimo en cantidad\apnota{23}{U iactitatur, W. A. fügt nach S. 312 Z. 34 zu in virtute maximum.}. Y bastante me admiro de que el egregio señor rehúya así las sagradas letras y huya de su rostro\apnota{24}{U eorum.}, desviándose a estas galerías inventadas por ingenio humano.

Por tanto, sigue aún en pie la autoridad a los Corintios, según la cual, para quitar el cisma, Pablo quita de en medio la persona de Pedro, de Pablo\apnota{26}{M Pauli Petri.}, de Apolo, etc., y deja solo a Cristo, quedándose para sí únicamente el ministerio de la palabra, que es verdaderamente el apostolado y la administración de la Iglesia. Y no vale lo que con aquella fría\apnota{29}{M fridiga (so!) ista, U frigida ista.} semejanza quiso disolver esta robustísima autoridad,

a saber, que una provincia sin cabeza queda expuesta a cismas y a muchos peligros, pues no faltará Iglesia sin cabeza cuando se den muchos obispos sin papa.

Como vemos que sucede en las federaciones, donde tanto más hermosamente se administran las repúblicas cuanto menos tienen una cabeza común, pues es evidente que el reino de Francia, el de España, el de Hungría y otros infinitos\apnota{33}{U Franciae, Angliae, Britaniae, Galliae et infinitorum.} en el mundo existen sin cabeza común y sin embargo no chocan entre sí\apnota{34}{M colliguntur, Mc colliduntur.}, contentándose cada reino con su rey\apnota{35a}{M regi.}. Y si en estas cosas temporales\apnota{35b}{M his temporalibus rebus, U iis temporalibus.} y fluctuantes, no\apnota{36}{M neque; et moribus fehlt in P.} teniendo los hombres, tan diversos en ingenios y costumbres, cabeza común ni en el cielo ni en la tierra, se mantienen la paz y la concordia de las tierras, cuánto\apnota{37}{M tanto.} [314] más podrá suceder esto en las cosas espirituales, constantísimas, esto es, en hombres dotados de una misma fe, de igual caridad y de las mismas costumbres, y que además tienen en el cielo una cabeza común, como dice el apóstol: Una fe, un Señor, a los Efesios [4, v. 5]\apnota{4a}{ad Ephes. fehlt in M.}. Por lo cual, aunque hubiese diez papas\apnota{4b}{P esset.}, aunque hubiese mil papas, no por ello se producirían cismas.

Por eso pido al D. D. que quiera resolver mejor aún\apnota{6}{M U adhuc d. d. velit.} la autoridad de Pablo, y asimismo aquello\apnota{7a}{M similiter illud, U similiter et illam.} a los Gálatas [2, v. 6]: Dios no hace acepción de personas\apnota{7b}{etc. statt non accipit.}, donde Pablo rechaza todo lo que en Pedro y en los apóstoles puede haber de personal o cualificado, como dicen. Aún no ha sido tocada\apnota{9}{M nondum est facta, U utraque nondum est tacta.}, menos aún resuelta.

A lo cual, para terminar la disputa, añado también aquello de que Pablo, en 1 Corintios 12 [v. 28], al describir el cuerpo eclesiástico, enumera en primer lugar a los apóstoles, en segundo a los profetas, en tercero a los doctores y casi en último lugar a los gobiernos\apnota{12}{M iubernationes.}, que el D. D. llama administraciones; luego, si el primado de Pedro fuese cosa tan grande y tan\apnota{13}{fehlt in M.} necesaria por\apnota{14}{fehlt in U.} derecho divino, Pablo habría descrito insuficientemente la Iglesia, al omitir lo principalísimo que hay en ella.

Además, sigue aún en pie esto: que ni Pedro ni todos los apóstoles juntos pudieron ordenar\apnota{17}{M ordinari.} a Matías, Hch 1, ni darle la administración. Por eso el pontífice romano se arroga más administración para conferir por todo el mundo\apnota{19}{M U orbem.} de la que no se encuentra en su predecesor el apóstol.

Igualmente pienso acerca de la ordenación de Pablo y Bernabé, Hch 13, a los cuales\apnota{20a}{fehlt in U, in Uv nachgetragen.} el Espíritu Santo, y no Pedro o\apnota{20b}{U autem, Uv aut.} la Iglesia, segregó. Por eso se gloría no sin razón de no haber recibido nada de los grandes apóstoles; más aún, escribiendo a los Corintios dice que en nada hizo menos que los demás apóstoles [2 Cor 12, v. 11].

Pero también aquella autoridad estrecha aún más a Pedro, en la que escribe a los Gálatas: El que obró en Pedro para el apostolado de la circuncisión, obró

también en mí entre los gentiles, y más abajo: Me dieron a mí y a Bernabé las diestras en señal de comunión, para que nosotros\apnota{27a}{M ipsi, gestrichen.} fuésemos a los gentiles y ellos a la circuncisión, con tal de que nos acordásemos de los pobres\apnota{27b}{M pauperes.}\apnota{28}{M erimus, U simus.} [2, vv. 8 ss.]. Si quisiera apoyarme en esta autoridad, ni el D. D. ni hombre alguno podría obligarme a que el apostolado de Pedro sea más amplio que a los judíos por\apnota{29}{M de iure.} derecho divino, y por ello todas las autoridades en que se apoya —Apacienta mis ovejas\apnota{30}{U meas etc. et.}, Tú eres Pedro, etc.— han de referirse\apnota{31}{M reverendas.} a los términos de su apostolado, como Cristo claramente en Jn 10 [v. 16] separa las ovejas de los gentiles\apnota{32}{P feceruit.} de las ovejas de Israel, diciendo: Tengo aún otras ovejas que no son de este redil.

Y al fin, como también ayer tocó el D. D., en Ap 21 [v. 14] se describe la nueva Jerusalén con doce fundamentos, además de aquel fundamento capital que es Cristo\apnota{36a}{U Christus, Uv Christum.}, y sin embargo allí no se pone ninguna diferencia\apnota{36b}{U ibi nihil differentiae.} entre los doce fundamentos, lo cual habría que hacer del todo si hubiera de pensarse algo tan grande acerca del primado de Pedro. A lo mismo valen muchísimas\apnota{38}{M plurimum.} figuras en las Escrituras, como los doce bueyes del mar de bronce, 3 Re 7\apnota{39a}{P 3. Reg. 4. M 1. (verbessert 3) regum 4 (a. R. 2 regum 4) duodecim boves mari sereni (so!), U duodecim boves 3. regum 4. maris enei; W. A. verbessert 3. reg. 7.}, asimismo los doce leoncillos\apnota{39b}{P M lenunculi.} del trono de Salomón, las doce piedras del Jordán y otras semejantes, en todas las cuales se describe una total igualdad de todos los [315] apóstoles, y ninguna desigualdad puede probarse en contrario por derecho divino\apnota{2}{M probari de iure divino; U in contrarium probari iure divino.}.

Sea esto dicho en favor de esta conclusión, que hasta hoy creo verísima\apnota{3}{M firmissimam.}, más aún, sé con certeza\apnota{4a}{U certe.} que no podrá ser nunca demolida\apnota{4b}{M efelli.}. Por eso ruego al egregio D. D.\apnota{5a}{U d. d. egre:; W. A. d. d. egregius.}, puesto que aún no ha opuesto ninguna autoridad\apnota{5b}{M adhuc nullam auctoritatem opposuit, U nullam adhuc posuit auctoritatem.} de derecho divino en favor de su sentencia\apnota{6}{M sententia sua.}, que al menos disuelva algunas clara y sin rodeos; de otro modo no lo tendré por teólogo que siente rectamente acerca de las Escrituras, y lo dejo, o más bien lo remito, al juicio no de la envidia ni del vulgo, sino de los jueces que han de designarse y de cualquier lector prudente y de todo oyente sincero.

El D. D. Eck.

No sé qué quiere el reverendo padre, que al inicio de su discurso dice que nada he respondido, y sin embargo en el curso del mismo se esfuerza por quebrantar las soluciones dadas por mí ayer, invictas y robustísimas; si esto es coherente, él mismo lo verá.

Después\apnota{14}{M U Deinde.}, en cuanto a que no quiere verse constreñido por la autoridad del concilio de Constanza, él verá, vean los jueces cuán justamente obre. El sínodo niceno

no lo rechazo en lo que\apnota{16a}{U iis.} es de fe, pero lo que\apnota{16b}{U atque (statt at quae).} mira a las costumbres, como las ordenaciones de los obispos, puede cambiarse según la calidad de los tiempos y lugares\apnota{17}{M posse.}.

Me pide que le aporte el derecho divino, lo cual, aunque ya lo he hecho muchas veces, sin embargo mi conclusión no lo contiene, sino solo esto: que el reverendo padre dijo en su resolutorio que la Iglesia romana antes de los tiempos de Silvestre no fue superior a las demás, lo cual niego en mi conclusión; con todo confieso que el primado del sumo\apnota{21}{U Romani.} pontífice es de derecho divino según lo alegado por mí.

Después\apnota{23}{U Deinde.}, para reprobar mi solución dice que mi distinción ha sido inventada por arbitrio humano, como si en teología no fuese lícito distinguir, igual que Arrio quiso burlarse de Atanasio, pues, demasiado aferrado a la letra\apnota{25}{M litera.}, dijo que el Hijo es criatura por aquello: El Padre es mayor que yo [Jn 14, v. 28]. Al disolverlo\apnota{26}{W. A. verbessert diluit.} Atanasio distinguiendo: El Padre es mayor que yo según la humanidad, pero no según la divinidad; como la Sagrada Escritura\apnota{28}{U scripturas sacras.} no había distinguido, Arrio no la aceptó y así permaneció pertinazmente en el error. Así en el presente caso, cuando los que esclarecen la Sagrada Escritura tendrán vida eterna, hay que dar distinciones cotejando los diversos pasajes\apnota{30}{M U passus scripturae.}, para que no se admita contrariedad en las sagradas letras, no sea que, por aferrarse demasiado a la letra, uno sea muerto por ella, sino más bien vivificado por el espíritu. Y como Pedro es antepuesto a los demás\apnota{32}{U praeferter, Uv praefertur.} muchas veces en las sagradas Escrituras y muchas veces se le encuentra igual a ellos, para evitar la contradicción es necesario poner una distinción conforme a la Sagrada Escritura.

Así también lo insinúan Jerónimo y Cipriano por mí\apnota{35a}{M a me.} alegados\apnota{35b}{U allegati et epistola 9, W. A. allegati lib. 4 epistola 9.}, quienes confiesan la igualdad del apostolado y sin embargo tributan a Pedro el primado, lo cual no puede hacerse si no se refieren a cosas diversas. [316] Me place decir que el apostolado es oficio, pero de modo que también diga honor. Pues estos nombres son de honor y de carga, como dije, según indica el apóstol a los Romanos, cap. 10 [v. 15]: ¿Cómo evangelizarán si no son enviados?

Sobre la administración, de nuevo doy mi asentimiento\apnota{4}{M assentio domino.} al reverendo padre, a saber, que aquella en lo temporal es muy inferior a aquella\apnota{5}{W. A. illam.} en lo espiritual. Roguemos al Dios óptimo que nuestros obispos reconozcan esto y administren lo espiritual no por vicarios y lo temporal por sí mismos\apnota{7}{M administrant.}. Llamé, pues, administración a la del régimen en lo espiritual de modo potestativo, y no obsta lo que dice el reverendo padre, que así él podría fingir que Juan es canciller, Marcos portador del sello, etc., porque, como dice Agustín, los filósofos usan palabras libres, pero a nosotros\apnota{11}{M autem.} solo nos es lícito hablar según una regla cierta; por eso conviene aportar cosas que reciban fundamento\apnota{12}{M recipiunt.} de la Sagrada Escritura,

como el primer miembro, el apostolado, en la igualdad de los apóstoles, y el otro miembro, el régimen, en el primado de Pedro. Por eso confieso que la autoridad de Pablo es robustísima, pero que en ni una sola letra me es contraria.

Despreció el reverendo padre el símil aducido por mí sobre los cismas en la provincia carente de gobernante, porque ya sucede así en los reinos fieles de España, Francia, Inglaterra y Hungría, que existen\apnota{18a}{U sunt.} sin cabeza, y sin embargo se mantienen entre ellos la paz y\apnota{18b}{U stat pax, quia.} la concordia. Me admiro de cómo el reverendo padre ha podido olvidar tantas matanzas, guerras y batallas, que no puede ignorar. No quiero decir\apnota{20}{M volo.} nada de tan diuturna desavenencia y perpetuo odio entre franceses e ingleses; ni del odio inveterado entre franceses y españoles, que tantas veces por el reino de Nápoles libraron guerras con máxima efusión de sangre cristiana\apnota{23}{M maximae Christianae.}, y ojalá al fin la paz fuese estable y permanente. Reconozco con el apóstol una sola fe, un solo Señor Cristo, pero venero al pontífice romano como vicario de Cristo. Queda, pues, la primera solución invicta y válida, al menos a mi exiguo juicio, del cual juzguen los más doctos.

Además, el padre digno de respeto dijo\apnota{25}{M U dicit.} que las palabras de Pablo no han sido tocadas por mí: Dios no hace acepción de personas\apnota{29}{U a. R. quia omnes sumus caeci, dixit Eccius (soll heissen Martinus?); W. A. fügt diese Randnote schon 2 Zeilen früher nach indicent ein.}, y por tanto menos aún resueltas. Quizá oyó mal, pues tanto por Ambrosio como por Jerónimo aduje el sentido de esas mismas palabras: Juan y Pedro habían sido idiotas e iletrados\apnota{31}{M alphabetici.} en la ley; Pablo, en cambio, había sido doctísimo\apnota{32a}{U fuit.}, aventajando a muchos de sus coetáneos, como de sí\apnota{32b}{M suos coetaneos, ut se (Mc de) se.} atestigua [Gá 1, v. 14]. No quiso, pues, Pablo estimar lo que habían sido Pedro y Juan, porque Dios no hace acepción de la persona del hombre\apnota{34a}{U accipit.}, pues así eligió al pescador como al docto\apnota{34b}{M doctam.} escriba. Pero con ello no se quita el primado, ni se imputa a Dios acepción de personas\apnota{35}{M personatur.}, aunque haya preelegido a uno sobre otro, como Pedro confesaba de Cornelio [Hch 10, v. 34]: En verdad he comprendido que no hay acepción de personas ante Dios.

Vengamos a lo aducido por el reverendo padre. [317] Cita al apóstol, 1 Cor 12 [v. 28], que al describir el cuerpo místico refiere apóstoles, profetas, doctores y en último lugar la administración. Digo que el apóstol describió egregiamente, pero que allí ni negó ni afirmó el primado; y es inválido el argumento\apnota{4}{U est ergo statt autem est.}: aquí Pablo no recuerda tal primado, luego el primado no existe en la Iglesia, tanto porque hasta los niños saben que el argumento por autoridad negativa no vale, como porque también lo no expresado en un lugar de la Sagrada Escritura debe tomarse, según la doctrina de Agustín en el libro sobre la doctrina cristiana, de otro lugar, como hemos hecho nosotros.

De igual modo respondemos a aquello de Ap 21. Confieso que son doce los fundamentos, y que estén figurados también por los doce bueyes, por los doce leoncillos\apnota{10}{P lenunculos.}, por las doce piedras del Jordán; pero en ellos\apnota{11}{M his.} se afirma el número duodenario de los apóstoles, y no por ello se niega el primado. Yo, en cambio, aduje aquel lugar porque es preciso mostrar concordes los pasajes\apnota{12}{M oportet fassuss.} de la Sagrada Escritura, y habiendo puesto el apóstol un solo fundamento, Cristo, y no pudiendo nadie poner otro, no debe negarse\apnota{15}{fehlt in M.} por ello que Pedro sea fundamento de la Iglesia, porque también el mismo san Juan puso doce fundamentos; lo cual, si no aceptas la distinción entre fundamento\apnota{16}{M susceperis de [fundamento fehlt], U funnamento.} de los fundamentos, esto es, principal, y fundamento sustituido o secundario\apnota{17}{M de [fehlt in P] fundamento subsecunda [Mc substituto] et secundario.}, no concordarás a Pablo con Juan.

En segundo lugar\apnota{19a}{M U secundo loco.} adujo que Pedro no pudo ordenar\apnota{19b}{M ordinari.} a Matías, y que así el Espíritu Santo segregó a Pablo y Bernabé. Respondo: como los apóstoles fueron obispos, no recuerdo haber encontrado en la Sagrada Escritura que fuesen ordenados obispos por Cristo, aunque en la última cena los ordenó primeramente sacerdotes. Por eso también puedo referir aquello a Pedro en cuanto a la ordenación episcopal; sin embargo, como admito que el oficio del apostolado solo es tributado por Dios, no es extraño que pidieran el auxilio divino, pues echaron suertes no al modo vulgar\apnota{25}{M modo, gestrichen, also vom Disputator widerrufen.}, sino que, como atestigua el sagrado Dionisio, aquella suerte fue una visible aparición del Espíritu Santo. Por eso aquel lugar habla en favor del primado de Pedro, porque ejecutó el primado que se le había concedido: En aquellos días, levantándose Pedro en medio de los hermanos\apnota{28}{M fratrum dicens.}, etc. Como también en lo de Safira y Ananías, como en la excusa de los apóstoles de que no estaban llenos de mosto, como en la respuesta y defensa ante el concilio\apnota{30}{M U consilio.} de los judíos [Hch 1-5] y en otras cosas contenidas en los Hechos de los Apóstoles, donde ejerció el régimen de la superioridad. Por eso creo que Matías fue hecho apóstol por el Señor y ordenado obispo por san Pedro. Así también hace tiempo concedí\apnota{33}{M concessit.} de Pablo que recibió\apnota{34}{M acceperat.} de Dios el apostolado, y confieso que trabajó más abundantemente que los otros.

Pero vengamos a lo tercero, que entre lo aducido por el reverendo padre parece hacer más\apnota{35}{U magis facere.} al propósito: El que cooperó con Pedro, a los Gálatas 2 [v. 8]. Y digo que Pablo no promulgó allí un derecho divino, sino que narró un hecho, pues también Pedro en los Hechos de los Apóstoles\apnota{38}{M dixit ad Judaeos: ecce convertimur ad gentes gestrichen, also vom Disputator widerrufen.} fue mandado instruir en la fe a Cornelio el gentil, y a los gentiles en Antioquía, donde presidió, e igualmente en Roma, donde al fin fue coronado con el martirio por Nerón el gentil, presidió y enseñó\apnota{40}{U a. R. non est verum, dixit Martinus.} juntamente a los judíos convertidos y a los que habían [318] de convertirse a la fe.

Por eso consta claramente que Pablo narró un hecho\apnota{2}{M U tunc narrasse.}. Paso por alto que san Pablo no recuerda el oficio del régimen o del primado de la presidencia, sino la ejecución del apostolado, según la cual los apóstoles se repartían por suerte diversas provincias, en lo cual también Pedro fue igual a los demás apóstoles según Anacleto, can. In novo, según Jerónimo contra Joviniano y sobre la carta a los Gálatas, según Cipriano en la carta 9 a Cornelio papa, lib. 4. Por esta razón la Sagrada Escritura ha de entenderse con efecto, de modo que advirtamos cuándo habla del oficio del apostolado y cuándo del régimen de la potestad o presidencia.

No importa que el reverendo padre niegue aún que yo haya aducido derecho divino en favor del primado y que no haya disuelto las autoridades aducidas por él, porque lo mismo diría yo de él. Sin embargo, ninguno de los santos padres entendió las autoridades aducidas por el reverendo padre para la negación del primado de Pedro sobre los demás, como él las adujo. En cambio, de múltiples maneras\apnota{14}{M plurivariam.} los santos padres entendieron que las autoridades aducidas por mí hablan del primado de Pedro. Por eso, en torno a esta conclusión, si alguna vez hablé algo con más acritud y dureza contra el reverendo padre, no quisiera\apnota{17a}{U nollem.} que pensarais que fue dicho\apnota{17b}{M dictam.} contra su persona, sino que la causa misma que se ventila me parecía tan preñada y necesaria que alguna vez hubieron de usarse\apnota{19}{M uterquam.} palabras más acres. En las demás conclusiones experimentaréis siempre\apnota{20}{M reperiemini.} la modestia eckiana.

Por eso también yo quiero terminar por ahora esta conclusión, porque reputo\apnota{21}{M repetens, U repetam.} que san Pedro tuvo de Cristo el primado de toda la Iglesia por la promesa a él hecha en Mt 16, como entienden Jerónimo, Cipriano, Agustín\apnota{23}{U Hieronymus, Ambrosius, Cyprianus.} y otros; por la nominación, por la que\apnota{24}{M U quia.} es nombrado primero en Mt 10, según la glosa allí mismo; por el pago del tributo, Mt 17, donde solo él fue equiparado a Cristo, según Agustín y Ambrosio; por que Cristo rogó por su fe indefectible y le mandó confirmar a los hermanos, Lc 22, según Crisóstomo y san León; porque fue constituido pastor de las ovejas de Cristo\apnota{28}{M Christi, Mc a Christo.}, Jn último [cap. 21], según Crisóstomo y Gregorio; porque Cristo dice a Pedro: Sígueme, no solo en el género del martirio, sino también en el orden del magisterio, según Teófilo; porque solo Pedro recibió el mandato de venir sobre el mar hacia Cristo [Mt 14], para que así el mundo, designado por el mar, le estuviese todo sometido\apnota{32}{fehlt in U, in Uv nachgetragen.} según Bernardo en su consideración a Eugenio, y por otras cosas aducidas arriba, reputo verdadera la conclusión comprobada\apnota{35}{U a concilio Constantiensi probatam.} por tantos santos padres y últimamente por el concilio de Constanza. No quiero, sin embargo, apoyarme en mi propia prudencia, dispuesto

a cautivar mi entendimiento en estas\apnota{36a}{U et in iis.} y en cualesquiera otras cosas al juicio de los jueces\apnota{36b}{fehlt in M.} que han de designarse, de la sede\apnota{37a}{U sedi.} apostólica y de otros varones prudentes y buenos\apnota{37b}{M vinorum.}.

Solo a Dios la gloria. Eck\apnota{38}{fehlt in U.}.

A las dos de la tarde se continuará la disputa.

[319] A las dos de la tarde del mismo día\apnota{1}{Die Zeitbestimmung [2 Zeilen] fehlt in M ganz, U hora secunda.}.

El D. D. Martín.

Nos vemos obligados a añadir un remate, porque el egregio D. D., al disolver mis objeciones, opuso a la vez muchísimas cosas y de nuevo me obligó a responder. Y aunque hubiera querido pasar esto por alto, sin embargo, para declarar plenamente mi sentencia a los jueces que han de designarse, respondo en pocas palabras.

Primero objetaba que Pedro presidía la elección de Matías\apnota{7}{M Matthei.}, exhortando a los apóstoles, etc. Creo, sin embargo, que ni al mismo\apnota{8}{M ipso.} D. D. le parece que por esto se pruebe el primado de Pedro, si Pedro predica, pues en lo anterior él mismo me respondió que una cosa es el apostolado y otra la administración del régimen. No dijo, pues, rectamente que Pedro al predicar así ejerció la administración del régimen\apnota{11}{non recte ergo dixit, Petrum sic praedicando exercuisse administrationem regiminis fehlt in P, offenbar durch ein Versehen (Homoioteleuton).}; de otro modo también Esteban será el primero y papa, porque\apnota{12a}{M U quia.} también él predicó e hizo muchos signos\apnota{12b}{fehlt in U, in Uv nachgetragen.} en el pueblo\apnota{13}{in populo fehlt in M.}.

Y en cuanto a que cree que Matías fue hecho apóstol y ordenado obispo por Pedro, me doy por contento con que crea lo que quiera; yo no lo creo si no lo prueba. Asimismo, que Pedro ejerció\apnota{15}{M excercuit, Mc exercuerit.} el primado en el castigo de Ananías y Safira puede decirse, pero no se prueba\apnota{16a}{M probari.}, pues también resucitó a muertos\apnota{16b}{et fehlt in U; M et resuscitavit et mortuos.}, lo cual también hicieron otros apóstoles. Asimismo, que excusó a los apóstoles para que no pareciesen ebrios tampoco es indicio de primado, pues, como dije, el egregio D. D. distinguió siempre el primado de todas las obras y administraciones comunes a los demás. Asimismo, que Pedro caminó sobre el mar y que Bernardo interpreta\apnota{20b}{U interpretatur.} por\apnota{20a}{in W. A. weggelassen.} el mar el mundo, lo concedo, pero nada atañe al primado, pues también nosotros debemos hollar el mundo.

Aquello de la indefectibilidad de la fe prometida por Cristo a ruego suyo\apnota{22b}{M probante.} lo paso de buen grado\apnota{22a}{M liber.}, pues concedo que la fe de Pedro nunca cayó, aunque él mismo cayó\apnota{23}{M ceciderat.}.

por la fe. Pues entonces el ladrón creyó con la fe de Pedro, cuando Pedro negó, como dice Agustín. Muy otra cosa es la fe que el primado\apnota{25}{M aliud fides et primatus.}.

Asimismo, en cuanto a aquello de que Pedro pagó el estáter por Cristo y por sí, en lo cual se dice que fue igualado a Cristo, lo concedo, pero nada aporta al primado, antes bien va contra el primado, pues enseguida sigue, como también interpreta Jerónimo\apnota{28}{M etiam interpretatur.}, que los demás apóstoles, ofendidos por sospechar que Pedro habría de ser mayor, comenzaron a discutir quién sería el mayor. Por eso Cristo, llamando junto a sí a un niño, reprimió su ambición, y así lo que sigue del texto combate el primado con más fuerza que lo que precede lo favorece.

Y en cuanto a aquello\apnota{32}{fehlt in U.} del último capítulo de Juan [c. 21], que lo dicho sólo a Pedro: Sígueme, según Teófilo haya de entenderse no sólo del género de martirio, sino también\apnota{33}{fehlt in P.} del orden del magisterio, rechazo a Teófilo, tanto porque\apnota{34}{U cum quod, W. A. cum.} Agustín parece sentir mejor al decir así: Sígueme padeciendo males temporales. Y aunque Agustín no lo dijera\apnota{35}{M dixerit.}, el propio texto del Evangelio lo probaría\apnota{36}{M convincit.}. Pues corrió la voz entre los discípulos de que aquel discípulo no moriría, al oír decir a Cristo: Si\apnota{37}{M sic modo.} quiero que él permanezca, ¿a ti qué? Tú sígueme. Por eso es manifiestísimo que aquella imitación ha de entenderse del género de martirio y de la pasión.

[320] Por lo cual ruego que el D. D. omita traer cosas nuevas, y sobre todo las que no hacen al caso, cuando también yo podría aducir aquello de Hch 8 [v. 14], donde los apóstoles enviaron a Pedro y a Juan como a un inferior, y Hch 15 [v. 13 ff.], donde Santiago confirmó y modificó el discurso de Pedro. Pero no he querido aducir estas cosas ni otras semejantes, pues hace ya tiempo que concedí a Pedro sólo el primado de honor y únicamente negué el primado de potestad sobre los demás apóstoles. Pues ésta es\apnota{6}{M enim est.} aquella acepción de persona, más aún, aquel primado de honor, que Dios no tiene en cuenta, como arriba aduje, y no aquella persona del pescador o\apnota{8}{U et.} de la pobreza, como expuso Ambrosio según refiere el doctor, pues Jerónimo interpreta mejor a Pablo por la magnífica cualidad\apnota{9}{M interpretavit.} por la que los pseudoapóstoles quisieron subvertir a los gálatas\apnota{10}{U quam pseudoapostoli ad Galatas.} como apartándolos de la autoridad de Pedro. Estoy pues\apnota{11}{U enim.} muy contento de que el D. D. haya confirmado\apnota{12}{M confirmavit etc.} su proposición con autoridades de Jerónimo, Ambrosio, Bernardo y León, sobre todo a partir de esta autoridad: Tú eres Pedro, con tal de que mi sentencia, como se ha dicho en lo precedente, se tenga por más confirmada con autoridades iguales y mayores y que hablan más a la intención, remitiendo también esto\apnota{15}{U quaeque.} al juicio, como arriba.

El D. D. Eck.

El reverendo padre vuelve a abrir hoy lo ya concluido y rebate como principales las cosas que yo había añadido como accesorias y adornos\apnota{18a}{M emplamata.} en la coronación de la conclusión\apnota{18b}{M adiecit, U in conclusione adiecit.}, y hasta tal punto procura la metamorfosis, que ignora\apnota{19}{M curam ut ignorem, U curat ut ignorem.} si se ha transformado en oponente o en respondiente. A la cosa, brevemente\apnota{21}{M brevius.}.

Nunca soñé aquella consecuencia: Pedro predicó, luego fue el primero. Luego reprueba mi creencia sobre Matías y, sin embargo, no disuelve mi argumento. Matías y los demás apóstoles fueron obispos y no fueron ordenados por Cristo, ni se ordenaron a sí mismos; por tanto, no fueron ordenados sino por Pedro, a quien Cristo constituyó pastor universal al apacentar las ovejas\apnota{27a}{universalem und pascendo oves fehlt in U, ersteres in Uv nachgetragen, M pasce oves.}. Por lo cual crea conmigo, o\apnota{27b}{W. A. vel aliud.} responda a ello.

En cuanto a la excusa de Pedro, me imputa\apnota{29a}{M iterum imponit.} cosas nunca pensadas\apnota{29b}{U cogitatum.}, que yo habría separado el primado de las obras de administración de los demás apóstoles, cosa que\apnota{30}{fehlt in U.} nunca fue de mi intención; antes bien, le alegué\apnota{31}{fehlt in U.} al bienaventurado Gregorio: que los demás obispos han sido llamados a participar en la solicitud, no en la plenitud de la potestad.

Sobre Bernardo, acerca del tránsito de Pedro sobre el mar, dice que nada hace al primado. De lo cual me admiro, cómo puede decirlo, si lee\apnota{34}{M U legit.} a Bernardo. Pues Bernardo, por su intención, quiso probar a partir de esto el primado de Eugenio\apnota{35}{M a. R. Petri.} sobre los demás, y que todo el mundo debe estarle\apnota{36}{M U debere ei esse (esse fehlt in M).} sujeto, no así a los demás obispos y apóstoles\apnota{37}{M apostolis, vel episcopis, U episcopis vel apostolis, Uv = P.}.

Sobre el estáter quiere sacar de la letra siguiente lo contrario, lo cual, sin embargo, no hizo\apnota{39}{M feci.}, sino que Cristo reprimió la murmuración de los apóstoles. Pero yo no tomo las palabras de la sagrada Escritura por mi propia cabeza, [321] sino según los santos padres, pues a partir de esto el bienaventurado Ambrosio quiso que se entendiera que Pedro habría de ser superior a los demás, lo cual también Agustín atestigua expresísimamente en las cuestiones sobre los Evangelios\apnota{3}{M a. R. August. q. 69.}. Por lo cual\apnota{4}{U quaestio mihi non occurrit quare.} no nos apoyemos en nuestro sentir, sino en los santos padres.

Sobre la indefectibilidad de la fe pasó de largo\apnota{5}{P M U defectibilitate transiit pater (pater fehlt in P).}, aduciendo, sin embargo, aquello de Agustín de que la fe\apnota{6}{M U fides.}, que Pedro negó, el ladrón la conservó en la cruz. Y yo paso de largo sobre aquel dicho. Sepa, sin embargo, el reverendo padre que san Cipriano, apoyado en las palabras de Cristo, quiso que ninguna herejía surgiera en Roma como\apnota{8}{U vel.} en las demás Iglesias, como ayer dije, que en la sede romana hubo doscientos

y doce papas y, sin embargo, antes de su remoción ninguno\apnota{10}{M nullus sensit.} erró en la fe de modo definitivo y dictando sentencia.

En cuanto a Teófilo, a quien niega por su propio arbitrio\apnota{11a}{M U pro suo nutu.} sin\apnota{11b}{U sive.} aducir autoridad ni razón alguna, me admiro de que desprecie así a los padres eclesiásticos quien es tan hostil a los sofistas. Adujó a Agustín, quien no se opone a Teófilo en que Cristo habló del género de martirio. Pero enseñe dónde negó Agustín\apnota{15}{P negagaverit (so!)} el orden del magisterio, cuando en el mismo tratado 124 sobre Pedro confiesa que él fue por gracia más abundante el uno y mismo primer apóstol. Y en la carta 162 dice: En la Iglesia romana, dice, siempre estuvo vigente el principado de la cátedra apostólica. Por lo cual no me\apnota{18}{M me non.} pida que no defienda nuestra conclusión, que reputo verdadera.

Lo que trajo\apnota{20}{M U Suaque.} sobre Pedro enviado a Samaría, y por ello arrebatado el primado, como si no supiéramos que alguien puede enviarse a sí mismo, como el Hijo de Dios enviado por el Padre y por sí mismo, no ha empezado a ventilarse\apnota{22}{M ventilare.} ahora por vez primera, sino que todo fue resuelto hace tiempo por Alfarus, Juan de Torquemada\apnota{23}{M per Johannem de Turre crematae.} y otros. Y si Pedro cedió ante Santiago en Jerusalén en el concilio, por razón del lugar y de la ancianidad, no le arrebata el primado a Pedro. Pero, para parecer que hace algo, el reverendo padre\apnota{26a}{reverendus pater fehlt in U.} atribuye a Pedro el primado de honor, acaso como\apnota{26b}{M velut forte velut.} a un ambicioso, aunque oyó de mí que, según Cipriano, Jerónimo y Anacleto, los apóstoles fueron de igual honor. Luego el reverendo padre atribuye a Pedro un primado de honor que los santos padres niegan, y en cambio niega a Pedro el primado de régimen que la concorde sentencia de los santos padres y el concilio de Constanza atribuye a Pedro.

Cese, pues, ruego, de gloriarse de estar apoyado en mayores\apnota{32}{M U se maioribus.} autoridades contra tantos santos padres, contra tan célebre concilio, pues ya mostré muchas veces que las autoridades aducidas por él\apnota{33}{M eam.} no niegan el primado según la sentencia de ningún santo padre. Por lo cual, más bien, con los griegos que al fin se hacen sabios, como se dice de los troyanos, aprenda por fin que el pontífice romano, verdadero vicario de Cristo, ocupa el primer lugar en el orbe, como los propios griegos, el emperador, el patriarca y los primados lo confesaron e hicieron obediencia a la Iglesia romana en el año 1439, el 22 de noviembre, y así se haga la paz en nuestros días. Todo lo cual, como arriba, encomiendo al juicio de aquellos a quienes toca y tocará. Solo a Dios la gloria. Eck\apnota{39}{fehlt in M U.}.

[322] El D. D. Martín.

Casi todo lo que ha dicho el egregio D. D. no me desagrada, en especial aquello de Agustín de que el principado de la cátedra apostólica\apnota{3}{M catholicae (Mc apostolicae) principatu.} estuvo siempre vigente

en la Iglesia romana, si hubiese añadido una palabra, a saber, el principado de potestad sobre\apnota{4}{U et unum verbum; M si addisset adhuc verbum unum principatus super.} todos los obispos, y creo que Bernardo retuerce el texto de Pedro caminando\apnota{5}{M U ambulante Petro.} sobre el mar, pues la consecuencia atribuye manifiestamente aquella caminata a la fe\apnota{6}{M U sequentia .... tribuant.}, no al primado\apnota{7}{M primatum.}, al decir Cristo, cuando Pedro se hundía: Hombre de poca fe, ¿por qué dudaste [Mt 14 v. 31]? Y en cuanto a lo que desea, que\apnota{8}{M quod.} yo crea con él que los demás apóstoles fueron ordenados obispos por el apóstol Pedro, de momento no me place obedecer, pues todos fueron\apnota{10a}{U fuerunt.} obispos igualmente como Pedro, según prueba el texto aducido por Pedro\apnota{10b}{M induxtus.} acerca de Judas\apnota{11a}{de Juda fehlt in M.}: Y su episcopado tómelo\apnota{11b}{M accipiet.} otro [Hch 1 v. 20; Sal 109 v. 8]. Esto sobre aquello\apnota{11c}{M ista.}.

El D. D. Eck.

Sobre el episcopado, en cuanto a que también Judas fue obispo según aquello del Salmo: Y su episcopado tómelo\apnota{14a}{M accipiet.} otro, digo que allí episcopado se usa\apnota{14b}{M ibi usurpari ibi.} por apostolado y no por el orden episcopal. Pues para la mayoría es dudoso si Judas llegó a ser\apnota{16}{M etiam Judas fuerat.} siquiera sacerdote. Pero esto nada hace al caso. Sin embargo, Judas salió tras tomar el bocado\apnota{17}{M porcella.}; por lo cual, cuando en la cena Cristo ordenó sacerdotes a los discípulos, Judas no fue\apnota{18a}{M fuerat.} consagrado obispo. Lo encomiendo para que se juzgue\apnota{18b}{U haec iudicanda.}, u otra cosa.

Fue renovada la disputa y enseguida se retomó la materia del purgatorio en torno a la novena posición.

El D. D. Eck comenzó como oponente\apnota{20}{Die Zwischenbemerkung lautet in M Eccius contra nonam conclusionem doctoris Martini sexta feria in die Kiliani. De animabus in purgatorio; in U De purgatorio. Contra nonam conclusionem argumentatur Eccius.}.

No sólo entre los teólogos, sino también en la sagrada Escritura y\apnota{21}{M U in (statt et).} en los santos padres se halla que esta vida presente es estado y término de merecer y desmerecer. Por lo cual los que están en el purgatorio ya no pueden merecer\apnota{23}{U posse, Uv possunt.} y así tampoco recibir\apnota{24}{etiam fehlt in U; M U non maiorem gratiam accipere.} mayor gracia. Que esto\apnota{25}{M U hoc.} es así lo pruebo por Jer 25 [v. 14]: Les daré según sus obras y según los hechos de sus manos; así también dice Pablo: Es preciso que todos nosotros seamos manifestados ante el tribunal de Cristo, para que cada uno reciba lo propio del cuerpo, según hubiere obrado, sea bueno, sea malo [2 Cor 5 v. 10]. Conociendo, pues, el temor del Señor, persuadimos a los hombres.

Añado a Jerónimo, para descender por grados, entendiendo allí al apóstol: Esto dice, digo, que mientras establecidos en el mundo peregrinamos

lejos del Señor, con buena conducta procuremos agradarle en el futuro, no, como algunos piensan, que después de salir del cuerpo\apnota{33}{U e corpore.} allí merezcamos algo obrando. Lo mismo confiesa la glosa ordinaria allí\apnota{34}{fehlt in U.}: En vano, pues, se promete el hombre después del cuerpo lo que descuidó conseguir aquí. Así también lo entendió Ambrosio: que sólo los hechos del cuerpo son premiados en otra parte.

De lo cual consta que, como la gracia aumenta\apnota{37}{M agatur.} según la correspondencia de los méritos y las almas en el purgatorio no podrán\apnota{38}{M poterint.} obrar nada meritoriamente por lo alegado, tampoco [323] podrá aumentar en ellas la gracia. Lo suscribe clarísimamente el bienaventurado Agustín en el Enquiridión cap. 109 y 110, y repite las mismas palabras principalmente en la cuestión 2 de las ocho cuestiones a Dulcicio\apnota{3}{M stulcicii.} después de muchas cosas. Por lo cual, dice, aquí se adquiere todo mérito por el que alguien pueda después de esta vida ser aliviado\apnota{4}{M revelari.} o agravado. Y nadie espere merecer ante Dios, una vez muerto, lo que aquí descuidó. Cesa\apnota{5}{M cesset.}, pues, después de la muerte el aumento de la obra meritoria; cese\apnota{6}{U cessat.} también el aumento de la gracia correspondiente a los méritos.

El D. D. Martín.

El egregio D. D. opuso dos autoridades de la divina Escritura y cuatro de los santos padres. La primera, de Jeremías, donde Dios dice: Les daré según sus obras. Me admiro bastante de que el egregio D. D. juzgue que éstas y otras semejantes pugnan contra mí, cuando yo en mi resolutorio aduje muchas más en esta forma y expuse cómo nada pugnan\apnota{14}{M pugnet.} contra mí o igualmente pugnan\apnota{15a}{M pugnet, U pugnant.} contra el purgatorio. Pues como la divina Escritura\apnota{15b}{U divina scriptura.} en su totalidad no contiene absolutamente nada sobre el purgatorio, sino que todo lo dice o del infierno o del cielo, quiero con esta\apnota{16}{M U hoc hac.} única respuesta haber respondido a todas las autoridades, que no hacen\apnota{17a}{U faciant.} al propósito, y en consecuencia\apnota{17b}{Uv sequentur.} tampoco todas\apnota{18}{fehlt in M U.} las exposiciones de todos los padres, a no ser que mencionen expresamente el purgatorio. Pues es creíble que, cuando tratan las santas Escrituras, sigan también el sentido de la santa Escritura\apnota{20}{fehlt in M.}. Y por eso con otras razones y autoridades conviene que se enseñe que las almas están ciertas y que no aumenta en ellas la caridad\apnota{21}{U eas in charitate.}.

El D. D. Eck.

Aunque el reverendo padre dice haber alegado\apnota{23}{U allegasse dicat.} más autoridades en el resolutorio y haberlas\apnota{24}{fehlt in U.} expuesto, ya las vimos. Vimos también que su glosa al texto nada aporta, lo cual ya mostraremos. Esto tampoco

aceptamos: que en la sagrada Escritura nada se halla sobre el purgatorio; dicho que ciertamente sería favorable a griegos y picardos, pero adverso a la fe cristiana, como rectamente recuerda en el resolutorio; sino opongamos su única solución al filo\apnota{28}{Löscher „in usto legitur arcem" beruht auf Irrtum.}, por la que dice que lo aducido nada hace al propósito. ¡Con qué elegancia respondió\apnota{29}{M respondeat.}! El Apóstol dice: Cada uno recibirá según lo que hubiere obrado en el cuerpo; el D. padre dice que el hombre recibe también según lo que hubiere obrado en el purgatorio. Agustín dice que todo mérito se adquiere aquí, y el D. padre dice en contra que también en el purgatorio se adquiere mérito. Jerónimo dice que, después de salir del cuerpo\apnota{34}{M excesseribus, Mc excesserimus.}, nada merecemos\apnota{35}{M promeremus (nos fehlt).} ante Dios; el D. padre dice en contra que, después de salir del cuerpo, aún merecemos en el purgatorio; así de Ambrosio. Juzguen ya vuestras ilustrísimas dominaciones, excelencias y prestancias\apnota{36}{fehlt in P.} si lo aducido por mí\apnota{37a}{M aducta, U inducta.} nada hace\apnota{37b}{M faciunt.} al propósito, cuando directamente contraría lo dicho por el padre\apnota{37c}{M directe dictis patrum, U directe nihil dictis patrum.}. Pido, pues, que realmente lo resuelva o muestre cómo no hace al pro-[324]pósito, y no se admire de que\apnota{1a}{U nihil faciat nec meretur (Walch XV Sp. 1198 moretur) quomodo; Uv verbessert faciant und miretur.} también yo le voy a aducir\apnota{1b}{fehlt in M.} más autoridades, para que\apnota{2}{M quod.} sea manifiesto que esto no es de teologastros, sino de verdaderos teólogos.

El D. D. Martín.

Yo, que creo firmemente, más aún, osaría\apnota{5}{M ausum.} decir que sé que hay purgatorio, fácilmente me persuado de que en las Escrituras se hace mención de él, como aquello de Mt [12 v. 32] que aduce\apnota{6}{M illud inducit illud Matthaei.} Gregorio en el Diálogo: No se le perdonará\apnota{7}{U remittittur.}, ni en este siglo ni en el futuro, queriendo\apnota{8}{M nolens.} que ciertos pecados se perdonen en el purgatorio.

Admito también aquello de 2 Mac [12 v. 45 f.]: Santo y saludable es el pensamiento de rogar por los difuntos\apnota{9}{M U sancta est et salubris cogitatio pro defunctis orare (U exorare etc.)}. Pero quiero esto: que en toda la Escritura no se halla memoria del purgatorio que pueda sostenerse en contienda y convencer, pues también el libro de los Macabeos, al no estar en el canon, es poderoso para los fieles, pero contra los pertinaces nada vale\apnota{12}{M U facit.}. Y el dicho de Gregorio facilísimamente se elude\apnota{13}{U eliditur.}, que ni aquí ni en el futuro se perdone el pecado\apnota{14}{M remittatur peccatum.}, esto es, nunca. Por eso no quiero que se me tenga por sospechoso de favorecer a bohemios y griegos. Esto por lo primero\apnota{15}{Hoc pro primo fehlt in M.}.

A la réplica del señor D. digo\apnota{16}{U dicam.} nada distinto de lo que ya antes dije. Pues repite lo mismo, y por eso también respondo lo mismo. Dije, pues, en

las sagradas letras nada se halla sobre el purgatorio, y por eso ni ellas ni sus exposiciones pueden aplicarse al negocio del purgatorio. Pues lo devuelto es\apnota{19}{U autem.} o bien a los buenos o mal a los malos después de la muerte. En el medio están los que están en el purgatorio, de quienes no hablan dichas autoridades, a quienes ni se ha devuelto bien ni\apnota{21a}{U neque.} mal, y\apnota{21b}{fehlt in M.} así el purgatorio pasa siempre por en medio de las autoridades\apnota{22a}{Mc 'illarum' auctoritatum.}\apnota{22b}{U transiit, Uv transit.}.

El D. D. Eck.

Que se excuse\apnota{24}{M U excusat} diciendo que cree y sabe que hay purgatorio, lo acepto bien y ya lo sabía antes. Pero que niegue que el purgatorio se prueba en la sagrada Escritura, lo rechazo como falso y grecánico, permaneciendo en el propósito, aunque ni aquello\apnota{27}{M U illud.} ni otras cosas mucho mayores puedan convencer contra protervos y pertinaces, como constó\apnota{28}{U patribus, Uv patuit.} en Arrio y otros. Pero que por eso se diga que no se halla en la Escritura, porque no pueda convencer al pertinaz, es del todo capcioso y adverso a nuestra religión.

Sobre los libros de los Macabeos, de los que dice que valen para los fieles pero que no están en el canon, también esto, digo, es falso, pues aunque entre los hebreos no estuvieron en el canon, la Iglesia, sin embargo, los recibió en el canon, como atestigua el padre del señor padre, Agustín, en el libro de la ciudad de Dios\apnota{34}{U testatur lib. 8. [W. A. lib. 18].}, y san Ifón inserta en sus decretos la constitución por la que la Iglesia recibió aquellos libros en el canon.

Que el reverendo padre quiera persistir en su propósito de que lo alegado por mí no hace al propósito, cuando, sin embargo, directamente le contraría. En primer lugar\apnota{38}{M U primo enim.}, lo que dice de que aquellas autoridades miran al cielo o al infierno no ayuda, porque directamente hacen mención del estado de merecer y desmerecer y definen\apnota{1}{U diffiniunt.} que este [325] estado termina con la muerte. En cambio, el reverendo padre, sin apoyarse en ningún santo doctor\apnota{2a}{M doctori.}, osa prorrogar a las almas que salen del cuerpo\apnota{2b}{U existentibus.} el término de merecer y afirmar que los méritos del purgatorio son remunerados en el cielo, y la sentencia apostólica clama que sólo se remuneran las obras que el hombre obró en el cuerpo. Por lo cual examiné lo dicho en su resolución\apnota{5}{U perpendi.}, pero no pude aceptarlo como insuficiente, como antes abrí\apnota{6}{M aperuit, U aperiat.} aquellas contrariedades.

El D. D. Martín.

En primer lugar\apnota{8}{U primum.}, en cuanto a lo que dice el egregio D. D. de que no por eso ha de negarse algo de la Escritura porque los pertinaces no puedan ser convencidos\apnota{9}{M possunt.}, lo dice óptima y verdaderísimamente, pero yo hablo\apnota{10a}{M loquar.} de aquellos pertinaces que con nuestra propia\apnota{10b}{M U de his ..... nos nostra.}

autoridad y con nuestro propio dardo pueden 'atravesarnos'. Pues es evidente\apnota{11}{M convincere (Mc a. R. confodere); enim fehlt.} que el libro de los Macabeos 'pertenece' al Antiguo 'Testamento'\apnota{12}{M ad vetus 'testamentum pertinere'.}. Cuando, pues, 'san Jerónimo' hubiere descrito\apnota{13a}{M scripserit, Mc descripserit, U conscripserit.} el canon hebreo y defina que sólo valen en\apnota{13b}{M de.} contienda los libros que son del canon, y haya sido recibido en esta sentencia suya\apnota{14}{U sua sententia.}, fácilmente seremos perforados por nuestro propio dardo, a no ser que persuadamos a los fieles\apnota{15}{M U telo verberabimur (Uv perforabuntur) nisi fidelibus (M videbimus).}.

En segundo lugar prueba\apnota{16}{M probat.} que el libro de los Macabeos ha sido recibido en el canon; contiende sobre una equivocación, y fácilmente nos concordaremos. Sé que la Iglesia recibe este libro, y esto\apnota{18}{fehlt in M.} lo dije, pero la Iglesia no puede atribuir más autoridad o\apnota{19}{M vel.} firmeza al libro que la que por sí mismo tenga, así como aprueba y rechaza\apnota{20a}{fehlt in M.}\apnota{20b}{M U recipit.} los opúsculos de los demás padres, pero no por eso confirma todo ni hace mejores las cosas que en ellos se dicen. Pase, pues, esa equivocación de canon y canon\apnota{21}{M approbat statt omnia confirmat; igitur statt ergo; et canonem 'et canon'; U non ideo confirmat aut meliora reddit. Transeo ergo ista, quae in multis dicuntur canon et canon; Uv = P, nur omnia fehlt.}.

En tercer lugar dice que las autoridades aducidas hacen mención directa\apnota{22}{M intentionem.} del estado de merecer y que éste termina con la muerte\apnota{23}{fehlt in M.}. Respondo como antes: Porque nada recuerdan del purgatorio, sino sólo del cielo o\apnota{24}{M U tantum de caelo et.} del infierno, por eso el que muere para el infierno no merece y el que muere para el cielo no merece\apnota{25}{M U haben die beiden mit moriens ad anfangenden Sätze umgestellt, M beidemal nihil statt non, et fehlt in M U.}, y así termina el estado de merecer de uno y otro; de otro modo, con esas mismas autoridades podría probarse evidentísimamente el purgatorio.

En cuarto lugar dice que yo no me apoyo en ningún doctor\apnota{27}{U auctore.} y que prorrogo el término de merecer, y que afirmo que los méritos del purgatorio son remunerados. Lo confieso, pues esto lo hice para disputar y oír cosas mejores que las que yo sepa. Pues yo nada sé del purgatorio, sino que las almas allí padecen y han de ser ayudadas con nuestras obras y oraciones, dispuesto a ser enseñado humildemente, si pueden transmitirse más cosas sobre él mismo\apnota{31}{M ipsis.}. Por lo cual el Apóstol aducido contra mí\apnota{32a}{U inductus.}, que sólo se remuneran las obras obradas en el cuerpo\apnota{32b}{M in corpore gesta opera remunerabuntur.}, ha sido tomado por el señor\apnota{33}{fehlt in M.} D. en relación al purgatorio, como ya dije que no puede tomarse, sino sólo al cielo o\apnota{34}{M et.} al infierno. Si se observa esa relación, es patente que no pugnan contra mí, o demostrarán\apnota{35}{M U pugnet contra me aut demonstrabit (M demonstabit) purgatorium.} el purgatorio.

Mañana a la hora séptima se continuará la disputa.

El sábado 9 del mes de julio a la hora séptima antes del mediodía\apnota{36}{Die Zeitbestimmung lautet in M Eccius sabbato die mane, Löscher liest falsch D. Marci, W. A. vermutet 8. Marie; U Nona Julii.}.

El D. D. Eck.

Primero, cuando el D. padre dice que las autoridades aducidas\apnota{38}{U inductas.} no se entienden del purgatorio, place en cuanto a esto, porque ponen\apnota{39}{fehlt in U, in Uv nachgetragen.} el mérito sólo en esta vida, luego excluyen el mérito [326] en el purgatorio\apnota{1}{U ergo purgatorium excludunt in hac vita, Uv = P.}. Además es cierto que Agustín en el Enquiridión habla también del purgatorio, porque allí mismo prueba que las almas son aliviadas por los sufragios de los vivos. Y Jerónimo increpa a los que piensan que las almas despojadas merecen, el cual afirma el purgatorio, sobre el último capítulo de Isaías [66 v. 24].

Pero en cuanto a que dice que su conclusión es disputación, como no teniendo gran confianza en ella, me admiro muchísimo, cuando ya después de un año la ventiló e inscribió la disputa con magnífico título contra los errores nuevos y viejos, y afirma que la sentencia de los teólogos más recientes\apnota{8a}{M recensiorum.} ni siquiera es verosímil para un necio\apnota{8b}{M stultam.} y con desprecio los llama teologastros\apnota{9}{M nuncupans.}. Pero en cuanto a que imagina una equivocación del canon, no lo soporto\apnota{10}{M patio.}, porque Agustín en aquel libro 18 de la ciudad de Dios no pudo equivocar el mismo término, máxime en el punto adversativo, al decir que no estuvo en el canon entre los hebreos, sino entre la Iglesia. Además es cosa explorada que, habiendo sido escritos\apnota{14}{U scripta.} más evangelios, por autoridad de la Iglesia cuatro fueron recibidos en el canon, y así el prólogo atestigua que los libros de los Macabeos fueron recibidos, pero que por la Iglesia son anotados entre las historias de los volúmenes divinos. — Pero mejor a la cosa.

Porque se funda en esto, que el purgatorio no está expreso en las sagradas letras, contra lo cual está el decreto del concilio\apnota{18}{U consilii.} florentino, que también los griegos asumieron abjurado el error, y después\apnota{19}{U deinde.} muchos lugares de la sagrada Escritura, como aquello: Pasamos por agua y fuego, y nos sacaste\apnota{20}{M reduxisti, U perduxisit.} al refrigerio [Sal 66 v. 12]; aquello del Eclesiastés\apnota{21}{U Ecclesiastes.} 4 [v. 14]: De la cárcel y de las cadenas sale para reinar; aquello de Mt 5 [v. 25]: Dijo el Señor: Ponte de acuerdo etc., no sea que seas echado en la cárcel. En verdad te digo, no saldrás de allí hasta que pagues el último cuadrante. Donde por cárcel entiende el lugar del purgatorio, del que el hombre no es liberado hasta haber satisfecho plenamente, como interpreta el bienaventurado Ambrosio\apnota{25}{U Gregorius, Uv = P.} en el lugar correspondiente sobre Lucas. Sin embargo, la razón de Gregorio, que retoma el bienaventurado Bernardo en el sermón 66 sobre el Cantar, es válida: que Cristo insinúa que algunos pecados se perdonan en el siglo futuro.

Pero el texto más abierto es 1 Cor 3 [v. 12—15]: Si la obra de alguno ardiere, sufrirá detrimento, pero él mismo será salvo, mas así como por fuego. Y el día del Señor lo declarará, porque será revelado en fuego, y el fuego probará cuál sea la obra de cada uno, y allí sobre la paja, el heno\apnota{31}{M veno.} etc., donde el bienaventurado Ambrosio dice que el Apóstol habló expresamente\apnota{32a}{M U expresse.} del purgatorio. Lo mismo\apnota{32b}{M U idem.} atestigua la glosa interlineal, y también la glosa ordinaria. Lo suscribe el bienaventurado Jerónimo en el libro 2 contra Joviniano, y también el bienaventurado Gregorio lo toma del purgatorio, lo que también hace el bienaventurado Bernardo sobre el Cantar. Y para que nadie me objete a Agustín en el Enquiridión, cap. acaso\apnota{36}{fehlt in M U.} 67, quien dice que aquellas palabras pueden entenderse también del fuego del presente siglo, lea al mismo más completo allí mismo y en la cuestión 2 de las ocho cuestiones a Dulcicio, y hallará que Agustín toma aquellas palabras del Apóstol, como son manifiestísimas, también\apnota{38}{M accipi.} del purgatorio.

[327] Por lo cual el purgatorio se afirma no sólo en los libros de los Macabeos, lo que ciertamente bastaría a la Iglesia, sino también en otras Escrituras, a no ser que Gregorio, Ambrosio, Agustín, Jerónimo y Bernardo sean del número de los teologastros; de otro modo no sé cómo el reverendo padre puede confesar que sabe que hay purgatorio, cuando quiere apoyarse tan válidamente en la letra de la sagrada Escritura, y cuando tan clarísimos son los testimonios de la Escritura en favor de los teólogos más nuevos. Él no podrá asignar ni un solo lugar en la Escritura\apnota{7a}{M scripturae statt in scriptura.} donde se afirme que los méritos o la gracia aumentan en los que se purifican\apnota{7b}{M purgantis.}, cuando en cambio dice el Sabio, Ecl\apnota{8}{U Ecclesiastes.} 11 [v. 3]: Y si el leño cayere hacia el austro o hacia el aquilón, en el lugar donde cayere, allí estará. Glosa ordinaria: Esto es, el lugar que aquí te hubieres preparado, entonces tendrás. Y\apnota{10}{fehlt in U.} porque en la casa de mi Padre, dice Cristo, muchas mansiones hay, al que muere se le asigna cierta mansión, más allá de la cual no puede ascender por sus méritos en el purgatorio. Así lo atestigua el Damasceno lib. 2 cap. 4: Conviene saber que lo que es en los hombres la muerte, eso es en los ángeles la caída, y lo declara en cuanto al término de merecer. Así el Sabio Eclesiastés\apnota{15}{U Ecclesiastes.} 9 [v. 10]: Todo lo que puede tu mano obrar al instante etc.\apnota{16}{fehlt in M U.}, como para aquel sentido lo alegó aquel devoto y digno ministro de la Iglesia, Juan Capistrano\apnota{17}{M cupristanus.}. Así lo enseñó el Apóstol a Gál 6 [v. 7 ff.]: No erréis, Dios no es burlado. Pues lo que el hombre sembrare, eso también segará. Porque los que siembran en su carne, de la carne segarán corrupción etc.\apnota{19}{M metet; etc. fehlt in M U.}, y haciendo el bien no desfallezcamos; luego, mientras tenemos tiempo, obremos el bien\apnota{20}{M oparemur; bonum fehlt in M, U tempus habentes operemur bonum.} a todos. Glosa: lo que el hombre sembrare\apnota{21}{M seminaveris.}, a saber, con el trabajo de la vida presente. Luego no en el purgatorio

…se siembran los méritos, sino en la vida presente. Lo cual confirma también Crisóstomo por aquello de Jn 9 [v. 4]: «Me conviene obrar las obras del que me envió, mientras es de día; pues viene la noche, cuando nadie puede obrar». El día, dice Crisóstomo, es la vida presente. Por eso Agustín nos exhorta a que debemos\apnota{26}{M U hortatur nos debere (M deberi).} obrar mientras vivamos, no sea que la noche nos sorprenda. Añádase como testigo también el bienaventurado Ambrosio, que entiende que el Apóstol habla del tiempo de la vida presente a nosotros concedido\apnota{28a}{M U concessae.}, para que vivamos justamente. Por lo cual dice el egregio salmista, Sal 104 [v. 23]\apnota{28b}{M 33 statt 103.}: «Saldrá\apnota{29}{P M U exivit.} el hombre a su obra y a su labor hasta la tarde».

Por último, expresa es aquella sentencia de Agustín, además de las arriba alegadas\apnota{30}{M allegatam.}, donde sin embargo hace mención del purgatorio; su sentencia es la del cap. 3 De fide ad Petrum: «El tiempo de alcanzar la vida eterna lo dio Dios a los hombres solo en esta vida». También Jerónimo, en el lib. 3 sobre aquello de Gal 6 [v. 5]: «Cada uno llevará su carga», muestra cómo los hombres son ayudados por sus méritos en esta vida, pero, llamados ante el tribunal, lo que sucede en la muerte, entonces ni Job ni Noé ni Daniel podrán interceder por nadie, sino que cada uno llevará su carga. Pues también\apnota{36}{M U ut.} bellamente enseña el bienaventurado Bernardo que\apnota{37a}{fehlt in M U.} en la muerte de cada hombre se realiza\apnota{37b}{M U fiat (U fiet) de eo.} un juicio particular, lo cual no podría darse si [el alma] tuviera todavía plazo para acumular\apnota{39}{U cumulandi.} sus méritos en el purgatorio y aumentar la gracia; porque, como dice por el profeta Zac [4, v. 7], igualará gracia con gracia, devolverá gracia de gloria en proporción a la gracia de la fe. [328] Y aunque aquellos testimonios no fuesen tan claros, siendo así que son clarísimos, esto solo debería apartar del propósito al reverendo padre: que un teólogo, sin fundamento teológico, a su propio arbitrio\apnota{3}{M U suo arbitrio.}, no siendo juez ni asesor del juez, quiera prolongar\apnota{5}{M promerendi vult prorogare.} a las almas que parten el plazo de merecer. Por lo cual, en cuanto a esta partícula, parece que debe concluirse que las almas en el purgatorio padecen bastante y son purificadas de sus delitos, pero no merecen mayor gracia.

Sobre lo segundo principal, porque\apnota{7}{M qua, U quo.} niega que las almas en el purgatorio estén ciertas de su salvación, esto de nuevo juzgo que no está fundado en ninguna autoridad de la sagrada Escritura. En cambio, de que estén ciertas de su salvación, aduzco primero\apnota{9}{fehlt in U.} aquello de Ap 5 [v. 1 ss.]: «Vi en la diestra del que está sentado sobre\apnota{10}{U supra.} el trono… Y nadie fue hallado digno de abrir el libro, ni en el cielo ni en la tierra ni debajo de la tierra». En el infierno indudablemente no hay dignidad alguna de abrir el libro; por lo cual habla del purgatorio, en el cual también hay a veces santos varones que resplandecen con milagros, como del santo Pascasio atestigua el bienaventurado Gregorio en el diálogo 4, y la historia del santo Severino, obispo de Colonia, de nación, sin embargo,

suevo, afirma\apnota{15}{P M U asserit, Uv accidit.} lo mismo. Y más abajo: «Y todos cantaban, en el cielo, en la tierra y debajo de la tierra\apnota{16}{M terra.}». Mas quienes están como desesperados de su salvación no cantan; por lo cual, por el hecho mismo de que el alma sagrada afirma que allí hay canto\apnota{17}{U tantum, Walch XV. 1206 und nach ihm W. A. richtig cantum.}, afirma que las almas están ciertas\apnota{18}{M certas esse.} de su salvación.

Además, en el canon de la misa, donde se ora por los difuntos en el purgatorio, decimos así: «Acuérdate, Señor, de tus siervos, que nos precedieron con el signo de la fe y duermen en el sueño de la paz», etc.; «a ellos y a todos los que descansan en Cristo», etc.\apnota{21}{fehlt in M U.} Si, pues, según el sacratísimo canon de la misa las almas descansan en Cristo, ¿cómo pueden estar en tal turbación y horror como de desesperación, según declara el reverendo padre en su resolutorio de la concl. 14 a la 20? Pues no veo cómo en tal horror, temblor, turbación y como desesperación, cosas que todas dicen grandísima inquietud, las almas que han de purificarse puedan decirse que duermen en el sueño de la paz, siendo así que todo lo dicho se opone a la paz. En cambio, con razón serán juzgadas por los verdaderos cristianos que descansan en paz, mientras aguardan seguras el término de su purgación. Estas son en parte las cosas que me movieron y hoy me mueven a disentir del reverendo padre, dispuesto sin embargo a estar al juicio e información\apnota{30}{M informatione.} de otros.

El D. D. Martín.

Tres puntos por orden ha desarrollado contra mí copiosísimamente el egregio D. D.: primero, que el canon de la Escritura\apnota{33}{Mc 'sacrae' scripturae.} comprende los libros de los Macabeos; segundo, que el purgatorio se prueba también por otros lugares de la Escritura; tercero, ha intentado mostrar que las almas están ciertas de su salvación.

Respondo: sobre los dos primeros no hay ninguna\apnota{36}{M U nulla est.} controversia entre mí y su señoría. Por eso no era necesario acumular tantas autoridades para aquello [329] que quizá afirmo con más constancia que él mismo, como quien he profesado saber\apnota{1}{fehlt in M.} que existe el purgatorio. Pero aquí se debate el punto capital: si puede\apnota{2}{U posset.} probarse que las almas en el purgatorio merecen y que la gracia en ellas\apnota{3}{M eius statt in eis.} aumenta; pasemos, sin embargo, por orden.

Primero dijo que, con las autoridades ayer aducidas, establecí que el mérito está solo\apnota{4}{M solum meritum.} en esta vida, luego se excluye en la vida futura, lo que él entiende del purgatorio, y así adapta las autoridades al purgatorio. Yo, en cambio, concedo que se trata de la vida futura, del infierno o del cielo. Adujó también a Agustín en el Ench. hablando del purgatorio, y a Jerónimo increpando a quienes afirman que las almas separadas merecen. Una y otra cosa concedo. Me reprende también por haber dicho

que yo disputo\apnota{10}{M quod disputem statt quod dixerim me disputare.} esta proposición como no teniendo confianza en ella, yo que con magnífico título inscribí\apnota{11}{M U scripserim.} que disputaría contra los errores nuevos y viejos, y afirmé que la sentencia de los más recientes\apnota{12}{veriorem.} ni siquiera es verosímil para un necio, y llamé con desprecio «teologistas» a esos\apnota{13}{U nominarim; M similem (Mc a. R. veri) affirmari ad contemptim 'eos' theologistas nominari (Mc a. R. appellari).}.

Digo, como\apnota{14a}{M quod et statt sicut.} ayer dije, que todavía nada sé «del negocio de las almas y del purgatorio»\apnota{14b}{Mc de negotiis animarum in purgatorio.}, y a esto llamo\apnota{15a}{M appello.} error: que algunos se atreven a afirmar\apnota{15b}{M assere.} lo que ignoran, y las opiniones de los padres, como del bienaventurado Tomás, de Buenaventura y semejantes\apnota{16}{et similium fehlt in M.}, que ellos mismos no afirmaron, ellos las proclaman como ciertos artículos de fe. A estos yo los llamo teologistas y no teólogos. Las opiniones deben tratarse en las escuelas; al pueblo deben predicarse\apnota{19a}{fehlt in U.} las palabras y obras de Dios. Sal 19 [v. 2]: «Los cielos narran la gloria de Dios, y las obras de sus manos\apnota{19b}{M psal. 108, U psal. 8; M U etc. statt et opera manuum eius.}». No condeno yo,\apnota{20}{fehlt in M.} pues, las opiniones de los óptimos padres, sino que resisto a aquellos artífices que de las opiniones de los hombres nos forjan artículos de fe, lo cual no es oficio de buen teólogo.

Que yo haya usado equívocamente el canon contra Agustín, lib. 18, cap. 36,\apnota{23}{M de trinit. c. 26, U c. 26.} me obligaron a ello el divino Jerónimo y\apnota{24}{M U item.} Eusebio, reseñando en la historia eclesiástica las autoridades de los antiguos\apnota{25a}{M U et antiquorum.}. Por eso subsiste la equivocidad, en tanto\apnota{25b}{U cum.} que de un modo siente Agustín y de otro Jerónimo con\apnota{26}{U de.} los más antiguos acerca del canon, y por consiguiente no queda ninguna fuerza\apnota{27}{U relictum est.} del argumento en la contienda. Si el prólogo de Jerónimo enumera\apnota{28}{U enumeret, M enumerat.} los libros de los Macabeos entre los volúmenes divinos, no lo recuerdo; paso por alto aquello de que los cuatro evangelios han sido recibidos\apnota{29}{M accepta, gestrichen, recepta.} por autoridad de la Iglesia. Pues esta será otra materia.

Después de esto dice que es contra el concilio Florentino que el purgatorio no\apnota{32a}{Mc quod non, letzteres Wort [rot] durchstrichen.} esté expreso en la Escritura. Respondo que\apnota{32b}{fehlt in U.} el concilio no puede hacer que sea de la Escritura\apnota{33}{M scripturae facere, U facere de scriptura esse.} lo que por su naturaleza no es de la Escritura, como tampoco la Iglesia pudo hacer los evangelios, aunque aprobó\apnota{34}{M approbaverit.} los evangelios.

Por lo cual veamos las autoridades. Primero aquella de Sal 66 [v. 12]\apnota{36}{M prima illa psal. 165.}: «Pasamos por fuego y agua». Respondo: no vale para el purgatorio; habla de las persecuciones de los santos, como en otros muchos lugares: Sal 17 [v. 3]: «Me examinaste con fuego»; asimismo 26 [v. 2]: «Abraza mis riñones

y mi corazón», y 1 Pe 1 [1, v. 6 s.]\apnota{39}{U 1. Petr. 1.}: «Un poco, si ahora es preciso entristeceros en varias [330] tentaciones, para que la prueba de vuestra fe sea mucho más preciosa que el oro, que por el fuego se prueba\apnota{2a}{M probetur.}». Y en breve, aquel es un tropo vulgarísimo\apnota{2b}{M 'iste' est vulgatissime, U est iste vulgatissimus.} de la Escritura: entender por fuego y agua\apnota{3}{M per aquam et ignem.} las tribulaciones. Por eso están demasiado apegados a letras y sílabas, como me imputa el egregio D. D., quienes toman el fuego por el purgatorio\apnota{4}{U ideo nimis haeret . . . . accipit.}.

Asimismo también aquello del Eclesiastés: que alguien sale de la cárcel y de las cadenas\apnota{5}{M catena.} al reino\apnota{6a}{M etc. statt egreditur ad regnum.}. Bastante claro es\apnota{6b}{U est clarus.} el texto que disputa sobre la vanidad de este mundo: que por giro del azar el que ahora es siervo se hace rey, y el que es rey\apnota{7}{M his statt is, qui rex statt et qui rex, U et qui rex est.} se hace siervo. Con estas y semejantes autoridades, si lucháramos contra quienes [el purgatorio] niegan, no haríamos de nosotros y de la Iglesia sino ludibrio ante los adversarios; aunque yo gustosísimamente admita\apnota{10}{M haec omnia libentissime admitto, U libentissime omnia haec admitto.} todas estas cosas, y cuantas haya semejantes.

Tal es también aquello de Mt 5 [v. 25]: «Ponte de acuerdo con tu adversario, no sea que te entregue al verdugo\apnota{12}{M U tortori etc.}», donde dice que por la cárcel se entiende, según Ambrosio, el lugar del purgatorio; lo admito de buen grado. Pero como otros padres lo expusieron en sentido diverso, sobre todo Agustín, que lo expone del infierno del que nunca se saldrá, la autoridad se hace dudosa: persuade a los fieles, pero no convence a los que resisten\apnota{15}{U resistens.}; por callar que la ilación del texto ni siquiera consiente que se entienda del purgatorio, pues habla del que consiente con el adversario y del que disiente, queriendo que el que disiente peque condenablemente y contra el precepto de Cristo, cosas que\apnota{18}{P qui.} pertenecen al infierno, no al purgatorio.

La razón de Gregorio, dice, consentida por Bernardo, es válida: que Dios perdona\apnota{21}{U remittit.} algunos pecados en el futuro. Respondo: es verdadera entre fieles, pero inválida ante los que resisten.

Después de esto dice que es texto clarísimo aquel de 1 Cor 3 [v. 15]: «él mismo\apnota{23}{U si cuius, ipse.} se salvará, pero como por fuego», y que Ambrosio y Jerónimo, la glosa\apnota{25}{U glossa.} ordinaria y la interlineal, Gregorio y Bernardo entendieron al Apóstol del purgatorio, aunque Agustín diga en algún lugar\apnota{26}{M alicui.} que las mismas palabras pueden entenderse del fuego del presente siglo. Respondo: tan poco clarísimo es este texto, que hoy confieso, tras haber examinado mucho, que todavía no sé\apnota{28}{M non scire.} el sentido genuino de Pablo, tan variantes son los intérpretes; aunque de buen grado admita para mí el del purgatorio. Pero como el Apóstol dice abiertísimamente que la obra de cada uno será probada por el fuego, fuego en el que dice que será revelado el día del Señor, que declarará la obra de cada uno

la obra, como son palabras abiertas, a mi tenue juicio parece hablar del fuego de la conflagración y del juicio extremo, o, como Agustín lo tropologiza\apnota{34}{M tropologisat, U tropologice.}, del fuego de la persecución temporal, por la que sobre todo se prueba la fe y la doctrina de la fe, y cuantas cosas sobre estas se edifican.

Por lo cual todavía nada manifiesto se tiene de las sagradas letras\apnota{36}{M ex sacra scriptura.} sobre el purgatorio que valga en la contienda. No tengo, pues, a los beatísimos padres por teologistas, porque con ellos confieso el purgatorio; ni ellos sancionaron como artículos sus opiniones e ignorancias sobre el estado de las almas, como hacen los teologistas.

[331] Se admira el D. D.\apnota{1}{fehlt in P.} de cómo puedo saber que existe el purgatorio, no teniendo yo ni un solo lugar\apnota{2}{M habeo.} en las Escrituras a mi favor, siendo así que estos, como dice, clarísimos testimonios\apnota{3}{M pro theonericis scripturis (so!).} de la Escritura juzga que favorecen a los teólogos neotéricos. Respondo: no es necesario confesar por qué vía haya conocido el purgatorio, o cualquier otra cosa. Además, estos\apnota{4}{M dehinc statt deinde haec.} clarísimos testimonios han sido aducidos a favor del purgatorio, no a favor del estado de las almas en el purgatorio.

Ya aquella autoridad de Ecl 11 [v. 3]\apnota{7a}{P M decimo.}: «En cualquier lugar\apnota{7b}{fehlt in P M.} donde cayere el leño, ya al austro ya al aquilón\apnota{8}{M sive austrum sive aquilonem.}, allí estará», donde la glosa ordinaria entiende: el lugar que aquí te hayas preparado, entonces\apnota{9}{M tibi hic praeparaveris tunc; tunc fehlt in U, in Uv nachgetragen.} tendrás, no sé con qué ingenio se arrastra al propósito; pues si por lugar preparado y que se ha de tener\apnota{10}{U et humidum statt et habendum (W. A. et habitum).} entiende el purgatorio, habrá que permanecer eternamente en el purgatorio; pero si entiende que lo que aquí\apnota{12}{fehlt in M, U quod si hic.} hayas preparado, esto es, merecido, entonces no hace a la autoridad de Ecl 11, que nada habla del mérito, sino de la muerte del hombre.

Paso por alto aquello de que Cristo dice que en la casa de su Padre hay muchas mansiones, y\apnota{15}{M in domo patris mei sunt mansiones multae; 'et' et (so!)} que al que muere se le asigna cierta mansión, más allá de la cual no puede ascender. Todo es a mi favor. Sé que a cada uno se le asigna una mansión después de la muerte, pero no llega enseguida a la mansión después de la muerte, a no ser que tome de nuevo la mansión por el purgatorio, y así excluya a los moribundos del reino eterno\apnota{18}{M accipit et sic morientes a regno coelorum (Mc 'ab regno aeterno') excludat; U excludit, Uv excludat.}.

Asimismo también aquello del Damasceno: lo que para los ángeles la caída, eso es para los hombres la muerte. Respondo: si, pues, después de la muerte cayeren\apnota{21a}{M U ceciderit.} en el purgatorio, se sigue que eternamente serán purgados\apnota{21b}{U purgabitur, Uv purgabuntur.}; o es preciso entender al Damasceno, como todo lo precedente, de las dos últimas mansiones, la condenación y la bienaventuranza, y no del purgatorio.

Así también aquello de Ecl 9 [v. 10]: «Todo lo que pudiere tu mano, óbralo al instante, porque ni razón ni sentido hay en los infiernos, adonde tú\apnota{25}{M quod tu, Mc quo tu.} te apresuras». Si lo entiende del purgatorio, de nuevo el purgatorio será infierno; luego habla solo del infierno, sin memoria del purgatorio.

A Juan Capistrano lo admito de buen grado, pero fuera de la contienda\apnota{28}{M conventionem.}.

Ya que lo que dice el Apóstol, Gal 6 [v. 7], habla contra nuestra sentencia: «Lo que el hombre sembrare en\apnota{30a}{fehlt in P.} la carne, eso también segará», y aquello de Jn 9 [v. 4]\apnota{30b}{U haec metet; P M U Johan. 8.}: «Viene la noche, cuando nadie podrá obrar\apnota{31}{M operare.}», lo que Ambrosio expone del tiempo de la vida presente, en el que\apnota{32a}{M de praesentis temporis vitae exponit in quo, U de praesenti temporis vita [Uv tempore vitae] exponit, quomodo.} merecemos, claro es que nada hace\apnota{32b}{M U facit.} al caso. Pues el Apóstol no habla del purgatorio, sino del juicio extremo. Y Cristo, por la noche —sea dicha con reverencia salvo\apnota{34}{M U salva tamen.} la de Crisóstomo— entiende la infidelidad, como es claro por el texto próximo: «mientras yo estoy\apnota{35}{M sunt, Mc a. R. 'sum'.} en el mundo, luz soy del mundo», queriendo decir que fuera de la fe de Cristo nadie puede\apnota{36}{M potest, Mc 'possit'.} obrar bien; sin embargo admito a Crisóstomo, porque mira al juicio extremo\apnota{37}{M indicii.}.

La autoridad de Sal 104 [v. 23]: «Saldrá el hombre a su obra hasta la tarde», se permite arrastrarla en sentido figurativo a la vida y muerte del hombre, pero en sentido genuino, y que luche en la contienda, habla de la admirable disposición de Dios, [332] que así dispuso al hombre: que por la mañana salga hasta\apnota{1}{U ita disposuerit quod homo exeat ad; usque fehlt in M.} la tarde del día natural a su obra\apnota{2}{M U operationes suas.}. Y Agustín\apnota{3}{P Augustinus.} rectamente siente, en el cap. 3 De fide ad Petrum, que el tiempo de alcanzar la vida\apnota{4a}{U viiam, Uv vitam.} eterna les ha sido dado a los hombres solo\apnota{4b}{U tamen.} en esta vida, porque, como dije, siempre miran a la vida futura\apnota{5}{M Lücke hinter ad, Mc 'praesentem (gestrichen) vitam futuram respiciunt'.}, no al purgatorio.

Y Jerónimo, sobre Gal 6 [v. 5]: «Cada uno llevará su carga», bien\apnota{6}{M Et quod Hieronymus, sonst = P; U Et Hieronymus super Galat. 6. Unusquisque portabit onus suum etc. et bene.} muestra que los hombres son ayudados por sus méritos en esta vida, pero, llamados ante el tribunal, lo que sucede en la muerte, entonces ni Job ni Noé\apnota{8}{nec Noe fehlt in U.} ni Daniel pueden interceder por nadie. Respondo que esta autoridad es robustísima\apnota{9}{M quod haec auctoritas est robustissima est (so!), U quod auctoritas est robustissima.} contra el egregio D. D., porque si es verdad que el hombre en la muerte es llamado ante el tribunal, y entonces ni Job ni Noé\apnota{11}{nec Noe fehlt in M U.} ni Daniel pueden interceder, en vano ora la Iglesia por los muertos. Y así se niega el purgatorio. Por lo cual Jerónimo habla del tribunal novísimo, después de terminado el purgatorio.

A Bernardo también lo adujo, en cuanto que en la muerte del hombre se celebra un juicio particular, lo cual no se daría si [el alma] tuviera todavía plazo para merecer. Respondo: sea lo que fuere del juicio particular, la consecuencia no vale.

Paso por alto la autoridad de Zac [4, v. 7]: «Igualará\apnota{17}{P adaequavit.} gracia con gracia», porque se expone en sentido tropológico y bueno, pero no propio.

Al fin de este artículo dice\apnota{19}{Mc fügt bei 'Eccius'.}: aunque estos testimonios no fuesen tan claros, esto, sin embargo, debería apartarme a mí, teólogo, a saber, que así\apnota{20}{M testimonia haec, hoc tamen me deberet abstinere (Mc 'absterrere') theologum, quod (sic fehlt, Mc 'si').} sin fundamento, a mi arbitrio, no habiendo sido yo ni juez ni asesor\apnota{21}{M fuerit, verbessert fuerim; U fuerim et tamen.}, habría prorrogado\apnota{22}{Mc 'prorogarem'.} a las almas que parten el término de merecer. Lo mismo devuelvo contra el D. D.\apnota{23a}{U ipsum D. D.}, no habiendo sido tampoco él ni juez ni asesor\apnota{23b}{M assessor nec iudex.}, y sin embargo define\apnota{24}{M diffiniat.} para las almas un estado cierto según sus opiniones, sin fundamento; sobre todo porque no quiere que esto sea solo opinión\apnota{25}{M opinionem esse.} (lo que toleraría gustosísimamente), sino ciencia cierta\apnota{26}{M Lücke, Mc 'tolleraret sed certam sententiam'.}.

A la hora segunda de la tarde continuará la disputa.

A la hora segunda de la tarde del mismo día.

El D. D. Martín continuó lo comenzado, como sigue\apnota{29}{M die Zeitbestimmung fehlt, U Hora secunda continua: disputatio per eundem d. Martinum, qui dixit se declaraturum recte et candide de scripturis sentire.}.

Y como mi sentencia no parece bastante entendida por el egregio D. D., en cuanto pueda me expondré\apnota{31}{M U expono.} más abiertamente: que las Escrituras y los santos padres antiguos en sus sentencias miran a la vida futura, en la que o se salvan\apnota{32}{M samnantur.} o se condenan las almas, sin pensar entretanto nada del purgatorio. Por eso aquellas muchas autoridades que niegan el estado de merecer después de esta vida no pertenecen al purgatorio. Que si aún no se entiende, fínjase que no hay purgatorio alguno, como tampoco estuvo su pensamiento en las mentes de ellos. Entonces todavía se sostendrán y serán verdaderas todas las autoridades aducidas: que antes de la muerte está la vida del merecer. Además, no hablo de los méritos de las almas en el purgatorio en el sentido de que obren algo, sino de que reciben gracia más amplia. Como entre todos consta [333] que la culpa no se perdona, ni siquiera la venial, sin aumento de la gracia, y el divino Gregorio escribe que en el purgatorio se perdonan las culpas, en el diálogo 4, y se recita en la dist. 25\apnota{2}{M dist. 24.}, cap. Qualis, así, sin embargo, afirmo estas cosas\apnota{3a}{U hoc.} de modo que las tengo por opinión,

más aún, para que lo confiese\apnota{3b}{M habeant imo ut fateor.}, por ignorancia. Creo que solo a Dios le es conocido el estado de las almas en el purgatorio.

Al tercer punto principal, sobre la ciencia de su salvación, donde adujo\apnota{5}{M adduxit.} primero Ap 5 [v. 3]\apnota{6}{M quintum capitulum, U quintum caput.}, que nadie fue hallado ni en el cielo ni en la tierra ni debajo de la tierra digno de abrir el libro, queriendo entender por «debajo de la tierra» el purgatorio, como en el que «alguna vez» también\apnota{8}{U et aliquando.} estuvieron santos varones, como Pascasio y Severino:

Digo que esta glosa no tiene autoridad; por eso con la misma facilidad se desprecia con que se prueba. Confieso, ciertamente, que los santos estuvieron en el purgatorio; más aún, que nadie está en el purgatorio\apnota{11}{M esse in purgatorio.} sino santo. Podría, pues, decir un porfiado que «debajo de la tierra» es lo mismo que el infierno o cualquier otra cosa, como que ni los demonios ni los hombres ni los ángeles pueden\apnota{13}{M quod quod nec daemones nec angeli nec homines possit (so!)} abrir el libro; igual que se dice que la triple máquina de las cosas honra a Dios, y según el Apóstol, tanto las celestes como\apnota{15a}{fehlt in P, in Uv gestrichen.} las terrestres como las infernales doblan la rodilla\apnota{15b}{M flectunt.}; si es así que también los demonios tiemblan; más aún, como es el modo en el mismo libro, «debajo de la tierra» significa también a los muertos\apnota{17}{U mortuus.}, como dice en otro cap. [20, v. 13]: «La tierra dio sus muertos, y el infierno dio sus muertos\apnota{18}{M infernus dedit . . . . et terra dedit.}».

En cuanto a lo que seguidamente aduce: todos cantaban\apnota{19}{M contabant, Mc cantabant.} en el cielo y en la tierra y debajo de la tierra, y que los desesperados no pueden cantar, respondo: nunca dije que las almas desesperen, sino que, siendo una misma la sentencia de la Iglesia de que es igual la pena de los infiernos y del purgatorio, dije que son semejantes\apnota{22}{U similes esse.} a los desesperados; del modo como se lee que también en esta vida algunos fueron tentados de desesperación, como en el Salterio\apnota{23}{M temptatos . . . ut psalterium.} en muchos lugares: «No apartes tu rostro de mí, y seré semejante\apnota{24}{M U assimilabor.} a los que descienden al lago [27, v. 9; 28, v. 1]». Además, cantar\apnota{25}{M cantare, Mc cantari.} al Señor no siempre es alegrarse y gozarse; más aún, el cántico nuevo\apnota{26}{U nonum.} es el cántico de la cruz, esto es, alabar y llevar a Dios en medio de las tribulaciones, y hasta en la muerte.

En cuanto a lo tercero, que adujo el canon de la misa, donde orando por los difuntos decimos: «Duermen en el sueño de la paz» y «descansan en Cristo», no ve, según dice, cómo puedan\apnota{30a}{M possunt.} decirse que duermen en el sueño de la paz, si están en tal turbación\apnota{30b}{fehlt in M.}, temblor e inquietud\apnota{31}{U inquitudine, Uv inquietudine} máxima. Por eso, interpretando este descanso de paz, dice que aguardan seguras el término de su purgación. Respondo: no me satisface esta glosa, y se prueba lo mismo por lo mismo. Pues que estén inquietos lo prueba lo que sigue en el canon:

«A ellos, Señor, y a todos los que descansan en Cristo, concédeles lugar de refrigerio, de luz y de paz», y aquello que todos oramos: «Dales, Señor, el descanso eterno». Asimismo: «Dales la paz», lo cual no puede entenderse como «dales aguardar seguras el término de su purgación», que es lo que el D. D. interpretó que es la paz. Más rectamente, pues, a mi juicio descansan en paz en cuanto al cuerpo. Pues dormir en las sagradas letras en paz significa descansar en el sepulcro. [334] Y así es patente que el egregio D. D. no probó la ciencia de su salvación; más aún, en cierto modo quita las penas del purgatorio, mientras\apnota{2}{M cum.} les atribuye paz.

El D. D. Eck.

Como nos ha sido fijado un tiempo estrecho para terminar hoy esta materia, de modo que por la calidad de la materia no puedo refutar las soluciones vanas\apnota{5}{M nihil non reputari.} ni tocar otros puntos del purgatorio, aduciré unas pocas cosas, para que fácilmente se reconozca que el reverendo padre ha evitado la médula de nuestra materia y ha buscado desvíos. De donde, primero, habiendo yo aducido a Agustín —que todo mérito aquí se adquiere\apnota{9}{M comparatum.}, y que nadie espere alcanzar ante Dios, una vez muerto, lo que aquí descuidó—, y asimismo a Jerónimo, dice que una y otra cosa concede. Pero ¿qué obstinación es entonces esta\apnota{11}{M est tunc; U obstinatia.}, si concede con Agustín que todo mérito aquí se adquiere, para decir que la gracia en las almas, aun en el purgatorio, aumenta?

Ni vale en este caso su colirio para todas las autoridades, de las que dice que no hablan del purgatorio, pues el bienaventurado Agustín dijo aquello mismo tratando de las almas en el purgatorio, tanto en el Enchiridion como en el libro De cura pro mortuis agenda, c. 1. Pues después de la muerte no podemos hacer lo que aproveche\apnota{16}{U prodest.}, sino recibir lo que hicimos. En esta dificultad había caído\apnota{17}{M inciderit.} Agustín: cómo podrían aprovechar a los muertos los sufragios, no mereciendo nada\apnota{18}{M possint quia nihil mererent, U possunt, quia nihil mererentur.}. Recurrió a esto: a que en vida merecieron lo que después de la muerte\apnota{19}{M pro mortem.} les aprovechase; por lo cual en vano se esforzaba en instruirme sobre el sentido\apnota{20}{M U de sensu suo.}, que bastante bien había yo comprendido. Ni su persuasioncilla tiene momento alguno ante un cristiano, porque las veniales no se perdonan sin aumento de la gracia. Pruébelo esto el reverendo padre. Con ninguna autoridad está esto cimentado; más aún, es falso: como el pecado venial no produce ofensa de Dios por misericordia de Dios, no es necesaria nueva gracia para su borradura, sino que basta que alguien padezca por él satisfacción.

En cuanto a lo que en segundo lugar se excusa del error, porque no veo aducida ni una letra por el D. padre a favor de su nueva doctrina: por eso con razón

debe reputarse\apnota{28a}{M suputari, Mc reputari.} sospechosa, siendo así que la sentencia de los teólogos más recientes\apnota{28b}{M neotericorum doctorum.} se apoya en tantas autoridades de la sagrada Escritura y de los santos padres. Ni le excusa esto: que a algunos imputa hacer\apnota{30}{P M facere se.} de las opiniones de Tomás o de Escoto como artículos de fe, cosa que a mí no me consta. Pues\apnota{31}{fehlt in U, in Uv nachgetragen.} yo no aduje a Tomás o a Escoto, sino a las clarísimas lumbreras de la Iglesia junto con la sagrada Escritura; vea él mismo a qué teologistas acusa\apnota{32}{M incusat.}.

Tercero, a Agustín opone, contra los libros de los Macabeos, con más fuerza al bienaventurado Jerónimo; pero Jerónimo en ninguna parte niega que los libros de los Macabeos estén en la Iglesia dentro del canon de la Biblia; es más, en el prólogo consta constantemente haber sido anotado entre las historias de los volúmenes divinos. Por lo cual no debió equipararse a los opúsculos de los santos padres en el can. Sancta\apnota{37}{fehlt in P M U.} Rom. 15 dist.

Y habiéndole yo opuesto lo del concilio Florentino, respondió\apnota{38}{U respondet.} que el concilio no puede [335] hacer que algo sea de la Escritura si no lo es. Esto, ciertamente, es verdad; pero ¿qué es esto: mancillar con tanta temeridad un concilio\apnota{2}{U consilium.} tan laudable, hasta decretar este absurdo? Siendo así que fueron varones doctísimos los de aquel concilio, prefiero creer al concilio, que es regido por el Espíritu Santo, que al señor Lutero; no porque el concilio haga que algo sea\apnota{5}{M concilium faciat esse aliquid, U consilium faciat aliquid (esse fehlt).} de la Escritura si no lo es, sino porque creo que el concilio\apnota{6}{M concilium melius, U consilium melius.} tiene el sentido y la inteligencia de las Escrituras al decretar que esto es de la Escritura, lo cual en la Escritura se halla.

Vanida\apnota{8}{M evanida, Mc evalida.} es su excusa de que siempre puede hallarse una exposición tal que el texto alegado no valga en la contienda contra los pertinaces, pues ese sería el resbaladero de todos los herejes, que siempre podrían aportar algún afeite de exposición con el que\apnota{11}{M U quibus.} contendieran que las verdades católicas no están expresas en la sagrada Escritura. Así hasta hoy duraría aún la pérfida herejía de Arrio: que el homousía\apnota{12}{M omusia.} no podría probarse tan expresamente por las sagradas letras que en la contienda no escapara\apnota{14}{M U non evaderet.} un pertinaz con cualquier afeite.

Igualmente bien sabemos que la Iglesia no puede hacer los evangelios; sin embargo, la Iglesia hace que, dejados los evangelios de Nicodemo, de Bartolomé, de Tomás y de otros, solo a cuatro prestemos fe indudable; en cuyo juicio de la Iglesia hay que estar en la elección de los evangelios\apnota{18}{M U standum est in evangeliorum acceptatione.}. Así también en la inteligencia y exposición de las sagradas Escrituras.

Además, asumiendo las autoridades por mí aducidas desde el principio, quiso que fuera innecesario lo hecho, porque también él creería que el purgatorio existe,

no recuerda que, frente a tantas\apnota{22}{M ex hoc.} autoridades que fijan en la vida presente el término de merecer, quiso zafarse porque el purgatorio no se probaría a partir de las Sagradas Escrituras; dicho lo cual, con razón debí yo refutarlo, no fuera que los picardos y otros cismáticos lo tomasen a su favor, ellos que, como el venerable padre, no saben que el purgatorio existe. Dejadas a un lado algunas, contemplemos al menos dos citas.

En Mt 5 [v. 25] no quiso entender por cárcel el purgatorio, porque también Agustín, lo cual en absoluto se me ocultaba, entiende por cárcel el infierno\apnota{28}{M U intelligit.}. Añado, más allá del venerable padre, que Crisóstomo entiende la vida presente, pero que es más recta la sentencia de Ambrosio, o al menos no despreciable, él que por cárcel entiende el purgatorio. Las palabras de Cristo esto indican: que aquél debe pagar hasta el último cuadrante\apnota{32}{P U dodrantem.}; pero en el infierno no hay pago alguno, como tampoco hay redención. Ni el subterfugio añadido puede rechazar\apnota{33}{U ad Ambrosii expositionem refellere.} la exposición de Ambrosio, a saber, que el que ha de ser arrojado a la cárcel pecó mortalmente, puesto que también por los pecados veniales y por los pecados mortales, aunque ya contritos, se sufre pena en el purgatorio, lo cual recibimos de los Macabeos, cuando dice: Santa y saludable es, pues, la resolución de orar por los muertos\apnota{37}{M U sancta (M sanctam) et salubris est ergo (ergo fehlt in M) cogitatio pro defunctis exorare (Mc 'orare pro defunctis', M Lücke).}, para que sean liberados de los pecados. Pues estos que habían sido muertos, y por quienes Judas Macabeo hizo\apnota{38}{M fecerat.} ofrendas, habían pecado mortalmente a causa de los despojos de los ídolos\apnota{39}{M spolia ablata.}, aunque se cree que se arrepintieron en [336] la misma matanza, según la glosa ordinaria allí mismo, y aquello del Sal [78. v. 34]: Cuando los hacía morir, lo buscaban.

En cuarto lugar, cuando hube aducido el clarísimo lugar del apóstol Pablo, 1 Cor 3 [v. 12ss.], que Agustín, Ambrosio, Jerónimo, Gregorio, Bernardo, Isidoro con la glosa ordinaria entienden del fuego del purgatorio, confiesa el reverendo padre que no tiene el sentido genuino de este lugar; por eso, despreciados tantos padres, aportó una sentencia nueva: que el apóstol habla del fuego de la conflagración, porque el apóstol hizo mención del día del Señor, y de que el fuego probará la obra de cada uno; yo, que siempre oí que debe creerse sobre todo a los santos más antiguos, abrazo con veneración la sentencia de los santos padres y no acepto aquella glosa nueva ni las palabras del D. Lutero, salvo que lo pruebe con autoridad de la Sagrada Escritura. Ni le ayudan las consideraciones que pondera\apnota{12}{M ponde dat, Mc ponderat.} en esta ocasión, que también los santos padres leyeron bien. Pues aunque peculiarmente el día del juicio final se diga día del Señor, como también recuerda Bernardo, sin embargo en la muerte de cada hombre, cuando se dicta el juicio\apnota{14}{U fit.}, puede decirse día del Señor, y no acepta una distribución tan exacta y capciosa, como suelen hacer los sofistas\apnota{16}{M solem, Mc solent.}, en la palabrita «de cada uno». Pero sobre esto, reciba a quienes sobreedificaron paja, madera y heno\apnota{17}{M venum.} sobre el fundamento, como también Agustín complica la distribución de Jn 1 [v. 9]: Ilumina\apnota{18}{P illuminet.} a todo hombre que viene a este mundo. De donde es invicta la sentencia apostólica: que quienes sobre el fundamento edifican madera, heno\apnota{20}{M venum.}, paja, se salvan, pero por el fuego de la purificación, de modo que ningún mal quede impune.

A aquello del Eclesiastés\apnota{22}{P U Ecclesiastes.} sobre la caída del leño responde: Si habla del purgatorio, así se probaría que el alma permanece perpetuamente en el purgatorio. Digo que sintió óptimamente la glosa: que por leño se entiende el hombre, como\apnota{24}{U hinc Uv sicut.} también en Ez 31 [v. 8] significa la criatura racional: Todo leño del paraíso etc.\apnota{25}{fehlt in M.}, y por caída la muerte; y sin embargo no se sigue que por eso permanezca perpetuamente en el purgatorio, sino que, como expone egregiamente el bienaventurado Jerónimo, por el austro se significan el bien, y por el aquilón el mal y la condenación\apnota{28}{M damnationem significari, U damnationem significare.}. Por eso el sabio no entendió aquellos lugares, sino el estado: que en la caída o es bueno, y así perseverará sin aumento de gracia por la supresión de los pecados, o si es malo, perdurará en la malicia\apnota{30}{U malignus statt in malitia.}.

Sobre la mansión en Jn [14 v. 2] dijo que, inmediatamente después de la muerte, se asigna al alma una mansión cierta. ¿Cómo, pues, será cierta, si se produce un aumento de gracia mayor? Pues entonces necesariamente sobrevendrá\apnota{33}{M non tunc necessario accedat.} una mansión más alta; también sería consecuencia que el que descendiera al purgatorio con muchos veniales estaría mejor que el que descendiera con pocos o el que muriera sin ninguno, porque se le haría mayor aumento de gracia para la supresión de los muchos veniales. Añado que sería nocivo y condenable orar por los muertos, y más les aprovecharía si perdurasen en las penas del purgatorio por el aumento de gracia mayor, como en algún lugar escribió el reverendo padre: que el alma liberada por sufragios ajenos\apnota{39}{M suffragii.} es menos beatificada que si por sí misma padeciese lo suficiente en el purgatorio. Lo cual es no sólo contra la Sagrada [337] Escritura: Santa, pues, y saludable es\apnota{1}{fehlt in U.} la resolución etc., sino también contra la observancia y la piedad de toda la Iglesia, que debemos tributar a los muertos, como el bienaventurado Agustín expone bellísimamente en el libro sobre el cuidado que ha de tenerse de los muertos\apnota{3}{M habendam.} y en el libro Enquiridión.

A lo de Damasceno respondió además que no habla del purgatorio, puesto que así perdurarían perpetuamente en el purgatorio. Y yo siento que ni Damasceno ni las demás autoridades hablan del purgatorio, pues de otro modo me dañarían y aprovecharían al venerable padre. Pero como hablan de la vida presente, en la que ponen\apnota{9a}{M ponant.} el término de merecer y desmerecer, no puede extender\apnota{9b}{U porrigi.} el tiempo de merecer más allá del término establecido, esto es, la muerte. Por eso no acepto otra exposición que la que el propio autor quiera, el cual habla en el lugar alegado del término de merecer.

Así también puedo decir de los demás, y no puedo ser increpado con semejante invectiva como el reverendo padre, porque\apnota{14}{M quo.} con la autoridad de los santos padres y de la Sagrada Escritura cierro el curso y el término de merecer con la muerte. Él, en cambio, sin apoyo alguno\apnota{15}{M finiant .... nullo.} de la autoridad de la Sagrada Escritura, les extiende y prorroga el término, contra Agustín y contra Jerónimo, aunque sea propio de su modestia decir que esta sentencia suya es su opinión; pero entonces no debió llamar error a lo opuesto.

En cuanto a la autoridad de Jerónimo\apnota{19}{M U De auctoritate Hieron.}, intentó volverla contra mí: que ni Noé ni Job ni Daniel interceden por nadie después de la muerte, porque allí se habla del juicio final, lo cual\apnota{21}{M U quod tamen.} no se prueba por la letra. Pues no sólo en el juicio final, sino también inmediatamente\apnota{22}{M etiam statim.} después de la muerte cada uno llevará su carga.

Finalmente, pasando a otra partícula, respondió a aquello de Ap [5 v. 3ss.] que por debajo de la tierra se entiende\apnota{25}{M U sub terram intelligatur.} el infierno. Pero entonces lo siguiente sobre el cántico no puede sostenerse, porque los demonios y los condenados no cantan, sino que blasfeman y aúllan\apnota{26}{M plasfemant et egulant.}. Por eso por debajo de la tierra\apnota{27}{M U sub terram.} necesariamente puede entenderse el purgatorio y no los infiernos. Aunque no niego que en\apnota{28}{M quin.} la Sagrada Escritura también se haga mención de los demonios en los infiernos, como en el apóstol: de los celestes, terrestres e infernales\apnota{29a}{M in inferis ut coelestia etc., U in inferis in apostolo coelestia terrestria, Uv = P.}. Pero esto nada tiene que ver con el asunto\apnota{29b}{M romphum.}.

Aportó además una glosa sobre aquella palabra «canta», según la cual sería llevar la cruz y alabar a Dios en las tribulaciones; glosa que, como forjada por el venerable padre en este lugar, desprecio con la misma\apnota{32}{P M U ea.} facilidad. Más aún, Juan no tolera ser expuesto así, él\apnota{33a}{P M U quia.} que dice: Cantaban en el cielo, no en la tierra\apnota{33b}{non fehlt in P, M cantabo in coelo et in terra; U cantabunt (Uv cantabat) in coelo in terra, wie oben zuerst bei Walch, 1223.}; pero en el cielo no llevan la cruz ni alaban en las tribulaciones, porque enjugó Dios toda\apnota{34}{M U laudant deum quoniam abstersit (U abstergit) deus omnem (U sic omnem).} lágrima de los ojos de los santos y ya no hay trabajo etc.\apnota{35}{M ab oculis eorum etc. (Die Worte von et iam an fehlen), U ab oculis eorum iam non est nullus labor et dolor etc.}.

Pero aun si se concediera su exposición, como si\apnota{36}{M U et.} fuera verdadera, aunque repugna mucho a la letra, todavía no quebrantaría este propósito, pues los que alaban a Dios en las tribulaciones después de la muerte tendrían certísimo signo de la futura bienaventuranza, mientras que los condenados, por el contrario, con odio consumado de Dios prorrumpirán\apnota{39}{U prorumpunt.} en perpetua blasfemia del Creador. [338] Por eso las almas que han de purificarse\apnota{1}{fehlt in P.}, seguras de su salvación, cantan y alaban a Dios, aguardando el término de la purificación.

Pero cuando hube aducido el canon de la misa, en el que se afirma que ellos duermen en el sueño de la paz, no tolera mi inferencia: si en el sueño de la paz, luego están seguros\apnota{4}{U securae.} de su salvación y no en horror, temblor, pavor y casi desesperación. Aportó una doble razón\apnota{6a}{U adduxi.}: primero, porque sigue que les concedas un lugar de refrigerio, de luz\apnota{6b}{P et lucis.} y de paz\apnota{7}{M indulgeas etc.}; segundo, porque frecuentemente oramos: Dales, Señor, el descanso eterno, y: Dales la paz. Por eso dio otra glosa: que duermen en el sueño\apnota{8}{U somnum.} de la paz, esto es, en cuanto al cuerpo. No me contenta\apnota{9}{M contento.} ni la objeción ni su nueva glosa\apnota{10}{M glossam, Mc glossa.}, que no tiene apariencia alguna ni fundamento alguno. Pues dice el canon: Acuérdate de ellos etc., que duermen en el sueño de la paz, y más abajo\apnota{11}{U eorum, qui dormiunt (in somno pacis fehlt); et infra fehlt in M.}: Y de todos los que descansan en Cristo. No dice: De los que descansan en el sepulcro. Además, ¿qué importa al alma que el cuerpo tenga paz o no, ya sea arrojado al mar o\apnota{13}{M sive mari iactetur ac, U sive a mari iactetur sive.} lacerado por las fieras, como atestigua el bienaventurado\apnota{14}{fehlt in M, U d. statt b.} Agustín? ¿Qué hombre soporta que el cuerpo exánime y exangüe duerma en paz? De donde no admite de buen grado las distinciones en los teólogos más recientes, y sin embargo aquí distingue egregiamente entre la paz del cuerpo y la del alma.

Por eso ha de abrazarse la sentencia más verdadera\apnota{17}{M esse.}: que las almas que han de purificarse duermen en paz y descansan en Cristo. De donde aquel descanso y aquella paz excluyen la turbación del alma, el horror y la casi desesperación, pues estarían inquietísimas\apnota{19}{U sunt.} las almas que fueran agitadas por esta turbación de la mente, hasta casi desesperar.

A lo aducido, en cambio, respondemos sin mucho gasto: Que oramos y suplicamos a Dios que les dé la paz, no la paz que\apnota{22}{M quem.} ya tienen ahora, esto es, la seguridad de su salvación, sino la paz eterna, puesto que la paz es, como atestigua Bernardo, tanto respecto de la culpa como de la miseria. Dé, pues, Dios a las almas en el purgatorio la paz frente a las penas, como les confirió la gracia contra el pecado\apnota{25}{U peccata.}.

Por eso sigue siendo verdadera e inconmovible la verdad: que las almas en el purgatorio no casi desesperan de su salvación.

El D. D. Martín.

Contra las soluciones replica el egregio D. D. desde Agustín, que dice\apnota{29}{M dom. egreg. doctor ex Aug. dicentem.}: Todo mérito se adquiere aquí\apnota{30}{M comparatum.}, y tratando del purgatorio, que después de la muerte sólo recibamos lo que hicimos. Respondo brevísimamente: Luego o no hay purgatorio, o el purgatorio es\apnota{32}{M U erit.} lo que se recibe por la vida pasada. Por eso es patente que Agustín no puede entenderse del purgatorio. Pues no hicieron esto en la vida pasada, para recibir el purgatorio. Aunque también esta autoridad me patrocina, en que el mérito por el que merecen ser ayudadas\apnota{35}{M iuventur iuvari.} lo merecieron aquí. Yo, en cambio, confieso que no puedo entender cómo son ayudadas las almas en el purgatorio y liberadas sin don alguno de gracia, por la sola quitación de las penas: entiéndalo quien pueda.

Luego dice que es falso\apnota{38a}{M falum, Mc falsum.} que los veniales no\apnota{38b}{M sine, Mc absque.} se remitan sin aumento de gracia.

Cuando, pues, quiere el D. D. que los veniales puedan\apnota{39}{M U possunt; quod venialia fehlt in M.} remitirse sin aumento de gracia, no lo concedo hasta que lo pruebe. Más aún, es contra el texto expreso\apnota{40}{fehlt in M.}, en Rom 7 [v. 14ss.]\apnota{41}{P Rom. V.}, Gal 5 [v. 17], donde el apóstol dice: Vendido estoy bajo el pecado, ¿quién me librará de este cuerpo de muerte\apnota{42}{M de corpore mundi, U de morte corporis.}? Responde: La gracia por Jesucristo [339]\apnota{1a}{M respondit: gratia dei Jesu Christi, U = P, nur gratia dei.}. Es cierto, en cambio, que el apóstol no estaba en pecado mortal\apnota{1b}{M esse, verbessert fuisse, M U in peccatis mortalibus.}. Lo demás lo dejo a la futura disputa\apnota{2}{M ad disputationem futuram, U futurae disputationi.} entre los egregios señores doctores Andrés y Juan Eck\apnota{3}{M Johannem et Andream.}. Y es erróneo esto: que el pecado venial no ofenda a Dios, pues a Dios desagrada toda inmundicia y reprueba al que quebrantare los mínimos, Mt 5 [v. 18s.]\apnota{5}{Matth. 5. fehlt in M.}. Ni probó el D. D. que baste con que padezcan lo suficiente.

En segundo lugar me acusa de no haber aducido ni una letra en favor de mi sentencia, cuando él en favor de la sentencia de los teólogos más recientes\apnota{7}{M recensiore; theologorum fehlt in M U.} adujo tantas autoridades de la Escritura y de los padres. Respondo: Tanto más sospechosa me es\apnota{8}{fehlt in U.} su sentencia cuanto más adujo, porque los trajo por los cabellos y retorciendo\apnota{9}{M optorto.} el cuello, como bastante dije en lo precedente. Menos peca quien duda en sus pensamientos que quien intenta establecer\apnota{11}{M statuere, verbessert constituere (construere?)} sus dudas con palabras divinas.

Dice también que no le consta quiénes establecieron como artículos de fe las opiniones de Tomás o de otros, porque adujo lumbreras de la Iglesia junto con la Escritura. Advierte, pues, que mire a quiénes llamo\apnota{14}{M U ut statt quod; M appellam.} teologistas. Dije hoy y de nuevo digo: Permito que las opiniones sean opiniones y a nadie imputo esto como vicio, pero arrastrar la Escritura que se resiste al sentido ajeno para afianzar\apnota{16}{M U in alienum sensum pro formandis (U confirmandis).} opiniones y permanecer pertinazmente en esto, esto digo que es oficio de teologistas.

\noindent\textit{[nota editorial sin llamada en el original, al inicio del aparato: coepit d. Martinus latine prosequi. W. A. 299, 3–11 fehlen in M, a. R. [rot] fecit sermonem vulgarem.]}

\noindent\textit{Notas del aparato sin llamada en el cuerpo:} [15b] M cephale, U Kephale, a. R. κεφαλη.; [15] M dici.

En cuanto a lo otro, sobre el canon de los libros, donde, apoyado en la autoridad de Jerónimo y del concilio de Florencia\apnota{19}{fehlt in M.}, prefiere creer al concilio, que es regido por el Espíritu Santo, antes que a mí, y\apnota{20}{W. A setzt dafür ei.} le doy las gracias. Pues sabe piadosamente, nunca quise que se me creyera. Pero respondo brevemente: concilie él primero\apnota{21a}{U primum.} a Jerónimo consigo, quien\apnota{21b}{M quia.} en el prólogo galeato\apnota{22a}{P Uv allegato.} enumera expresamente los libros de los Macabeos y algunos otros entre los apócrifos\apnota{22b}{M apographa.}, por cuya autoridad resulta que el libro de los Macabeos me es grato y aprobado, pero queda abierto a los contenciosos para\apnota{24}{M pati ad. Mc schreibt über 'fateor'.} el rechazo. Bastante he dicho hoy sobre el concilio, más aún, como dice el propio D. D.: el concilio no yerra, pero si yerra, no es concilio, y para hablar según mi sentir, creo que el concilio y la Iglesia nunca yerran en lo que es de fe, en lo demás no es necesario que no yerren.

Y en cuanto a que exclama que este es el escondrijo de los herejes, quienes con esta confianza refutarían o aceptarían\apnota{29}{M quemlibet librum negarent, gestrichen vor quaslibet (vom Disputator widerrufen); M U vel reciperent, Mc a. R. vel acceptarent.} cualesquiera exposiciones de las autoridades, respondo: ¿por qué no se quejaron\apnota{30a}{M non sunt conquesti hoc.} de esto san Agustín, Jerónimo y los demás triunfadores\apnota{30b}{M victoriosissimi.} victoriosísimos sobre los herejes, sino que, dejados de lado los lugares ambiguos, se esforzaron en combatir con lugares ciertos y claros? Hagamos esto también acerca de las almas en el purgatorio.

A la cuestión y a las autoridades\apnota{33}{et auctoritates fehlt in M U.}.

A las soluciones de las autoridades\apnota{34}{U ad auctoritatum solutiones.} dijo que por mí han sido despreciados Agustín, Jerónimo, Ambrosio, Bernardo, Gregorio, Isidoro\apnota{35}{M Ambrosium Hieronymum; Isidorum fehlt; U Aug. Ambros. Bernh. Hieron. Gregor. Isid.} con la glosa ordinaria\apnota{36}{M ordinariis.}. Esto lo dijo por modestia, o más bien por molestia eckiana, demasiado deseoso de mover envidia\apnota{37}{M invidii, Mc invidiae.}. Dije hoy que no han sido despreciados por mí, y de nuevo lo digo.

Que no acepte mi glosa, no me importa. Pruebe él, en cambio, también la suya: siendo el texto claro acerca del día del Señor y del fuego, en\apnota{39}{fehlt in M.} el que será revelado el día del Señor. Lo cual, aunque puede aplicarse al purgatorio, como dije y así no reprobé la sentencia del señor doctor [340], a saber, que el purgatorio existe, sin embargo con esta inteligencia no puede taparse la boca al contencioso.

Asimismo tampoco desprecié aquello de Mt 5 [v. 25] sobre la cárcel, según la sentencia de Ambrosio. Pero\apnota{4}{fehlt in U.} lo que añade el D. D., que las palabras de Cristo indican que debe pagarse hasta el último cuadrante y que por eso no debe entenderse del infierno: esto\apnota{5}{M posse, übergeschrieben debere, beides gestrichen, daneben deberi; M U respondeo hoc.} queda suficientemente diluido por el divino Jerónimo, quien contra

a Helvidio\apnota{6}{P Uv Helpidium.} expone rectamente esta dicción «hasta que», que\apnota{7}{M U ut.} no está obligada a significar lo que Ambrosio quiso aquí con el D. D. Pues así también en Mt 1 [v. 25] se escribe\apnota{8}{U hic ponit nam et sic . . ., M U scribitur de Joseph.}: Y no la conoció, hasta que dio a luz a su hijo, y sin embargo no se sigue que después del parto la conociera; así aquí no se sigue que\apnota{10}{M U quod.} haya de salir después del pago, sino que pagará y no saldrá.

Y en cuanto a que refutó que no pecó mortalmente quien disiente del adversario\apnota{12}{M dissentio ab adversario.}, y que también los mortales, con tal\apnota{13a}{U non.} que sean contritos, y los veniales se purgan\apnota{13b}{M U purgantur.} en el purgatorio: digo que este texto no habla de mortales contritos, pues un mortal contrito ya\apnota{15}{M nam, Mc iam.} no es mortal, y el penitente ya no es quien disiente de su adversario\apnota{16a}{M ab adversario.}, luego Cristo habla\apnota{16b}{M deus.} del que disiente y así permanece.

Paso por alto aquello del caso del leño en el Eclesiastés, donde aduce a Jerónimo\apnota{17}{M U dicit.} entendiendo por el austro el estado bueno y por el aquilón el estado malo; concedo ambas cosas, y nada para el purgatorio.

Sobre la mansión cierta en Juan, donde arguye de este modo: ¿cómo se hará cierta, si se produce un aumento de gracia, al que se deba una mansión más alta?; lo mismo puede argüirse de cualquier fiel después del bautismo\apnota{22}{M baptisma.}, para quien hay una mansión cierta predestinada desde la eternidad.

Además, los que están en el purgatorio están ordenados de tal modo que de este modo lleguen a una mansión cierta. También aquel argumento humano nada concluye, a saber, que se seguiría que el que desciende con muchos pecados veniales\apnota{25}{fehlt in M U.} estaría mejor que el que desciende con pocos. Si esta razón es buena, hay que temer que la meretriz llegue a ser mejor\apnota{27}{P fieret.} que la bienaventurada Virgen, como si no hubiera distintos grados de las almas en el purgatorio, cuando él mismo hoy afirmó que los santos están en el purgatorio como más excelentes que los demás.

Después aduce el argumento silvestrino, que les aprovecharía si perdurasen en las penas, pues así también convendría a los mártires morir y padecer hasta el día del juicio, como si ignorase el D. D. que para las almas están constituidas las penas hasta una medida cierta. Pues no\apnota{33}{M U ergo.} es condenable orar por los muertos, como tampoco fue condenable que el apóstol orase por sí y pidiese que se orase, aunque su virtud crecía en la enfermedad. Así cualquier fiel debe orar por cualquier necesidad de cualquier fiel, debe ayudar, sin que obste que aquel por esa necesidad merezca más y más.

Por lo mismo quiero que quede suelto lo que contra la observancia de la Iglesia y\apnota{38}{M et ecclesiae} la piedad sería, a saber, que dije que el alma liberada por sufragios ajenos es menos bienaventurada que si por [341] sí misma padeciese\apnota{1}{fehlt in M.} bastante en el purgatorio. Esto lo tomó como a modo de corolario de mis dichos, donde yo sentí que debe socorrerse a las almas, por mucho que fuese género más perfecto de mérito satisfacer en todo a la justicia divina.

Al Damasceno lo paso por alto con el término de mérito y lo remito a lo precedente.

Asimismo, de nuevo me reprende porque sin estar apoyado en ningún testimonio prorrogo el término del mérito, mientras él se apoya en muchos y lo niega. Respondo: tampoco él se apoya en ningunos sino entendidos con violencia\apnota{7}{M intellexisset.}, como bastante se ha dicho.

Queriendo también cogerme en mis palabras, dice que llamé opinión a mi sentencia, luego mal habría llamado error a lo opuesto\apnota{9}{M U appellaverim errorem.}. Digo como antes. No sólo llamo opinión, sino también ignorancia a mi sentencia; llamé error no a la opinión contraria, sino a que establezcan\apnota{12}{M statuunt proferunt.} la opinión como verdad.

Después dice que la autoridad de Jerónimo sobre Noé, Job y Daniel\apnota{13a}{Hieronymi fehlt in M U, de fehlt in M; M U Daniel.} vale no\apnota{13b}{U non valere.} sólo para el juicio extremo, y que yo no puedo demostrar esto por la letra. Lo dejo al juicio del que sienta mejor, pues el texto\apnota{15a}{U textus manifeste.} tiene manifiestamente que ni Noé ni Daniel ni Job\apnota{16}{M raptum statt rapto, sonst = P; U rapto nec Job nec Daniel.} oran por el hombre arrebatado al tribunal\apnota{15b}{M tribunal.}, que se produce en la muerte, lo cual debe entenderse en absoluto del juicio extremo, o al menos no del purgatorio, porque por los muertos en el purgatorio oran Noé, Job, Daniel y toda la Iglesia\apnota{18}{M U Daniel et (M etc. et) omnes ecclesiae.}.

A la autoridad del Ap responde rectamente que los demonios bajo tierra\apnota{19}{M ad tertium (?), Mc sub terra, U in inferno, Uv = P.} «no cantan, sino que blasfeman etc.»\apnota{20}{mit sub terra non cantant beginnt in M der Passus, der aus dem Notariatsprotokoll abgeschrieben ist, am Rande Mc [rot] verba notarium iuratorum. M liess eine Seite dafür frei; da sie nicht reichte, sind 4 Blätter (Signatur post o ante p) später eingeheftet. U non laudent (Uv non cantant), sed blasphemant deum.}.

El señor D. Martín, excluido por el tiempo, cuando ya habían transcurrido dos horas, se ofreció a entregar por escrito a los notarios lo restante y a exhibirlo primero al señor doctor Eck. Y ofreció\apnota{22}{M autem sequenti luce.} lo infrascrito por medio de su notario\apnota{23}{in U lautet die Zwischenbemerkung: Reliqua exclusus tempore doctor Martinus vult signare in scheda et ostendere domino doctori et dominis notariis. Obtulit autem sequenti loco (Lö sequenti luce) infra scripta.}.

Que los demonios en el infierno no alaban, sino que blasfeman a Dios, lo sé, pero que por eso «debajo de la tierra» signifique el purgatorio, en el que canten las almas, no se sigue. Primero, porque están en tribulaciones y penas y aún no ha enjugado Dios toda lágrima de los ojos de ellas\apnota{27}{M U eorum.}. El D. D., en cambio, niega que cantar sea en las tribulaciones alabar a Dios. Por eso

habla contra sí mismo, al poner a la vez en el purgatorio bajo tierra\apnota{30}{M subtus terram (U motus terra, Uv = P) in purgatorio subtus terram.} la alabanza y las penas, en las que no alabarían. Niego además también yo que las almas canten por eso, porque tienen un signo certísimo de bienaventuranza. Pues este signo certísimo habría que probarlo. Además, yo no dije que «bajo la tierra» signifique sólo el infierno, sino también otras cosas, cualesquiera que están bajo la tierra, en todas las cuales abunda la alabanza de Dios, pues todas están llenas de la gloria y alabanza de Dios, que es laudable en todas sus obras. De otro modo conviene que el señor doctor no sólo entienda el purgatorio por «debajo de la tierra», sino también alguna cuarta cosa por «debajo del mar», pues en el mismo Ap 5 también se recita «debajo del mar» y todas las cosas que hay en él. Más rectamente, pues, entendemos significadas por este cuaternario todas las criaturas que lo que el señor doctor\apnota{38}{fehlt in M U.} Eck imagina.

[342] Que\apnota{1}{hier beginnt die eingeheftete Lage von 4 Blättern, s. o.; Me (rot) quia per schedam sunt haec notariis praesentata.} alabar a Dios en las tribulaciones sea cantar, pruebo que no es fingido por mí, como dice el D. D., porque en Rom 5 [v. 3] dice: Nos gloriamos\apnota{2}{M gloriamini.} en las tribulaciones, y en Sal 42 [v. 9]: De día mandó el Señor su misericordia y de noche (esto es, en la tribulación), su cántico, y en Sant 5 [v. 13]: Si alguno está triste, ore con ánimo sereno y salmodie; asimismo: Bendeciré al Señor\apnota{5}{P domino.} en todo tiempo [Sal 34 v. 2]. En todo, dice, tiempo, también adverso. En cambio, es vituperado aquel de quien dice Sal [49 v. 19]: Te confesará, cuando le hicieres bien. Confieso, sin embargo, que esto no lo dije de los condenados bajo tierra, sino para excluir la sentencia del señor doctor, que quiere que cantar sea sólo de la seguridad de la bienaventuranza en el purgatorio.

Sobre el canon de la misa no está contento con que referí el dormir en paz al cuerpo, porque dice el canon: A los que descansan en Cristo, no a los que descansan en el sepulcro. Además, ¿qué importa para el alma que el cuerpo se encuentre en el mar o en el aire y no descanse? Asimismo, ¿quién (dice) llamaría al cuerpo exánime dormir en paz? Respondo: nada vale esta argucia, pues los que descansan en el sepulcro o en cualquier lugar, agitados según el cuerpo, descansan verdaderamente en Cristo, esto es, como en Ap 14 [v. 13]: descansan\apnota{16}{U requiescant.} de sus trabajos, a saber, porque el alma separada del cuerpo ya no trabaja en el cuerpo entre las diversas molestias de este mundo. De otro modo, ¿cómo negará que el cuerpo de Cristo en el sepulcro duerme en paz y descansa\apnota{18}{M U quiescere.} en Dios, cuando en tantos lugares así habla\apnota{19}{M loquitur.} de él la Escritura y canta la Iglesia: En paz en esto mismo dormiré y descansaré [Sal 4 v. 9]; y Sal 16 [v. 9]: Y mi carne descansará en esperanza? Asimismo Ap 14 [v. 13]: Bienaventurados los que mueren en el Señor. Asimismo Esteban en Hch 7 [v. 59]\apnota{21}{P M U act 8.}: Durmió en el Señor, lo cual sin duda se dice de la muerte corporal por tropo de la Escritura. Si a Eck

no le duerme el cuerpo exánime, deje, por favor, que duerma para Dios que lo resucita y para la autoridad de la Escritura\apnota{24}{M U suscitaturo deo: M auctoritatem.}.

Al final dice que oramos\apnota{25}{M orare.} por la paz, no la que tienen, sino la que les toque tener una vez purgados de la miseria de las penas. Respondo: esta\apnota{26}{M U hoc.} es una petición de principio, porque aquella paz de seguridad, que tendrían, aún no está probada. Queda, pues, que están en pena, y que oremos para que sean librados de ella y tengan paz.

Todo esto quiero haberlo dicho así para mostrar que ignoro lo que Dios hace en el purgatorio, y que se atreven demasiado\apnota{30}{U nimis, Uv et nimis.} los que, impacientes de esta ignorancia, prefieren fingir lo suyo antes que confesar que ellos\apnota{31}{U se.} lo ignoran. Lo mío no puedo probarlo, pero ni aquellos lo suyo. Por lo cual rectamente mi conclusión sólo quiso que no está probado ni por las Escrituras ni por razones que las almas estén ciertas al menos todas de su salvación y que la gracia no aumente en ellas; baste\apnota{34}{M sufficimus, U sufficit.} que sabemos que padecen y que debemos socorrerlas, lo demás ha de dejarse sólo a Dios.

Esto se actuó en presencia de los honorables varones Juan Teuber de Weysmann y Bartolomé Schaller de Ernfriedsdorf, de las diócesis de Bamberg y de Meissen, testigos laicos especialmente requeridos y llamados para lo premiso. El domingo diez del mes de julio, en el lugar acostumbrado de la disputa, a la hora del mediodía o casi\apnota{36}{M Wassmann statt Weissmann, ad praemissis statt ad praemissa, am Ende fehlt vel quasi: U hat statt der ganzen Angabe nur Acta sunt haec 10. Julii praesentibus etc.}.

El D. D. Eck ofreció lo siguiente en cierta cédula el lunes once del mes de julio por la mañana a la hora séptima\apnota{37}{U Eccius per schedam respondit 11. Julii. Löschers Anm. p. 436 „in M Sto additur hora prima" beruht auf Irrtum.}.

Las soluciones del reverendo padre me parecen insuficientes, sobre todo porque dice que Agustín no habla del purgatorio, cuando trata máximamente esto\apnota{39}{U id.}, a saber, [343] del purgatorio. Ni ha de negarse que las almas de los difuntos son aliviadas por la piedad de los suyos que viven, cuando por ellos se ofrece el sacrificio del mediador o se hacen limosnas en la Iglesia, pero esto aprovecha a los que, cuando vivían, merecieron\apnota{4}{U possint.} que esto les aprovechase después. Pues hay un modo de vivir ni tan bueno que no requiera aquello\apnota{5}{U ista.} después de la muerte, ni tan malo que no le aprovechen después de la muerte. Por lo cual aquí se compara todo mérito por el que alguien pueda después de esta vida ser aliviado o agravado. Ve ahora, Lutero, y di que Agustín no habló aquí del purgatorio.

Además no prueba que para la supresión del pecado venial se requiera aumento de la gracia. Confieso que el aumento de la gracia quita alguna vez lo venial\apnota{10}{U venialia.}, pero también lo venial se quita por la satisfacción. Por eso el apóstol alegado en Rom 7 [v. 14 ss.] no habla del pecado venial, y además si hablase, no excluiría la satisfacción. Pero sobre su consentimiento, como impertinente, nada digo de presente. Lo venial, confieso, ofende, pero venialmente, porque no constituye al hombre en odio de Dios ni lo hace enemigo de Dios, lo cual por lo demás con la gracia no\apnota{15a}{fehlt in U, in Uv nachgetragen.} estaría, por eso en el purgatorio tienen\apnota{15b}{M U habent quidem.} ciertamente pecados veniales sin enemistad con Dios.

Pero en cuanto a que se queja de que tuerzo las autoridades por los cabellos hacia el propósito, perezca yo si en toda la disputa adujo siquiera una autoridad tan pertinente como yo aduje aquí. Pero esto júzguenlo los jueces, no Lutero.

Se indigna de por qué tampoco Agustín y Jerónimo se quejaron de este escondrijo de herejes que dije; buen Dios, ¡si\apnota{21}{W. A. verbessert quasi.} no clama en todas partes Agustín\apnota{22}{M U Augustinus et Hieronymus.} contra los pelagianos, Cipriano y Ambrosio contra los arrianos y novacianos, que interpretan a su arbitrio la sagrada Escritura!

Pide que pruebe mi sentencia, cuando no puede probar la suya; pero la mía, ya probada, la dejo al juicio de los jueces.

Sobre el «hasta que» otra es la razón de Jerónimo a Helvidio, porque por otra parte se prueba por la Escritura que María permaneció siempre virgen, por tanto el «hasta que» allí no dice consumación, lo que aquí no puede hacer.

Además de mis razones por brevedad de gracia no me detengo. Pero en cuanto a que dijo\apnota{29}{U dicit.} que yo le hice cierto por mis dichos de que las almas liberadas por sufragios de los vivos son menos bienaventuradas, me imputa mal, pues no\apnota{31a}{M nisi.} me finjo sueños. Refiere\apnota{31b}{U reperietar, Uv refert.}, en efecto, en la dilución de mis anotaciones, que Pascual no quiso que se le disminuyesen sus premios, por eso prefirió arder, pero dejo esto y lo encomiendo a los jueces.

Además, no niego que en las tribulaciones alguien alabe a Dios\apnota{34}{M eum}, pero en este lugar no puede aceptarse, porque también dice que cantan en el cielo, cuando de ellos haya enjugado\apnota{36a}{M et coelo cantare eos dicit, ab illis abstergcrit (so!), U = P, nur et statt cum.} Dios toda lágrima\apnota{36b}{U lacrimam etc.}. De donde consta al lector cómo invierte mi sentencia\apnota{37}{M suam, verbessert meam.}. Por eso las almas en el purgatorio cantan, según Juan, a Dios, y pregunto, explique la causa de la incertidumbre\apnota{38}{M U nisi certitudinis.} de la gloria; por lo cual falsísimamente finge\apnota{39}{M comminiscatur.} que hablo contra mí mismo. Por eso no niego que los buenos se gloríen en las tribulaciones y canten, sino que negué aquella glosa que él mismo fingió, que [344] cantar sea en las tribulaciones alabar a Dios, de otro modo los ángeles no cantarían a Dios ni los bienaventurados.

Que intenta defender un error patentísimo, a saber, que las almas duermen con sueño de paz en cuanto\apnota{4}{U scilicet quoad.} al cuerpo etc., porque el alma separada del cuerpo

no trabaja\apnota{5a}{M U laborat.} en el cuerpo entre varias molestias: mira la astucia, porque\apnota{5b}{M vabriciem; M U qui.} habiendo referido al cuerpo, al resolver, lo que descansa en el sepulcro\apnota{6}{M quiescente, U quiescentes.}, ya se traslada al descanso del alma respecto de los trabajos del cuerpo. Pero oye, buen varón, ¿es esto descansar de los trabajos, porque a saber no padece en el cuerpo las molestias, pero siente fuera del cuerpo en el alma molestias incomparablemente mayores? ¿De qué sirve si descanso de la fiebre y soy infestado por el cálculo?

Sobre el descanso de Cristo\apnota{11}{M quete.} otra cosa es, al esperar la resurrección sin corrupción. Bienaventurados, pues, los que mueren en el Señor. Pero ¿quién dirá que el cuerpo es bienaventurado sino por redundancia del alma en el cuerpo? Así también\apnota{13}{fehlt in U.} del bienaventurado Esteban, que durmió en el Señor.

Después\apnota{15}{M U demum.} dice que pido el principio: Están en paz, y sin embargo les pedimos la paz. Y\apnota{16}{W. A. verbessert sed.} porque es ridículo referir la paz tenida al cuerpo, se sigue que aquella es en el alma la ciencia\apnota{17}{M iam statt illam; P U scientia.} de la seguridad.

Que dice que no puede probar lo suyo, lo reputo verdaderísimo, en cambio otros prueban lo suyo, aunque el pertinaz y testarudo ni siquiera asienta a las demostraciones, porque siempre contendería\apnota{20}{M U quin statt quoniam; M contenderit.} que no es demostración, como hacen Gregorio de Rímini y Pedro de Ailly contra Aristóteles en las razones sobre el primer motor en el 7 y 8 de la Física. Por lo cual concluyo con la común sentencia, tanto que en el purgatorio la gracia no aumenta, como que ellas\apnota{24}{W. A. verbessert animas.} están ciertas de la salvación.

El lunes once del mes de julio a la hora séptima se empezó a disputar la materia de las indulgencias\apnota{26}{U De indulgentiis. Die undecima Julii, quae fuit dies lunae. Lö. p. 438 Anm. „In M Sto additur hora septima".}, siendo el D. D. Eck\apnota{27}{M Egreg. D. D Eccius Damit endet der aus dem Notariatsprotokoll abgeschriebene Passus in M. Es folgt eine leere Seite mit der Notiz vacat. Die nun folgende Lage von 8 Blättern müsste die Signatur „P" tragen; da aber das 1. (leergelassene) Blatt fehlt, so ist auch die Signatur weggefallen. Das 2. Blatt beginnt: De indulgentiis Eccius.} quien disputaba.

En primer lugar, antes de descender a esta disputa de las indulgencias, protesto que no es de mi mente ni propósito\apnota{29}{U propositum.} querer con esta mi disputa contravenir los mandatos del sumo pontífice, por los que manda en el decreto que empieza: Cum postquam, que no deba defenderse ni predicarse lo contrario a ciertos puntos sobre las indulgencias bajo pena de excomunión, sino que más bien es de mi intención defender las verdades aprobadas\apnota{33}{M approbare, U veritatem . . . . approbatas (so!).} en el mismo decreto. Prefijado esto así, accedo a impugnar\apnota{34}{M U impugnanda m.} la conclusión undécima.

Decir que las indulgencias son útiles a los cristianos es verdadero y piadoso, ni las indulgencias son vicio de la buena obra, por lo cual no parece bien dicho que [345] quienes esto afirman estén locos. Esto se prueba porque la Iglesia en aquello\apnota{1}{M his.} que es de fe no yerra, y que mira a la salvación de las almas. Más aún, como dice san Cipriano, Dios no permite que yerre la mayor parte del clero, sino que en los concilios generales y por la universidad del clero ya desde hace trescientos años las indulgencias han sido reputadas útiles y piadosas para los cristianos. Lo cual constó en el concilio de Vienne\apnota{5}{M binensi, U Wienensi.}, donde fueron aprobadas las indulgencias\apnota{6a}{fehlt in U, in Uv nachgetragen.} dadas por Urbano IV para la veneración de la venerable eucaristía, pues el concilio\apnota{6b}{U consilium.} añadió esta razón motiva, para que los fieles de Cristo estuviesen más dispuestos a los obsequios de veneración y honor que han de tributarse al sacramento.

Así Inocencio III en el concilio general, donde promulgó la utilísima constitución: Omnis utriusque, recibida\apnota{10}{M recepta.} por toda la Iglesia, dispuso también sobre las indulgencias de los hospitales, diciendo: Os imponemos en remisión de los pecados que de los bienes a vosotros confiados les erogéis gratos subsidios\apnota{12}{M de bonis collatis grata charitate subsidiis, U = P, nur subsidia charitatis.} de caridad, para que por vuestra subvención\apnota{13a}{M subfensionem.} se socorra la indigencia de ellos y vosotros por estos bienes\apnota{13b}{U et vos praeponat.} y otros que hiciereis inspirándoos el Señor podáis llegar a los gozos eternos. En el mismo concilio fue también limitada la potestad de dar\apnota{16a}{M danda, verbessert dandi.} indulgencias de los prelados inferiores, la cual en el concilio de Lugdun\apnota{16b}{U in Lugduniensi consilio.} fue después aprobada. Pero si las indulgencias fuesen vicio de la buena obra e inútiles a los cristianos, ¿qué necesidad habría de disminuir en el prelado inferior la colación del vicio de la obra y conferirla a la plenitud de la potestad?

Se añade que el laudable concilio de Constanza\apnota{20}{U consilium.}, que también condenó entre los errores el desprecio\apnota{21}{M contemptor, Mc contemtum.} de las indulgencias, confirió indulgencias a todos los que estaban en el concilio y ayunaban el viernes, lo cual habría sido manchar e infectar las buenas obras de tantos padres. Y Gersón, a quien el reverendo padre llama ilustre teólogo, como fue estudiosísimo\apnota{25}{M honesti studii honestissimus.} de lo verdadero y honesto, decidió que la concesión de indulgencias no ha de tenerse en poco ni\apnota{26}{M seu, U pervidendam seu.} despreciarse, sino abrazarse devotamente en la fe, esperanza y caridad de nuestro Señor Jesucristo, que dio tal potestad de las llaves eclesiásticas a los hombres. Pues consta que es más fructuosa y más acepta a Dios la operación apoyada en tales indulgencias que la otra, y esto\apnota{29}{M et ob, U et.} en igualdad de condiciones no apoyada. Esto Gersón.

Así también otros santos y varones de óptima\apnota{31}{M optimi.} reputación en tiempo de las indulgencias dadas quisieron, las cuales también fueron dadas por el bienaventurado Gregorio hace ya

novecientos años y por Pascasio\apnota{33}{M Paschasium.} hace seiscientos años. Por eso la sola autoridad de tan gran padre debe mover a cualquier cristiano a creer que las indulgencias son útiles\apnota{34}{M utiles esse.} a los cristianos, aunque no ignoro\apnota{35a}{M ignoro.} que Guillermo de Auxerre\apnota{35b}{M Altioridensis, U Altisidorensis.} refiere que no faltaron quienes aun viviendo el bienaventurado Gregorio contradijeron al santo padre en las indulgencias.

Además ha de añadirse\apnota{38}{M addendo, U addam.} esto: si las indulgencias son vicio de la buena obra, esto sobre todo a mi arbitrio podría sospecharse\apnota{39}{U possem.} de que son satisfactorias, pero [346] por la obra satisfactoria no merecemos menos que si no fuese satisfactoria, de otro modo sería mejor que no hiciésemos ninguna obra satisfactoria, para no disminuir nuestros méritos para este estado. Pero sembrando en bendiciones también segaremos\apnota{4a}{U metemus.} de bendiciones. Pero aquí oiré\apnota{4b}{M audeam.} al reverendo padre. Añádase tanto consenso de toda la cristiandad en los jubileos celebrados por los sumos pontífices, por Bonifacio\apnota{6}{M celebratos statt celebratis; W. A. per Bonifacium 8.}, Clemente VI, Urbano VI, Nicolás V, pontífice integérrimo y doctísimo, Sixto IV\apnota{7}{M Sixtum IV etc.}, y que con común consenso de los cristianísimos\apnota{8}{M christianorum.} reyes y príncipes para los pasajes y cruzadas el sumo pontífice dio muchas veces plenísimas\apnota{9}{fehlt in M.} indulgencias.

Siendo, pues, indefectible la fe de la Iglesia, por la que rogó Cristo para que no desfalleciese, en Lc 22, y prometió estar con nosotros todos los días hasta la consumación del siglo en Mt 28 [v. 20], no ha de admitirse que la Iglesia haya errado con perdición de las almas por tanto tiempo, y porque los decretos de los sumos pontífices han de ser aceptados por todos los que tienen la comunión de la santa Iglesia, como dice Gregorio en el can. Praeceptis, 12. dist., asimismo en el can. Omnia decretalia, 25.\apnota{16}{P M U 24.} q. 1. con muchos concordantes, y el pontífice moderno definió bajo pena de excomunión latae sententiae que las indulgencias son útiles para borrar la pena debida por los pecados a aquellos\apnota{18}{M U poenam peccatis debitam his (U iis).} que unidos por la caridad son miembros de Cristo, y que tal remisión se compensa como por el mérito de Cristo y de los santos, por autoridad del sumo pontífice: por todo esto quiero persuadido al fiel cristiano de que las indulgencias no son vicio de la buena obra, ni que decir que es útil a los cristianos tener indulgencias\apnota{22}{M U esse utiles esse.} sea estar loco, lo cual no es otra cosa que decir que todo el clero en la Iglesia estuvo loco por tanto tiempo.

El D. D. Martín\apnota{24}{P D. D. Eccius.}.

Con doce argumentos impugna el egregio D. D. mi conclusión undécima. A los cuales antes de responder\apnota{26}{M U respondeam.}, primero diluiré aquello,

que dije que es locura tener\apnota{27}{M U haberi.} por buenas para el cristiano las indulgencias y que sean vicio de la buena obra. Digo, pues, que cuando el profeta en Sal 40 [v. 5] se atreve a llamar locuras falsas, y en Sal 119 [v. 85] fabulaciones de inicuos a las mismas tradiciones de los hombres, que de suyo no eran malas, sino porque prestaban a los insipientes una mala confianza contra la justicia de Dios: cuánto más es locura tener por bien de los cristianos las indulgencias, que ni están preceptuadas ni aconsejadas\apnota{32}{M U consultae sunt.} ni son necesarias para la salvación, más aún, son remisiones de muchas buenas obras, cuando todo nuestro bien y necesario para la salvación está comprendido en los preceptos y consejos. Más aún digo: cuando la salubérrima ley de Dios, al decir de Ezequiel y del apóstol Pablo en Rom 7 [v. 11 ss.], no es buena para el hombre, más aún, es ocasión de pecado, y sólo la gracia es bien para el cristiano, cuánto menos\apnota{37}{P M U magis, in W. A. verbessert.} deben tenerse por bien las indulgencias, que ni con la bondad de la gracia, ni con la de ninguna ley pueden en modo alguno compararse\apnota{38}{U possint.}, [347] siendo nada y cierta privación de buenas obras. No porque las juzgue\apnota{2}{M arbitremur.} nocivas o perniciosas, como diremos.

Cuando, pues, el egregio D. D. dice primero que la Iglesia «en aquello»\apnota{3}{U iis.} que es de fe no yerra, ni en aquello\apnota{4}{U iis.} que mira a la salvación de las almas, dice rectamente, pero tales no son las indulgencias, como se ha dicho.

Y en cuanto a que Cipriano dice que Dios no permite que\apnota{6}{M U quod.} yerre la mayor parte del clero, vea él mismo qué dijo. Cierto es que en tiempo de la pérfida arriana\apnota{7}{M Arriani.} erró la mayor parte y de los obispos más elocuentes y de los varones más doctos en tan gran artículo de fe, hasta el punto de que casi ningún obispo católico residía en su cátedra.

En cuanto a que dice que en los concilios generales y desde hace trescientos años se han tenido\apnota{11}{U iam statt et a; habitis.} por útiles y piadosas para los cristianos las indulgencias y aprobadas en el concilio de Vienne\apnota{12}{M binensi, U Wienensi.}, las que Urbano IV largamente concedió para excitar en los fieles\apnota{13}{M infidelibus statt in fid.} la reverencia a la venerable eucaristía, respondo: nunca negué que las indulgencias sean útiles, sino no a los cristianos, esto es, a los que obran fervorosamente, aman y buscan a Cristo según su nombre; a estos\apnota{16}{U iis.}, en efecto, les es ingrata la remisión de las obras, grata en cambio la imposición de las obras. Segundo digo que en estas\apnota{17}{U iis.} cosas no ha de mirarse a persona alguna\apnota{18}{M unius.}, sea concilio sea papa, sino qué se dice, sobre todo cuando este error sobre las indulgencias, como en cosa no necesaria, no es peligroso, en igualdad de condiciones por lo demás, esto es, si saben que no son necesarias ni confían\apnota{20b}{M perfidant.} en ellas que son\apnota{20a}{U eas esse.} no necesarias.

Al segundo, sobre el concilio general y sobre Inocencio III\apnota{21}{M de concilio generali Innocentii.}, que impuso para la disposición de los hospitales que confiriesen en remisión de los pecados

los gratos subsidios etc., digo que aún no entiendo esta palabra del papa, pues la remisión de los pecados me parece cosa distinta de las indulgencias\apnota{24}{U indulgentias.}, y si fuese lo mismo, digo como antes.

En cuanto a lo tercero, sobre el concilio Lugdunense\apnota{26}{M luculenso, gestrichen, a. R. luetnensi.}, en el que se dice aprobada también la limitación de la potestad de conferir indulgencias de los prelados inferiores, lo admito todo. Mas a lo que sigue: si fuesen inútiles, no habría sido necesario disminuir en los prelados inferiores\apnota{29}{M in inferiori praelato.} la colación del vicio de la obra; digo que las indulgencias fueron llamadas por mí vicio de la obra\apnota{30}{Mc 'boni' operis.}, según fui impugnado por el D. D. en mi conclusión 42 en el Resolutorio, donde dije así: no que las venias sean malas y nocivas, sino que el abuso perverso daña, cuando no harían tal obra si no hubiese venias. Pues así el fin de una obra de este tipo se convierte en la venia misma, y es clarísimo que raramente contribuirían los hombres si no se prometiesen indulgencias. Por eso siempre hay allí\apnota{35}{M U est ibi.} peligro, al menos de obra viciosa.

En cuarto lugar adujo el laudable concilio Constanciense, que entre los demás errores condenó el desprecio de las indulgencias. Respondo: nunca yo las desprecié ni enseñé que deban despreciarse, a no ser que se entienda el desprecio en este sentido: que con los mismos gastos con que se redimen las indul-[348]gencias podemos\apnota{39}{M U possumus.} hacer cosas incomparablemente mejores, o mejores que las indulgencias mismas. Por lo cual la prerrogativa de lo mejor no es desprecio de lo peor, así como el oro no es desprecio de la madera o del heno\apnota{2}{M veni, verbessert feni.}, aunque sea más excelente.

Por lo mismo, a lo quinto, donde aduce a Gersón, que la concesión de las indulgencias no ha de tenerse en poco, añado: en su género, o como él mismo dijo\apnota{5}{M dixerit.}, en igualdad de condiciones. Ni creo que esta sentencia agrade bastante al señor doctor, a saber, que las indulgencias han de recibirse en la fe, la esperanza y la caridad, habiendo él sostenido que la satisfacción puede cumplirse y remitirse también a quienes\apnota{8}{U iis.} están fuera de la gracia. Admitiría que es más fructuosa la acción apoyada en indulgencias que la otra no apoyada, en igualdad de condiciones, con tal de que aquellas demás condiciones iguales se comparen rectamente; así como no niego que puede ser mejor una sola oración dominical de un laico que todas las horas canónicas de un sacerdote.

En sexto lugar\apnota{13}{U Sexo, Uv Sexto.} aduce que el bienaventurado Gregorio hace novecientos años, y Pascasio hace seiscientos\apnota{14}{M sexcentum.} años dieron indulgencias, aunque, según refiere el Altisiodorense\apnota{15}{M 'referente' Wilhelmo Altiodorisiensi, U recitante Wilhelmo Altisidiorensi.}, algunos contradijeron a san Gregorio. Respondo: aún no tengo historia digna de fe\apnota{16}{M vide, verbessert fide.} sobre las indulgencias gregorianas. Pero sea lo que fuere, no se sigue por eso\apnota{17}{fehlt in M.} que las indulgencias sean otra cosa que indulgencias, esto es, remisiones de buenas obras, a las que nadie podrá llamar

buenas para los cristianos, al menos\apnota{18}{fehlt in M U.}, por muy santos que sean los concilios o quienesquiera\apnota{19}{U quantumlibet sancta concilia vel quamlibet sancti.} que las hayan dado o aprobado. Por eso, que el D. D. no actúe contra mí solo\apnota{20}{M solis, U non solum contra me agat solis.} con nombres de autoridades, sino también con la verdad de la cosa misma.

En séptimo lugar sospecha que son por eso\apnota{22}{M ideo esse, U eas esse ideo.} vicio de la buena obra, porque son satisfactorias. Digo que no son satisfactorias, sino remisiones de la satisfacción, y no han de tenerse\apnota{24}{M satisfactorium nec habendi.} como obra satisfactoria. Por eso mal se comparan con este argumento a las obras satisfactorias, como si por ello se disminuyeran nuestros méritos si se abandonan las indulgencias; antes es más verdadero lo contrario\apnota{26}{fehlt in M.}: los méritos aumentan cuando se abandonan las indulgencias.

En octavo lugar dice que por consenso de los fieles fue recibido el año jubilar de Bonifacio VIII, Clemente VI, Urbano VI, Nicolás V, Sixto IV etc.\apnota{29}{Urbani VI., Nicolai V. fehlt in M U, Sixti IV. fehlt in M; U Sixti statt Sixti IV.}. Respondo: ¿quién sabe si el consenso de los fieles lo recibió? Ni es Bonifacio de tanta reputación que cuanto hizo haya de tenerse al punto por bien hecho, sobre todo en esta materia\apnota{31}{fehlt in U.} no buena ni necesaria para la salvación, él que\apnota{32}{M quia.} también perpetró otras monstruosidades que en cosas necesarias dieron mal ejemplo.

En noveno lugar, que los Romanos pontífices dieron indulgencias por el pasaje etc.\apnota{35}{M passaio, verbessert passagio; etc. fehlt in U.} con consenso de los cristianísimos\apnota{34}{M roma, gestrichen, christianorum, verbessert christianissimorum.} reyes y príncipes, digo como antes: no es milagro que Dios desprecie esto y permita que se haga y obre lo que es impertinente para la salvación, él que sobre nosotros vela en esto, para enseñar lo útil, como dice el apóstol a\apnota{38}{fehlt in U.} Tito 3 [v. 8], esto es, sus preceptos y consejos.

En décimo lugar arguye que la fe de la iglesia es indefectible y que Cristo está con nosotros, y que por eso la iglesia no pudo errar tanto tiempo con perdición de las almas. [349] Digo: no erró toda la iglesia. Además, aunque errase en estas\apnota{1}{U iis.} cosas de nada (hablo siempre de los cristianos), no hay peligro, salva la fe. Variar las opiniones y errar en lo temporal no suprime la iglesia de Cristo.

En undécimo lugar sostiene que los decretos del Romano pontífice\apnota{4}{W. A. Romanorum pontificum (U Ro: pon:).} han de aceptarse. Respondo: sin duda han de aceptarse, pero con juicio, como se tiene en la dist. 19 cap. Anastasius\apnota{6}{P dis. XXII. vel XXI. ca. Anastasius, a. R. Eccius invenit esse XIX. dis.; M U dist. 22. ca. Anastasio, Mc verbessert im Text 22 in 19, a. R. 19 vel ut alii 22.}; porque los Romanos pontífices fueron hombres y tuvieron hombres a su alrededor, y pudieron pecar\apnota{7}{M U errare.}. Además, no hay ningún decreto que preceptúe redimir indulgencias, ni hasta hoy lo hay que haya declarado abiertamente el valor de las indulgencias.

En último lugar, sobre la definición del pontífice moderno, a saber, que las indulgencias son útiles y tomadas del mérito de Cristo y de los santos, respondo: ni ella\apnota{11}{U ipse.} misma lo expresa bastante ni lo prueba ni con una\apnota{12a}{M pro una (statt vel una), Mc per unam.} sola sílaba, sobre lo cual dije más copiosamente en mis actas\apnota{12b}{Mc actibus, verbessert actis.} augustanas. No se sigue, pues, que todo el clero en la iglesia haya enloquecido tanto tiempo\apnota{14}{M tanto tempore in ecclesia.}, pues entretanto siempre hubo quienes también les desagradaron y contradijeron al menos privadamente las indulgencias, como atestigua el vulgatísimo\apnota{16}{Mc 'pro' vulgatissimum.} proverbio: las indulgencias son piadosos engaños.

El D. D. Eck.

He recibido la sentencia del reverendo padre, en muchas cosas más mitigada y ecuánime de lo que las palabras de las conclusiones pretendían o lo que sus otros escritos me parecían significar\apnota{21}{M significabantur, verbessert significare videbantur.}. Ni es de mi instituto defender las indiscretas proclamaciones de las indulgencias o sus abusos, sino tutelar\apnota{23}{U possunt statt possint, M putari statt tutari.} la verdad y cuánto puedan aprovechar las indulgencias razonablemente dadas. Y para ser breve, no discuto la excusa de su conclusión. Pero cuando asume que las indulgencias no están preceptuadas\apnota{24}{M praeceptas fore.}, que no son necesarias, lo confesamos. Mas en todo este mi exiguo juicio aparece que la dificultad de nuestra disensión está principalmente en que el reverendo padre dice que las indulgencias no son sino remisiones de buenas obras, interpretación que no recuerdo haber hallado en ninguno que dé o admita indulgencias, pues así las indulgencias serían verdaderamente nocivas, peligrosas y condenables, ya que así los concilios y los sumos pontífices remitirían aquí\apnota{30}{U haec.} las buenas obras, de modo que o bien el hombre no merecería tanta bienaventuranza, lo cual es contra la salvación de las almas, o bien los hombres, engañados por la remisión de las buenas obras, serían enviados al purgatorio\apnota{32}{M remitterentur.} para sufrir allí las penas aquí no expiadas, lo cual es verdaderamente condenable\apnota{33}{M dandabile esset.}. De donde el consenso unánime de todos los que dan o admiten indulgencias es que las indulgencias no son remisiones de buenas obras, sino remisión de la pena temporal aún debida y no pagada por el pecado contrito. Como León, el papa moderno, expresamente definió en el citado decreto, lo que Sixto papa\apnota{37}{M diffinit, U definit; papa fehlt in M U.} IV hizo antes que él expresísimamente en tiempo de su jubileo.

De donde no puedo concordar lo que el reverendo padre dice en sus resoluciones, que las indulgencias solo se dieron a los perezosos, lo que [350] también afirma en el sermón vulgar, a los que no quieren cumplir los cánones penitenciales. Y sin embargo en el mismo resolutorio dice que los cánones penitenciales fueron impuestos a los perezosos, pues varones por lo demás esforzados harían penitencia toda la vida\apnota{3}{P totam vitam vel omnem, M totam vitam, U omnem vitam.}. De donde serían infelicísimas las indulgencias, si fuesen remisiones de buenas obras.

Para resolver lo primero respondió que las indulgencias no atañen a la salvación de las almas, cuando por cierto interesa muchísimo a la salvación de las almas si promueven o impiden la salvación. Interesa también a nuestra fe si el papa, como vicario de Cristo, hecha la recompensación por los méritos de Cristo y de los santos, puede\apnota{9}{M posset.} remitir y borrar las penas debidas\apnota{8}{M debitae.} al pecado.

Sobre Cipriano me objetó lo que se le ocurría al leer a Cipriano, aunque no me consta plenamente si, como en Grecia, también en otras regiones de la iglesia la mayor parte del clero\apnota{12}{P ecclesiae.} estuvo infectada.

Luego a los concilios\apnota{13}{M dein; U consilia.} respondió que las indulgencias no son buenas para el cristiano, esto es, para quien con fe y caridad sigue a Cristo. Esto ciertamente destruye los dichos de los concilios y la largición de las indulgencias, porque no se conceden sino a los fieles de Cristo contritos y confesados bajo la forma común. Por lo cual, según la sentencia común, las indulgencias no son buenas para los malos, porque no las consiguen, sino para los buenos y verdaderos\apnota{18}{M veris.} cristianos.

No acepto que el reverendo padre\apnota{19}{U quod 2. p[ater].} usurpe aquí el dicho de Séneca: no ha de atenderse quién, sino qué se dice, pues sobre todo en estas\apnota{20}{M quoniam quam maxime in his; U in iis.} cosas que son de la fe conviene atender también quién lo dice, y el concilio\apnota{21}{M U consilium.} legítimamente congregado ha de ser escuchado por todo cristiano.

Dice que el error no es peligroso si\apnota{23}{M sic.} aquí se yerra sobre las indulgencias. Mas a mí me parece lo contrario, pues fueron juzgados\apnota{24}{U fuerunt.} errores peligrosos los que desde el principio detrajeron a vuestras sagradas órdenes de frailes mendicantes.

A la constitución de Inocencio\apnota{26}{U Innocentii, Uv Innocentii.} en el concilio general dice que no entiende cómo valen, o que las indulgencias sean remisión de los pecados. Digo, como siempre: no concederé que tan solemne y legítimo concilio haya podido errar\apnota{28}{M U concilium (U consilium) et legitimum potuisse errare.}, pues no es abandonado por el Espíritu Santo, y la remisión de los pecados ha de entenderse en cuanto a la pena, porque el nombre de pecado no siempre se usa por la culpa, sino también por la pena debida a la culpa, como en los Macabeos, para que sean liberados de los pecados, lo cual no puede entenderse de la culpa mortal, sino de la pena debida a la culpa. Y si se dice que las indulgencias se dan para que sean absueltos de la pena y de la culpa, nadie debe estimar que las indulgencias remiten la culpa, sino que, porque el papa concede que de la culpa sean absueltos por comisarios diputados mediante el sacramento de la penitencia, y después\apnota{36}{U deinde.} consigan las indulgencias, esto es, las remisiones de las penas.

[351] En cuanto\apnota{1}{Vor quod hat W. A. die in U a. R. stehende Zeitangabe Hic continuavit Eccius hora secunda die 11. Julii 1519. Uv deutet die richtige Stellung dieser Angabe hinter facerent an; vgl. unten Z. 5.} a que se excusa de haber dicho\apnota{2}{M dicit.} que las indulgencias son vicio de la buena obra, porque también él quiso que, en igualdad de condiciones, la obra buena con indulgencias sea mejor que sin indulgencias, pero como suele hacerse, porque así los hombres redimen las indulgencias de modo que no harían aquellas buenas obras sin las indulgencias.

Transcurridas ya dos horas, lo restante se difirió a la hora segunda de la tarde; llegada esta, prosiguió como sigue\apnota{5}{Die Zeitangabe fehlt in M, doch ist der neue Anfang an der Schrift deutlich zu erkennen; über U und Uv vgl. oben Z. 1.}.

Aunque sobre aquel dicho podrían aportarse\apnota{6}{M illo verbo 'inferri'.} muchas cosas, y en caso semejante, cuando los vecinos por amistad de los vecinos acuden a las oblaciones y ofertorios más que por respeto a los muertos, sin embargo me parece, salvo siempre el mejor juicio, que esto no es vicio de la buena obra. Primero\apnota{9a}{M primo.}, porque el sacro concilio\apnota{9b}{U consilium.} quiere por las indulgencias excitar a los hombres a las buenas obras. Segundo, porque aunque hagan o exhiban tal cosa por conseguir las indulgencias, sin embargo esas obras no se inficionan por tal fin, pues este fin de las indulgencias no discrepa del fin último, de modo que la remisión de las penas se ordena congruentemente a Dios; de otro modo, prescindiendo de la subordinación de los fines, poquísimas o ningunas serían las buenas obras.

Sobre el desprecio de las indulgencias reprobado\apnota{15}{M concilium improbatum, U consilium improbatum.} por el concilio Constanciense, dice el reverendo padre que él no las despreció; lo confieso ingenuamente: si siempre hubiese usado de esta moderación, como hoy se expuso, y quizá haya de estarse a su interpretación\apnota{18a}{M interpretatione.} en conciencia sin fuero, sin embargo su sermón escrito\apnota{18b}{M scripta.} vulgar exhibe otra inteligencia a los laicos. Pero esto nada me importa.

A Gersón respondió admitiéndolo, pero juzga que Gersón me es adverso, quien reputa útiles las indulgencias\apnota{21}{M gratiam, gestrichen, indulgentias.} cuando se perciben en fe, esperanza y caridad. Que yo en otro tiempo haya sostenido que puede hacerse satisfacción en pecado mortal, digo que también el cristianísimo canciller sostiene esto: que alguien puede satisfacer de la penitencia impuesta en pecado mortal; de otro modo obrarían imprudentísimamente los oyentes de confesiones imponiendo penitencia duradera. Pero si alguien\apnota{26}{M quaespiam (Verwechslung der Abbreviatur).} consigue las indulgencias en pecado mortal, nada importa al caso; la razón de la diversidad es obvia, pues las indulgencias no suelen concederse sino a los contritos y confesados,

por eso no es extraño si los pecadores\apnota{28}{M peccatores existentes.} no son capaces de ellas, porque las indulgencias tanto valen cuanto suenan\apnota{29}{M sunt, Mc sonant.}.

Añadió el reverendo padre que la oración dominical\apnota{30}{M dominici.} del laico vale tanto como las horas canónicas del clérigo, si lo demás es igual; esto no lo admito.

Sobre el bienaventurado Gregorio dice que aún no ha visto recibida la historia. Mas la constantísima fama, como dice el bienaventurado Agustín en el libro sobre la Trinidad hablando de Alejandría, no permite que esto se niegue\apnota{34}{U negare.}. Niega en cambio que las indulgencias sean satisfactorias, sino que solo sean remisiones de buenas obras. En esto creo que principalmente disentimos y que aquí se funda gran parte de la controversia. Pero antes recuerdo que las indulgencias no son remisiones de buenas obras, como cree el reverendo padre, [352] sino remisiones de las penas debidas\apnota{1}{M debita.} a los pecados. De donde el hombre verdaderamente cristiano que recibe las indulgencias no hace menores, sino más buenas obras que si no consiguiese las indulgencias. Así pues entiendo que las indulgencias\apnota{3}{U indulgentis.} son satisfactorias, como declaran\apnota{4}{IV. und VI. fehlen in P, M U Sixtus Clemens (U Clem. 8) papa modernus declararunt.} el papa Sixto IV, Clemente VI y el papa moderno, no porque el que ha conseguido las indulgencias satisfaga, sino porque ningún mal queda impune y ningún pecado queda sin venganza (can. Sicut primi de poenit. dist. 1). Así el vicario de Cristo, dispensador de su tesoro, donde el hombre no devuelve el equivalente por la pena del pecado, él mismo lo recompensa del tesoro a él encomendado con el suficientísimo mérito de Cristo y por él satisface.

Por lo cual concedí en lo anterior que las indulgencias de suyo no borran la culpa, sino la pena, y que tal remisión de las penas en\apnota{11}{fehlt in M.} las indulgencias y no remisión de buenas obras se hace patente por aquello de Isaías: el Espíritu del Señor sobre mí, porque me ungió, para anunciar a los mansos me envió, para curar a los contritos y predicar a los cautivos la indulgencia y a los encerrados la apertura etc.\apnota{14}{fehlt in U.}. Cristo, asumiendo de sí las mismas palabras, por la palabra indulgencia\apnota{15}{U indulgentias.} tomó el nombre de remisión. Así dice Teófilo, como lo enarra\apnota{16}{M eum renarrat, U enim renarrat, Uv eum renarrat.} el bienaventurado Tomás: pues antes las almas de todos eran retenidas por los infiernos, hasta que vino quien predicara a los cautivos la remisión.

Luego sostiene el reverendo padre que los méritos aumentan abandonadas las indulgencias. Esto yo no lo veo\apnota{19a}{M vidi.}, pues, como dije, los que abandonan\apnota{19b}{M non, gestrichen, relinquentes, U non redimentes vel reliquentes (so!)} las indulgencias quieren a veces consumir\apnota{20}{P absumere} ese dinero en buenos convites, no aumentar los méritos. Confieso que a veces podría\apnota{21}{M non merita (Mc 'indul.') redimere fateor quod aliquando possit.} alguien hacer mayor

obra meritoria con aquel dinero\apnota{22}{U meritoritorium, in Uv verbessert; M peccata, Mc pecunia.} que si redimiese indulgencias; no conviene especificar los casos\apnota{23}{M specificare, U specificari.}.

Sobre los jubileos celebrados no quiere admitir que accedió el consenso de los fieles, y que la sospechosa vida de Bonifacio VIII podría hacer sospechoso\apnota{25a}{M Bonif. octava vita posset suspecta, U = P, nur potest statt posset.} su jubileo. Digo\apnota{25b}{M dicam (?).}, que es bastante grande el consenso del pueblo, pues de todo el orbe cristiano\apnota{26}{M Christiani.} confluyen los fieles de Cristo en el jubileo; cómo podría probarse\apnota{27}{M probari.} de otro modo el consenso de los fieles, no lo veo.

A lo que dije que con consenso de los príncipes fueron dadas por el sumo pontífice indulgencias por el pasaje y para los cruzados, responde que poco concluye, si aquí erraron en cosa de nada, como si fuese nada engañar a los fieles de Cristo y que el unánime consenso de los príncipes y reyes cristianos descendiese a tan falaz engaño de los fieles.

Además dice que no hay ningún decreto expreso y claro sobre las indulgencias. Le opongo la declaración sixtina y ya novísimamente la declaración del pontífice moderno, donde dice que la iglesia Romana, a la que las restantes están obligadas a seguir como a madre, transmitió que el Romano pontífice, sucesor de Pedro el clavero\apnota{37}{M clafigeri.} y vicario en la tierra de Jesucristo, con la potestad de las llaves, cuya función es abrir quitando en los fieles de Cristo sus impedimentos, a saber, la culpa y [353] la pena debida por los pecados actuales. La culpa ciertamente\apnota{1}{U deletam statt debitam, M scilicet statt quidem.} mediante el sacramento de la penitencia, y la pena temporal debida según la divina justicia por los pecados actuales, mediante la indulgencia eclesiástica: que por causas razonables pueden concederse a los mismos\apnota{4}{M U concedi fehlt; M eiusdem.} fieles de Cristo, que por la caridad que une son miembros de Cristo, sea en esta vida, sea en el purgatorio\apnota{5}{M purgatoria, Mc purgatorio.}, indulgencias de la superabundancia de los méritos de Cristo y de los santos, y concediendo la indulgencia por autoridad apostólica tanto por los vivos como por los difuntos\apnota{6}{M = P, nur et pro statt quam, U et tam pro vivis quam pro defunctis.}, dispensar\apnota{8}{M U dispensare etc.} el tesoro de los méritos de Cristo y de los santos. Donde determina claramente con lo siguiente que las indulgencias no son remisiones de buenas obras, sino remisiones de las penas, que se hacen\apnota{9}{U fiat.} hecha la compensación de los méritos de Cristo.

Y a lo que yo razonaba\apnota{11}{P rationabar; M quo statt quod.}, que la iglesia en tantos años no erró toda, respondió que no erró toda la iglesia, porque siempre entretanto\apnota{12}{M interim.} hubo quienes tuvieron por sospechoso este negocio. De ahí aquello vulgado: las indulgencias son piadosos engaños de los fieles. A lo cual digo: cuando ha de darse razón de alguna cosa, no ha de escucharse lo que cada uno diga, sino lo que estatuyen los buenos, doctos, honestos y óptimos. Tantos\apnota{16a}{U quot.} en efecto fueron

varones probadísimos y honestísimos por todo el orbe cristiano\apnota{16b}{M probatissimi viri per totam Christianam.} en trescientos años, y algunos de ellos referidos en el catálogo de los santos, que sin embargo no improbaron\apnota{18}{U probaverunt.} de este modo las indulgencias, aunque les desagradase el abuso de ellas como en toda cosa. Santo es Tomás, Buenaventura, Alberto el Carmelita, Bernardino\apnota{20a}{M Bernhardinus Antonius de Badna (so!)} y otros grandes varones, Alejandro\apnota{20b}{fehlt in M.}, Gersón, Capistrano etc., que sin embargo no reprobaron la concesión de las indulgencias. Por lo cual nada importa que a Juan Wiclef y a Ulrico Kalteisen en Inglaterra, al señor Juan de Wesalia\apnota{23a}{M Johanni bicleff (Mc 'bicler') et Ulrico Calte in Anglia, d. Joh. de Besalia.}, que sin embargo retractó esto en Maguncia, y a los Bohemios\apnota{23b}{et fehlt in U; P Boiemis.} no les agradasen las concesiones de las indulgencias.

Por lo cual concluyamos que las indulgencias razonablemente dadas para el honor de Dios y la exaltación de la fe no son inútiles a los fieles de Cristo devotos, que no presumen, sino que las reciben en humildad y acción de gracias\apnota{27}{M U cum humilitate et (M et cum) gratiarum actione.}, resecados y repelidos los abusos, de modo que demos más bien gracias a Dios por tan gran don, que si soberbiamente lo\apnota{29}{M U illud.} despreciásemos, teniendo siempre en el ánimo que la mejor indulgencia\apnota{30}{M contritionem gestrichen, indulgentiam.} es la verdadera contrición. Sin embargo, en todo caso, si la sentencia hoy expuesta fue la mente del reverendo padre, habría preferido que lo hubiese dicho claramente y no hubiese dado a los simples ocasión de pensar otra cosa, porque no hay ninguno de todos los oyentes que haya entendido en la conclusión undécima el sentido hoy dado: decir que las indulgencias son un bien para el cristiano es delirar, pues son verdaderísimamente vicio de la buena obra. Mas estoy bien contento de su exposición, en cuanto a mí atañe.

El D. D. Martín.

Contra mis soluciones opone el egregio D. D. primero que las indulgencias no son remisiones de buenas obras, ni ha visto quien esto afirme. Seguirse [354] también de esto que las indulgencias son verdaderamente peligrosas, nocivas\apnota{1}{M obnoxias.} y condenables, sino que por unánime consenso de los que las dan y admiten\apnota{2}{M Lücke, Mc 'admittentium esse'.} son remisiones de las penas debidas\apnota{3}{pro peccatis debitarum fehlt in M.} por los pecados. Respondo: bastante me admiro de qué quieran estas palabras del señor doctor, pues sin duda por penas debidas por los pecados entienden todos los trabajos de la satisfacción, que es la tercera parte de la penitencia. Entre los cuales ciertamente están las buenas obras, las limosnas, las oraciones, los ayunos, las vigilias\apnota{6}{M et vigilhae.} y semejantes. Para lo cual hay texto expreso en el lib. 5\apnota{7}{Mc 'lib. 6.'} cap. Cum ex eo, donde el papa dice que por las indiscretas indulgencias se enerva la satisfacción penitencial, y casi todas las bulas de indulgencias tienen esta cláusula: de las penitencias impuestas relajamos.

Pero ea, dado que fuesen penas y no obras, digo que es peor remitir las penas que las obras, pues la vida pasiva y penosa perfecciona y aprovecha más que la meramente\apnota{11}{U aut prosit quam mere; M mera.} activa, como dice el apóstol en 2 Cor 12 [v. 10]: cuando\apnota{12a}{P 2. Corinth. XI. forte dum, M 2. ad Corinth. 11 dum (Mc 'cum'), U 2 Corinth. 11 cum.} estoy enfermo, entonces soy más fuerte. Pues es más pura\apnota{12b}{Mc a. R. durius.} la obra que, obrando solo Dios y padeciendo nosotros, se perfecciona, que la que se efectúa cooperando nosotros, según aquello de Dt 32 [v. 39]\apnota{14}{P M U Deut. 22.}: yo heriré y sanaré, mataré y vivificaré\apnota{15}{M occidam etc.}. Por lo cual, según el sentir del egregio señor doctor, se siguen de las indulgencias más absurdos que según mi sentir, y en verdad no sé si diga\apnota{16}{M dicam, Mc dicat.} la verdad.

En segundo lugar no puede concordar que dije que las indulgencias solo fueron dadas a los perezosos y que los cánones igualmente solo\apnota{18}{M tantum similiter.} fueron impuestos a los perezosos. Digo: no recuerdo haber dicho que los cánones fueron impuestos a los perezosos; pero si lo dije, mi sentido es\apnota{19}{M esset.} que los esforzados y fervientes pecadores\apnota{20}{U ferventes sive poenitentes, Uv = P.} hacen más de lo que los cánones exigen.

En tercer lugar\apnota{21}{M U Tertio contendit eas esse necessarias ad salutem, quia expedit nosse an promoveant, necne ad salutem, Mc streicht die Worte von expedit an, Uv streicht das Ganze; in P fehlt alles ausser tertio; vgl. Einleitung p. 7.}, aunque confiesa que las indulgencias no son necesarias para la salvación, conviene sin embargo saber si promueven o no\apnota{23}{U Etsi non fateatur necessarias esse ad salutem expedit tamen scire, promoveant salutem necne.} la salvación. Respondo: ni esto es necesario saberlo, sino si dañasen a la salvación; pues no es necesario saber lo no necesario para la salvación\apnota{25}{M necessaria saluti.}.

Asimismo, ¿qué interesa si\apnota{26}{M U quod interest an (M quod, Mc 'an').} el papa puede remitir por los méritos de Cristo? No sé si interese\apnota{27}{M an id 'intersit', U an inersit (Uv intersit).}. Sobre los méritos de Cristo, después\apnota{28}{M = P, Mc 'quae merita Christi postulant' (a. R. wiederholt).}.

Después de esto dice: mi sentencia destruye los dichos de los concilios y la largición de las indulgencias, que con forma común se conceden a los contritos y confesados. Respondo: no por eso destruye mi sentencia los dichos de los concilios, porque son contritos y confesados los que son perezosos para pagar las satisfacciones, a los que la Escritura\apnota{32}{M in scriptura.} también llama enfermos, y por estos no han de condenarse las indulgencias.

Me reprende también\apnota{34}{M U in me.} por haber usurpado el dicho de Séneca: atiende no quién sino qué se dice\apnota{35}{M attendi, fehlt in U.}, que en estas cosas que son de la fe conviene atender también quién lo dice. Respondo: esto no viene al caso, porque la materia de las indulgencias no es materia de fe. Además, en materia de fe muchísimo más necesario es atender no quién, sino qué se dice,

porque no ha de hacerse acepción de personas, como manifiestamente previno Cristo\apnota{39}{fehlt in M U, von Mc Uv nachgetragen.} en Mt 7 [v. 15] y 24 [v. 5]: guardaos de los falsos profetas, que vienen a vosotros etc. Asimismo: muchos vendrán\apnota{40}{U veniunt.} en mi nombre diciendo: yo soy [355] Cristo; y Jn 10 [v. 27]: mis ovejas oyen mi voz. Por tanto\apnota{1}{M audient, P Deinde.}, tampoco al sumo pontífice ha de creerse temerariamente, para que se mantenga la regla del apóstol Juan [1 Jn 4 v. 1]: carísimos, probad los espíritus si son de Dios.

Dice luego que hay peligro en el error sobre las indulgencias, porque fue tenido por error peligroso el de quienes desde el principio detrajeron a las órdenes mendicantes. Sea lo que fuere de los detractores\apnota{6}{U detractatoribus.}, por esto no se prueba error peligroso en las indulgencias, ni es al punto error lo que las órdenes mendicantes o los sacerdotes de las iglesias juzgaron error. Por lo mismo, a lo que dice que el concilio no yerra porque es regido por el Espíritu Santo, añado yo: a saber, en aquello en que es regido por el Espíritu Santo, esto es, en las cosas de la fe. Y para decirlo alguna vez: no dije que el concilio erró en dar indulgencias, sino que puede errar\apnota{12}{U errasse in concilio in dandis indulgentiis, sed postea (Uv. a. R. sed posse errare).}.

Paso aquello de que entendió\apnota{13}{MU intelligit.} la remisión de los pecados como remisión de las penas, pues a menudo el pecado se toma por la pena, en 2 Mac 12 [v. 46]\apnota{14}{P M Macc. 2, U Macc. 20.}. Yo no hallo a menudo que el pecado se tome por la pena, salvo quizá donde se escribe de Cristo que él mismo llevó nuestros pecados, lo cual sin embargo tampoco me atrevería\apnota{16}{M ausum.} a decir que simplemente se tome por la pena. Pero esto es ajeno.

Para confutar aquello de que llamé a las venias\apnota{18}{U dixi esse.} vicio de la obra, dice que hay subordinación de los fines en las obras\apnota{19}{M in bonis operibus.}; por eso, si las venias se refieren a Dios, no vician la obra. Respondo: da un ejemplo\apnota{20}{M da exemplum 'et da', die letzten beiden Worte gestrichen, U da exemplum et.} de demostración. Pues dije hoy que raros se hallan que tanto diesen gratis y por Dios cuanto dan por las indulgencias; por sus frutos, pues, los conoceréis\apnota{22}{U cognosceris.}; más aún, si oyesen que pueden hacer algo mejor que redimir indulgencias, es probable que ningunos\apnota{23b}{U nullas.} redimirían indulgencias, si se\apnota{23a}{M si audeam melius posse se, U si audiant melius se posse.} puede hacer mejor.

Sobre el desprecio de las indulgencias dice: si hubiese usado de esta moderación\apnota{25}{M moderamini.}, habría de estarse a mi\apnota{26}{M esset eius, U est eius.} declaración; pero el sermón vulgar indica otra inteligencia a los laicos. Respondo: vea con más diligencia el mismo sermón el egregio señor doctor y hallará que dije con expresas palabras que las indulgencias no han de despreciarse, sino dejarse libres; que a nadie\apnota{29}{M U neminem esse.} ha de disuadirse, como a nadie ha de urgirse, a saber, que son de los perezosos

consuelo\apnota{30}{M esset videlicet 'pigrorum' (ursprünglich virorum?) solatia, U = P, nur solatia.}. Por eso, sin culpa mía, no entienden esta declaración y esta sentencia bastante moderada.

Paso por alto aquello de Gerson, donde dice: una cosa es satisfacer en los pecados y otra rescatar indulgencias en los pecados: porque en estas últimas\apnota{33}{U iis.} obsta la forma de la concesión, que solo permite las indulgencias a los confesados y contritos, mientras que allí el mismo Juan Gerson piense consigo mismo. Digo: dejo\apnota{35}{U relinquendo, Uv relinquo.} esto como algo fuera de la capacidad de mi entendimiento, pues en verdad no comprendo cómo el pecador puede satisfacer en los pecados y no\apnota{37a}{W. A. lässt dieses non fälschlich aus.} puede recibir también la licencia\apnota{37b}{M possit etiam recipere (dicentium, gestrichen, Mc a. R.) 'licentiam'.} de no satisfacer, cuando parece requerirse más para la perfección de las obras que para la\apnota{38}{fehlt in M.} omisión.

[356] Sobre Gregorio, aunque la historia no se tenga por recibida\apnota{1}{M etiam recepta.}, dice que la fama no permite negarlo. Digo: nadie crea temerariamente a la fama, y paso por alto.

Finalmente vuelvo a esto: que las indulgencias son satisfactorias y no remisiones de las buenas obras, sino de las penas; quizá contendemos en un equívoco, pues lo satisfactorio el D. D. lo llama así porque se permite, en virtud de las indulgencias, no satisfacer, y esta permisión\apnota{6}{U promissio.} de no satisfacer se tiene por satisfactoria. Yo llamo satisfactorio a lo que cumple la satisfacción, no a lo que la omite\apnota{7}{in M durchstrichen, daher unleserlich, Mc. a. R. omissionem.}.

Adujo también la autoridad de Is 61 [v. 1]\apnota{8}{P U Esai. 66, M Esai. 60.}, donde Cristo, leyendo de Isaías en Lc 4 [v. 18]\apnota{9}{M Esaie.}, dice: El Espíritu del Señor sobre mí, porque me ungió para predicar a los cautivos la indulgencia, donde Cristo dijo remisión\apnota{10}{fehlt in P M.}. Añadió a Teófilo, que expone que las almas estaban retenidas etc., hasta que vino\apnota{11}{etc. fehlt in M; M veniat.} quien predicase a los cautivos la redención. Que aquella autoridad\apnota{12}{M dico quod illas auctoritas.} nada hace al caso, lo declaran las propias palabras y la consecuencia del texto\apnota{13}{M verba textus et sequentia.}, porque allí mismo dice Cristo: Hoy se ha cumplido esta Escritura\apnota{14}{M haec scriptura.} en vuestros oídos; además, porque habla del verdadero año del jubileo, esto es, de la plenitud del tiempo, de la corona del año benigno, y no de las remisiones de las penas, sino de la remisión de los pecados, de la que habla toda la Escritura, habiendo comenzado el jubileo bonifaciano tanto tiempo después. Asimismo, que niegue que se aumenten los méritos dejadas de lado las venias, lo paso por alto, porque con ello no se refuta que alguna vez quieran gastar el mismo dinero en convites.

Sobre los jubileos celebrados dice que es bastante grande el consenso del pueblo, pues confluyen a las indulgencias tantos\apnota{22}{fehlt in M U.} fieles en todo el orbe, exceptuada

sin embargo Italia y Roma, y no ve cómo de otro modo pueda yo probar el consenso\apnota{23}{M Nec videat quomodo alias possit probari consensus; U possim probare consensum, sonst = P.} de los fieles. Respondo: yo probaría que es consenso de los fieles\apnota{24}{M ego fehlt; esse fidelium probarem.} el que, revelada la cosa y expuesta la verdad de las indulgencias, moviese a los fieles a confluir. Ahora, en cambio, confluyen ignorando qué y por qué, pues la mayor parte cree que hace\apnota{26}{fehlt in M.} una cosa necesaria y meritoria.

Después de esto, el egregio D.\apnota{28}{fehlt in M.} D. juzga que no es cosa de poco momento que sean engañados tantos fieles, sobre todo magnates. Respondo: tampoco es cosa de gran momento, puesto que también el divino Jerónimo, sobre Mt 23 [v. 35], acerca de la sangre del profeta Zacarías, dice: No condenamos el error que desciende de la piedad de la fe\apnota{32}{M descendunt, verbessert descendit.}. Así también este error es sin daño del alma, e incluso sin lucro de dinero el que aportan por honor de la santa iglesia a los rectores de las iglesias.

Cuando yo\apnota{34}{fehlt in M.} dije que no hay ningún decreto sobre las indulgencias, opuso la Sixtina y la novísima Leonina, donde se declaran las virtudes de las indulgencias. Digo: esta es otra cuestión. Yo dije que no hay ningún decreto que preceptúe rescatar indulgencias, y por eso que no son necesarias.

Al final, al proverbio aquel: las indulgencias son piadosos engaños de los fieles, con el que quise decir que no erró toda la Iglesia, dice que no ha de escucharse lo que cada uno diga, sino lo que establezcan los varones buenos y honestos. Ahora bien, también algunos\apnota{40}{M statuunt nunc etiam aliqui (sancti, gestrichen) etiam. U = P, nur statuunt.} [:357] inscritos en el catálogo de los santos no reprueban las indulgencias. Respondo: ni yo tampoco, ciertamente. Esto añado, sin embargo: que la Iglesia, al recibir opiniones, no hace por eso de las opiniones verdades. Por eso, a aquella declaración de León X\apnota{3}{M U istam Leonis X. (X. fehlt in M).} digo, como dije en las actas, que no se prueba con meras palabras, sobre todo de hombres, que las indulgencias estén tomadas del tesoro de Cristo y liberen de las penas requeridas por la justicia divina\apnota{5}{U iustitia divina.}, aunque no condeno la opinión. Pues, por decir lo que siento, como los méritos de Cristo, ya se tomen\apnota{7}{M accipientur.} como sufragios ya de cualquier otro modo y nombre, no dejan de ser méritos de Cristo, ni por el uso o la aplicación sufren metamorfosis. Pero si son méritos de Cristo, gracia y verdad son, según aquello del Sal 25 [v. 10]: Todos\apnota{10}{M universi.} los caminos del Señor son misericordia y verdad, y Jn 1 [v. 14ss.]: La gracia y la verdad fueron hechas por Jesucristo; allí mismo: Lo vimos lleno de gracia\apnota{12}{P M gratia.} y de verdad. Por lo cual, aunque un ángel del cielo persuadiese otra cosa, no creeré que esté en mano de hombre\apnota{13}{M unius.} alguno dar la gracia y la verdad, esto es, dispensar los méritos de Cristo; sin embargo no resisto al pontífice. Sino que me remito a la declaración en las actas.

El D. D. Eck.

Porque, como dije, gran parte de la disensión reside en si las indulgencias son remisiones de las buenas obras o de las penas, donde dije que a mi parte descienden los doctores de la Iglesia y los pontífices romanos\apnota{19}{W. A. Romani.} que hasta ahora escribieron sobre las indulgencias. Pero esto lo rechaza\apnota{20}{M rescripserunt statt scripserunt; U refellit.} aguda y doctamente el reverendo padre: porque remitir la pena «sería remitir los trabajos de la satisfacción»\apnota{21}{Mc 'satisfactionum'.}, entre los cuales están las buenas obras, contra el decreto de Inocencio, cap. Cum ex eo, de poenitentiis et remis.\apnota{22}{et remis. fehlt in M U.}, «no sea enervada la satisfacción penitencial». Digo que bien se remiten los trabajos de la satisfacción, pero no por eso las buenas obras, pues aunque la oración, el ayuno etc. sean satisfacciones, sin embargo son buenas obras aun si con ellas\apnota{25}{U iis.} no satisfacemos. Se añade que la satisfacción no solo se hace\apnota{26a}{M fit, Mc 'satisfit'.} con el trabajo de las buenas obras, sino también por satispasión\apnota{26b}{P M Uv satispassio.}. Por lo cual la mayoría quiere que no deba obligarse al hijo de confesión a recibir en sí la satisfacción; según aquello común: mejor\apnota{28}{U melior.} es que el confesor envíe al hijo con pequeña penitencia al purgatorio que con grande al infierno. Y no es extraño que el pontífice recortase\apnota{30}{M resequit, Mc 'resecuit'.} las plenarias indulgencias para eludir las satisfacciones penitenciales, si por doquier se hiciesen\apnota{31}{P Uv fieret; M fiat, U fiant; M U eludens statt eludendas, Uv = P.} por los prelados. Así la partícula de la bula sobre las penitencias impuestas, esto es, las que según la justicia divina deberían imponerse\apnota{32}{M deberet.}, como declaran los doctores eclesiásticos\apnota{33}{M U ecclesiae.} que escriben sobre las indulgencias.

Por lo demás, que argumente que es peor remitir las penas que las obras, no lo acepto, porque las penas aprovechan en tanto que el hombre sea instruido a obrar bien, y así la virtud se perfecciona en la enfermedad, pero esto a los jueces\apnota{36}{W. A. indices.}. Y porque errar en las indulgencias sería «peligroso» para las almas, como hoy aduje, por eso la Iglesia no yerra si solo quiere\apnota{38}{M U ideo ecclesiae non errant, si . . . . volunt (M vult).} que los perezosos contritos reciban las indulgencias; ciertamente a los cartujos no deberían venir\apnota{39}{M deberet, U debent.}.

[358] En materia de fe dice\apnota{1}{M dixit.} que ha de atenderse sobre todo qué se dice. Yo, en cambio, reputo la autoridad potentísima en teología, y que se nos manda escuchar a los sentados sobre la cátedra de Moisés etc.\apnota{3}{M sedentem statt sedentes, U supra statt super; etc. fehlt in M U.}. Por eso debemos atender quién lo dice, para cautivar el entendimiento, según el precepto del apóstol, en obsequio de la fe; por eso\apnota{5}{M proptera.} se celebraron concilios, para que con su autoridad fuesen exterminados los errores.

Sobre los mendicantes paso por alto.

Pero lo que dice, que el concilio en aquello en que es regido\apnota{7}{M concilium in his, quae regantur, U consilium in iis in quibus regitur.} por el Espíritu

santo no yerra, digo yo, siguiendo la sentencia de los antiguos padres\apnota{8}{fehlt in M.}, que siempre debe presumirse que el concilio\apnota{9}{U consilium.} es regido por el Espíritu Santo mientras no conste lo contrario, a saber, la ilegitimidad de la congregación.

Pero todavía insiste fuertemente el reverendo padre en el vicio de la buena obra, que el que da dinero\apnota{12a}{M pecunias.} por las indulgencias no lo daría de otro modo\apnota{12b}{M daturas.}, como si esto fuese vicio de la buena obra. Preguntaría yo al «reverendo padre» si él mismo «con sus hermanos» tiene el aniversario de algún «príncipe» o de algún noble con treinta misas, porque se dan buenas presencias, no habiendo de tener de otro modo el aniversario\apnota{15a}{Mc 'cum suis fratribus habet anniversarium' (a. R., die beiden letzten Worte auch im Text) 'cuiusdam principis' (Lücke ausgefüllt) eum 30 vel 20 missis non 'habiturus etc.' (statt anniversarium); U = P, nur habet statt habeat, nobilis statt alicuius nobilis.}: ¿acaso esto es vicio de la buena obra? [Martín duda. Eck: Y así no será salva toda carne, tanto en la cogulla como fuera\apnota{15b}{Die Zwischenbemerkung in M U a. R.; M dubitet statt dubitat, Eccius und tam in cuculla quam extra fehlen.}.]

Sobre las indulgencias no despreciadas por él, juzguen quienes hayan leído el sermón vulgar\apnota{17}{M vulgarem suum sermonem legent.}.

Por lo demás, el reverendo padre, hecho intérprete mío, dice que quizá las indulgencias son así satisfactorias\apnota{19a}{U factas esse.} porque por ellas\apnota{19b}{U istas.} alguien «no esté obligado a satisfacer». Este no es mi sentido, sino que, debiendo el hombre satisfacer de lo suyo también por los pecados contritos\apnota{20}{M etiam pro peccatis.}, por las indulgencias satisface de lo ajeno, porque el sumo pontífice le da del tesoro de la Iglesia con qué pagar, de modo que así, según la común sentencia de los padres, tampoco en las indulgencias el pecado quede\apnota{23}{P U manet.} impune, quedando a salvo que tampoco podría satisfacer de lo propio sin que accedan los méritos de Cristo.

Por lo demás, sobre el consenso de los fieles, creo que nadie es tan necio que crea que es cosa necesaria las indulgencias en el jubileo, y si estima que el viaje a Roma le es meritorio, tampoco aquí yerra\apnota{28}{M esse sibi meritorium, nec hic erret, U esse meritorium nec hic errat.}, pues subsiste que la obra por la que uno consigue las indulgencias es meritoria y sin embargo las indulgencias mismas no son meritorias.

Por lo demás, en las palabras sixtinas\apnota{30}{M Sixtini, U Sixtinae.} y del pontífice moderno no hago fuerza, porque sigo la cátedra de Pedro y su sentencia, el que en ella se sienta, mientras no haya caído en herejía, lo que no suceda. Sé que él no aprueba\apnota{32}{M Scio enim eum probare.}, sino que define, al cual, como vicario de Cristo, al definir\apnota{33}{M diffinire und diffiniendo.} juzgo que es de fe indefectible. Pero todo esto lo encomiendo al juicio de aquellos a quienes atañe, dispuesto a tachar los errores, si se mostraren\apnota{35}{M demonstrentur.}.

Añadió el D. D. Martín: Y yo\apnota{36}{M Addidit D. D. fehlt, U Jussit Martinus addere: Et ego.}.

[359] Mañana a la hora séptima antes del mediodía se continuará la disputa entre los mismos E. D. doctores y se tratará la materia de la penitencia.

El día martes doce del mes de julio, a la hora séptima de la mañana.

El D. D. Eck\apnota{5}{M Die Zeitbestimmung von crastina an fehlt, Überschrift De poenitentia; U duodecima Julii argumentatus est Eccius respondente Martino de poenitentia. In nomine tuo, dulcis Jesu.}.

Contra la conclusión del reverendo padre y parte de su sermón sobre la penitencia, en defensa de los óptimos padres predicadores intento probar que la verdadera penitencia comienza por el temor aun de la pena, y que bien se prepara uno\apnota{9}{M se quis praeparet.} discutiendo, recordando etc.

Primero, porque tal modo de predicar lo observa\apnota{10}{M U observavit; U a. R. Eccius.} el Señor Jesús y su precursor san Juan. Pues en Lc 15 el hijo pródigo representa al penitente según Agustín en las cuestiones evangélicas, Ambrosio en el libro 2 sobre la penitencia cap. 3, Crisóstomo, Jerónimo y otros, pero a este\apnota{13}{M eam.} nos lo propone Cristo así: vuelto en sí dijo: Cuántos\apnota{14}{U quod in se conversus dixit: quantum.} jornaleros en la casa de mi padre abundan en panes, y yo aquí perezco de hambre; me levantaré e iré a mi padre y le diré: Padre, pequé contra el cielo y delante de ti etc. Aquí el Señor Cristo, describiendo el modo del penitente, expone primero el movimiento que lo hace penitente por la magnitud de los premios, a saber, «abundan en panes», y por el temor de la pena, a saber, «aquí perezco de hambre», por cuyos grados elevado comenzó a meditar la verdadera penitencia, a saber, «y diré: Padre, pequé» etc. Y como\apnota{20}{U peccavi et hoc.} explica Basilio: tres son allí los grados de la penitencia, la esperanza del premio, el temor de los suplicios y el sincero amor de la bondad paterna. Y así, antes de hacerse amigo de Dios quien fue enemigo, primero se hace siervo.

Así Juan comenzó por el temor: Raza o linaje de víboras, ¿quién os mostró\apnota{24}{M demonstrabit, U monstrabit.} huir etc.?

Esto\apnota{25a}{U sic.} también se hizo en todo el Antiguo Testamento\apnota{25b}{U instrumento.}, de modo que los hombres fuesen atraídos por el temor de las penas a la obediencia de los mandamientos, como en los libros legales, las historias y\apnota{26}{fehlt in U.} los profetas se halla por doquier. Baste aducir\apnota{27}{M repetitur statt reperitur; M U adducere, Uv = P.} un lugar, Sal 89 [v. 31ss.]: Si dejaren sus hijos mi ley y no anduvieren\apnota{29}{U ambulaverit, Uv ambulaverint.} en mis juicios, si profanaren mis justicias y no guardaren mis mandamientos, visitaré con vara sus iniquidades de ellos y con azotes sus pecados; con las concordancias\apnota{31}{M concordanti, U concordantibus, Uv = P.}. Por lo cual los predicadores, imitando hasta ahora a la sagrada Escritura, tuvieron un buen modo de preparar la penitencia, recordando la gravedad de los pecados, del suplicio eterno etc.

Se añade la razón dionisiana: porque Dios administra así las cosas, que conduce lo ínfimo a lo superior por lo intermedio. Por lo cual, como el pecador está en lo ínfimo, porque no es digno del pan que come, queriendo ascender a lo superior de la gracia lo hará\apnota{37}{U faciat.} por el medio del temor.

Bien sería de perfección que alguien inmerso en el lodo [360] del pecado pudiese alzarse con el mero contemplar el amor de Dios y el amor de la justicia\apnota{1}{M erire, Mc erigere.} a alcanzar la gracia; pero ¿quién es este y lo alabaremos?

Más aún, el bienaventurado Agustín aprueba el modo de predicar de nuestro tiempo, enseñando expresamente que nadie llega a la caridad y a la verdadera gracia sino precediendo el temor, no filial, sino aun servil. Agustín lo atestigua en el tratado 9 sobre Jn\apnota{5}{P M U Aug. est tract. 9; M Test., verbessert Augustinus; super Joh. fehlt in U, in Uv nachgetragen.}: comience, pues, el temor, porque el inicio de la sabiduría es el temor del Señor. Pero cuando la caridad comience a habitar, es expulsado el temor que le preparó lugar. Pues cuanto aquella crece, este decrece, y cuanto aquella se hace más interior, el temor es arrojado fuera; mayor caridad, menor temor; menor caridad, mayor temor. Pero si no hay ningún temor, no hay por donde entre\apnota{9}{U intrat.} la caridad (da una comparación), como vemos\apnota{10}{M vitemus, Mc videmus.} que por la cerda se introduce el lino con que se cose: la cerda entra primero, pero si no sale, no sigue el lino: así el temor ocupa primero la mente, porque entró para introducir la caridad. Alegra aquello del salmista [30 v. 12]\apnota{13a}{M U psalmi.}: Convertiste mi luto en gozo\apnota{13b}{M gaudium etc.}. Y abajo: dicha sentencia está tomada de las Escrituras, pues quien está sin temor no podrá justificarse. Es necesario, pues, que entre primero el temor, por el que venga\apnota{15}{M veniet.} la caridad. El temor es medicamento, la caridad es salud. De esto se sigue abiertísimamente que, siendo la penitencia medicinal y siendo el temor el medicamento según Agustín, y no entrando la caridad sino por el temor, bien se predica que la penitencia comienza por el temor.

Y no acepto lo que el reverendo padre refiere al inicio del resolutorio, la voz del reverendo padre Staupitz\apnota{20}{M stabicii.} como bajada del cielo, que la penitencia comienza por el amor y la dilección. Pues a esta se opone expresamente el bienaventurado Ambrosio, buen médico de las almas, en la epístola a Estudio: donde se prescribe pena, debe haber penitencia de los pecados; donde se dona remisión, gracia es; precede la penitencia, sigue la gracia. Ni hay, pues, penitencia sin gracia, ni gracia sin penitencia. Pues la penitencia debe primero condenar el pecado, para que la gracia pueda abolirlo\apnota{26}{M U illud.}. Aquello

mismo atestigua Crisóstomo en todo el libro sobre la compunción del corazón y en la Homilía 80 sobre la penitencia y en el Sermón 29.

Añado también a Isidoro en el libro 2 sobre el sumo bien, cap. 12: La compunción del corazón es la humildad de la mente con lágrimas, que nace del recuerdo del pecado y del temor del juicio, y en todo el libro sobre la contrición del corazón y en el espejo de los pecadores Agustín no hace otra cosa, más aún exhorta al pecador que quiere hacer penitencia a que recuerde los tres abismos\apnota{32}{P abyssus.}, a saber, de sus pecados, de las penas y de los juicios de Dios. Así el bienaventurado Bernardo sobre el Cantar serm. 16 y en otros muchos lugares, Gregorio en los Morales lib. 2 y 5, e igualmente sobre el cap. 29 de Job, y también en el libro 2 sobre Ezequiel, donde en la Homilía 19 dice: Escrito está: El inicio de la sabiduría es el temor del Señor. Consta sin duda que del temor se asciende al Señor, pero no se vuelve de la sabiduría al temor. El profeta, pues, hablaba de lo celestial a lo ínfimo, comenzó más bien por la sabiduría y descendió al temor. Pero nosotros, que tendemos de lo terreno a lo celestial, contamos los mismos [361] grados ascendiendo, para poder llegar del temor a la sabiduría. El mismo en el lib. 1 sobre Ezequiel\apnota{2}{M Ezechielem.} dijo que hay dos alas, a saber, el temor y la penitencia, que se dice que cubren los cuerpos, esto es, los pecados.

De todo lo cual, en cuanto a una partícula sobre el temor, quiero tener por concluido que bien predicaron los predicadores de nuestro tiempo y del anterior que la penitencia comienza por el temor\apnota{6}{P poenitentiam . . . . incipiunt.}: y que con sus sermones se esfuerzan en sembrar en el pueblo el temor de Dios, para que así sea extirpada la semilla del diablo, según aquello de Orígenes en el lib. 3 sobre Job: bueno es el temor de la pena y de los juicios, al cual si el demonio no lo expulsa, no podrá sembrar la semilla de los pecados. De donde después\apnota{9}{U deinde.} es consecuencia que la penitencia no comienza por el amor y la dilección de la justicia\apnota{10}{M incipere a dilectione et iustitia, a. R. 'ab amore'.}, aunque si así comenzase, confesaría\apnota{11}{M inciperer faterer, U inciperet fateor.} que sería más laudable y más perfecta que la que comienza por el temor de la pena\apnota{12}{M U inciperet.}. Pero nuestra fragilidad no lo tolera, a la cual condescendiendo el Señor Jesús y los predicadores predican el temor como grado por el que consigamos el verdadero amor de la justicia.

El D. D. Martín.

No es este el modo de entender o interpretar felizmente las divinas Escrituras, si de diversos lugares se entresacan diversos dichos sin tener en cuenta ni la consecuencia ni la comparación. Más aún, este es el canon vulgadísimo de errar en las sagradas letras. Conviene, pues, al teólogo, si no quiere errar, poner ante los ojos toda la Escritura y comparar lo contrario con lo contrario, y como

los dos querubines de rostros opuestos\apnota{22}{U Cherobin statt Cherubim; M fultibus, verbessert vultibus.} hallar el consenso de una y otra diversidad en medio del propiciatorio; de otro modo, el rostro de cualquier querubín desviará lejos el ojo que lo sigue del propiciatorio, esto es, de la verdadera inteligencia de Cristo.

Por eso nunca me pareció el egregio D. D. más alejado de las sagradas letras que hoy, y tanto más cuanto que al concluir confiesa que sería más laudable y más perfecta la penitencia si comenzase por el amor de la justicia, como yo sostengo, que si comenzase por el temor de la pena\apnota{27}{fehlt in M.}. Como si no hubiese que esforzarse por hacer laudablemente y, como dice Juan, dignos\apnota{28}{M dignus.} frutos de penitencia. Pues esto no lo acepto en absoluto: que dice que nuestra fragilidad\apnota{29}{P fragilitarem.} impide que comencemos la penitencia por el amor de la justicia. Si ha de obrarse según nuestra fragilidad y exponerse la Escritura, nunca haremos penitencia, sino que de día en día seremos peores. Por lo cual, antes de responder a sus objeciones, primero me declaro\apnota{33}{M U primo declaro (M declaram).} yo mismo.

Toda buena vida es necesario que se instituya por alguna ley; por eso la ley es principio de la penitencia y\apnota{34}{fehlt in M U.} de cualquier buena obra. Por lo cual en\apnota{35}{M U et in.} el penitente ante todo es preciso que o se revele o se sugiera la ley contra la que obró y según la cual debe obrar. Pero manifestada la ley o revocada a la memoria, pronto sigue el aumento del pecado si falta la gracia, porque naturalmente la voluntad odia la ley, como son expresos testimonios de Pablo a Rom y a Gál. Rom\apnota{39a}{fehlt in U.} 5 [v. 20]: La ley entró para que abundase el delito, Gál 3 [v. 19]\apnota{39b}{P M U Gal. 5.}: La ley [362] fue puesta por causa de las transgresiones. Por eso el bienaventurado Agustín en el libro\apnota{1}{fehlt in U.} sobre el espíritu y la letra c. 3: pues el libre albedrío antes de la gracia para nada vale sino para pecar, no para hacer penitencia, como dice el D. D. (pues esto es pelagiano). Sigue Agustín: Y cuando comience a no estar oculto\apnota{4}{M eum coeperit cognoscere, Mc 'non latere', U cum coeperit cognosci se non latere.} lo que ha de hacerse, si el Espíritu Santo no difunde la caridad en nuestros corazones, no se ama, no se recibe, no se vive bien. Allí dice clarísimamente Agustín que la ley de Dios no puede amarse si no se ha recibido la gracia del Espíritu\apnota{7}{U spiritus sancti.}. Pero si la ley no se ama, su contrario el pecado no se odia. Luego es imposible hacer penitencia antes de la dilección de la ley. Esto es lo que quiere el apóstol en Rom 4 [v. 15]: La ley obra ira, esto es, muestra el pecado, pero no da la gracia para que el pecado sea odiado. Por eso permanece el odio de la ley y el amor del pecado, por mucho que el hombre sea sacudido por increpaciones extrínsecas\apnota{12a}{U fornisecas, Uv forinsecas.} o intrínsecas con temor servil

de esclavo\apnota{12b}{fehlt in M.}. Pues aunque se abstenga de la obra del pecado, no puede sin embargo abstenerse del amor del pecado. Y esto lo enseña también Cristo en Jn 6 [v. 44], donde dice: Nadie viene a mí si mi Padre no lo atrae. Concedo, pues, que la ley, el recuerdo de los pecados, la contemplación de las penas pueden aterrar al pecador, pero nunca harán\apnota{16}{M U faciunt.} un penitente.

Respondo, pues, a la primera objeción, sobre el hijo pródigo\apnota{17}{M prodago.} en Lc 15, donde comenzó la penitencia por el recuerdo de la magnitud del premio\apnota{18}{M a 'recogitatione' magnitudinis praemii, U recordatione magn. praem.} diciendo: Abundan en panes\apnota{19a}{M U panibus mercenarii.} etc. Digo que este hijo pródigo\apnota{19b}{U filius prodigus.} comenzó verdaderamente por el amor de la justicia, porque vuelto en sí conoció primero el bien y, conocido el bien, entendió su mal. Pero esta conversión suya en sí mismo no la tuvo por su fragilidad ni por el temor de la pena, como el mismo D. D. aquí\apnota{22}{M hic D. D., in U fehlt hic.} dice que el temor de la pena movió después de la magnitud de los premios. Luego no comenzó por la pena y el temor\apnota{24}{M a poeno timoris.}. La tuvo, en cambio, por el Padre que interiormente atrae e inspira la dilección de la casa paterna, cuando dijo: Cuántos\apnota{25}{M quanto.} jornaleros en la casa de mi padre. Pues habiendo vivido\apnota{26}{M dixisset.} antes en pecados, ni conoció ni amó el bien, ni odió el pecado, aunque no podía ignorar el pecado. Luego fue necesaria otra afección, esto es, el amor del bien.

En segundo lugar adujo a Juan Bautista, Lc 3 [v. 7], que habría comenzado por el temor diciendo: ¿Quién os mostró\apnota{30}{M monstrabit.} huir etc.? Respondo: una cosa es predicar la penitencia y otra comenzar la penitencia; una cosa es predicar la buena obra y otra comenzar la buena obra. El predicador amonesta, aterra\apnota{32}{P movet terret, M monet terrat.}, atrae etc., pero nada sigue si la gracia no mueve la voluntad.

Lo mismo digo al Sal 89 [v. 33]\apnota{34a}{U psal. 58.}: Visitaré con vara sus iniquidades\apnota{34b}{iniquitates eorum fehlt in M.} etc. Puede ser golpeado el pecador, pero si la gracia no coopera, nada se aprovecha\apnota{35}{M proficit.}, como en Jer 5 [v. 3]: [363] Los golpeaste y no les dolió\apnota{1}{P M U percussi eos et non doluerunt (M U voluerunt, Uv = P).}; y en Is 1 [v. 6]: Desde la planta del pie hasta la coronilla no hay en él sanidad: ¿sobre qué os golpearé\apnota{2}{M U percutiam vos.} más? porque nada aprovechó\apnota{3}{M perfecit.} golpeando.

Y bastante me admiro de que el egregio D. D., olvidado de la luz de la naturaleza, Aristóteles, ose establecer el inicio\apnota{5}{M in ipsum, gestrichen, initium.} de alguna virtud en el temor que apremia, cuando aquel con tantas palabras ya recibidísimas en el lib. 3 de la Ética se esfuerza en persuadir: conviene que la obra buena sea voluntaria y se haga con libre voluntad. Pero la voluntad ciertamente o es amor o es hermana del amor. De ahí también se dice que conviene

que la obra buena se haga por el libre albedrío\apnota{9}{U arbitrum.}. Añado también aquello: que Cristo nunca obligó con temor a los pecadores a la penitencia\apnota{11}{M poenitentias.}, sino que atrajo suavemente a cuantos llamó, como a Zaqueo, a la Magdalena, a los apóstoles y a todos. Como también dice en Jer cap. 31 [v. 3]\apnota{12}{cap 31 fehlt in P M; U cap. 2.}: Con perpetua caridad te amé, por eso te atraje compadeciéndote\apnota{13}{M U miserans tui.}.

Digo, pues, que el temor del Señor ciertamente es necesario, pero filial, porque sin amor es imposible soportar la conversión de sí en que es aterrado, contrito y humillado el pecador, según aquello de 1 Re 2 [1 Sam v. 6]: El Señor conduce a los infiernos y saca de nuevo. Creo, además, que la sentencia del mismo egregio D. D. y de todos los doctores escolásticos está conmigo y pugna contra sus objeciones, pues todos convienen en que la contrición debe hacerse en caridad, si ha de ser buena y meritoria la penitencia. Lo cual ciertamente entiendo yo que la contrición se hace moviendo e imperando la caridad, para que\apnota{20}{M quod.} así sea voluntaria, gozosa, amorosa la penitencia. Por eso, aunque\apnota{21}{M licet, Mc a. R. dicit.} Juan increpase a los judíos y los aterrase, no por eso se sigue que los penitentes comiencen también\apnota{22}{M etiam incipere.} por el terror. O si comienzan por el terror, si no accede la gracia, son más verdaderamente hipócritas que penitentes.

Añadió el D. D. la razón de que Dios administra así\apnota{25}{M sic, gestrichen, deus res sic.} las cosas, que conduce lo ínfimo por lo intermedio a lo superior, lo cual aplicó\apnota{26}{M appavit, Mc aptavit.} al temor, queriendo que el pecado sea lo ínfimo, el temor lo intermedio, el amor lo superior, lo cual paso por alto y no acepto.

Dice también que es de perfección si el hombre puede alzarse con el mero\apnota{29}{M U posset se mero (U mere).} contemplar a Dios y el amor de la justicia a alcanzar la gracia. Pero ¿quién es este, y lo alabaremos? Respondo: ni por temor ni por amor puede el hombre alzarse\apnota{31}{M U homo erigere.} a alcanzar la gracia, sino que la gracia previene y mueve al mero contemplar a Dios y al amor de la justicia.

A Agustín, que enseña que el temor precede antes de la gracia y, entrando la caridad, es expulsado\apnota{35}{M expellit.} el temor, digo: si rectamente se entiende, lo admito, esto es, que la penitencia aún no está comenzada cuando el temor precede a la caridad. Pero entrando la caridad se comienza la penitencia, esto es, el amor de la justicia y el odio del pecado. Pero si la caridad no entrase\apnota{38}{M intret, Mc 'intraret'.}, el temor no obraría sino mayores pecados. Por lo cual aquella semejanza de la cerda\apnota{39}{M cetae.} y el lino me parece probada, con tal de que no se entienda que el temor induce la caridad, pues por las palabras del mismo D. D. entiendo que si no sale la cerda\apnota{40}{M ceta.}, [364] no sigue el lino, esto es, si el temor, que impide la verdadera penitencia, no es expulsado\apnota{2a}{M pellat.} por la caridad que entra, nunca verdaderamente se hace penitencia, por así decir\apnota{2b}{ut sic dixerim fehlt in M.}.

A Ambrosio, aducido en la carta a Studium —que la penitencia precede y la gracia sigue— y lo demás, igualmente a Isidoro, Acerca del sumo bien, etc., igualmente los tres abismos\apnota{4}{etc. fehlt in U; P abyssus statt abyssos.} —los pensamientos de los pecados, de las penas, de los juicios\apnota{5a}{M U recogitandos statt recogitationes; poenarum, iudiciorum fehlt in U.}—, a Bernardo sobre el Cantar, a Gregorio en los Morales\apnota{5b}{M Greg. in moral Bernh. super cantica (so!)} y a los demás padres, que enseñan que desde el temor\apnota{6}{M a timore, Mc 'per timorem'.} se ha de ascender a la sabiduría y comienzan la penitencia por el temor, los admito de buen grado, pero no los entiendo contra el apóstol Pablo, que enseña acerca de la ley y del temor de la ley. Digo también yo que, tenida la caridad, el hombre es movido a la vez al temor de Dios, y así comienza\apnota{9}{M incipit.} la penitencia por el temor en la caridad; de otro modo permanece firme la sentencia de que el temor tiene pena y\apnota{10}{fehlt in U.} no obra el bien, sino que odia la ley.

No concluye, pues, el egregio D. D. que la penitencia no comience por el amor de la justicia, por bueno que sea el temor de la pena por su origen\apnota{13}{P M ex Origene.}. Pues no ha de hacerse penitencia por temor de la pena, sino por temor de Dios, porque aquél es siervo, que no permanece en la casa, y éste, en cambio, hijo y heredero. Por eso tampoco admito, sobre aquello de Prov 1 [v. 7]: El principio de la sabiduría es el temor del Señor, la interpretación referida al temor de la pena, que antes de la gracia atormenta al hombre infructuosamente, pues dice expresamente: temor del Señor, no temor de la pena. El temor de la pena es más bien principio de la insipiencia. Vea, pues, el egregio D. D. de no confundir en un solo caos el temor servil y el filial\apnota{19}{U timorem filialem.}, no sea que se cierre para sí mismo\apnota{20}{U sibi ipse.} la inteligencia de las Escrituras y de los padres.

El D. D. Eck.

El reverendo padre intenta escabullirse con afeite de palabras de tan expresas autoridades de la sagrada Escritura y de los santos padres. Y para persuadiros de esto, se ha atrevido a decir que yo estaba muy lejos de la inteligencia de la sagrada Escritura, buscando aquí un subterfugio acerca de los querubines que se miran mutuamente. Pero juzguen, a quienes compete, cuál de los dos siente\apnota{26}{M recte.} más rectamente acerca de la sagrada Escritura. Sin embargo, para excluir sus soluciones del todo impertinentes, aduce con frecuencia dos cosas, de las que en el sermón recitado no recordó\apnota{28}{U secundum viam comm. non memini.} el camino común, ni ningún predicador, que yo recuerde\apnota{29}{M meminerit.}, las negó, y también este temor, en cuanto dispone a la verdadera penitencia, es prevenido por la inspiración divina. Pues es indudable entre los cristianos, contra la perfidia de Pelagio\apnota{31a}{U Christianum statt Christianos; M Pelagii, Mc Pelagianorum.}, que el inicio de nuestra salvación\apnota{31b}{M initium nostrae salutis.} lo tenemos por inspiración de Dios. Por lo cual no fue necesario aducir esto ni por ello reprender el modo de predicar.

En una cosa, sin embargo, me parece que se engaña por equivocación: en que cree que esta gracia, con la que Dios previene los corazones de los hombres moviéndolos, es la caridad, siendo así que es otro don gratuito de Dios. Y con esto queda resuelta la sentencia del bienaventurado Agustín, en El espíritu y la letra, contra los malvados pelagianos.

Lo otro, lo que dice de que el temor es infructuoso si no se añade la caridad, ¿quién, ni entre los escolásticos ni entre los predicadores, lo negó jamás, cuando todos siguen la sentencia del apóstol Pablo a los corintios acerca de la caridad?

[365] Además, en cuanto a que dije que es más laudable la penitencia que comienza por el amor, asumió que nosotros debemos hacer ya\apnota{2}{M U illam.} conforme a aquello de Juan: Haced dignos frutos de penitencia [Mt 3,8]. Pues si hubiera de obrarse según nuestra fragilidad, nunca haríamos penitencia. Digo que también nosotros podemos hacer dignos frutos de penitencia, aun si comenzamos por el temor y llegamos a la caridad.

Y me extraña que el reverendo\apnota{5}{fehlt in M U.} D. padre quiera hacernos ángeles y olvidados de nuestra fragilidad\apnota{6}{U oblitus.}, cuando el bienaventurado Gregorio, en el lib. 2 sobre Ezequiel, declaró expresamente\apnota{7}{M declaverit.} que el profeta descendió de la sabiduría al temor, mientras que nosotros por los mismos grados, ascendiendo, llegamos del temor a la caridad, del mismo modo que muchas otras cosas se perdonan a nuestra fragilidad.

Para declarar su pensamiento muestra cómo la dilección de la ley debe preceder a la penitencia: porque por la ley abundó el delito [Rom 5,20], y a Gal 3 [v. 19] etc.\apnota{11}{P M U ad Gal. 5 etc. (etc fehlt in M).}. Digo: es verdad, como atestigua san Agustín contra Fausto, lib. 22, que el pecado hecho o dicho\apnota{13a}{U dictum, Uv delictum.} contra la ley no sería\apnota{13b}{M esset.} pecado, si no hubiera ley que lo prohibiera, pero así como por la transgresión de la ley abundó\apnota{14}{M U abundat.} el delito, así por la observancia de la ley aumenta el mérito. Por lo cual todos esos subterfugios no contribuyen en nada al propósito: y que las soluciones dadas son contra la inteligencia de la sagrada Escritura\apnota{26}{(Sin definición en el aparato impreso; la llamada 26) del cuerpo carece de correlato y la nota 16 carece de llamada — probable errata de numeración del editor.)}, lo manifestaremos después de comer.

A la hora segunda de la tarde se continuará la disputa.

A la hora segunda, pues, del mismo día

El D. D. Eck prosiguió su réplica comenzada hoy\apnota{20}{Die Zeitbestimmung fehlt in M, dafür nur Prosequar replicam; U hora secunda continuavit Eccius semiplenam orationem, a. R. Prosequamur hodie incepta.}.

El reverendo padre respondió a lo del hijo pródigo\apnota{21}{M prodogo.} que había precedido el amor de la justicia, conforme a aquello: «Y vuelto en sí», porque había tenido memoria de la casa paterna. Pero esto no deshace la objeción, porque aquella conversión

fue hecha por consideración de la pena, pues no tenía a nadie que le saciara con algarrobas, ya que mientras duró el dinero no se volvió en sí, sino que, hambriento, vuelto en sí etc.\apnota{25}{U est conversus; etc. fehlt in M U.}. Se añade que el bienaventurado Agustín dice que sus palabras fueron palabras de quien meditaba la penitencia, no de quien la hacía aún. Pues si vuelto en sí lo hubiera sido por amor de la justicia, ya habría comenzado la penitencia, contra Agustín. Por lo cual su solución no quebranta el propósito y se opone manifiestamente a Basilio.

A Juan respondió que una cosa es comenzar y otra predicar, y en otro lugar dijo que los predicadores aterrorizan, pero no hacen penitente por el tormento de las penas\apnota{31}{M cruciatu statt per cruciatum.}. De nuevo la solución es vana, porque si Juan predicó, ciertamente pretendía el fruto de su predicación, y porque les insinuaba\apnota{33}{M insinuebat, U insinuavit.} el temor, es señal de que de ese modo se inicia la penitencia. Y nuestros predicadores predican cómo comienzan, Dios lo sabe.

Además, en cuanto a que atribuye el inicio de la penitencia\apnota{36a}{P nuntium.} a la gracia, también estos\apnota{36b}{U nostri.} predicadores y doctores nunca lo negaron, porque Dios previene aspirando\apnota{37}{M U praevenit, Uv = P.}.

Además, no quiere aceptar a Aristóteles en las escuelas teológicas, y sin embargo se atreve a oponérmelo. Pero digo que es sentencia constante de los paganos, [366] que deteniéndose en el temor de la pena no hay virtud perfecta\apnota{1}{M in timore (so!) . . . . . perfectum retuli.}, conforme a aquello: Los malos odian pecar por miedo de la pena, los buenos odian pecar por amor de la virtud\apnota{2}{U amore etc.}. Ni la obra de la virtud debe ser coaccionada, sino libre. Pero una cosa es que el libre albedrío sea inducido, y otra\apnota{4}{M U aliud est.} que el libre albedrío sea coaccionado. Confieso que Cristo llamó con palabras benévolas a los apóstoles, a Zaqueo, a la Magdalena. Pero que él mismo llama alguna vez con más dureza, nos lo indica el hecho de Pablo, acerca del cual dice Agustín, 23 q. 4 can. Quis nos\apnota{7}{fehlt in U.}: ¿A quién hizo Cristo violencia? ¿A quién coaccionó? Y en el evangelio, acerca de la llamada a la cena, dice: Obliga a entrar. Como lo prosigue con amplias palabras el bienaventurado Gregorio en la homilía. Calle, pues, ruego, el reverendo padre, que aduce solo\apnota{10}{W. A. verbessert asserens.} la llamada benévola.

Además, juzga que solo el temor filial es necesario al penitente, de lo cual me admiro sobremanera, cuando hoy me amonestó a la verdadera\apnota{12}{M verarum.} inteligencia de las Escrituras, cosa que él mismo primero no se corrigió, pues también el bienaventurado Agustín en el lugar hoy alegado sobre Juan habla del temor servil, que la caridad echa fuera\apnota{14}{U super Johanne; P M foris.}, y de él dice que habla el Sabio: El inicio de la sabiduría es el temor del Señor [Prov 1,7]. Y así entiende\apnota{16a}{M U intellexit.} del temor servil la glosa sobre Sal [111,10]: El inicio de la sabiduría etc.\apnota{16b}{fehlt in M.}: El miedo del juicio es la puerta de la conversión a Dios, y es glosa de Casiodoro. Así la glosa sobre aquello del apóstol a Rom 8 [v. 15]: Pues no recibisteis

espíritu de servidumbre en temor, dice que el temor servil es bueno y viene de Dios. Así Pablo atestigua que aquél es bueno, 1 Tim 5 [v. 20]\apnota{20}{P M U 1. Tim. 4.}: Al que peca repréndelo delante de todos, para que también los demás tengan temor. Así dice Agustín sobre los Salmos, y lo refiere\apnota{22a}{M enarrat.} el Lombardo en el 3: Temor servil es cuando por temor de la gehena\apnota{22b}{M jehennae.} el hombre se contiene del pecado, porque teme la presencia del juez y las penas. Y más abajo: Bueno es este temor y útil, aunque insuficiente, por el cual poco a poco se hace costumbre de la justicia. También Beda expone: El inicio de la sabiduría es el temor del Señor, a saber, el servil.

Por lo cual busque otro patrocinio y dé otra solución, pues esa\apnota{27}{M U quoniam; M istam, verbessert ista.} no concuerda con los santos padres. Y si no quiere atenerse a los dichos de los santos, tendrá en contra suya los pasajes de la sagrada Escritura. Pues mientras el temor del Señor\apnota{28}{M U timor domini.} es inicio de la sabiduría, todavía la caridad echa fuera\apnota{29}{P M foris.} el temor, pero el temor filial según David permanece para siempre y por los siglos de los siglos. Por lo cual uno es el temor con el que se inicia la sabiduría, y otro el que permanece al crecer. De donde, cuando el reverendo padre niega que el temor introduzca la caridad, contradice expresamente a san Agustín, que en el lugar alegado dice todo esto\apnota{33}{M agit.}: que el temor servil es inicio de la sabiduría y que por la caridad es de nuevo expulsado. Suscribe la sentencia agustiniana el melifluo Bernardo, Serm. 52 sobre el Cantar, que por brevedad no aduzco. Pero aún digo más: que la caridad no expulsa el temor como [367] una cosa incompatible\apnota{1a}{M incompatibile.} a otra, sino que con el progreso, al aumentar la gracia, disminuye\apnota{1b}{M minuit.} el temor servil, y aumentada la gracia se quita del todo el temor.

Ni importa lo que hoy ponderaba el D. padre: que el temor del Señor, no el temor de la pena, es inicio de la sabiduría, puesto que tanto el temor de la pena como el temor de reverencia son ambos del Señor y ambos vienen del Señor. Solo son diversos los medios de temer, como atestiguan Beda y otros, aducidos por el Lombardo, además de Agustín antes citado. Y está bastante claro por las palabras de Cristo, si nada más tuviéramos, Mt 10 [v. 28]: No temáis a los que matan el cuerpo, pero no pueden matar el alma, sino temed más bien al que puede perder alma y cuerpo en la gehena\apnota{9}{M jehennam.}. Pues si el temor servil fuera condenable, como dice el reverendo padre, ¿por qué nos invitaría Cristo a él? Y que así lo entendieron Agustín y los demás, lo someto al juicio de aquellos a quienes competa.

Además dice que todos los escolásticos concuerdan en que la contrición no es de ninguna utilidad si no se hace en caridad. Esto es ciertamente verdad, pero según Agustín a la caridad no se llega sino por el temor. El temor es medicina, la caridad salud. De donde, no añadiéndose la caridad, todos

bien saben que aquel temor es insuficiente, por eso lo ponen en el camino\apnota{17}{M illum in via constitutum illum (letzteres gestrichen).}, no en el término.

Pero\apnota{19}{M U ad.} la razón tomada del bienaventurado Dionisio la pasó por alto, mas no la abrió el D. padre: siendo el pecado lo ínfimo y la caridad lo superior, cuál sea el medio para llegar\apnota{21a}{M perveniendo.} del pecado a la caridad, y ciertamente no puede dar\apnota{21b}{M nequid.} otro que Basilio, Beda, Agustín, Bernardo y los demás\apnota{22}{et ceteri fehlt in M U.}: el temor.

Después\apnota{23}{M Dein.} asume que la gracia previene al temor y al amor; si habla de la gracia de la moción divina, con la que Dios nos previene aspirando, ciertamente lo confieso\apnota{24}{M fateor (Mc 'accipio') tunc, U fateor quidem tunc.}; pero si habla de la gracia de la caridad, aquello no lo acepto, porque de aquella\apnota{25}{M tunc illius.}, que rectamente\apnota{26}{M recta.} se llama sabiduría en las Escrituras, el inicio es el temor del Señor, aunque ese temor sin caridad sea infructuoso.

Además, en cuanto a que admite a Ambrosio, Gregorio y los demás doctores, pero no contra el apóstol acerca del temor de la ley, no sé qué nieblas difunde\apnota{29}{M oppugnat.}. Diga claramente si\apnota{30}{U aut.} esos santos padres por mí alegados repugnan o no a los dichos del apóstol. Si repugnan\apnota{31}{M repugnat.} al dicho del apóstol, muéstrelo, cosa que juzgo que nunca podrá hacerse, pues fueron ejercitadísimos en las sagradas Escrituras y llenos del Espíritu Santo, y entendían al apóstol Pablo tan bien como nosotros. Si no repugnan a los dichos de Pablo, aténgase a la sentencia de ellos y no repruebe a los predicadores que enseñan el susodicho modo de la contrición y de la penitencia, ni excluya de los hombres pecadores ese temor servil, útil y como medio necesario.

No le vale que estime que, tenida la caridad, la mente es movida al temor, puesto que esto sería retroceder a modo de cangrejo\apnota{38}{M retrogradi.} contra la sentencia del bienaventurado Gregorio, lib. 2, homilía 19 sobre Ezequiel\apnota{39}{U super Ezechiele.}. Consta sin duda alguna que del temor se asciende a la sabiduría, pero no se vuelve de la sabiduría al temor. Tiene, pues, el temor pena, que al crecer la caridad disminuye y es totalmente absorbida.

[368] Paso por alto, por la estrechez del tiempo, otras cosas contenidas en el mismo sermón, que sin embargo habrían de discutirse y examinarse por extenso, si no estuviera prefijado el día de hoy para terminar esta materia, como es el punto de la confesión sellada de los pecados en especie\apnota{4}{M de sigillato, sonst = P; U de sigillata peccatorum confessione in speeu.}, según recuerda Crisóstomo sobre el cap. 12 de Mt; el punto de que los escolásticos pongan tres partes de la penitencia, a saber, contrición, confesión y satisfacción, las cuales puso antes Crisóstomo en el sermón 29 sobre la penitencia; el del doble modo de confesar, lo grave al sacerdote y lo demás a Dios, con otros puntos, todo lo cual

puedan\apnota{8}{M possint.} tratarse y examinarse por extenso. Pero, echado ya el fundamento, someto todo el sermón a la vez y lo que aduje por mi parte al juicio de los jueces que han de elegirse.

El D. D. Martín.

Esperaba que el egregio D. D. refutara mis soluciones, sobre todo lo que había aducido como fundamentos tomados de Pablo\apnota{13}{U induxeram; M ea quae pro fundamento ex Paulo aduxeram (so!) 'quae pro fundamentis' (Mc 'eam quam' am Anfang).} acerca de la ley, que antes de la caridad no obra sino ira y aumenta el pecado —callo que disponga a la gracia—, del mismo modo\apnota{15}{fehlt in U.} que el temor servil, fruto de la ley, obra ira y aumenta el pecado. Pero él, saltando valientemente estas cosas, recantó la misma cantilena, y pretende persuadirnos de que esto no es afeite\apnota{17}{U succum.}, sino la médula de la Escritura; discurriré punto por punto.

Primero dice\apnota{19a}{fehlt in M, U proinde dicit.} que no fue necesario aducir que el inicio de nuestra salvación\apnota{19b}{M U nostrae, Mc 'nobis'.} nos viene de Dios que inspira, donde entiende que también se da el temor que dispone. Respondo: todo esto es claramente erróneo por la palabra de Pablo, que dice que es imposible que la ley se cumpla, más aún, que los pecados no aumenten, si el Espíritu Santo no difunde la caridad en nuestros corazones [Rom 8,3 ss.]. Convenía que esas autoridades tan expresas fueran refutadas, y aquella del bienaventurado Agustín, cap. 3 de El espíritu y la letra: «Cuando\apnota{24}{M in libro de spir. et lit. c. 3.; M U cum statt dum.} comience a no estar oculto cómo se ha de vivir, si no se da la gracia, no se recibe, no se vive bien». Vaya, pues, con su temor servil, que no obra sino odio de la ley y de Dios, y con injuria es llamado «el que dispone a la gracia».

Encuentra también una equivocación de la gracia: que una cosa es la caridad y otra el don con que primeramente\apnota{29a}{U primum.} somos movidos. Pase esa\apnota{29b}{M ista, U haec.} distinción, nada aporta al propósito, es elusión\apnota{30}{M delusio.} de las palabras de Pablo. Clarísimo es el texto de Pablo: a no ser que la gracia y la caridad nos hagan amar la ley, la ley obra siempre ira. Pero ese don con que primeramente somos movidos no amará\apnota{32}{M primo movemur, non diluet.} la ley, sino la caridad del Espíritu.

En tercer lugar concedió que el temor es infructuoso sin la caridad, y que esto nadie lo negó. Respondo: ¿por qué, pues, enseñan penitencias infructuosas y resisten a mí, que enseño la fructuosa? Dejo constancia de que esta confutación está del todo subvertida\apnota{35}{M docentem Relinquo hanc omnino esse superstem, U docenti Relinquo hanc esse omnino subversam.}, pues en tantos opúsculos Agustín\apnota{36a}{W. A. Augustini (unnötige Änderung).} —y aunque Augusto no lanzara tantos rayos\apnota{36b}{M voluminibus, Mc 'fulminibus'.}—, sin embargo Pablo solo concluye que cualesquiera obras antes de la caridad son pecados y condenables [369] y

no disponen a la gracia. Y por estos dichos puede responderse a casi todas sus réplicas. Él, según su costumbre, atempera los dichos de la divina Escritura a los dichos de los padres, más aún los arrastra a la inteligencia que él tiene en los padres, cuando más bien los escritos de los padres deben cotejarse y juzgarse según los dichos de la Escritura.

En cuarto lugar no acepta lo que asumí: que debemos hacer aquella penitencia más laudable, sino que dice que también podemos hacer dignos frutos de penitencia si comenzamos por el temor. Esto, según Agustín\apnota{7}{M U beatum Aug.}, lo expongo así: si añadiéremos pecados a pecados, puesto que, como ya se ha dicho muchas veces, todo lo que se hace antes de que la voluntad sea sanada\apnota{9}{M hec gratia (?), gestrichen, sanatam.} por la gracia es fruto del árbol malo, que no puede hacer frutos buenos\apnota{10}{M U bonos fructus.}. Luego por el temor servil nunca se llega a la caridad.

En quinto lugar, apoyado en Gregorio sobre Ezequiel, que enseña que desciende de la sabiduría al temor, mientras que nosotros al contrario ascendemos\apnota{13}{U descendere, Uv ascendere.} del temor a la caridad, respondo así: que el divino Gregorio no debe excluir ni excluyó la caridad del temor que inicia la penitencia, como\apnota{15}{P sic.} entiende el D. D.; como tampoco de la caridad se excluye todo temor servil, sobre todo en esta vida, pues oficio de la caridad es expulsar el temor servil, y ello durante toda la vida, e introducir el temor filial.

En sexto lugar, a las autoridades de Pablo —que la ley fue puesta por causa de las transgresiones y aumenta el delito— dice que nada aportan al propósito y que yo busqué subterfugios, aduciendo a Agustín contra Fausto, lib. 22: que el pecado es lo dicho o hecho contra la ley de Dios etc.\apnota{21}{fehlt in U.}; no sé qué quiere el D. D. Aquí no se disputa qué es el pecado, por eso con palabras superfluas se ha perdido este tiempo. De esto trata el objeto: que por la autoridad de Pablo el temor servil no puede ser sino pecado y aumentar el pecado por la virtud\apnota{24}{M virtutum.} de la ley, si falta la gracia. Y así las autoridades de Pablo no son subterfugios, sino rayos que trituran la sentencia del D. doctor hasta el polvo. No soy, pues, yo quien de hombres hace ángeles, olvidado de la fragilidad, sino que enseño a no hacer dioses de los pecadores, cuando se olvidan de su fragilidad, por la cual antes de la gracia no pueden sino hacer el mal.

En cuanto a que dice que la conversión del hijo pródigo fue hecha por consideración de la pena, lo niego. A la prueba de que no tenía a nadie que le saciara con algarrobas, digo: si dentro no hubiera sido atraído y vuelto en sí, antes habría muerto de hambre que haber vuelto\apnota{32a}{M rediisse.}, para que permanezca la sentencia: Nadie viene a mí si mi Padre no lo atrajere [Jn 6,44]\apnota{32b}{M eum traxerit.}.

En cuanto a que Agustín dice\apnota{33}{fehlt in M.} que sus palabras fueron palabras de quien meditaba la penitencia, no de quien la hacía aún, espero que eso favorezca mi causa, pues meditar

la penitencia, sobre todo desde lo hondo del corazón, es ciertamente comenzar la penitencia. Luego aquel «aún no la hacía» debe entenderse de la perfección hacia fuera.

Lo mismo se dirá a Basilio\apnota{37}{fehlt in M.}.

Acerca de la palabra de Juan Bautista dice que es vana la solución de que yo dijera que una cosa es enseñar la penitencia y otra comenzarla. Respondo: salva [370] la reverencia del egregio D. D.\apnota{1a}{U egre: d. doct.;, W. A. egregius dom. doctor.}, no parece entender a Pablo ni haber conocido la virtud\apnota{1b}{M virtutis.} de la ley. Pues la ley enseña ciertamente lo santo\apnota{2}{P facta.}, justo y bueno, pero solo la gracia comienza, hace y perfecciona esas cosas. Luego, aunque\apnota{3}{etiam fehlt in M, von Mc hinter si nachgetragen.} Juan hubiera enseñado que el temor es inicio de la penitencia, no por eso se sigue que la penitencia comience por el temor. Del mismo modo que si insinúo cualquier obra buena por terrores y amenazas, no por eso la obra buena\apnota{6}{U opus bonum.} comienza por el terror y las amenazas, sino por la caridad.

Se indigna de que le haya objetado a Aristóteles en las escuelas teológicas, al que sin embargo yo no acepto. Confieso que es menor de lo que valdría en la escuela teológica, pues fue seductor de los doctores escolásticos. Pero quise satisfacer a mi conclusión: que contra su santo Aristóteles desvarían quienes comienzan la penitencia por el temor y no por la libre voluntad. Pues la distinción\apnota{11}{M distinctio.} del libre albedrío coaccionado e «inducido» la paso por alto. Pues no es verdad que el libre albedrío sea jamás coaccionado al bien ni pueda siquiera ser coaccionado. Pero si es coaccionado, es arrastrado a lo diverso\apnota{14}{M raptum.} y odia su coacción. Solo por la gracia es atraído, esto es, se hace verdaderamente libre, como enseña Augusto contra Juliano, lib. 2\apnota{15}{M contra Pelagianum (lib. 2 fehlt); U contra Jul. lib. 2 et in multis locis.}.

Donde dije que Cristo llamó a los apóstoles y a la Magdalena con blando atractivo, contrapone que Pablo fue llamado con más dureza, y que en el evangelio se mandó al siervo que obligara a entrar. Primero, me admiro bastante de que, cuando solemos aducir el ejemplo de Pablo para defender la gracia, ellos se nos escabullen diciendo que aquello fue hecho milagroso y no hace regla. Aquí, sin embargo, el egregio D. D., como teniéndolo no por milagro, lo aduce como regla. Pero esto lo dejo. Digo que Pablo no pudo convertirse de corazón por dentro si no lo atraía la gracia, como enseña el bienaventurado Augusto contra las cartas de los pelagianos: que la gracia de Dios hace de los que no quieren y resisten, quienes quieren y siguen. Así también se resuelve aquello del evangelio: que el siervo puede obligar con la palabra, pero si el Señor no silba, nada se sigue\apnota{26}{U a. R. Ut Esaias loquitur.}.

A aquello de Agustín, que entiende el temor servil en la palabra de Salomón: El inicio de la sabiduría es el temor del Señor, y aquello de la glosa al Sal 111: El miedo del juicio es puerta de la conversión a Dios, respondo: si está presente\apnota{29}{M assit.}

la gracia; de otro modo el temor servil sin la gracia (pues tampoco Augusto la excluye) no obra sino ira. Conviene, en efecto, que\apnota{31}{M U quod.} Augusto no luche, como en verdad no lucha, con Pablo, que condena todo lo que está fuera de la gracia. La glosa sobre aquello de\apnota{32}{M id.} Rom 8 [v. 15]: No recibisteis espíritu de servidumbre en temor, que dice que el temor servil es\apnota{33}{M servilem timorem esse, U esse servilem timorem.} bueno, más bien la rechazaría, como que habla expresamente contra el texto. Pues el apóstol dice: No recibisteis espíritu de servidumbre en temor, condenándolo. O digo que la glosa\apnota{36}{fehlt in M.} no expone el texto.

Aquello del apóstol a Timoteo [1, 5,20]: Reprende delante de todos, para que también los demás\apnota{37}{U ut de, Uv ut et.} tengan temor, el D. D. lo adaptó al temor servil, cosa que dejo. Yo lo tomo del temor filial, hasta que pruebe otra cosa. Pero en cuanto a que Augusto, recitando el maestro, dice en el lib. 3: Temor servil es cuando por temor de la gehena\apnota{1}{M pro timorem jehennae.} se contiene del pecado, porque teme la presencia del juez y las penas etc.\apnota{2}{M poenae metuat (etc. fehlt); U quo poenitentiam iudicis, quo poenas metuit etc.}: se contiene, digo, solo del pecado externo, pero dentro aumenta el odio de la justicia que amenaza con la gehena. Y más abajo: Bueno es el temor y útil, aunque\apnota{4}{M dicit, U licet et.} insuficiente, por el cual poco a poco se hace costumbre de la justicia. Esto es, a mi juicio, costumbre de desesperar y de odiar a Dios, si se excluye la gracia; pero es verdad si se incluye la gracia.

No es, pues, necesario\apnota{7}{M ergo est.} que yo dé otra solución, a no ser que el egregio D. D. pruebe primero\apnota{8}{M U primum.} que los santos padres hablan del temor servil excluida la gracia, o que la caridad no tiene nada que ver en expulsar el temor servil. Por lo cual lo que en consecuencia adujo acerca de la expulsión del temor, del inicio de la sabiduría y de cómo expulsa la caridad al temor por progreso mediante el aumento de la gracia, se entiende bastante\apnota{12}{P intelligitur.} por lo precedente.

Resta la robustísima autoridad de Mt 10 [v. 28], que juzga que ella sola basta: No temáis a los que matan el cuerpo, sino temed etc.\apnota{14}{sed timete eum fehlt in M, U etc. statt qui occidunt corpus.}; [que] por eso el temor servil, al que Cristo nos invita, no es condenable. Respondo primero: si esto quiere el D. D., se contradirá con lo precedente, donde el temor servil se dice insuficiente; por eso es igualmente absurdo\apnota{17a}{M absordum.} decir que Cristo nos enseñó cosas insuficientes\apnota{17b}{U insufficientiam.}. Digo, sin embargo, yo que allí no se enseña\apnota{18}{U domini.} el temor servil, porque también el temor filial teme ofender a Dios y separarse de él. Y aun si hablara del mero temor servil, todavía no se entiende excluida la gracia, sino más bien incluida, puesto que, con el testimonio del apóstol y de Agustín, toda ley y doctrina es letra que requiere espíritu.

Convino también en que los escolásticos dicen con verdad: que la contrición no vale fuera de la caridad, pero no lo refutó. Permanece, pues, la autoridad de ellos mismos contra ellos: a no ser que refute que la contrición en caridad es lo que comienza a hacerse\apnota{24}{P M incipi et fieri.} por la caridad.

Aquello de Agustín place bien: que el temor es medicina y la caridad salud, el temor a saber en la caridad imperfecta, y la caridad perfecta.

Aquella razón tomada de Dionisio\apnota{27}{M U ex Dionysio sumptam (M sumpta).}, sobre lo ínfimo, lo medio y lo sumo, dice que la pasé por alto\apnota{28}{M U transiisse.} y no abrí cuál sea el medio entre el pecado ínfimo y la caridad superior. Confieso que la pasé por alto de buen grado\apnota{29}{M U transiise.}, por creer que al propio D. D. le era bastante evidente\apnota{30}{M U perspectam.} que esta razón nada aporta al propósito. Pues Dionisio habla de los órdenes ínfimos, medios, sumos\apnota{31}{M U infimis (U infinitis) mediis et summis.}. Pero yo no tengo en absoluto ningún medio entre el pecado y la gracia, como tampoco Cristo, cuando dice: Quien no está conmigo, está contra mí. Y de nuevo: O haced bueno el árbol o hacedlo malo. Pero creo que el propio D. D. afirma eso mismo: que la gracia y el pecado se oponen inmediatamente entre los escolásticos.

No acepta el D. D. que yo dije que la gracia previene al amor y al temor, a no ser que hable\apnota{37}{M loquor.} de la gracia de la primera moción. Respondo: yo siento con el apóstol y Agustín que, si la ley no es amada\apnota{38}{M dirigatur.} (lo cual es propio de la caridad y no de la primera moción), no se vive bien. Luego ni Dios es temido o\apnota{39}{M colitur seu timetur, U timetur nec colitur.} venerado.

[372] Al final instruye contra mí un silogismo cornuto: o Ambrosio, Gregorio y los demás repugnan a los dichos del apóstol o no; si así, que\apnota{2}{M sic sic ut, U si sic quod.} lo muestre; si no, que se atenga a la sentencia de ellos. Respondo y paso por en medio: no repugnan a los dichos del apóstol, si\apnota{4}{fehlt in M U.}, y me atengo a la sentencia de ellos, pero no a la inteligencia eckiana, más aún al error. Pues no excluyen la caridad del temor, sea servil o filial.

Dice el D. D. que es\apnota{7}{fehlt in M.} retroceder a modo de cangrejo lo que dije: tenida la caridad, la mente es movida al temor. Me admiro bastante, sea de aquel escorpión o de aquel cangrejo\apnota{8}{M U istum.}, pues también aquel poeta pagano dijo\apnota{9}{M dixerim.}: Cosa llena de solícito temor es el amor. Como si ignorásemos que el amor es fuente y cabeza de todos los afectos. Pues por eso la pena y el infierno se temen servilmente, porque la vida y el placer se aman puerilmente. Y para extirpar este amor y temor\apnota{12}{M timorem et amorem; U pueriliter et serviliter. Ad hunc timorem et amorem.} se difunde la caridad de Dios, por la que amemos otra vida y temamos otra muerte, esto es, la separación de Dios.

Pasó por alto el D. D. los puntos señalados en mi sermón y todo

discutiendo el sermón\apnota{15}{M et totum sermonem confutandum, U et de toto sermone confutando.}. Digo brevemente: Tiene cálamo y papel, acométalo con confianza; se verá si ha refutado el sermón o si se ha burlado de sí mismo. Esto\apnota{16}{M U vel se ipsum irriserit. Hoc (M haec).} lo dejo al juicio de los que han de juzgar.

El D. D. Eck.

Puesto que por la estrechez del tiempo no puedo responder a esto\apnota{19}{fehlt in M; U iis.} que el reverendo padre ha aducido, me remito sin embargo a los jueces\apnota{20}{fehlt in P, in M gestrichen wegen des folgenden me.}: que nada he pasado por alto, ni fue sentir de los predicadores o de los doctores que la ley se cumpla sin caridad, ni enseñaron algunos una penitencia infructuosa, sino de qué modo por el temor servil se llega a la fructuosa. Y en cuanto a que con magnífico aparato de palabras se ha apartado de la solución de hoy, [dijo el D. D. Martín: no me he apartado\apnota{21a}{(En el aparato impreso numerada 24a — errata de numeración del editor.) U a. R. Martinus non recessi respondit. [M a. R. nur Mar.]}] en la que\apnota{21b}{(En el aparato impreso numerada 24b — errata de numeración del editor.) U quia.} había dicho que el Sabio habla del temor filial, y ahora en cambio admite que habla del temor servil: pero sin excluir la gracia, lo cual ni el texto ni los santos doctores permiten. Pues por sabiduría entiende la caridad, y así\apnota{27a}{M U ita.} el inicio sería antes del inicio. Todo\apnota{27b}{M initium et Bernhardus et Augustinus ita intelligit sic et totus.} Agustín tampoco se afana en cómo el temor servil con caridad sea inicio de la caridad, sino en cómo el temor servil ocupa primero la mente y entra primero y así introduce la caridad, lo cual es también el sentir de Gregorio al reprobar el retorno canceroso, donde el reverendo padre en lugar de Gregorio me trae un escorpión: Es cosa del solícito etc.\apnota{31}{M solliciti certe, U sollicita etc.}. Todos los santos doctores hoy citados quieren, pues, que el temor servil sea inicio de la caridad, según la interpretación dada muchas veces y que suelen dar los predicadores, sobre lo cual me remito\apnota{34}{M U refero me.} a los jueces.

Mañana a la hora segunda, por la fiesta de santa Margarita, disputarán los mismos señores doctores sobre la materia de la absolución de la pena y de la culpa por cualquier simple sacerdote.

[378] Miércoles 13 de julio a la hora segunda de la tarde\apnota{2}{statt Zwischenbemerkung und Zeitangabe M nur In die Margarethae}.

El D. D. Eck.

Acerca de la materia de las conclusiones cuarta y quinta, reverendo padre, puesto que el tiempo que se nos ha fijado es demasiado estrecho, tocando sin embargo el fondo

del asunto tocaremos de paso algunas cosas, a saber, aquello de que cualquier sacerdote absuelve de la pena y de la culpa, contra el uso común de toda la Iglesia, y lo que decís en el sermón vulgar y en la conclusión, que por ninguna Escritura se prueba que la justicia divina exija del pecador alguna pena o satisfacción. Y me apoyo en que en la absolución sacramental la culpa se remite simultáneamente, pero la pena eterna debida al pecado se conmuta en temporal. Esto consta por\apnota{10}{P et.} la tradición de la sagrada Escritura y el uso de los santos padres. Y para ser más breve, esto\apnota{12}{fehlt in M.} lo quiso expresamente Ambrosio sobre Lucas, Jerónimo en el libro 1 contra Joviniano, Agustín en la cuestión 1 de las ocho cuestiones a Dulcicio, y oportunamente Ambrosio sobre Lc 5 atestigua que con la satisfacción se disuelve la pena del pecado.

La Escritura está a la mano: pues el pecado de Adán, aun remitida la culpa, es castigado en toda la posteridad. Por eso, pasando la culpa, permanece la pena. Esta razón la toca Agustín en el tratado 124 sobre Mt y la glosa a 2 Re [Sam.] 12. Igualmente de David, 2 Re 24 [v. 10], donde David se arrepintió y su corazón le golpeó: he pecado, dijo, gravemente en este hecho, pero te pido que traspases\apnota{19}{M U ut transferas.} la iniquidad de tu siervo. Ahora bien, el pecado no se traspasa en cuanto a la culpa, porque, como dice el profeta, la misma alma que pecare morirá\apnota{21a}{M U morte morietur.}. Resta, pues, que el pecado se traspase en cuanto\apnota{21b}{M quantum ad.} a la pena. Así por el profeta Natán se dice a David: El Señor ha traspasado tu pecado. Donde la glosa: Dios borra el delito, pero no lo deja sin castigo. O el hombre penitente se castiga en sí mismo, o Dios, vindicando, hiere junto con el hombre\apnota{24}{Mc ac enim homo . . . ponit.}.

Y para venir a los santos padres, de quienes el uso y la práctica de la Iglesia se ha derivado principalmente desde los mismos\apnota{27}{fehlt in M U.} apóstoles. Pues primero ningún pecado queda impune, y la pena es ordenadora de la culpa. Estas razones las toca Agustín y las refiere Graciano en el can.\apnota{29a}{fehlt in U.} Como los primeros, sobre la penitencia, distinción 1\apnota{29b}{U dist. 7.}. Pues Dios no permitiría que se hiciera el mal, si no hiciera que por la justicia, ordenándolo, estuviera\apnota{30}{M melius dare.} mejor en el mundo. Dice, pues, Agustín en el libro sobre la medicina\apnota{31}{M poenitentia.} de la penitencia: No basta cambiar las costumbres a mejor y apartarse de los hechos malos, si no se satisface también a Dios por lo\apnota{32}{U iis.} que se ha hecho, con el dolor de la penitencia, con el gemido de la humildad, con el sacrificio del corazón contrito, cooperando las limosnas etc. Pues no se dijo que solo os abstengáis de los pecados, sino también de los pasados\apnota{35}{U de peccatis.}, dice, rogad al Señor para que os sean perdonados etc.

[374] Pero el reverendo\apnota{1}{M U Sed si reverendus.} padre, como en el resolutorio y en otras partes, a mi juicio, se esfuerza en hacer que aquella satisfacción se haga a la Iglesia, pero que Dios

[kalligraphisch geschrieben]; de absolutione sacerdotis (sacerdotali?); a. R. circa conclusionem 4 et 5; U Tertia decima Julii Mcccc. XIX. hora secunda. Circa materiam conclusionum quartae et quintae; a. R. Eccius.

no requiera tal pena, y que la que Dios requiera el hombre no la quite; contra ello opongo yo: Primero, que Agustín dice en el Enquiridión\apnota{4}{fehlt in P.}, cap. 71 (habla de la oración dominical): Esta oración borra por completo los pecados mínimos y cotidianos, borra también aquellos de los que se aparta la vida de los fieles llevada criminalmente, pero cambiada a mejor arrepintiéndose. Donde consta que la oración dominical borra los veniales y también los mortales, no en cuanto a la culpa, sino en cuanto a la pena, porque cambiada arrepintiéndose, cuyo cambio en Agustín no puede entenderse sino del cambio de la pena eterna en temporal.

A los que fueron hace cuatrocientos años no los aduzco, como a Guillermo de París, Guillermo de Auxerre\apnota{11}{M Altiodorensem et alii (so!); U non adduco huc Wilhel: Paris: Altisidio:}, que lo afirman a plena boca. Añádase Cipriano en la epístola a Fido, donde dice así\apnota{12}{P Fidium; ubi sic ait fehlt in M.}: Hemos leído tus letras, carísimo hermano, en las que significas acerca de cierto presbítero Víctor, que Terapio, colega temerario, le dio precipitadamente la paz antes de que hubiera hecho plena penitencia y hubiera satisfecho al Señor Dios contra quien había delinquido\apnota{15}{U dederat.}. Aquí san Cipriano no dice que Víctor no hubiera satisfecho a la Iglesia, sino al Señor Dios.

El mismo Cipriano al papa Cornelio dice de los herejes: Se afanan\apnota{17}{M elaboram.} por consumar\apnota{18}{M consumen.} con la malicia del diablo su obra, para que la divina clemencia no cure en su Iglesia a los heridos. Corrompen con fraude de mentiras la penitencia de los míseros, para que no se satisfaga a Dios indignado. Clarísimamente el santo\apnota{20}{M U sacer, Mc 'sanctus'.} mártir increpa a los herejes, que estimaban que la penitencia impuesta por la Iglesia no era satisfacción de Dios.

Y aunque\apnota{22a}{U quamvis (et fehlt).} el reverendo padre sin modestia me reproche\apnota{22b}{M improberet.} muchas veces, como si no tratara idóneamente la inteligencia de las santas Escrituras, sobre lo cual juzguen hoy los jueces designados cuál siente más rectamente acerca de la fe y del sentido de las sagradas Escrituras\apnota{25}{M literarum,}: sin embargo, por cuanto yo veo, todos los tratadistas de las sagradas Escrituras descienden a esta sentencia. ¿Para qué querría la Iglesia imponer tan grave obra\apnota{27}{M U onus.} a los penitentes, si el Señor y cabeza de la Iglesia no la exigiera? Esta razón la toqué en las anotaciones, pero no fue mía, sino de Crisóstomo, como refiere Graciano en la cuestión 26, can. Alligant\apnota{29}{M P alligat, U allegat.}: hombre, a quien impones grave peso de penitencia, o rechazará la penitencia, o, no pudiendo soportarla\apnota{31}{M nequid.}, escandalizado pecará más. Además, aunque erremos imponiendo una penitencia modesta, ¿no es mejor dar razón por misericordia que por crueldad? Donde el padre de familia es generoso, el dispensador no debe ser tacaño. Si Dios es benigno, ¿por qué quiere el sacerdote ser austero? Esto en favor de los confesores.

El mismo en la Homilía 31 dice sobre la epístola a los Hebreos: El pecado con la confesión se minora, pero no se minora en cuanto a la culpa, porque es impío, como dice Jerónimo, esperar de Dios media venia, luego se minora en cuanto a la pena.

[375] Añado a Gregorio, 4 de los Morales, § 49, donde después de muchas cosas dice así\apnota{1}{post multa ita dicit fehlt in M; U cap. 49 post multa ita dicit.}: Pero puesto que Dios no deja sin castigo ningún pecado, o nosotros lo perseguimos llorándolo, o él mismo juzgando, resta que para su enmienda la mente vigile siempre solícita etc.\apnota{4}{fehlt in M U.}.

Esto sintió también el bienaventurado Jerónimo, cuyas palabras están transcritas en el can. Mensuram, sobre la penitencia, dist. 1: Los cánones no prefijan suficientemente la medida del tiempo en hacer penitencia para cada crimen, para decir\apnota{8}{M dicam.} de cada uno cómo deba enmendarse, sino que establecen dejarlo más bien al arbitrio del sacerdote inteligente\apnota{9}{M relinquentis.}, porque ante Dios no vale tanto la medida del tiempo como la del dolor.

A Agustín\apnota{11}{W. A. Augustinus.} lo refiere también en el can. Nullus: Ninguno recibe venia de la pena grave debida, si no ha pagado alguna pena, aunque mucho menor de la que deba; así en efecto la largueza de la misericordia es impartida\apnota{13}{M U impertitur a deo.} por Dios de modo que no se abandone la disciplina de la justicia. Aquí tiene el egregio señor D. Carlstadt a Agustín aducido en mi defensa, cuya sentencia no es de nuestro propósito proseguir al presente.

Añádase por último Isidoro, lib. 2 del Sumo Bien, cap. 13: Aunque por la penitencia haya propiciación de los pecados, sin embargo no debe ser sin temor, porque la satisfacción del penitente solo se pesa en el juicio divino, no en el humano. Por lo cual la penitencia impuesta no lo es porque se satisfaga\apnota{20}{U satisfacit.} solo a la Iglesia, sino también porque se satisface a Dios. Pues también ante Dios es más duradera la pena que la culpa, como atestigua Agustín en el tratado 124\apnota{22}{M tract. 134.} sobre Juan.

Habidas todas estas cosas por verdad indudable, fácilmente se hace patente la impugnación de la conclusión quinta, donde dice el reverendo padre que cualquier sacerdote debe absolver\apnota{25}{M U absolvere poenitentem.} de la pena y de la culpa, y así cualquier sacerdote aldeano sería en su parroquia obispo, arzobispo y papa. Que esto es falsísimo con toda evidencia\apnota{27}{M manifestissimum.}, además del uso de toda la Iglesia, no solo de los aduladores, se prueba primero por lo\apnota{28}{U iis.} que se ha dicho, porque por la absolución sacerdotal se borra la culpa, pero permanece la pena, aunque conmutada; después\apnota{29}{M dein.}, porque el sacramento de la penitencia es\apnota{30a}{M est quasi.} judicial, como describe la forma de este juicio el bienaventurado Agustín en el cap. quizá\apnota{30b}{fehlt in M U.} 2 de la medicina de la penitencia. Y a dictar juicio pertenece la jurisdicción, a no ser que quiera hacer en la Iglesia de Dios un caos anaxagórico y una confusión máxima.

Y como la jurisdicción se extiende a dos cosas materiales tenidas\apnota{33}{M habito.} en este sacramento, pues tanto el pecador penitente como los pecados mismos son de la materia de la penitencia, y nadie niega que la jurisdicción en los prelados inferiores, para quitar la confusión, esté coartada en cuanto a la materia de los que pecan, pues de ahí se tiene la determinación del sacerdote propio, y porque la sentencia no dada por su juez es nula, por eso el que absuelve a un no súbdito nada hace.

Así de modo semejante en la materia de los pecados puede coar- [376] tarse la jurisdicción por la misma razón y para detestación de los crímenes\apnota{1}{U detestatione.}. Y que tampoco pueda absolver plenamente de los que le han sido encomendados consta por que tan gran potestad, si no se funda en autoridad, contra el consenso de toda la Iglesia, cuya práctica debe ser para el hombre\apnota{4}{M omni.} cristiano como regla, no debe atribuirse. Pero en lo anterior se mostró por los clérigos mayores, Cipriano, Crisóstomo y Agustín, que el débito de la pena permanece después de la culpa. Por eso ceda el reverendo padre a la autoridad de los santos padres y obedezca a la costumbre\apnota{7}{M consuetudine.} de toda la Iglesia, o exponga por qué no deba hacerse así, por autoridad de la sagrada Escritura, de los concilios o de los santos padres.

Dijo al final el D. D. Eck que quería replicar no obstante, si le pareciera necesario, después de la respuesta del señor doctor Martín\apnota{9}{Diese Schlussbemerkung Ecks fehlt in M U.}.

Respondió el D. D. Martín\apnota{10}{U nur a. R. Martinus.}.

Concedo al señor doctor que tenga la última palabra, porque así lo quiere\apnota{11}{U a. R. Adduxit tamen se velle replicare, si dominus d. Eccius plus iusto responderet.}. Pero el egregio señor doctor argumenta contra mi conclusión queriendo probar que la pena es requerida por Dios, y por lo primero aduce la Escritura de Gn 3, donde el pecado de Adán aun remitida la culpa es castigado en la posteridad, luego permanece la pena pasando la culpa. Quiero al señor doctor ligado a este ejemplo, y que no replique si no confirma, o ceda en la sentencia\apnota{16}{P M U sententiam.}. Si esa pena es requerida por el pecado e igualmente ha de sentirse de cualquier pena, como infiere, tengo lo propuesto, que la pena no puede ser remitida por el papa ni por ningún sacerdote, porque ninguno remitió aún la muerte, los infinitos géneros de penas, enfermedades y miserias semejantes, que todas fueron infligidas por el primer pecado. Y así la prueba del señor doctor habrá probado\apnota{20}{M U probavit.} mi conclusión contra sí mismo\apnota{21}{U se ipsam, Uv se ipsum.}.

Segundo, que, traspasado el pecado de David, no fue traspasada sin embargo\apnota{22}{M tamen translata.} la pena, también hace por mí, como en el mismo sermón vulgar dije seguidamente\apnota{23}{M in eodem sermone vulgari frequenter.}: Porque era pena que Dios requirió, por eso nadie

la puede quitar\apnota{24}{M poterit, U potuit}. Si el señor doctor no me ha entendido, lea con más diligencia, pues yo dije y digo que esta pena que el papa o el hombre puede soltar Dios no la requiere\apnota{26}{M deus hanc poenam non requirit.}; pero la que él mismo requiere, a saber, inmediatamente\apnota{27a}{scil. immediate fehlt in M; U quam aut ipse requirit (scil. immediate loquitur).}, el hombre no la puede\apnota{27b}{P posset.} soltar.

Tercero, adujo la glosa: Dios borra el delito, pero no lo deja sin castigo; nada pugna contra mí. Pues se venga, ya\apnota{29}{U sine, Uv sive.} por el mismo hombre contritándolo, ya por la Iglesia enmendándolo, ya por sí mismo juzgando. Y este último y el primer género de pena no están en el arbitrio de ningún hombre, como dice el apóstol en 1 Cor 11 [v. 31 s.]: Si nos juzgáramos a nosotros mismos, ciertamente no seríamos juzgados por el Señor, pero cuando somos juzgados por el Señor somos corregidos, para no ser condenados con este mundo. Ahí está claro de qué modo Dios requiere la pena y no la requiere.

Cuarto, a Agustín referido por Graciano, que Dios no permitiría\apnota{35}{M pateretur.} que se hiciera el mal\apnota{36}{P mala.}, si por la justicia, ordenándolo, no hiciera que estuviera mejor: me admiro del egregio\apnota{37a}{fehlt in M U.} señor doctor, que zurce así estas y semejantes autoridades, cuando nadie las niega\apnota{37b}{M dum . . . . . neget; U negat.} ni siente lo contrario. Pues yo solo impugné siempre esto, que en virtud de las llaves se jactan de soltar las penas requeridas por la justicia divina. Pues esto no creo que sea verdad, ni se probará. Pues no dijo a Pedro: Lo que yo ligo, tú lo sueltas, sino: Todo lo que tú soltares, soltado será [Mt 16 v. 19].

Quinto, aquello de la medicina de la penitencia: No basta cambiar las costumbres a mejor\apnota{5a}{M U commutari.}, si no se satisface también a Dios por lo\apnota{5b}{U iis.} que se ha hecho, con el dolor de la penitencia, con el gemido de la humildad, con el sacrificio del corazón contrito. Todo esto yo siempre lo quise muchísimo, ¿para qué pues se jactan por las indulgencias de que esto se remite, si no basta cambiar las costumbres a mejor y la justicia divina requiere\apnota{8}{M U requirat.} estas satisfacciones? En vano se gloría pues el señor doctor de que no se dijo que\apnota{9}{U quod.} solo os abstengáis de los pecados, sino de los pasados rogad al Señor\apnota{10}{M Dominum etc.}. Y así consta que aún ninguna Escritura se ha aducido contra mí, aunque yo podría aducir muchas más fuertes en mi favor\apnota{12}{P U contra me; M multa fortiora introducere contra me.}.

Después de esto aduce a Agustín en el Enquiridión: Borra por completo esta oración los pecados mínimos y cotidianos, borra también aquellos de los que se aparta la vida de los fieles llevada criminalmente, pero cambiada a mejor arrepintiéndose\apnota{15}{M commutata discedit, U mutata etc.; discedit fehlt in P.}. Esto es lo que

dije, que el pecador después de la conversión está tenido a la cruz y a las pasiones de la vida, según Dios las infligiere, las cuales no puede soltar el hombre, y de nuevo esta autoridad hace por mí\apnota{17}{fehlt in M.}.

A los que fueron hace cuatrocientos años\apnota{18}{U qui a ecce (M quadragentis) annis fuerunt.} no los aduce, y me place.

A Cipriano, en la epístola a Fido, donde condena a Terapio\apnota{19}{P Fidium; M in (fehlt in U) epistola ad filum, ubi damnat Herapium.} porque dio precipitadamente\apnota{20a}{U dederat.} la paz a Víctor antes de que hubiera hecho\apnota{20b}{M fecisset.} plena penitencia y hubiera satisfecho al Señor: He aquí, dice, que no a la Iglesia, sino al Señor dice que Víctor no\apnota{22}{fehlt in M.} había satisfecho. Respondo: Lea y confronte bien Cipriano el señor doctor y hallará que incluso aquellos a quienes habían dado\apnota{23}{M dederam.} a la paz, por eso fueron dados a la paz aun precipitadamente, para que soportaran más expeditamente las cruces y los martirios. Los cuales él mismo por muchas epístolas expone que son penas y azotes por los pecados\apnota{25}{M flagellas, M U pro peccatis.} infligidos por Dios. Por lo cual Víctor de este modo aún no satisfizo a Dios; y sin embargo satisfizo, porque satisfizo a la Iglesia, a la cual Dios nos quiere oír. Pues esto tienen las palabras de Cristo, que la Iglesia debe imponer penas, cuando dice: Todo lo que ligares. Y de este modo podría admitir que Dios requiere penas que la Iglesia pueda soltar, porque hizo pacto con ella.

La otra autoridad de Cipriano a Cornelio, en la que de nuevo escribe que los herejes impiden a los que pecan que satisfagan a Dios indignado: Aunque Cipriano allí no habla de los penitentes\apnota{32}{U loquitur.}, sino de los herejes, que como justos y como si hubieran obrado bien se excusaron y defendieron en sus pecados, respondo sin embargo como a lo anterior\apnota{34}{M U priorem.}.

Después de esto añadió la razón: ¿Para qué querría la Iglesia imponer carga a los penitentes\apnota{36}{M poenitenti.}, si el Señor no la exigiera, e induce una larga autoridad después del cap. Alligant, todo lo cual paso por alto\apnota{37a}{M U non exigit et longam post (fehlt in U.) auctoritatem capituli allegat (U allegat capitulum) inducit quae (U quae inducit) omnia transeo.}, y digo que a la Iglesia le está mandado que\apnota{37b}{U ut.} castigue y juzgue a los pecadores, y si no lo hiciere, Dios lo hará irremisiblemente según la sentencia de Pablo arriba aducida de 1 Cor 11 (v. 31 s.\apnota{39}{M U adductam 1. Cor. 11 (U 1. Cor. 13, Uv 1. Cor. 11).}), y así no puede soltarse.

[378] Igualmente lo mismo de la Homilía 31 sobre los Hebreos\apnota{1}{de fehlt in M U; M homilia Chrysostomi, U = P.}: El pecado con la confesión se minora. De cuya palabra colige el señor doctor de este modo: No se minora en cuanto a la culpa, porque sería impío esperar de

Dios\apnota{3}{M sperare a deo.} media venia: luego en cuanto a la pena. Admito todo, según lo predicho.

Ya aquella autoridad de Gregorio\apnota{5}{fehlt in M.} 4 de los Morales: O nosotros lo perseguimos llorándolo, o él mismo juzgando; apenas pudo aducirse otra más apta en mi favor. Igualmente también aquel texto siguiente de Jerónimo, cap. Mensuram, del que sin embargo dudo que sea de Jerónimo, lo admito todo\apnota{8a}{M sequens Hieronymi c. mensura admitto totum, U Hieronymi c. mensuram admitto, quod . . . . . dubito. [Absatz.] Admitto totum.}, que los cánones no prefijan suficientemente, y por eso se deja\apnota{8b}{U relinquo.} al arbitrio del sacerdote. Y añado yo\apnota{9}{M at nego statt addo et ego.}: mucho más al arbitrio de Dios, que solo es ponderador de los espíritus y juzga no ignorante ni inicuo\apnota{10}{M inique indicat, U iniquus indicat.}.

Place también aquella autoridad de Agustín en\apnota{11}{fehlt in U.} el cap. Nullus: Así la largueza de la misericordia es impartida por Dios de modo que no\apnota{12}{M quod nulli (letzteres gestrichen, Mc 'non').} se abandone la disciplina de la justicia. Puede valer por mí y por el señor doctor\apnota{13}{M pro me et domino doct. esse, ideo transeo, U pro me et pro d. d. valere transeo.}.

E Isidoro, aunque\apnota{14}{M U quamvis.} no es autor bastante grave en estas cosas, place sin embargo que diga que la satisfacción del penitente solo se pesa en el juicio divino, no en el humano. Luego mucho menos puede remitirse en juicio humano, pues la llave de la potestad no debe\apnota{17}{M debet.} obrar si no es anterior la llave de la ciencia, que sepa qué y cuánto suelta\apnota{18}{M solveat.}.

Que ante Dios es más duradera la pena que la culpa, lo concedo por Agustín, salva sin embargo la autoridad del apóstol\apnota{20}{M U Pauli.} en Rom. 7 [v. 18] que dice: No hallo en mí, esto es\apnota{21}{M id est.}, en mi carne, el bien. Cuya sentencia es que la pena y el pecado se acaban a la vez.

Esto de la primera conclusión. — Contra\apnota{24}{U Sed contra.} la otra dice que es falsísimo con toda evidencia y fuera del uso de toda la Iglesia que cualquier sacerdote deba\apnota{25}{M U debet.} absolver de la pena y de la culpa al penitente. Y esto lo prueba primero por lo dicho, es decir, de la nada. Segundo por la razón de que el sacramento de la penitencia sea como un juicio, y a dictar juicio pertenece la jurisdicción. Y que la jurisdicción está coartada en los prelados inferiores para quitar la confusión, tanto en los que pecan como en los pecados, segundo para detestación de los crímenes, pues de otro modo cualquier sacerdote aldeano sería obispo, arzobispo y papa.

Respondo y digo dos cosas, primero\apnota{31}{U primum.} que no sé hasta el día de hoy si esta coartación de la jurisdicción\apnota{32}{M illius.} ha hecho lo que se pretende, a saber, la detestación de los crímenes\apnota{33}{M criminis.} y la remoción de la confusión. Lo cierto es que ha sucedido muy de otro modo, pues los pecados más crasos aun se ríen en las curias mayores, los cuales en las propias parroquias podrían castigarse muy bien, si se hubiera guardado el modo que instituyeron los apóstoles y guardaron los santos padres hasta después del concilio Niceno, donde se definió y mucho después se guardó que las diócesis no se mezclaran y cada uno hiciera penitencia en su diócesis. De [379] lo cual existe con otras la clarísima epístola de Cipriano, al romano pontífice Cornelio epístola III\apnota{2}{U pontificem 3.}: Pues habiéndose estatuido por todos nosotros y siendo igualmente equitativo y justo que la causa de cada uno se oiga allí\apnota{3}{M et illic.} donde el crimen fue admitido, y a cada pastor se haya adscrito una porción de la grey, que cada uno rija y gobierne, habiendo de dar razón de su\apnota{5}{M suam, verbessert sui.} acto al Señor, conviene en todo caso que aquellos a quienes presidimos no anden vagando ni choquen con su subdola y falaz temeridad la concordia coherente de los obispos\apnota{7}{U confidere, Uv collidere.}, sino que defiendan allí su causa donde puedan tenerse acusadores y testigos de su crimen\apnota{8a}{M U habere et testes . . . . . possint etc.}. Habla en efecto de los\apnota{8b}{U iis.} que habían pecado en África y habían corrido\apnota{9}{U peccaverunt; ad fehlt in M.} al romano pontífice Cornelio. Y así consta el uso de la primitiva Iglesia acerca de ligar y soltar a los pecadores.

Y como, según muestra el apóstol en Hch 20 [v. 17 y 28], lo mismo es lo que\apnota{11}{fehlt in U.} obispo y presbítero, y según Tit. 1 [v. 5] cualquier ciudad debe tener su obispo por derecho divino, mucho más útil sería\apnota{13}{fehlt in M.} para corregir los pecados el modo de que cualquier sacerdote en su parroquia ligara y soltara al penitente. Cuyo ejemplo mostró el apóstol en 1 Cor 5, donde estando con los corintios presente en espíritu entregó al fornicario a Satanás, reprendiéndoles que ellos no lo hubieran hecho.

Con qué derecho o con qué felicidad de la Iglesia haya sido suprimido este modo prescrito por derecho divino y robustecido por tanto tiempo, véanlo otros. Yo ciertamente no puedo negar\apnota{19a}{M negavi.} que vemos a simple vista que así se hace, así se estatuye, que\apnota{19b}{U quod.} tanto los que pecan como los pecados se reservan y a una sola alma se le imponen seis o siete pastores según, bajo y sobre. Pero si así deba hacerse o convenga, no lo defino\apnota{22}{M U diffinio.}. Sé que el inferior está tenido a obedecer al que restringe y veja, aunque por ningún derecho divino. El superior, sin embargo, como dice mi conclusión, peca gravísimamente si reserva los pecados ocultos sin causa razonabilísima, antes aún dudo, y cuanto alcanzo, creo sin temeridad, que ningún pecado oculto debe\apnota{26}{M deberi.} reservarse ni puede, deseando oír la prueba contraria.

Segundo digo que la Iglesia\apnota{27}{U ecclesiam, Uv ecclesia.} no se arruinaría si el mismo cura fuera obispo, arzobispo y papa y solo con la concordia coherente\apnota{28}{P hac sola statt ac sola; M adhaerente, Mc 'cohaerente'}, como dice Cipriano, se juntaran, como fue el uso de la Iglesia anterior.

Por tanto, que también se diga condenado el artículo\apnota{29}{M Proinde volo damnatos ostendere articulos.} en el concilio Constanciense acerca de estas\apnota{30}{M his.} reservaciones, no me importa bastante. Esto sé, que fue probado y reprobada esa reservación\apnota{31}{M reservatio huiusmodi.} en la primitiva Iglesia y en la institución de los apóstoles, y también ahora, como enseña la miserable experiencia de la Iglesia\apnota{32}{M experientia ecclesiae.}, sería utilísimo y saludabilísimo para cohibir los pecados y quitar la detestable confusión de todos los obispados que hoy vemos. Por qué causa haya sido mudado en su tiempo, lo paso por alto. Mudable lo hallo, mudable lo dejo. Esto\apnota{35}{W. A. verbessert Mutabilem invenio. Relinquo haec.} al arbitrio de los jueces\apnota{36}{U iudicum. Hora transit.}.

[380] El D. D. Eck\apnota{3}{U hat hier die Zeitangabe Quarta decima die Julii mane hora septima continuavit Eccius praeter pactum; vgl. Z. 8.}.

El reverendo padre desde el inicio con bastante jactancia tiene en poco lo por mí aducido, como si en nada le obstara, lo que máximamente se le opone. Pues él mismo quiso en el sermón vulgar que no se hiciera conmutación de la pena eterna en\apnota{6}{U in poenam; M aeternae mutatae in poenam.} temporal, y contradice la opinión común, como si no pudiera probarse\apnota{7}{M mutari.} que Dios exija alguna satisfacción fuera de la portación de la cruz. Y en el sermón latino alaba

Esto fue comenzado solo por el D. D. Eck hacia la hora cuarta vespertina, pero interrumpido y reasumido en el exordio de la disputa del día siguiente, a saber

a la hora séptima antes del mediodía del jueves 14 de julio\apnota{8}{Die Zwischenbemerkung fehlt in M; an der Schrift ist der neue Ansatz zu merken, allerdings schon hinter sermone vulgari. U setzt die Zeitangabe an den Anfang der Rede, vgl Z. 3; continuavit wird dadurch sinnlos.}.

el dicho de los vulgares más allá de toda doctrina de los doctores escolásticos dada sobre la penitencia: NVMMER THVN DIE HOCHSTE BVES\apnota{10a}{P CHVN statt THVN; M nummer thun die hochste busse; U nymmer thuen die hochste puess.}. La óptima penitencia es la vida nueva\apnota{10b}{M secundum vitam, Mc 'nova vita'.}, según su glosa. Añadieron tal persuasión él mismo y el propugnador: si se apartare el impío etc.\apnota{12}{U quia siaverterit se impius a iustitia (so!) sua etc.}. Si imputa a nueva pena, ¿cómo se dice\apnota{13}{M dicit.} que no se acuerda? y consta su conclusión que dice manifiestamente que peca el sacerdote que no ababsuelva de la culpa y de la pena. A este error me opuse yo con las autoridades de los santos padres, de las que él quiso burlarse y con las que quiso fascinar a los oyentes, como si solo en él residiese la potestad de interpretar la Sagrada Escritura.

Por eso aduje con razón contra él a Agustín: No basta cambiar las costumbres a mejor y apartarse de las malas acciones\apnota{19}{M U recedere etc.}. Donde resulta evidentísimo que la doctrina alegada de su sermón\apnota{20}{M ex 'eius' sermone.} es falsa, pues la vida nueva no es la mejor penitencia, ya que no basta según Agustín. Y con razón también aquella doctrina\apnota{22}{M U doctrina illa.} es refutada por el bienaventurado Ambrosio, libro 2 sobre la penitencia, cap. 5: Los apóstoles enseñaron la penitencia según el magisterio de Cristo. Y más abajo: Quien en efecto hace\apnota{23}{M ait.} penitencia, no solo debe lavar con lágrimas su pecado, sino también cubrir y tapar sus pecados con hechos más enmendados. Clarísimas palabras del santo padre, a saber, que haciendo penitencia debemos tapar los pecados con hechos más enmendados.

Mas como las autoridades eran expresísimas, a saber, que Dios no perdona el pecado sin castigo, se refugió en una admirable distinción de las penas, él que sin embargo suele reprobar a los escolásticos por el uso de distinciones, y dijo que la pena con que Dios quiere castigar el pecado no puede\apnota{30}{M possit, U posset.} ser quitada por el hombre ni por el papa, lo cual es falsísimo y anula la potestad de las llaves. Pues en la penitencia, hecha la conmutación de la pena, ciertamente el hombre puede pagar aquella pena, según el mismo Apóstol, aducido por el reverendo padre, 1 Cor 11: Si nos juzgásemos a nosotros mismos, ciertamente no seríamos juzgados por el Señor. Por lo cual, si satisfacemos por esta pena, Dios no exige de nosotros otra por el pecado\apnota{35}{M pro hoc peccato aliam non exigat.}; de lo contrario, contra el profeta, lo que nosotros castigáramos, si Dios quisiera castigarlo, ya lo castigaría dos veces.

[381] Además fueron claras las palabras de Cipriano, Crisóstomo, Gregorio y\apnota{1}{U Grego: 4 statt Greg. et.} Jerónimo, a saber, que la pena impuesta por el sacerdote en la satisfacción\apnota{2}{M iuncta a sacerdoti; M U in satisfactionem.} es pena debida a Dios, y así de Víctor, por no haber cumplido la penitencia impuesta, se dice según Cipriano que aún no había satisfecho a Dios. Lo cual también atestigua Teodoro en su penitencial, al que\apnota{4}{M quam.} sigue casi al pie de la letra Beda, aunque quizá tampoco estos\apnota{5}{M quamvis forte hi.} autores, como tampoco Isidoro, le satisfacen\apnota{6}{M U satisfacit.}.

De donde, diga lo que diga de palabra, que todas las autoridades están a su favor, es sin embargo necesario que en su corazón disienta de ellas: pues estas\apnota{8}{M haec.} autoridades subvierten su doctrina y su conclusión. Pues si el sacerdote peca no absolviendo de la pena y de la culpa, entonces los obispos habrían pecado no absolviendo a Víctor de la pena y de la culpa, y todos los sacerdotes pecarían por todo el orbe cristiano no absolviendo\apnota{11}{non fehlt in P M U; M absolvendo.} fuera de las indulgencias.

Ni se libra con ello, a saber, que deba llevar la cruz y que Dios exija esta pena, pues ese llevar la cruz no es otra cosa que la vida cristiana, como el mismo reverendo padre doctamente lo expuso; sino que con esto conviene satisfacer por lo pasado y suplicar al Señor por lo pasado.

Sobre la reserva de los casos mi conclusión nada contiene; sin embargo creo que las reservas moderadas son útiles, como también lo experimentan los prelados en los monasterios. Confieso ingenuamente que estoy en el voto de Gerson, que en el concilio\apnota{19}{U consilio.} de Constanza interpeló a favor de suprimir la inmoderada reserva de los casos, y sobre todo me desagrada aquella reserva, como también al señor padre, cuando lleva por compañera la avaricia, esto es, la pena pecuniaria\apnota{22}{U pecuniarum.} aneja.

No dice que la Iglesia vaya a arruinarse si el sacerdote fuese en su parroquia obispo y papa; ciertamente me parece que, si hubiese corrupción de aquel hermosísimo orden jerárquico, también habría\apnota{24}{M esset etiam.} subversión de la Iglesia.

Pero permanezca yo en lo principal, pues el pecado no queda sin castigo, según Agustín, Gregorio etc.\apnota{26}{fehlt in U.}. Por lo cual con mérito se afirma que la satisfacción es la tercera parte de la penitencia. Y Agustín lo probó plenamente en el Enchiridion, a saber, que por la oración dominical satisfacemos por lo cometido criminalmente. Y como ayer dijo con verdad el reverendo padre, Dios hizo un pacto con la Iglesia: si la Iglesia no lo hace, Dios lo hace. Tengo, pues, el parecer de los escolásticos y de los predicadores, a saber, que conviene o que nosotros satisfagamos, o que Dios lo exija. Y si por nuestras oraciones o buenas obras satisfacemos, ¿por qué no podría hacerse esto accediendo la virtud de las llaves, que Dios no en vano confirió a su esposa la Iglesia?\apnota{33}{M tradidit.} Y puesto que, según la sentencia de Gregorio, cap. Decreto 2, q. 6,\apnota{34}{P M U de cetero 2 q. 7.} los demás obispos han sido llamados a participar de la solicitud, teniendo el papa la plenitud de la potestad, por las indulgencias dadas por él se satisface a la pena debida por los pecados con el pago hecho del tesoro de la Iglesia\apnota{36}{M factae.}, como después de Sixto declaró el papa actual, mandando bajo pena de excomunión que así se enseñe\apnota{37}{fehlt in M.}, [382] se sostenga y se predique. Por lo cual, si el reverendo padre enseñare, predicare o disputare en contra, ya ha sido alcanzado por el anatema. Me plugo, sin embargo, que además de la llave de la potestad pusiera ayer la llave de la discreción, aunque niega varias llaves contra los doctores escolásticos en la conclusión 7 del Resolutorio.

Esto quise añadir en todo caso, para que los futuros jueces entendieran más abiertamente cuál fuese nuestra controversia en este punto, pues, si el reverendo padre está por la sentencia de los doctores por mí alegados, no se opondrá ni a los escolásticos ni a los predicadores ni a mí. Por lo cual, si quisiere, también

podrá expresar más claramente su sentencia para información de los señores jueces.

El D. D. Martín
sin subir a la cátedra, sino desde el lugar acostumbrado de asiento
levantándose respondió\apnota{10}{M U Martinus (die Bemerkung fehlt).}.

A estas necedades e ineptitudes del señor doctor de ayer ya respondí bastante; pues repite lo mismo y, como ridículo citarista, yerra siempre en la misma cuerda.

En segundo lugar, no ha alcanzado el objeto de la controversia. Pues la cuestión no es si Dios deja el pecado sin venganza, lo cual probó\apnota{14}{U probavi, M probatum est.} con harta abundancia, sino si el papa o la Iglesia remiten las penas que Dios exige. Sobre esto nada probó, lo cual dejo al juicio de los jueces y de todos los oyentes.

En tercer lugar, hoy ha callado las Santas Escrituras; por eso me atengo a su primera prueba\apnota{17}{M probatione eius.} de ayer, del cap. Gén 3,\apnota{18}{M U ex Genes. 3 (U 30) ca.} donde probó que Dios requiere penas, las cuales muestra allí la Escritura irremisibles. Lamento que el D. D. penetre tan profundamente las Escrituras como el mosquito las aguas; más aún, parece huir de su faz como el diablo de la cruz. Por lo cual, salvas las reverencias de los padres, prefiero yo la autoridad\apnota{21}{M auctoritates} de la Escritura, lo cual encomiendo a los futuros jueces\apnota{22}{'M iudicio 'iudicum'.}.

El D. D. Eck.

Porque el impaciente monje añadió\apnota{24}{M scurrilia quaedam addit.} ciertas burlas aparte de la gravedad teológica, sobre lo cual juzguen los varones íntegros\apnota{25}{M in quae (Mc 'de quo') 'interim' iudicem (verbessert iudicent, folgt ein gestrichenes Wort).} si rectamente aduje contra él, los jueces juzgarán. Pero que esta fue nuestra materia consta por la conclusión 4: decir que Dios, al remitir la culpa, remite la pena, etc. Esta era la piedra que había que remover\apnota{28}{volvendum fehlt in M, Lücke ausgefüllt Mc per nos solvendam; U per nos volvendum.}, y puesto que prefiere la autoridad de la Sagrada Escritura a los padres, como si él solo, cual otro oráculo de Apolo, tuviese la inteligencia de las Escrituras\apnota{29}{M solus ipse habeat scripturam.} por encima de los santos padres, y aduce en apariencia la autoridad ayer citada por mí, digo dos cosas.

Primero, que la aduje con el fin de mostrar que, cuando el impío se hubiese apartado de su injusticia, Dios aún se acuerda, no en cuanto a la culpa, sino en cuanto a la pena, lo cual está a favor de los escolásticos y de los predicadores.

Segundo, cuando el reverendo padre vuelve contra mí aquella autoridad, porque las penas impuestas por Dios al pecado de Adán sean irremisibles por el papa y el hombre, es verdad, y lo confieso, pues estas son penas que

no solo alcanzan a la per-[383]sona, sino que alcanzan también a la naturaleza, y por eso no es extraño que estas penas no sean remisibles\apnota{2}{U sunt irremissibiles.} por el hombre; pero con esto no se prueba que las penas personales debidas por el pecado no sean remisibles por el sacerdote o el papa\apnota{3}{U a papa vel sacerdote.}. Mas en esto me remito\apnota{4}{U remitto.} a los jueces, dispuesto a cambiar de sentencia si de otro modo y mejor\apnota{5}{M aliter in melius, U aliter ad melius.} me enseñaren\apnota{*}{In P folgt hier die Bemerkung: Et sic imposita fuit extrema manus disputationi inter E. D. Doctores praenominatos, inchoato deinceps evestigio, imo instaurata rursus disceptatione inter praestantissimos D. Doctores, videlicet eundem Johannem Eccium et Andream Carolostadium in prima fronte supra memoratos etc. Den dann folgenden actenmässigen Schluss des Protokolls s. unten. — U hat Finita sunt haec 14. die Julii hora octava praesente frequenti concione auditorum. Soli Deo honor et gloria. Anno M. D. xix.}.

\noindent\textit{Notas del aparato sin llamada en el cuerpo:} [16] M sententias.; [27] M quodam.

\end{document}
